% Generated by build_chapter.py from the GY1974 files. Do not edit by hand; edit the files and rebuild.
\chapter{The academic year 1974--75}\label{ch:1974}
\noindent
This chapter indexes the 30 catalogue records assigned to 1974--75 in the supplied
package, with 265 mathematical occurrences, 10 definite problem entries,
78 bibliography entries and 92 critical corrections or qualifications.
Records retain their catalogue order; items follow the surviving text within
each record. Status statements refer to November 2026. The classifications
describe historical fate; the separate assessments describe literal correctness,
publication matching and verification.

\begin{center}
\begin{tabular}{@{}lrrrrrr@{}}
\toprule
Class & A & B & C & D & E & F\\
\midrule
Occurrences & 29 & 10 & 0 & 0 & 203 & 23\\
\bottomrule
\end{tabular}
\end{center}

\paragraph{The main mathematical lines.}
The finite-gap sequence connects Floquet theory, Riccati densities, stationary
KdV equations, hyperelliptic curves and Jacobians. Its chronology matters:
Novikov's paper was received in February and published in July 1974, before
the September meetings. The joint Dubrovin--Novikov announcement, received in
June and published later in 1974, gives the converse and algebraic-geometric
synthesis. The later survey supplies fuller hypotheses, not the first
circulation of all its results
\citep{novikov1974periodic,dubrovinnovikov1974periodicity,dubrovinmatveevnovikov1976nonlinear}.
Kushnirenko's local Milnor-number and global critical-point questions have
different inequality directions and nondegeneracy conditions
\citep[Theorems I--II]{kushnirenko1975newton}.

Giusti's lectures juxtapose scalar regularity, vectorial counterexamples and
relaxed variational problems. Segal moves from configurations and cancellation
to nerves, buildings and classifying spaces; the categorical material continues
1973-308 and 1973-309. Graev's main construction explains work already published
in 1973. Tulovskii's sequence is anchored in the 1971--72 FIO theory
\citep{hormander1971fourier,duistermaathormander1972fourierii}.
Sinai's phase-diagram lecture has a close subsequent Pirogov--Sinai publication,
whereas Dobrushin's conditional-law formulation predates the talk
\citep{pirogovsinai1975phase,dobrushin1968conditional}.
King's lectures explain divisor currents, counting functions and curvature
inequalities in relation to existing holomorphic-mapping theory
\citep{griffithsking1973nevanlinna,griffiths1971canonical}.

\paragraph{What the critical checks change.}
The notebook is often correct shorthand, but hypotheses and conventions cannot
be discarded. On P000555 the Hamiltonian quartic coefficient needs \(5/8\);
the copied \(5/2\) gives a nonzero Poisson bracket. Its spectral quintic has
a further normalization problem. Giusti's copied tensor on P000565 fails
the claimed radial-solution substitution. The two logarithmic integrals on
P000564 behave differently: the first has the wrong sign, while the symmetric
golden-ratio integral is correct. The Steinberg module requires reduced
homology at \(n=2\). Fourier signs, radial characteristics, Haar measures and
the divisor-counting boundary term provide further convention-sensitive checks.

\paragraph{Questions and the scope of development.}
No item is assigned C or D, and no Stage 2 research problem meets the selection
rule. This does not imply that the surrounding subjects contain no open problems.
Definite questions have established answers in specified settings; several
prompts lack enough data for an exact formulation. Entire functions of order
one, an elliptic analogue of resolution and the \(p\)-adic analogies are kept
separate from later thematic work. The verification section gives complete
local arguments, counterexamples to overbroad wording and independently checked
computations. These are editorial audits of old mathematics, without novelty claims.

\paragraph{Source boundaries and earlier years.}
P000551 is headed 25 March 1974, although records 331--332 are catalogued as
25 March 1975. Their five entries duplicate 1973-331.1--4 and 1973-332.1.
The requested identifiers and catalogue dates are retained as compatibility
aliases with explicit warnings. Removing these exact source duplicates leaves
260 occurrences newly indexed by this package. P000550 belongs to the preceding
finite-gap talk; P000552--553 belong to the 1 April 1974 spectral lecture.
Their earlier-index coverage is audited separately. Further links include
1967-132.7 on the ray space, 1973-326 on finite-gap inverse data and monotonicity,
and the current-group questions of 1969. Related problems are not merged when
the earlier index supplies only titles and cannot establish identical hypotheses.

\GStalk{331}{1975-03-25}{A. A. Kirillov}{P000551}{Diffeomorphism groups, gauge groups and exponential Hilbert space}
The catalogue date is retained as supplied, but the manuscript heading on P000551 reads \(25/\mathrm{III}\text{-}74\). These notes repeat the material already indexed as 1973-331.1--4; the chronology of the catalogue entry is unresolved. The mathematical content consists of a representation-theoretic programme, the automorphism group of a principal bundle, and symmetric Fock space. P000550 and P000552 belong to adjacent talks, documented separately in the boundary audit. The current-year identifiers are compatibility aliases for the five existing source occurrences 1973-331.1--4 and 1973-332.1, not newly dated mathematical occurrences. The catalogue all-links page independently retains the 1975 date; the manuscript and prior-index evidence still leave the catalogue chronology inconsistent.

\begin{GSitem}{1974-331.1}{F}{Representation programme for diffeomorphism and current groups}
\GSfield{Kind; source}{remark; \GSpage{P000551}}
\GSfield{The notes}{The groups \(G_1=\operatorname{Diff}M\) and \(G_2=H^M\), with smooth maps to a Lie group \(H\), are introduced in connection with describing representations.}
\GSfield{Statement}{For a smooth manifold \(M\) and Lie group \(H\), the current group \(C^\infty(M,H)\) has pointwise multiplication. No representation class, topology, irreducibility criterion or classification theorem is recorded.}
\GSfield{Reconstruction}{This item is a compatibility alias for the identically numbered 1973 occurrence; it is counted once when counting distinct source occurrences.}
\GSfield{Fate}{\citet[title and abstract]{vershikgelfandgraev1975diffeomorphisms} is a relevant later programme, received 15 May 1975 and published in 1975; it does not identify an unrecorded Kirillov theorem. Earlier occurrence: 1973-331.1.}
\GSfield{Present status}{As of November 2026, the source does not determine a sufficiently precise claim for an open-problem status.}
\GSfield{Assessment}{literal statement: not determinable; publication match: thematic only; verification: publisher metadata checked; development: none; confidence: medium}
\end{GSitem}
\begin{GSitem}{1974-331.2}{E}{Automorphisms of a principal bundle and the gauge extension}
\GSfield{Kind; source}{construction; \GSpage{P000551}}
\GSfield{The notes}{For a principal \(H\)-bundle \(P\to M\), its automorphisms are said to project onto \(G_1\) with fibre \(G_2\).}
\GSfield{Statement}{There is an exact sequence \(1\to\operatorname{Gau}(P)\to\operatorname{Aut}(P)\to\operatorname{Diff}(M)_P\to1\), where \(\operatorname{Gau}(P)=\Gamma(M,\operatorname{Ad}P)\) and \(\operatorname{Diff}(M)_P=\{f:f^*P\simeq P\}\). For the trivial bundle, the kernel is \(C^\infty(M,H)\) and all base diffeomorphisms lift.}
\GSfield{Reconstruction}{The bundle-preserving subgroup and the adjoint bundle are necessary additions; neither is specified in the shorthand. This item is a compatibility alias for the identically numbered 1973 occurrence; it is counted once when counting distinct source occurrences.}
\GSfield{Fate}{This is the standard bundle-automorphism construction; \citet[abstract and extension sequence]{wockel2007symmetry} formulates it for smooth symmetry groups (preprint December 2006; publication 2007). The construction itself is elementary and predates the talk. Earlier occurrence: 1973-331.2.}
\GSfield{Present status}{As of November 2026, the stated reconstruction is a standard established result.}
\GSfield{Assessment}{literal statement: false as written; publication match: exact; verification: read the result; development: annotated; confidence: high}
\end{GSitem}
\begin{GSitem}{1974-331.3}{E}{Symmetric Fock space}
\GSfield{Kind; source}{construction; \GSpage{P000551}}
\GSfield{The notes}{The exponential space is displayed as \(e^H=H_0\oplus H\oplus S^2H\oplus\cdots\), with one-dimensional \(H_0\).}
\GSfield{Statement}{For a complex Hilbert space \(H\), the bosonic Fock space is \(\mathcal F_s(H)=\bigoplus_{n=0}^{\infty}\operatorname{Sym}^nH\), the Hilbert direct sum with \(\operatorname{Sym}^0H=\mathbb C\).}
\GSfield{Reconstruction}{This item is a compatibility alias for the identically numbered 1973 occurrence; it is counted once when counting distinct source occurrences.}
\GSfield{Fate}{\citet[\S4, pp.~370--373]{araki1970factorizable} uses the exponential Hilbert space in current-group representations; received June 1969, revised November 1969, published 1969/70. Earlier occurrence: 1973-331.3.}
\GSfield{Present status}{As of November 2026, the stated reconstruction is a standard established result.}
\GSfield{Assessment}{literal statement: correct; publication match: exact; verification: read the result; development: annotated; confidence: high}
\end{GSitem}
\begin{GSitem}{1974-331.4}{E}{Exponential vectors and factorial normalization}
\GSfield{Kind; source}{definition; \GSpage{P000551}}
\GSfield{The notes}{Tensor coordinates containing \(\xi^{\otimes n}/n!\) occur beside a norm with weights \(1/n!\). The tensor-product normalization is not stated.}
\GSfield{Statement}{With the usual Hilbert norms on symmetric tensor powers, \(\varepsilon(\xi)=(\xi^{\otimes n}/\sqrt{n!})_{n\ge0}\) satisfies \(\langle\varepsilon(\xi),\varepsilon(\eta)\rangle=e^{\langle\xi,\eta\rangle}\). An equivalent convention places \(1/n!\) in the sector norm and uses coordinates \(\xi^{\otimes n}\).}
\GSfield{Reconstruction}{The two consistent conventions above are an editorial normalization, not an assertion that one specific tensor norm was intended. This item is a compatibility alias for the identically numbered 1973 occurrence; it is counted once when counting distinct source occurrences.}
\GSfield{Fate}{\citet[\S4]{araki1970factorizable} supplies the established exponential-space convention. The displayed factors cannot be combined without specifying a rescaling of each symmetric power. Earlier occurrence: 1973-331.4.}
\GSfield{Present status}{As of November 2026, the stated reconstruction is a standard established result.}
\GSfield{Assessment}{literal statement: not determinable; publication match: partial; verification: read the result; development: annotated; confidence: medium}
\end{GSitem}

\GStalk{332}{1975-03-25}{B. M. Gurevich (as catalogued)}{P000551}{KdV by Whitham’s method: heading only}
Only the heading \ru{Гуревич. О КдФ по методу Уизема} is recorded. The catalogue date 25 March 1975 conflicts with the manuscript date 25 March 1974, as in record 331. P000552--553 instead contain a separately dated spectral talk. The catalogue initials B.~M. are not silently replaced by A.~V., the author of the well-known Gurevich--Pitaevskii paper. The current-year identifiers are compatibility aliases for the five existing source occurrences 1973-331.1--4 and 1973-332.1, not newly dated mathematical occurrences. The catalogue all-links page independently retains the 1975 date; the manuscript and prior-index evidence still leave the catalogue chronology inconsistent.

\begin{GSitem}{1974-332.1}{F}{An unrecorded account of Whitham’s method for KdV}
\GSfield{Kind; source}{remark; \GSpage{P000551}}
\GSfield{The notes}{The speaker heading announces KdV by Whitham\textquotesingle s method; no equation, theorem or argument follows.}
\GSfield{Statement}{No mathematically definite assertion can be recovered from this heading alone.}
\GSfield{Reconstruction}{This item is a compatibility alias for the identically numbered 1973 occurrence; it is counted once when counting distinct source occurrences.}
\GSfield{Fate}{\citet[\S1]{gurevichpitaevskii1974shock} is a thematic comparator: submitted 3 April 1973, Russian publication August 1973, English translation February 1974. Its author is A.~V.~Gurevich, whereas this catalogue gives B.~M.; identity and contents are therefore not inferred. Earlier occurrence: 1973-332.1.}
\GSfield{Present status}{As of November 2026, the source does not determine a sufficiently precise claim for an open-problem status.}
\GSfield{Assessment}{literal statement: not determinable; publication match: thematic only; verification: read the result; development: none; confidence: low}
\end{GSitem}

\GStalk{334}{}{}{P000554}{Academic-year heading: 1974}
The year heading introduces the seminar\textquotesingle s thirty-second year. It has no separate mathematical assertion. The first dated meeting below it is 2 September 1974. The catalogue\textquotesingle s two-page range is a scan grouping, not a second lecture.


\GStalk{335}{1974-09-02}{Enrico Giusti}{P000554}{Scalar regularity, elliptic systems and mean curvature}
The introductory discussion distinguishes scalar regularity from the possible singularities of elliptic systems, then turns to prescribed mean curvature. The hypotheses are highly abbreviated. The detailed nonlinear-system example is deferred to record 341; the mean-curvature existence discussion continues in record 338.

\begin{GSitem}{1974-335.1}{E}{Scalar Hölder regularity for measurable uniformly elliptic coefficients}
\GSfield{Kind; source}{theorem; \GSpage{P000554}}
\GSfield{The notes}{The names Nash, De Giorgi and Moser accompany the assertion of regularity for one equation; the parenthesis mentions Hölder coefficients.}
\GSfield{Statement}{Let \(u\in W^{1,2}_{\mathrm{loc}}(\Omega)\) solve \(\operatorname{div}(A(x)Du)=0\), where \(A\) is bounded measurable and \(\nu|\xi|^2\le \xi^TA(x)\xi\), \(|A(x)|\le\Lambda\). Then \(u\) is locally Hölder continuous. Hölder continuity of \(A\) is unnecessary for this conclusion.}
\GSfield{Reconstruction}{The energy class, local conclusion and uniform bounds have been supplied explicitly.}
\GSfield{Fate}{\citet[\S2, scalar regularity discussion]{mingione2006regularity} records the classical De Giorgi--Nash theory; the central results date from the 1950s, before this talk.}
\GSfield{Present status}{As of November 2026, the stated reconstruction is a standard established result.}
\GSfield{Assessment}{literal statement: correct with missing hypotheses; publication match: exact; verification: read the result; development: annotated; confidence: high}
\end{GSitem}
\begin{GSitem}{1974-335.2}{E}{Using linear theory in scalar nonlinear regularity}
\GSfield{Kind; source}{claim; \GSpage{P000554}}
\GSfield{The notes}{The scalar nonlinear situation is described as following the same regularity scheme.}
\GSfield{Statement}{For a uniformly elliptic quasilinear scalar equation, regularity can be bootstrapped once the solution and its gradient lie in a range where the coefficients and their required derivatives are controlled. Smoothness of the coefficient formula alone does not impose those bounds.}
\GSfield{Reconstruction}{No single nonlinear equation is specified. The statement records the valid method and its missing conditions, rather than choosing an unrecorded theorem.}
\GSfield{Fate}{\citet[pp.~111--113]{giusti1982basic} explicitly distinguishes growth hypotheses, minimizers and merely stationary solutions. This is an expository methodological statement, not an unrestricted smoothness theorem.}
\GSfield{Present status}{As of November 2026, the stated reconstruction is a standard established result.}
\GSfield{Assessment}{literal statement: correct with missing hypotheses; publication match: partial; verification: read the result; development: annotated; confidence: medium}
\end{GSitem}
\begin{GSitem}{1974-335.3}{E}{Failure of scalar regularity for elliptic systems}
\GSfield{Kind; source}{claim; \GSpage{P000554}}
\GSfield{The notes}{Giusti and Miranda are associated with an example of a nonregular solution despite analytic coefficients.}
\GSfield{Statement}{There are uniformly elliptic quasilinear systems with real-analytic dependence on the unknown and discontinuous energy-class weak solutions. One corrected concrete example for \(n=N\ge3\) is verified in the companion\textquotesingle s coefficient audit for record 341.}
\GSfield{Fate}{\citet[\S3.2, equations (3.3)--(3.5)]{mingione2006regularity} attributes the analytic-coefficient example to Giusti--Miranda (1968). The original 1968 paper was not retrieved; the later formula and the residual calculation are independently available.}
\GSfield{Present status}{As of November 2026, the phenomenon is established; the literal tensor later copied on P000565 needs a coefficient emendation.}
\GSfield{Assessment}{literal statement: correct with missing hypotheses; publication match: partial; verification: indirect; development: annotated; confidence: high}
\end{GSitem}
\begin{GSitem}{1974-335.4}{E}{Partial regularity and the dimension of the singular set}
\GSfield{Kind; source}{theorem; \GSpage{P000554}}
\GSfield{The notes}{The singular set is stated to be closed and to have dimension less than \(n-2\).}
\GSfield{Statement}{For \(n\ge3\), bounded uniformly strongly elliptic \(A(x,v)\), uniformly continuous in \((x,v)\), and an energy weak solution of \(\operatorname{div}(A(x,u)Du)=0\), there is an open regular set on which \(u\) is locally Hölder of every exponent below one; its relatively closed complement has Hausdorff dimension strictly below \(n-2\). The assertion concerns regularity of \(u\), not an unrestricted gradient singular set.}
\GSfield{Reconstruction}{The target is partial regularity of the weak solution, not a universal bound for every nonlinear elliptic system.}
\GSfield{Fate}{\citet[\S1.1, equations (1.1)--(1.2)]{defilippismingionenowak2026partial} states the exact classical local bound and identifies its 1968 foundations. The original arbitrary-order paper \citet{giusti1969partial} was opened but its formula extraction did not permit full verification. The strict inequality is not deduced merely from zero \((n-2)\)-dimensional measure.}
\GSfield{Present status}{As of November 2026, the reconstructed second-order theorem is established; the coefficient and energy hypotheses must be retained.}
\GSfield{Assessment}{literal statement: correct with missing hypotheses; publication match: exact; verification: indirect; development: annotated; confidence: high}
\end{GSitem}
\begin{GSitem}{1974-335.5}{E}{Mean curvature of a graph and a size obstruction}
\GSfield{Kind; source}{definition; \GSpage{P000554}}
\GSfield{The notes}{The graph \(x\mapsto u(x)\) is assigned mean curvature \(n^{-1}\operatorname{div}(Du/\sqrt{1+|Du|^2})\); prescribed curvature cannot be arbitrarily large.}
\GSfield{Statement}{For an oriented \(C^2\) graph over \(\Omega\subset\mathbb R^n\), the signed average principal curvature is \(n^{-1}\operatorname{div}(Du/\sqrt{1+|Du|^2})\), up to orientation. Integrating the divergence over a subdomain bounds total prescribed curvature by boundary area.}
\GSfield{Fate}{\citet[\S1B]{giusti1976boundary} develops the associated perimeter condition. The flux obstruction is elementary and already part of the pre-1974 mean-curvature existence theory; record 338 gives its explicit form.}
\GSfield{Present status}{As of November 2026, the stated reconstruction is a standard established result.}
\GSfield{Assessment}{literal statement: correct with missing hypotheses; publication match: exact; verification: read the result; development: annotated; confidence: high}
\end{GSitem}

\GStalk{336}{1974-09-02}{S. P. Novikov; B. A. Dubrovin}{P000554-P000556}{Algebraic tori, commuting Hamiltonians and finite-gap spectra}
The talk links Jacobians of hyperelliptic curves to the stationary KdV hierarchy. It contains an explicit genus-two Hamiltonian construction, the Abel map, and the spectral interpretation of Bloch functions. The algebraic statements require a polarization, a nonsingular spectral curve and compactification of affine levels. Two polynomial coefficients and the claim that the Abel map is globally an isomorphism require separate scrutiny. The broad question about describing all tori does not specify a modern moduli problem.

\begin{GSitem}{1974-336.1}{E}{Complex tori and algebraic tori}
\GSfield{Kind; source}{definition; \GSpage{P000554}}
\GSfield{The notes}{A torus is written \(\mathbb C^n/\Lambda\), where \(\Lambda\) has a basis of \(2n\) vectors; algebraic tori are called abelian varieties.}
\GSfield{Statement}{A complex torus is \(\mathbb C^g/\Lambda\) for a discrete lattice of rank \(2g\) spanning the real vector space. It is an abelian variety precisely when it is projective, equivalently when it admits a positive Riemann form. Here “algebraic torus” means a projective complex torus, not the affine group \((\mathbb C^*)^g\).}
\GSfield{Fate}{\citet[complex tori and Riemann forms]{milne2008abelian} gives the classical criterion. The terminology in these notes is historical and should not be confused with an algebraic group torus.}
\GSfield{Present status}{As of November 2026, the stated reconstruction is a standard established result.}
\GSfield{Assessment}{literal statement: correct with missing hypotheses; publication match: exact; verification: read the result; development: annotated; confidence: high}
\end{GSitem}
\begin{GSitem}{1974-336.2}{F}{The unspecified problem of describing all tori}
\GSfield{Kind; source}{question; \GSpage{P000554}}
\GSfield{The notes}{Novikov asks how to describe all tori, noting that dimension two is already nontrivial.}
\GSfield{Statement}{The notebook does not say whether tori are polarized, whether isomorphism or isogeny is intended, or whether the aim is a parameter space, a Jacobian criterion or explicit functions. No unique definite problem follows.}
\GSfield{Fate}{The marked hyperelliptic-Jacobian construction of \citet[\S4]{dubrovinnovikov1974periodicity} is a precise nearby result, submitted 3 June 1974 and published later in 1974. It does not settle every possible interpretation of the question.}
\GSfield{Present status}{As of November 2026, the source does not determine a sufficiently precise claim for an open-problem status.}
\GSfield{Assessment}{literal statement: not determinable; publication match: thematic only; verification: read the result; development: none; confidence: low}
\end{GSitem}
\begin{GSitem}{1974-336.3}{E}{Hyperelliptic curves and their holomorphic differentials}
\GSfield{Kind; source}{construction; \GSpage{P000554}}
\GSfield{The notes}{The curve \(y^2=P_{2g+1}(x)\) is used with the forms \(x^kdx/y\).}
\GSfield{Statement}{If \(P_{2g+1}\in\mathbb C[x]\) has distinct roots, the smooth projective compactification of \(y^2=P_{2g+1}(x)\) has genus \(g\); \(x^kdx/y\), \(0\le k\le g-1\), is a basis of holomorphic differentials. Repeated roots require normalization and can lower the genus.}
\GSfield{Reconstruction}{Square-freeness, compactification and the index range have been made explicit.}
\GSfield{Fate}{This classical construction is the input to \citet[\S3]{dubrovinnovikov1974periodicity}, not a new assertion of that paper.}
\GSfield{Present status}{As of November 2026, the stated reconstruction is a standard established result.}
\GSfield{Assessment}{literal statement: correct with missing hypotheses; publication match: exact; verification: read the result; development: annotated; confidence: high}
\end{GSitem}
\begin{GSitem}{1974-336.4}{E}{The Jacobian and the missing Hodge quotient}
\GSfield{Kind; source}{construction; \GSpage{P000554}}
\GSfield{The notes}{The Jacobian is written \(H^1(R,\mathbb C)/H^1(R,\mathbb Z)\).}
\GSfield{Statement}{For a compact Riemann surface \(R\) of genus \(g\), \(J(R)=H^0(R,\Omega_R^1)^*/H_1(R,\mathbb Z)\), with the lattice embedded by integration. Equivalently,
\[J(R)=H^1(R,\mathbb C)/(F^1H^1(R,\mathbb C)+H^1(R,\mathbb Z)).\] The literal quotient omits a complex \(g\)-dimensional Hodge subspace.}
\GSfield{Fate}{\citet[\S2, analytic description]{milne1986jacobian} gives the period construction. The omission is mathematical: the literal numerator has complex dimension \(2g\) and its rank-\(2g\) lattice does not give the compact Jacobian.}
\GSfield{Present status}{As of November 2026, the stated reconstruction is a standard established result.}
\GSfield{Assessment}{literal statement: false as written; publication match: exact; verification: read the result; development: annotated; confidence: high}
\end{GSitem}
\begin{GSitem}{1974-336.5}{E}{Period matrices and the moduli diagram}
\GSfield{Kind; source}{construction; \GSpage{P000554}}
\GSfield{The notes}{A diagram relates curves, tori, periods and a domain of matrices.}
\GSfield{Statement}{After a symplectic marking and normalization, a principally polarized complex abelian variety is represented by a symmetric matrix \(\tau\) with \(\operatorname{Im}\tau>0\). Changing the marking acts by \(\operatorname{Sp}_{2g}(\mathbb Z)\). The resulting quotient describes principally polarized abelian varieties; the period map from curves lands in the Jacobian locus.}
\GSfield{Reconstruction}{Principal polarization and marking are editorial additions required to identify the matrix domain.}
\GSfield{Fate}{\citet[polarizations and moduli]{milne2008abelian} supplies the standard setting. The diagram does not specify a level structure or a fine moduli space, and its arrows are not treated as unqualified isomorphisms.}
\GSfield{Present status}{As of November 2026, the stated reconstruction is a standard established result.}
\GSfield{Assessment}{literal statement: correct with missing hypotheses; publication match: partial; verification: read the result; development: annotated; confidence: medium}
\end{GSitem}
\begin{GSitem}{1974-336.6}{E}{A genus-two polynomial Hamiltonian pair and its corrected quartic coefficient}
\GSfield{Kind; source}{construction; \GSpage{P000555}}
\GSfield{The notes}{The notes give \(J_1=p_1p_2+V\), \(J_2=p_1^2-2q_1p_1p_2+2(q_2-c)p_2^2+D\), claiming \(\{J_1,J_2\}=0\), where the quartic term of \(V\) is \(-5q_1^4/2\).}
\GSfield{Statement}{Identify \(c=c_1,d=d_1\) and put \[V_a=-\tfrac12q_2^2-\tfrac52q_1^2q_2-aq_1^4+cq_1^2-dq_1,\quad D=q_1^5+2cq_1^3+dq_1^2-4q_1q_2^2+4cq_1q_2-2dq_2.\] For \(\{f,g\}=\sum_i(f_{q_i}g_{p_i}-f_{p_i}g_{q_i})\), direct calculation gives \(\{J_1,J_2\}=(8a-5)q_1^3(q_1p_2-p_1)\). Thus the copied coefficient \(a=5/2\) fails, while \(a=5/8\) works.}
\GSfield{Reconstruction}{The identification of the differently subscripted parameters is explicit. The replacement concerns the last quartic coefficient only.}
\GSfield{Fate}{\citet[\S2, equations (18)--(19), p.~64]{novikov1974periodic} already publishes the genus-two Hamiltonian construction (received 18 February 1974; publication July 1974), before this September talk. The coefficient is also \(5/8\) after expanding \citet[\S4]{dubrovinnovikov1974periodicity}. The exact bracket audit below verifies the mathematical correction in the notebook convention.}
\GSfield{Present status}{As of November 2026, the corrected classical commuting pair is established; the printed notebook pair is not in involution.}
\GSfield{Assessment}{literal statement: false as written; publication match: exact; verification: read the result; development: annotated; confidence: high}
\end{GSitem}
\begin{GSitem}{1974-336.7}{E}{Commuting Hamiltonian vector fields}
\GSfield{Kind; source}{claim; \GSpage{P000555}}
\GSfield{The notes}{The two Hamiltonian fields are asserted to commute.}
\GSfield{Statement}{On a symplectic manifold, \([X_f,X_g]=\pm X_{\{f,g\}}\), with sign determined by convention. Consequently the corrected pair in item 1974-336.6 has commuting Hamiltonian fields. The uncorrected pair generally does not.}
\GSfield{Fate}{This is the standard Hamiltonian identity used in \citet[\S2]{dubrovinnovikov1974periodicity}. The application is contingent on the coefficient correction, not a second independent assertion of involution.}
\GSfield{Present status}{As of November 2026, the stated reconstruction is a standard established result.}
\GSfield{Assessment}{literal statement: false as written; publication match: exact; verification: read the result; development: annotated; confidence: high}
\end{GSitem}
\begin{GSitem}{1974-336.8}{A}{Common levels as open parts of abelian varieties}
\GSfield{Kind; source}{theorem; \GSpage{P000555}}
\GSfield{The notes}{The common levels \(J_1,J_2=\mathrm{constant}\) and flow orbits are called tori or abelian varieties.}
\GSfield{Statement}{For the corrected genus-two stationary-KdV system and generic constants giving a nonsingular spectral curve, the affine common level is identified with a Zariski-open part of its Jacobian. The complete abelian surface is obtained by compactification. Compact nonsingular real components support the real torus dynamics; arbitrary complex affine levels are not compact tori.}
\GSfield{Reconstruction}{Genericity and compactification are necessary. Singular spectral curves require a separate degeneration theory.}
\GSfield{Fate}{\citet[\S\S2--4]{dubrovinnovikov1974periodicity} and \citet[Ch.~II, \S4]{dubrovinmatveevnovikov1976nonlinear} give the finite-gap algebraic-integrability construction (1974 result; 1976 synthesis).}
\GSfield{Present status}{As of November 2026, the stated reconstruction is a standard established result.}
\GSfield{Assessment}{literal statement: correct with missing hypotheses; publication match: exact; verification: read the result; development: annotated; confidence: high}
\end{GSitem}
\begin{GSitem}{1974-336.9}{A}{Rational parametrization of the marked universal hyperelliptic Jacobian}
\GSfield{Kind; source}{theorem; \GSpage{P000555}}
\GSfield{The notes}{The collection of all the preceding tori is said to be unirational.}
\GSfield{Statement}{The precise nearby construction concerns a total family of hyperelliptic Jacobians, with a distinguished branch point and the coordinate/marking normalizations used in its polynomial model. It is a rational total parameter variety in that formulation, hence gives a unirational family after forgetting the finite marking. This does not make any fixed positive-dimensional complex abelian variety unirational.}
\GSfield{Reconstruction}{The family and its marking have been supplied from the matching result; an arbitrary fine-moduli interpretation is not asserted.}
\GSfield{Fate}{\citet[\S4, theorem]{dubrovinnovikov1974periodicity} states rationality of the marked total variety; the manuscript is its abbreviated geometric consequence.}
\GSfield{Present status}{As of November 2026, the stated reconstruction is a standard established result.}
\GSfield{Assessment}{literal statement: correct with missing hypotheses; publication match: partial; verification: read the result; development: annotated; confidence: medium}
\end{GSitem}
\begin{GSitem}{1974-336.10}{E}{The spectral quintic and a parameter-normalization discrepancy}
\GSfield{Kind; source}{construction; \GSpage{P000555}}
\GSfield{The notes}{The curve is written \[y^2=E^5+\tfrac c4E^3-\tfrac d{16}E^2+(\tfrac{J_1}{32}+\tfrac{c^2}{4})E+\tfrac{J_2}{256}+\tfrac{cd}{128}.\]}
\GSfield{Statement}{The spectral construction is established, but its coefficients depend on the Hamiltonian and spectral-coordinate normalization. With the repaired \(J_1\), canonical flow \(D_x=\{\cdot,J_1\}\), \(u=q_1\), and \(P=E^2+q_1E/2+(q_1^2-2q_2+2c)/16\), exact substitution into the squared-eigenfunction identity gives \[R=E^5+\tfrac c4E^3-\tfrac d{16}E^2+(\tfrac{J_1}{32}+\tfrac{c^2}{64})E-\tfrac{J_2}{256}-\tfrac{cd}{128}.\] This is a conditional normalization check; it does not establish a literal corrected reading of the notebook.}
\GSfield{Reconstruction}{The author-hosted Dubrovin--Novikov parameter change yields \(c^2/64\), but that text also prints a positive constant term and has apparent typographic defects elsewhere. The notebook does not specify its Lax matrix. The conditional computation fixes the coefficients only after the stated choices of potential, flow and divisor polynomial.}
\GSfield{Fate}{\citet[\S2, equation (18)]{novikov1974periodic} already publishes the genus-two spectral polynomial in July 1974. Compare \citet[\S4]{dubrovinnovikov1974periodicity} and the coordinate interpretation in \citet[Ch.~II, \S4, p.~105]{dubrovinmatveevnovikov1976nonlinear}. The companion gives an exact conditional polynomial audit, while preserving the unresolved comparison of all source normalizations.}
\GSfield{Present status}{As of November 2026, the spectral construction is established, and the stated normalization has been checked exactly; the literal coefficient and constant-term comparison remains an explicitly unresolved normalization discrepancy.}
\GSfield{Assessment}{literal statement: not determinable; publication match: partial; verification: read the result; development: annotated; confidence: medium}
\end{GSitem}
\begin{GSitem}{1974-336.11}{E}{Abel’s map from a symmetric power}
\GSfield{Kind; source}{construction; \GSpage{P000555}}
\GSfield{The notes}{The map \(S^gR\to J(R)\), given by sums of integrals of holomorphic differentials from a base point, is called an isomorphism.}
\GSfield{Statement}{For a smooth projective complex curve of genus \(g\), \(D\mapsto[D-gQ]\) maps \(\operatorname{Sym}^gR\) surjectively and birationally to \(J(R)\). Its fibre over \([D-gQ]\) is the complete linear system \(|D|\). Thus it is an isomorphism on the nonspecial locus, but generally not globally: in genus two the hyperelliptic pencil gives a positive-dimensional fibre.}
\GSfield{Fate}{\citet[Theorem~5.1 and the preceding fibre description]{milne1986jacobian} states the classical result. The source is a 2021 corrected version of the 1986 chapter; the mathematical theorem long predates this lecture.}
\GSfield{Present status}{As of November 2026, the stated reconstruction is a standard established result.}
\GSfield{Assessment}{literal statement: false as written; publication match: exact; verification: read the result; development: annotated; confidence: high}
\end{GSitem}
\begin{GSitem}{1974-336.12}{E}{Normalized Bloch solutions and the unit multiplier condition}
\GSfield{Kind; source}{definition; \GSpage{P000555}}
\GSfield{The notes}{For \(L=-d^2/dx^2+u\), the notes impose \(\Psi(x+T)=\lambda\Psi(x)\), \(|\lambda|=1\), \(\lambda=e^{ip}\), and \(\Psi(x_0)=1\).}
\GSfield{Statement}{For a real smooth \(T\)-periodic potential, Floquet solutions satisfy \(\Psi(x+T,E)=\lambda(E)\Psi(x,E)\). The condition \(|\lambda|=1\) characterizes real spectral bands; it is not true in open gaps. Normalization at \(x_0\) is meromorphic in the spectral point and fails as an ordinary finite value at its pole divisor.}
\GSfield{Fate}{\citet[Ch.~II, \S1]{dubrovinmatveevnovikov1976nonlinear} sets out the standard Floquet normalization used in the new finite-gap construction.}
\GSfield{Present status}{As of November 2026, the stated reconstruction is a standard established result.}
\GSfield{Assessment}{literal statement: correct with missing hypotheses; publication match: exact; verification: read the result; development: annotated; confidence: high}
\end{GSitem}
\begin{GSitem}{1974-336.13}{F}{The almost-periodic completeness warning}
\GSfield{Kind; source}{remark; \GSpage{P000555}}
\GSfield{The notes}{For almost-periodic potentials the collection of such eigenfunctions is said not to form a complete system.}
\GSfield{Statement}{An almost-periodic function need not have a period \(T\), so the preceding monodromy eigenvalue equation does not define a general Bloch basis. No function space, spectral measure or notion of completeness is specified here.}
\GSfield{Fate}{\citet[\S3]{dubrovinnovikov1974periodicity} distinguishes the periodic case from an imposed meromorphic finite-gap condition in the conditionally periodic case. This does not identify an exact completeness counterexample.}
\GSfield{Present status}{As of November 2026, the source does not determine a sufficiently precise claim for an open-problem status.}
\GSfield{Assessment}{literal statement: not determinable; publication match: thematic only; verification: read the result; development: none; confidence: low}
\end{GSitem}
\begin{GSitem}{1974-336.14}{E}{The two-sheeted spectral surface and its genus}
\GSfield{Kind; source}{construction; \GSpage{P000555}}
\GSfield{The notes}{Two eigenfunctions lead to a two-sheeted surface branched at spectral endpoints, described as having infinite genus.}
\GSfield{Statement}{The normalized Bloch functions for a real periodic Schrödinger operator are naturally considered on a two-sheeted spectral surface. Open-gap endpoints are branch points; closed gaps give removable double points after normalization. Infinitely many open gaps give an infinite-genus surface. With \(g\) open gaps, the compact finite-gap model is \(y^2=\prod_{j=0}^{2g}(E-E_j)\), of genus \(g\).}
\GSfield{Reconstruction}{The qualifier “infinitely many open gaps” is necessary; periodicity alone does not imply infinite genus.}
\GSfield{Fate}{\citet[Ch.~II, \S1]{dubrovinmatveevnovikov1976nonlinear} distinguishes the general periodic surface from its finite-gap algebraic specialization.}
\GSfield{Present status}{As of November 2026, the stated reconstruction is a standard established result.}
\GSfield{Assessment}{literal statement: correct with missing hypotheses; publication match: exact; verification: read the result; development: annotated; confidence: high}
\end{GSitem}
\begin{GSitem}{1974-336.15}{A}{The finite-gap isospectral family and its Jacobian}
\GSfield{Kind; source}{theorem; \GSpage{P000556}}
\GSfield{The notes}{The set of potentials with the same spectrum is identified with the Jacobian of the spectral surface.}
\GSfield{Statement}{Fix a real nonsingular finite-gap spectrum with \(g\) open gaps and admissible real divisor data. The corresponding smooth real isospectral family is a real \(g\)-torus in the complex Jacobian. The complex algebraic construction is described on an open Jacobian chart and its compactification. Infinite-gap spectra require additional analytic hypotheses and do not give this finite-dimensional statement.}
\GSfield{Reconstruction}{The finite-gap and real-admissibility conditions are made explicit.}
\GSfield{Fate}{\citet[\S3]{dubrovinnovikov1974periodicity} gives the finite-gap identification and linearized flows; publication followed the September lecture in 1974, with receipt on 3 June 1974.}
\GSfield{Present status}{As of November 2026, the stated reconstruction is a standard established result.}
\GSfield{Assessment}{literal statement: correct with missing hypotheses; publication match: exact; verification: read the result; development: annotated; confidence: high; problem \hyperref[prob:GP-1973-51]{GP-1973-51}}
\end{GSitem}

\GStalk{337}{1974-09-09}{Czesław Bessaga; I. M. Gelfand (intervention)}{P000556}{Z-sets, infinite-dimensional spaces and the space of rays}
The notes introduce negligible subsets of the Hilbert cube and the extension theorem for their homeomorphisms. They also mention the topological classification of separable infinite-dimensional Fréchet spaces. Gelfand\textquotesingle s space-of-rays example has an exact published antecedent in 1967. Two literal assertions need correction: the pseudoboundary is a countable union of Z-sets, and a maximal such union cannot mean maximal by ordinary inclusion.

\begin{GSitem}{1974-337.1}{E}{Z-sets defined by approximation of Hilbert-cube maps}
\GSfield{Kind; source}{definition; \GSpage{P000556}}
\GSfield{The notes}{A closed subset \(A\) of a complete metric space \(X\) is called a Z-set if continuous maps of the Hilbert cube can be approximated by maps missing \(A\).}
\GSfield{Statement}{For a metric space \(X\), a closed subset \(A\subset X\) is a Z-set in the stated sense if for every continuous \(f:Q\to X\) and every \(\varepsilon>0\) there is continuous \(g:Q\to X\setminus A\) with \(\sup_Qd(f,g)<\varepsilon\), where \(Q=[-1,1]^{\mathbb N}\).}
\GSfield{Fate}{\citet[introduction, pp.~281--282]{andersonchapman1971extending} recalls the established negligible-set theory and its equivalent formulations in Hilbert-cube manifolds. The ambient-space hypotheses matter when moving between those formulations.}
\GSfield{Present status}{As of November 2026, the stated reconstruction is a standard established result.}
\GSfield{Assessment}{literal statement: correct; publication match: exact; verification: read the result; development: annotated; confidence: high}
\end{GSitem}
\begin{GSitem}{1974-337.2}{E}{The pseudoboundary of the Hilbert cube}
\GSfield{Kind; source}{claim; \GSpage{P000556}}
\GSfield{The notes}{For \(Q=[-1,1]^{\mathbb N}\) and \(P=(-1,1)^{\mathbb N}\), the notes call \(Q\setminus P\) a Z-set.}
\GSfield{Statement}{The pseudoboundary is \(\bigcup_{j\ge1}\{x\in Q:|x_j|=1\}\), a countable union of closed Z-sets. It is dense and not closed in \(Q\), so it is not itself a Z-set under the preceding definition. Each coordinate face can be avoided by an arbitrarily small inward change in that coordinate.}
\GSfield{Fate}{The distinction is part of the Hilbert-cube setting of \citet[pp.~281--282]{andersonchapman1971extending}; the correction also follows immediately from the definition copied in the notes.}
\GSfield{Present status}{As of November 2026, the pseudoboundary is a sigma-Z-set, and the literal Z-set claim is false.}
\GSfield{Assessment}{literal statement: false as written; publication match: partial; verification: read the result; development: annotated; confidence: high}
\end{GSitem}
\begin{GSitem}{1974-337.3}{E}{Z-set unknotting in the Hilbert cube}
\GSfield{Kind; source}{theorem; \GSpage{P000556}}
\GSfield{The notes}{A homeomorphism between two Z-sets extends to a homeomorphism of the surrounding cube.}
\GSfield{Statement}{If \(A,B\) are closed Z-sets in the Hilbert cube \(Q\), every homeomorphism \(h:A\to B\) extends to a homeomorphism \(H:Q\to Q\).}
\GSfield{Reconstruction}{The ambient Hilbert cube is essential; the preceding definition for arbitrary complete metric spaces does not make the extension assertion true there.}
\GSfield{Fate}{\citet[Theorem~4 and corollary, pp.~37--38]{anderson1967special} gives the extension theorem in its equivalent deficient-set formulation; \citet[introduction]{andersonchapman1971extending} explicitly recalls Z-set unknotting. Both precede the talk.}
\GSfield{Present status}{As of November 2026, the stated reconstruction is a standard established result.}
\GSfield{Assessment}{literal statement: correct with missing hypotheses; publication match: exact; verification: read the result; development: annotated; confidence: high}
\end{GSitem}
\begin{GSitem}{1974-337.4}{E}{Countable unions of Z-sets}
\GSfield{Kind; source}{definition; \GSpage{P000556}}
\GSfield{The notes}{A \(Z_\sigma\)-set is introduced as a countable union of Z-sets.}
\GSfield{Statement}{A sigma-Z-set in \(X\) is a set \(A=\bigcup_{m\ge1}A_m\) with each \(A_m\) a closed Z-set in \(X\). The union need not be closed.}
\GSfield{Fate}{This is standard negligible-set terminology in the infinite-dimensional topology discussed by \citet[pp.~31--38]{anderson1967special}. The pseudoboundary is the preceding concrete example.}
\GSfield{Present status}{As of November 2026, the stated reconstruction is a standard established result.}
\GSfield{Assessment}{literal statement: correct; publication match: exact; verification: read the result; development: annotated; confidence: high}
\end{GSitem}
\begin{GSitem}{1974-337.5}{F}{The ambiguous maximal sigma-Z-set assertion}
\GSfield{Kind; source}{claim; \GSpage{P000556}}
\GSfield{The notes}{The notes say that a maximal \(Z_\sigma\)-set exists.}
\GSfield{Statement}{If maximal means maximal by inclusion among sigma-Z-subsets of \(Q\), the claim is false. Every Z-set is nowhere dense, so a sigma-Z-set is meagre and is not \(Q\). Any missing point is itself a Z-set, and adjoining it gives a strictly larger sigma-Z-set. A universal or absorbing-set meaning would require an additional definition.}
\GSfield{Fate}{The literal obstruction follows from the Baire theorem and the preceding definition. Searches in the opened Anderson and Anderson--Chapman sources did not identify a theorem with the unqualified maximality stated here.}
\GSfield{Present status}{As of November 2026, inclusion-maximality is refuted; an intended absorbing or universal notion is not recoverable from this sentence.}
\GSfield{Assessment}{literal statement: false as written; publication match: none found; verification: unresolved lead; development: annotated; confidence: low}
\end{GSitem}
\begin{GSitem}{1974-337.6}{E}{Every continuous real function on the ray space is constant}
\GSfield{Kind; source}{claim; \GSpage{P000556}; I. M. Gelfand}
\GSfield{The notes}{Gelfand considers rays in the space of real sequences and says that every continuous function on the quotient is constant.}
\GSfield{Statement}{Let \(X=\mathbb R^{\mathbb N}\) have its product topology and let \(S=(X\setminus\{0\})/\mathbb R_{>0}\) be the quotient by positive scalar multiplication. Every continuous \(S\to\mathbb R\) is constant. Equivalently, every continuous degree-zero homogeneous real function on \(X\setminus\{0\}\) is constant.}
\GSfield{Reconstruction}{The product topology and positive scaling are essential. A normed-space unit sphere gives a different quotient.}
\GSfield{Fate}{\citet[introduction, Russian p.~34]{gelfandfuchs1967topology} states this exact example; received July 1967, published 1967. The space is described there as Hausdorff but nonregular, not as a non-Hausdorff quotient. Earlier occurrence: 1967-132.7.}
\GSfield{Present status}{As of November 2026, the stated reconstruction is a standard established result.}
\GSfield{Assessment}{literal statement: correct with missing hypotheses; publication match: exact; verification: read the result; development: annotated; confidence: high; problem \hyperref[prob:GP-1974-01]{GP-1974-01}}
\end{GSitem}
\begin{GSitem}{1974-337.7}{E}{Homeomorphism type of infinite-dimensional separable Fréchet spaces}
\GSfield{Kind; source}{theorem; \GSpage{P000556}}
\GSfield{The notes}{All separable Fréchet spaces are said to be homeomorphic.}
\GSfield{Statement}{All infinite-dimensional separable real Fréchet spaces are homeomorphic to separable Hilbert space. Finite-dimensional spaces are excluded: their topological dimensions distinguish them from one another and from the infinite-dimensional case.}
\GSfield{Fate}{\citet[p.~31]{anderson1967special} explicitly summarizes the already established Anderson--Kadec and related results. The 1966 symposium paper appeared in 1967, before this talk.}
\GSfield{Present status}{As of November 2026, the stated reconstruction is a standard established result.}
\GSfield{Assessment}{literal statement: false as written; publication match: exact; verification: read the result; development: annotated; confidence: high}
\end{GSitem}

\GStalk{338}{1974-09-09}{Enrico Giusti}{P000556-P000558}{Prescribed mean curvature, boundary obstructions and BV relaxation}
The lecture first gives elementary failures of classical solvability, including a radial annulus and a one-dimensional solution with infinite slope. It then builds the variational setting in bounded variation: dual definitions of total variation and area, lower semicontinuity, compactness and a boundary mismatch term. The sign of the volume potential conflicts with the earlier PDE convention. The cutoff-minimization question is meaningful only after coercivity and the boundary class are specified.

\begin{GSitem}{1974-338.1}{E}{The prescribed-mean-curvature Dirichlet equation}
\GSfield{Kind; source}{problem; \GSpage{P000556}}
\GSfield{The notes}{The equation is \(\operatorname{div}(Du/\sqrt{1+|Du|^2})=-H(x,u)\), with \(u=\phi\) on the boundary; \(H=0\) is the minimal-surface case.}
\GSfield{Statement}{Given a bounded domain \(\Omega\), boundary data \(\phi\) and prescribed function \(H\), seek an appropriately defined solution of the displayed graph equation with the stated boundary trace. Classical, weak Sobolev and relaxed BV solutions are different formulations; none exists for arbitrary data without additional conditions.}
\GSfield{Fate}{\citet[introduction and \S1]{giusti1976boundary} treats the generalized Dirichlet formulation and recalls the earlier work on existence. The lecture\textquotesingle s examples establish why the solution class must be stated.}
\GSfield{Present status}{As of November 2026, the stated reconstruction is a standard established result.}
\GSfield{Assessment}{literal statement: correct with missing hypotheses; publication match: exact; verification: read the result; development: annotated; confidence: high}
\end{GSitem}
\begin{GSitem}{1974-338.2}{E}{A radial annulus and the maximal catenoid height}
\GSfield{Kind; source}{construction; \GSpage{P000556-P000557}}
\GSfield{The notes}{On \(1<|x|<2\subset\mathbb R^2\), radial minimal graphs have inner value \(M\) and outer value zero; beyond \(M_0\) no solution exists.}
\GSfield{Statement}{For \(0\le c\le1\), the radial solutions are \[u_c(r)=c\log\frac{2+\sqrt{4-c^2}}{r+\sqrt{r^2-c^2}}.\] Their heights increase to \(M_0=\log(2+\sqrt3)\). There is no \(C^2(\Omega)\cap C^0(\overline\Omega)\) minimal graph with these boundary values if \(M>M_0\): comparison gives uniqueness, rotation invariance forces radiality, and \(ru^{\prime}/\sqrt{1+(u^{\prime})^2}=-c\) forces \(c\le1\). At equality the inner-boundary slope is infinite.}
\GSfield{Reconstruction}{The solution class and the equality case are editorially explicit; \(C^1\) regularity up to the inner boundary fails at \(c=1\).}
\GSfield{Fate}{This is an explicit catenoid calculation illustrating the classical-versus-generalized distinction in \citet[introduction]{giusti1976boundary}. Its exact formula and threshold follow directly by integrating the radial equation.}
\GSfield{Present status}{As of November 2026, the stated reconstruction is a standard established result.}
\GSfield{Assessment}{literal statement: correct with missing hypotheses; publication match: thematic only; verification: read the result; development: annotated; confidence: high}
\end{GSitem}
\begin{GSitem}{1974-338.3}{F}{Soap-film diagrams as a geometric illustration}
\GSfield{Kind; source}{remark; \GSpage{P000557}}
\GSfield{The notes}{Drawings show a film lifted with a rod and a film spanning wire rings, accompanying an experimental check of loss of graph solvability.}
\GSfield{Statement}{The drawings illustrate the distinction between a single-valued graph over an annulus and a general spanning surface. They specify no surface functional, stability criterion, contact-angle condition or exact transition height for the physical experiment.}
\GSfield{Fate}{The neighbouring radial calculation is identifiable, but no precise theorem about the drawn experiment is recorded. The general variational background in \citet[introduction]{giusti1976boundary} is thematic only.}
\GSfield{Present status}{As of November 2026, the source does not determine a sufficiently precise claim for an open-problem status.}
\GSfield{Assessment}{literal statement: not determinable; publication match: thematic only; verification: read the result; development: none; confidence: low}
\end{GSitem}
\begin{GSitem}{1974-338.4}{E}{A continuous prescribed curvature with an infinite-slope solution}
\GSfield{Kind; source}{construction; \GSpage{P000557}}
\GSfield{The notes}{The example \(u(x)=3x^{1/3}\), \(-1<x<1\), is nondifferentiable at zero although the derivative of its curvature flux is continuous.}
\GSfield{Statement}{Using the real cube root, \(u\in W^{1,1}(-1,1)\) and \(u^{\prime}=|x|^{-2/3}\) almost everywhere. Its flux is \(q=(1+|x|^{4/3})^{-1/2}\), with \(q^{\prime}=-(2/3)\sqrt[3]{x}(1+|x|^{4/3})^{-3/2}\), a continuous function. Thus a weak graph solution may have continuous right-hand side without finite slope. It is not in \(W^{1,2}\).}
\GSfield{Fate}{The assertion is verified by direct differentiation and integrability at zero; it illustrates the degeneracy explained by \citet[\S1]{giusti1976boundary}. No exact published formula match was found.}
\GSfield{Present status}{As of November 2026, the stated reconstruction is a standard established result.}
\GSfield{Assessment}{literal statement: correct; publication match: thematic only; verification: read the result; development: annotated; confidence: high}
\end{GSitem}
\begin{GSitem}{1974-338.5}{E}{Raising one endpoint beyond the possible flux}
\GSfield{Kind; source}{claim; \GSpage{P000557}}
\GSfield{The notes}{After \(u(-1)=-3,u(1)=3\), the notes change to \(u(-1)=3,u(1)=3+A\) but integrate from the value \(-3\).}
\GSfield{Statement}{For the fixed right-hand side of item 1974-338.4, retain \(v(-1)=-3\). Any admissible flux is \(q_v=q+C\). Since \(q(0)=1\) and \(|q_v|\le1\), necessarily \(C\le0\); monotonicity of \(s/\sqrt{1+s^2}\) gives \(v^{\prime}\le |x|^{-2/3}\). Therefore \(v(1)\le3\), and data \(v(1)=3+A\), \(A>0\), are impossible.}
\GSfield{Reconstruction}{The left endpoint is restored to \(-3\), and the omitted integration constant is included. The equality flux corresponds to a weak solution with infinite slope.}
\GSfield{Fate}{This is the elementary one-dimensional obstruction adjacent to the explicit example. The general perimeter obstruction is treated by \citet[\S1B]{giusti1976boundary}.}
\GSfield{Present status}{As of November 2026, the stated reconstruction is a standard established result.}
\GSfield{Assessment}{literal statement: false as written; publication match: thematic only; verification: read the result; development: annotated; confidence: high}
\end{GSitem}
\begin{GSitem}{1974-338.6}{E}{The area-plus-potential functional and its Euler equation}
\GSfield{Kind; source}{method; \GSpage{P000557}}
\GSfield{The notes}{The functional is \(\mathcal F(u)=\int_\Omega\sqrt{1+|Du|^2}+\int_\Omega\lambda(x,u)\), with \(\lambda_t=H\).}
\GSfield{Statement}{For smooth admissible variations with fixed boundary values, the displayed plus sign gives \(-\operatorname{div}(Du/\sqrt{1+|Du|^2})+H(x,u)=0\), hence \(\operatorname{div}(Du/\sqrt{1+|Du|^2})=H(x,u)\). To match the earlier equation with \(-H\), use \(-\lambda\) or redefine the sign of \(H\).}
\GSfield{Fate}{The first variation is immediate. \citet[\S1]{giusti1976boundary} uses a consistent area-plus-volume convention; the notebook switches conventions without warning.}
\GSfield{Present status}{As of November 2026, the stated reconstruction is a standard established result.}
\GSfield{Assessment}{literal statement: false as written; publication match: exact; verification: read the result; development: annotated; confidence: high}
\end{GSitem}
\begin{GSitem}{1974-338.7}{E}{Removing the gradient cutoff in variational minimization}
\GSfield{Kind; source}{question; \GSpage{P000557}}
\GSfield{The notes}{A minimum is asserted under \(|Du|\le K\); then the notes ask whether the minimizers \(u_K\) converge as \(K\to\infty\).}
\GSfield{Statement}{On a bounded Lipschitz domain with nonempty fixed-trace admissible class, a fixed uniform gradient bound and suitable continuity of the potential give a direct-method minimum. Removing the bound requires a uniform BV and \(L^1\) estimate. It is false without coercivity: if \(H=h>0\) is sufficiently large and \(w\ge0\) has zero trace with \(h\int w>\int|Dw|\), then \(\mathcal F(-tw)\to-\infty\). Under a coercive relaxed functional, subsequences converge in \(L^1\) to BV minimizers, which need not retain the classical boundary values.}
\GSfield{Reconstruction}{The coercivity and boundary assumptions are not present in the question. The elementary large-force counterexample shows why they cannot be omitted.}
\GSfield{Fate}{\citet[introduction and existence hypotheses]{giusti1976boundary} explains the BV setting and necessary size restrictions. The unqualified convergence question has no affirmative answer; the positive compactness argument is supplied by the subsequent items.}
\GSfield{Present status}{As of November 2026, the conditional BV approach is classical; unconditional convergence is false and no separate open problem is inferred.}
\GSfield{Assessment}{literal statement: correct with missing hypotheses; publication match: partial; verification: read the result; development: annotated; confidence: medium}
\end{GSitem}
\begin{GSitem}{1974-338.8}{E}{A necessary flux-perimeter inequality}
\GSfield{Kind; source}{claim; \GSpage{P000557}}
\GSfield{The notes}{For a Lipschitz solution and \(A\subset\Omega\), the notes bound \(|\int_AH|\) by \((1-\varepsilon_0)\mathcal H^{n-1}(\partial A)\).}
\GSfield{Statement}{If \(u\) is a solution with \(\|Du\|_\infty\le K<\infty\), then for every smooth relatively compact subdomain \(A\), \[\left|\int_AH(x,u(x))\,dx\right|\le\frac{K}{\sqrt{1+K^2}}\mathcal H^{n-1}(\partial A).\] The same statement has a finite-perimeter formulation using normal traces. The sign convention disappears after taking absolute values.}
\GSfield{Fate}{\citet[\S1B]{giusti1976boundary} treats the same geometric obstruction. The constant in the notebook is precisely the supremum norm of the flux.}
\GSfield{Present status}{As of November 2026, the stated reconstruction is a standard established result.}
\GSfield{Assessment}{literal statement: correct with missing hypotheses; publication match: exact; verification: read the result; development: annotated; confidence: high}
\end{GSitem}
\begin{GSitem}{1974-338.9}{E}{Bounded variation and the dual definition of total variation}
\GSfield{Kind; source}{definition; \GSpage{P000558}}
\GSfield{The notes}{BV functions are \(L^1\) functions whose first distributional derivatives are Radon measures, with a dual supremum over bounded test vector fields.}
\GSfield{Statement}{For an open \(\Omega\), \(u\in BV(\Omega)\) means \(u\in L^1(\Omega)\) and \[|Du|(\Omega)=\sup_{g\in C_c^1(\Omega;\mathbb R^n),\ |g|\le1}\int_\Omega u\operatorname{div}g\,dx<\infty.\] This is equivalent to the distributional gradient being a finite vector Radon measure.}
\GSfield{Fate}{\citet[\S1A]{giusti1976boundary} recalls the standard definition used in the earlier area-relaxation theory.}
\GSfield{Present status}{As of November 2026, the stated reconstruction is a standard established result.}
\GSfield{Assessment}{literal statement: correct; publication match: exact; verification: read the result; development: annotated; confidence: high}
\end{GSitem}
\begin{GSitem}{1974-338.10}{E}{Relaxed area for a BV graph}
\GSfield{Kind; source}{definition; \GSpage{P000558}}
\GSfield{The notes}{The notation \(\int\sqrt{1+|Du|^2}\) is defined through a supremum over \((g_0,g)\) satisfying \(g_0^2+|g|^2\le1\).}
\GSfield{Statement}{If \(Du=\nabla u\,dx+D^su\), the relaxed area is \[\mathcal A(u)=\int_\Omega\sqrt{1+|\nabla u|^2}\,dx+|D^su|(\Omega).\] It equals \(\sup\int_\Omega(g_0+u\operatorname{div}g)\,dx\) over compactly supported \(C^1\) test fields with \(g_0^2+|g|^2\le1\). It is not obtained by squaring a singular measure as a density.}
\GSfield{Fate}{\citet[\S1A]{giusti1976boundary} gives this generalized-area formulation; the notebook\textquotesingle s dual definition is the appropriate interpretation.}
\GSfield{Present status}{As of November 2026, the stated reconstruction is a standard established result.}
\GSfield{Assessment}{literal statement: correct; publication match: exact; verification: read the result; development: annotated; confidence: high}
\end{GSitem}
\begin{GSitem}{1974-338.11}{E}{Lower semicontinuity of BV area}
\GSfield{Kind; source}{theorem; \GSpage{P000558}}
\GSfield{The notes}{The area is lower semicontinuous for \(L^1\) convergence.}
\GSfield{Statement}{If \(u_j\to u\) in \(L^1(\Omega)\), then \(\mathcal A(u)\le\liminf_j\mathcal A(u_j)\), with the area extended by \(+\infty\) outside \(BV\). Each functional in the dual supremum is \(L^1\)-continuous, and a supremum of such functions is lower semicontinuous.}
\GSfield{Fate}{\citet[\S1A]{giusti1976boundary} uses this standard direct-method fact. The dual formula in the notes supplies its proof.}
\GSfield{Present status}{As of November 2026, the stated reconstruction is a standard established result.}
\GSfield{Assessment}{literal statement: correct; publication match: exact; verification: read the result; development: annotated; confidence: high}
\end{GSitem}
\begin{GSitem}{1974-338.12}{E}{Compactness of bounded BV sets}
\GSfield{Kind; source}{theorem; \GSpage{P000558}}
\GSfield{The notes}{Bounded sets in BV are said to be relatively compact in \(L^1\).}
\GSfield{Statement}{If \(\Omega\) is bounded with Lipschitz boundary and \(\sup_j(\|u_j\|_{L^1}+|Du_j|(\Omega))<\infty\), there is an \(L^1(\Omega)\)-convergent subsequence. On a general open set only local compactness follows without control at infinity and the boundary. Bounding total variation alone does not control additive constants.}
\GSfield{Reconstruction}{The domain condition and the full BV norm are essential additions.}
\GSfield{Fate}{\citet[\S1A]{giusti1976boundary} states the bounded-domain framework in which the BV compactness argument is applied.}
\GSfield{Present status}{As of November 2026, the stated reconstruction is a standard established result.}
\GSfield{Assessment}{literal statement: correct with missing hypotheses; publication match: exact; verification: read the result; development: annotated; confidence: high}
\end{GSitem}
\begin{GSitem}{1974-338.13}{E}{Boundary mismatch in the relaxed Dirichlet functional}
\GSfield{Kind; source}{construction; \GSpage{P000558}}
\GSfield{The notes}{The area-plus-potential expression acquires \(\int_{\partial\Omega}|u-\phi|\,d\mathcal H^{n-1}\).}
\GSfield{Statement}{For a bounded Lipschitz domain, BV trace \(Tu\) and integrable boundary datum \(\phi\), the generalized Dirichlet energy contains \[\mathcal A(u)+\int_\Omega\lambda(x,u)\,dx+\int_{\partial\Omega}|Tu-\phi|\,d\mathcal H^{n-1}.\] This boundary term permits detachment from the prescribed trace while charging its vertical area. Existence still requires lower semicontinuity and coercivity of the entire functional.}
\GSfield{Fate}{\citet[pp.~501--505]{giusti1976boundary} explicitly recalls generalized-solution constructions from Giusti (1970), Giaquinta and Miranda (1973--1974). This framework was already available at the talk; the 1976 paper studies its boundary behaviour.}
\GSfield{Present status}{As of November 2026, the stated reconstruction is a standard established result.}
\GSfield{Assessment}{literal statement: correct with missing hypotheses; publication match: exact; verification: read the result; development: annotated; confidence: high}
\end{GSitem}

\GStalk{339}{1974-09-09}{B. A. Dubrovin}{P000558}{Finite-gap Bloch functions and their divisor}
This short continuation constructs the hyperelliptic spectral surface and locates the poles and zeros of a normalized Bloch solution. The real gap ovals are visible in the drawing. The divisor statement must allow cancellations, especially at the normalization point where the function equals one. The next dated lecture starts on P000559.

\begin{GSitem}{1974-339.1}{A}{The finite-gap hyperelliptic spectral curve}
\GSfield{Kind; source}{construction; \GSpage{P000558}}
\GSfield{The notes}{For a finite-gap potential, \(R(E)=\prod(E-E_i)\) and \(\Gamma_N\) is the surface of \(\sqrt{R(E)}\); \(N\) is its genus.}
\GSfield{Statement}{For a real smooth periodic Schrödinger potential with exactly \(N\) open gaps, list the distinct finite band endpoints \(E_1<\cdots<E_{2N+1}\). The smooth compactification of \(y^2=\prod_{j=1}^{2N+1}(E-E_j)\) has genus \(N\). Closed gaps are omitted after normalization.}
\GSfield{Fate}{\citet[\S3]{dubrovinnovikov1974periodicity} gives the algebraic spectral construction (received June 1974, published 1974). The hyperelliptic genus count itself is classical.}
\GSfield{Present status}{As of November 2026, the stated reconstruction is a standard established result.}
\GSfield{Assessment}{literal statement: correct with missing hypotheses; publication match: exact; verification: read the result; development: annotated; confidence: high}
\end{GSitem}
\begin{GSitem}{1974-339.2}{A}{Meromorphic continuation of the normalized Bloch solution}
\GSfield{Kind; source}{theorem; \GSpage{P000558}}
\GSfield{The notes}{Bloch functions continue analytically, corrected to meromorphically, on \(\Gamma_N\setminus\{\infty\}\).}
\GSfield{Statement}{With \(\Psi(x_0,E,x_0)=1\), the finite-gap Bloch solution at fixed \(x\) is meromorphic on the spectral curve away from the point above infinity. In the local parameter \(k=\sqrt E\), its exponential factor is \(e^{\pm ik(x-x_0)}\); infinity is generally an essential singularity.}
\GSfield{Fate}{\citet[\S3]{dubrovinnovikov1974periodicity} states the continuation result; \citet[\S\S1--2]{dubrovinnovikov1974periodic} develops the eigenfunction construction (submitted July 1974, Russian publication December 1974).}
\GSfield{Present status}{As of November 2026, the stated reconstruction is a standard established result.}
\GSfield{Assessment}{literal statement: correct with missing hypotheses; publication match: exact; verification: read the result; development: annotated; confidence: high}
\end{GSitem}
\begin{GSitem}{1974-339.3}{A}{Pole and zero divisors of degree N with possible cancellation}
\GSfield{Kind; source}{claim; \GSpage{P000558}}
\GSfield{The notes}{The function is said to have \(N\) poles \(\gamma_i(x_0)\) and \(N\) zeros.}
\GSfield{Statement}{There are effective degree-\(N\) divisors \(D(x_0)\) and \(D(x)\) such that the finite divisor of the normalized solution is \(D(x)-D(x_0)\), with common points cancelling. Generically the poles and zeros each number \(N\), with multiplicity; at \(x=x_0\), \(\Psi\equiv1\) and all cancel.}
\GSfield{Fate}{The divisor construction is in \citet[\S3]{dubrovinnovikov1974periodicity}. The cancellation qualification follows already from the normalization and prevents a false universal exact count.}
\GSfield{Present status}{As of November 2026, the stated reconstruction is a standard established result.}
\GSfield{Assessment}{literal statement: correct with missing hypotheses; publication match: exact; verification: read the result; development: annotated; confidence: high; problem \hyperref[prob:GP-1973-51]{GP-1973-51}}
\end{GSitem}
\begin{GSitem}{1974-339.4}{A}{One divisor point on each real gap oval}
\GSfield{Kind; source}{claim; \GSpage{P000558}}
\GSfield{The notes}{Each oval projects to a forbidden gap and carries one pole and one zero; poles depend on \(x_0\).}
\GSfield{Statement}{For the real nonsingular finite-gap class, each admissible Dirichlet divisor has one point on each real oval above a closed finite gap interval, with its sheet specified and the two sheets meeting at endpoints. \(D(x_0)\) is the pole divisor and \(D(x)\) the moving zero divisor before cancellation.}
\GSfield{Reconstruction}{The assertion concerns admissible real data, not arbitrary complex divisors.}
\GSfield{Fate}{\citet[\S3, the points \(\gamma_j\)]{dubrovinnovikov1974periodicity} gives the gap and sheet description; \citet[Ch.~II, \S1]{dubrovinmatveevnovikov1976nonlinear} supplies its detailed spectral formulation.}
\GSfield{Present status}{As of November 2026, the stated reconstruction is a standard established result.}
\GSfield{Assessment}{literal statement: correct with missing hypotheses; publication match: exact; verification: read the result; development: annotated; confidence: high; problem \hyperref[prob:GP-1973-51]{GP-1973-51}}
\end{GSitem}

\GStalk{340}{1974-09-16}{B. A. Dubrovin; I. M. Gelfand (questions)}{P000559-P000563}{Floquet theory and the polynomial finite-gap representation}
The lecture derives the monodromy criterion for bands, states the correspondence between finite-gap potentials and stationary KdV equations and Jacobian description, and expresses the logarithmic derivative of a Bloch function through a polynomial divisor. The proof sketch needs several repairs: a missing subtraction in the monodromy formula, a Wronskian factor, an infinity exponent, and a global complex-growth step in the polynomiality argument. The final variational equation describes stationary members of the hierarchy.

\begin{GSitem}{1974-340.1}{E}{Periodic Schrödinger bands and finite-gap potentials}
\GSfield{Kind; source}{definition; \GSpage{P000559}}
\GSfield{The notes}{The spectrum of \(L=-d^2/dx^2+u(x)\), \(u(x+T)=u(x)\), is a union of intervals; finitely many bands is called finite-zone.}
\GSfield{Statement}{For real smooth \(T\)-periodic \(u\), the self-adjoint operator on \(L^2(\mathbb R)\) has a closed band spectrum. A finite-gap potential has finitely many open bounded gaps, so its spectrum consists of finitely many compact bands and a final half-line, with closed gaps filled in.}
\GSfield{Fate}{\citet[Ch.~II, \S1]{dubrovinmatveevnovikov1976nonlinear} recalls this classical spectral setting before describing the 1974 finite-gap theory.}
\GSfield{Present status}{As of November 2026, the stated reconstruction is a standard established result.}
\GSfield{Assessment}{literal statement: correct with missing hypotheses; publication match: exact; verification: read the result; development: annotated; confidence: high}
\end{GSitem}
\begin{GSitem}{1974-340.2}{E}{Unimodular monodromy and its characteristic equation}
\GSfield{Kind; source}{construction; \GSpage{P000559}}
\GSfield{The notes}{Translation by \(T\) acts on a solution basis through \(\widehat T=(\alpha_{ij})\); the Wronskian gives determinant one.}
\GSfield{Statement}{For \(-y^{\prime\prime}+uy=Ey\), a fundamental solution matrix normalized at \(x_0\) has monodromy \(M(E,x_0)\in\operatorname{SL}_2(\mathbb C)\). If \(\tau=\tfrac12\operatorname{tr}M\), its multipliers satisfy \(\lambda^2-2\tau\lambda+1=0\). The determinant is one because the Wronskian is constant.}
\GSfield{Fate}{\citet[Ch.~II, \S1, monodromy construction]{dubrovinmatveevnovikov1976nonlinear} records these classical identities.}
\GSfield{Present status}{As of November 2026, the stated reconstruction is a standard established result.}
\GSfield{Assessment}{literal statement: correct; publication match: exact; verification: read the result; development: annotated; confidence: high}
\end{GSitem}
\begin{GSitem}{1974-340.3}{E}{The bounded-solution and discriminant criterion for the spectrum}
\GSfield{Kind; source}{theorem; \GSpage{P000559-P000560}}
\GSfield{The notes}{A spectral point has a bounded solution, hence a monodromy eigenvalue of modulus one; \(|\operatorname{tr}M|<2\) gives complex roots and \(>2\) real roots.}
\GSfield{Statement}{For real \(E\), \(E\in\sigma(L)\) if and only if \(|\tau(E)|\le1\), equivalently there is a nonzero solution bounded on all of \(\mathbb R\). In band interiors the multipliers are conjugate on the unit circle; in open gaps they are distinct real reciprocals. At an endpoint a bounded periodic or antiperiodic solution exists even when a second solution grows linearly.}
\GSfield{Fate}{\citet[Ch.~II, \S1]{dubrovinmatveevnovikov1976nonlinear} gives the Floquet spectral criterion. The endpoint qualification distinguishes one bounded eigenvector from boundedness of every monodromy iterate.}
\GSfield{Present status}{As of November 2026, the stated reconstruction is a standard established result.}
\GSfield{Assessment}{literal statement: correct with missing hypotheses; publication match: exact; verification: read the result; development: annotated; confidence: high}
\end{GSitem}
\begin{GSitem}{1974-340.4}{E}{Quasimomentum and the Bloch multiplier}
\GSfield{Kind; source}{definition; \GSpage{P000560}}
\GSfield{The notes}{The Bloch condition is written using \(u\), with \(\nu=e^{ip}\) and \(p\) called quasimomentum.}
\GSfield{Statement}{For a Bloch solution \(\psi\), write \(\psi(x+T,E)=e^{ip(E)}\psi(x,E)\). The phase \(p\) is defined modulo \(2\pi\) and a choice of the reciprocal multiplier; on spectral bands it can be chosen real. The symbol \(u\) in that line denotes the eigenfunction, although elsewhere it denotes the potential.}
\GSfield{Fate}{\citet[Ch.~II, \S1]{dubrovinmatveevnovikov1976nonlinear} fixes the same Floquet phase convention, subject to a possible factor \(T\).}
\GSfield{Present status}{As of November 2026, the stated reconstruction is a standard established result.}
\GSfield{Assessment}{literal statement: correct with missing hypotheses; publication match: exact; verification: read the result; development: annotated; confidence: high}
\end{GSitem}
\begin{GSitem}{1974-340.5}{E}{Base-point independence of the discriminant}
\GSfield{Kind; source}{claim; \GSpage{P000560}}
\GSfield{The notes}{The half-trace \(\tau(E)\) does not depend on the normalization point \(x_0\).}
\GSfield{Statement}{Monodromy matrices at different base points are conjugate by the intervening fundamental solution matrix. Their trace and multipliers are therefore independent of \(x_0\), although the individual matrix entries and normalized eigenvectors are not.}
\GSfield{Fate}{This standard conjugacy argument underlies \citet[Ch.~II, \S1]{dubrovinmatveevnovikov1976nonlinear}.}
\GSfield{Present status}{As of November 2026, the stated reconstruction is a standard established result.}
\GSfield{Assessment}{literal statement: correct; publication match: exact; verification: read the result; development: annotated; confidence: high}
\end{GSitem}
\begin{GSitem}{1974-340.6}{A}{Bloch continuation to the finite-gap curve}
\GSfield{Kind; source}{theorem; \GSpage{P000560}}
\GSfield{The notes}{Theorem 1 says that Bloch eigenfunctions analytically continue to \(\Gamma_N\), with normalization one at \(x_0\).}
\GSfield{Statement}{For a real nonsingular finite-gap potential, the normalized Bloch function extends meromorphically to \(\Gamma_N\setminus\{\infty\}\), with fixed pole divisor and an essential exponential singularity at infinity. ``Analytic continuation is not a claim of holomorphy at its poles or at infinity.}
\GSfield{Fate}{\citet[\S3]{dubrovinnovikov1974periodicity} is the matching 1974 result; this repeats items 1974-339.2--4 in a longer proof-oriented lecture.}
\GSfield{Present status}{As of November 2026, the stated reconstruction is a standard established result.}
\GSfield{Assessment}{literal statement: correct with missing hypotheses; publication match: exact; verification: read the result; development: annotated; confidence: high}
\end{GSitem}
\begin{GSitem}{1974-340.7}{A}{Finite-gap potentials are stationary for a higher KdV equation}
\GSfield{Kind; source}{theorem; \GSpage{P000560}}
\GSfield{The notes}{Theorem 0 states that every finite-zone potential is a stationary solution of a higher KdV equation.}
\GSfield{Statement}{For a real smooth periodic Schrödinger potential, having finitely many open gaps is equivalent to satisfying a nontrivial stationary equation in the KdV hierarchy, allowing a finite linear combination of the flows and integration constants. Equivalently, there is an odd-order differential operator commuting with \(L\) of the corresponding finite-gap type.}
\GSfield{Fate}{\citet[\S2 and the converse at the end of \S3]{dubrovinnovikov1974periodicity} gives both directions, referring the stationary-to-finite-gap direction to \citet{novikov1974periodic}. Novikov\textquotesingle s paper was received February 1974; the joint paper June 1974; both appeared in 1974. Earlier related occurrence: 1973-326.15.}
\GSfield{Present status}{As of November 2026, the stated reconstruction is a standard established result.}
\GSfield{Assessment}{literal statement: correct with missing hypotheses; publication match: exact; verification: read the result; development: annotated; confidence: high; problem \hyperref[prob:GP-1974-10]{GP-1974-10}}
\end{GSitem}
\begin{GSitem}{1974-340.8}{A}{The Jacobian description of the isospectral manifold}
\GSfield{Kind; source}{theorem; \GSpage{P000560}}
\GSfield{The notes}{Theorem 2 calls the set of potentials with fixed spectrum \(J(\Gamma_N)\).}
\GSfield{Statement}{For a fixed nonsingular finite-gap curve, admissible real divisor data identify the smooth real isospectral family with a real \(N\)-torus. Complex divisor coordinates identify an open part of the complexified construction with an open part of \(J(\Gamma_N)\), whose compactification is the Jacobian.}
\GSfield{Fate}{\citet[\S3]{dubrovinnovikov1974periodicity} supplies the algebraic-geometric identification. The real/complex and affine/compact distinctions refine the lecture\textquotesingle s shorthand.}
\GSfield{Present status}{As of November 2026, the stated reconstruction is a standard established result.}
\GSfield{Assessment}{literal statement: correct with missing hypotheses; publication match: exact; verification: read the result; development: annotated; confidence: high; problem \hyperref[prob:GP-1973-51]{GP-1973-51}}
\end{GSitem}
\begin{GSitem}{1974-340.9}{A}{Higher KdV flows on the isospectral Jacobian}
\GSfield{Kind; source}{claim; \GSpage{P000560}}
\GSfield{The notes}{KdV and its higher analogues act on the fixed-spectrum manifold.}
\GSfield{Statement}{The commuting hierarchy flows preserve the spectral curve and become constant-velocity translations in Abel coordinates on its Jacobian. On the admissible real component these are commuting flows on a real torus; periodicity in a particular time requires additional rationality of the velocity vector.}
\GSfield{Fate}{\citet[\S3, Abel coordinates]{dubrovinnovikov1974periodicity} and \citet[\S\S2--3]{dubrovinnovikov1974periodic} identify these flows (1974 submissions and publications).}
\GSfield{Present status}{As of November 2026, the stated reconstruction is a standard established result.}
\GSfield{Assessment}{literal statement: correct with missing hypotheses; publication match: exact; verification: read the result; development: annotated; confidence: high}
\end{GSitem}
\begin{GSitem}{1974-340.10}{E}{Monotonicity of a chosen band phase}
\GSfield{Kind; source}{claim; \GSpage{P000560}}
\GSfield{The notes}{The roots are said to run in one direction, with a parenthetical remark that this can be proved.}
\GSfield{Statement}{On the interior of a spectral band, choose the Bloch branch with positive conserved current \(j=\operatorname{Im}(\overline\psi\psi^{\prime})\). Its phase satisfies \(p^{\prime}(E)=(2j)^{-1}\int_{x_0}^{x_0+T}|\psi|^2dx>0\). The reciprocal multiplier travels in the opposite angular direction. Thus the two roots cannot both be said to rotate in the same direction without a branch convention.}
\GSfield{Reconstruction}{The specific monotonicity statement and current orientation are a reconstruction of the terse sentence.}
\GSfield{Fate}{The discriminant and phase theory is classical; compare \citet[Ch.~II, \S1]{dubrovinmatveevnovikov1976nonlinear}. The phase derivative identity is obtained by differentiating the differential equation and its Floquet boundary condition.}
\GSfield{Present status}{As of November 2026, the stated reconstruction is a standard established result.}
\GSfield{Assessment}{literal statement: correct with missing hypotheses; publication match: partial; verification: read the result; development: annotated; confidence: medium; problem \hyperref[prob:GP-1973-50]{GP-1973-50}}
\end{GSitem}
\begin{GSitem}{1974-340.11}{F}{The physical trapping and exponential-decay shorthand}
\GSfield{Kind; source}{remark; \GSpage{P000560}}
\GSfield{The notes}{The physical sketch contrasts transmission with trapping; a nearby fragment says that forbidden zones decay exponentially.}
\GSfield{Statement}{In an open gap the two Bloch solutions grow and decay exponentially in opposite spatial directions; neither is a nonzero solution decaying at both infinities for the unperturbed periodic operator. If the fragment instead refers to the widths of high-energy gaps, an exponential estimate needs regularity assumptions beyond the written periodicity. The notes do not identify which decay is intended.}
\GSfield{Fate}{\citet[Ch.~II, \S1]{dubrovinmatveevnovikov1976nonlinear} supplies the first Floquet interpretation. No exact assertion about gap-width asymptotics or a localized particle state can be recovered from this fragment.}
\GSfield{Present status}{As of November 2026, the source does not determine a sufficiently precise claim for an open-problem status.}
\GSfield{Assessment}{literal statement: not determinable; publication match: thematic only; verification: read the result; development: none; confidence: low}
\end{GSitem}
\begin{GSitem}{1974-340.12}{E}{Extrema of the discriminant and the sign of the lower bound}
\GSfield{Kind; source}{claim; \GSpage{P000560-P000561}}
\GSfield{The notes}{The graph is said to have infinitely many maxima and minima, with both written as \(\ge1\).}
\GSfield{Statement}{For a real periodic scalar Schrödinger operator, real critical points of the discriminant lie outside open band interiors. Its successive maxima satisfy \(\tau\ge1\) and minima satisfy \(\tau\le-1\), allowing equality at closed gaps. High-energy asymptotics yield infinitely many such extrema.}
\GSfield{Fate}{\citet[Ch.~II, \S1]{dubrovinmatveevnovikov1976nonlinear} gives the standard discriminant picture. The source drawing itself puts minima below \(-1\), confirming that the verbal inequality needs a sign or absolute value.}
\GSfield{Present status}{As of November 2026, the stated reconstruction is a standard established result.}
\GSfield{Assessment}{literal statement: false as written; publication match: exact; verification: read the result; development: annotated; confidence: high}
\end{GSitem}
\begin{GSitem}{1974-340.13}{E}{Open band edges, closed gaps and monodromy Jordan form}
\GSfield{Kind; source}{claim; \GSpage{P000561-P000562}}
\GSfield{The notes}{One selects simple roots of \(\tau(E)=\pm1\) for band endpoints; these give Jordan cells, while a diagonal matrix is degenerate.}
\GSfield{Statement}{At an endpoint of an open gap, \(\tau^2-1\) has a simple zero and monodromy has a nontrivial Jordan block with eigenvalue \(+1\) or \(-1\). At a closed gap, the zero is double and \(M=\pm I\). Removing these double points gives the finite-gap curve. This concerns a real self-adjoint scalar periodic equation.}
\GSfield{Fate}{\citet[Ch.~II, \S1]{dubrovinmatveevnovikov1976nonlinear} explains the scalar Floquet alternative. The notebook\textquotesingle s reference to Gelfand--Lidskii (1955) is broader; no exact claim from that paper is substituted here.}
\GSfield{Present status}{As of November 2026, the stated reconstruction is a standard established result.}
\GSfield{Assessment}{literal statement: correct with missing hypotheses; publication match: exact; verification: read the result; development: annotated; confidence: high}
\end{GSitem}
\begin{GSitem}{1974-340.14}{E}{The normalization-independent Riccati variable}
\GSfield{Kind; source}{construction; \GSpage{P000561}}
\GSfield{The notes}{The logarithmic derivative is \(\chi=-i\Psi^{\prime}/\Psi\); it does not depend on \(x_0\), is periodic in \(x\), and satisfies a Riccati equation.}
\GSfield{Statement}{On a fixed Bloch branch, away from zeros of \(\Psi\), multiplication by an \(x\)-independent normalization cancels in \(\chi\). Floquet periodicity gives \(\chi(x+T,E)=\chi(x,E)\), and substitution in \(L\Psi=E\Psi\) gives \(-i\chi_x+\chi^2+u-E=0\). The identities extend meromorphically across the divisor.}
\GSfield{Fate}{The Riccati construction and its role in conserved densities are recalled in \citet[\S2]{dubrovinnovikov1974periodicity}; the substitution itself is elementary.}
\GSfield{Present status}{As of November 2026, the stated reconstruction is a standard established result.}
\GSfield{Assessment}{literal statement: correct; publication match: exact; verification: read the result; development: annotated; confidence: high}
\end{GSitem}
\begin{GSitem}{1974-340.15}{E}{Amplitude, phase and the real part of the Riccati variable}
\GSfield{Kind; source}{claim; \GSpage{P000561}}
\GSfield{The notes}{The formulas are \(\chi_I=\tfrac12\chi_{R,x}/\chi_R\), \(|\Psi|^2=\chi_R(x_0)/\chi_R(x)\), and \(p=\int_0^T\chi_Rdx\) up to a phase integer.}
\GSfield{Statement}{For real \(E\) in a band interior, choose a branch with nonzero current. Separating the Riccati equation yields \(\chi_I=\tfrac12(\log|\chi_R|)_x\). Consequently \[\Psi(x)=\sqrt{\frac{\chi_R(x_0)}{\chi_R(x)}}\exp\!\left(i\int_{x_0}^x\chi_R(s)ds\right).\] The positive amplitude ratio is periodic, so \(p=\int_{x_0}^{x_0+T}\chi_Rdx\pmod{2\pi}\).}
\GSfield{Reconstruction}{The relation between real and imaginary parts applies on real bands, not literally to arbitrary complex \(E\).}
\GSfield{Fate}{\citet[\S3]{dubrovinnovikov1974periodicity} uses the amplitude-phase factorization. The source\textquotesingle s overwritten factor before \(\pi n\) is not recoverable; the multiplier fixes the modulus \(2\pi\).}
\GSfield{Present status}{As of November 2026, the stated reconstruction is a standard established result.}
\GSfield{Assessment}{literal statement: correct with missing hypotheses; publication match: exact; verification: read the result; development: annotated; confidence: high}
\end{GSitem}
\begin{GSitem}{1974-340.16}{E}{The high-energy Riccati expansion and local differential polynomials}
\GSfield{Kind; source}{method; \GSpage{P000561-P000563}}
\GSfield{The notes}{An expansion \(\chi\sim k+\sum_{m\ge1}\chi_m(x)/(2k)^m\), \(k=\sqrt E\), is obtained from the Riccati equation; the real part contains odd-index terms.}
\GSfield{Statement}{For a smooth potential, the formal recursion is \(\chi_1=-u\), \(\chi_{m+1}=i\partial_x\chi_m-\sum_{a+b=m}\chi_a\chi_b\). Thus \(\chi_2=-iu_x\) and \(\chi_3=u_{xx}-u^2\). Odd coefficients are real and even coefficients imaginary for real \(u\). The corresponding analytic high-energy expansion is asymptotic in suitable complex sectors; convergence of the infinite series is not asserted.}
\GSfield{Fate}{\citet[\S2]{dubrovinnovikov1974periodicity} records the recursion framework underlying the higher KdV integrals. The displayed first coefficients are checked directly from the notebook\textquotesingle s convention.}
\GSfield{Present status}{As of November 2026, the stated reconstruction is a standard established result.}
\GSfield{Assessment}{literal statement: correct with missing hypotheses; publication match: exact; verification: read the result; development: annotated; confidence: high}
\end{GSitem}
\begin{GSitem}{1974-340.17}{E}{The monodromy formula and the missing subtraction}
\GSfield{Kind; source}{claim; \GSpage{P000561}}
\GSfield{The notes}{Formula (1) has \(\sqrt{1-\tau^2}/\alpha_{21}\) plus \((i/2)\alpha_{11}\alpha_{22}/\alpha_{21}\), with no visible operator between the two diagonal entries.}
\GSfield{Statement}{With the row convention \(c(x+T)=a_{11}c(x)+a_{12}s(x)\), \(s(x+T)=a_{21}c(x)+a_{22}s(x)\), and multiplier \(\mu=\tau+i\sqrt{1-\tau^2}\), one has \[\chi(x_0,E)=\frac{\sqrt{1-\tau^2}}{a_{21}}+\frac i2\frac{a_{11}-a_{22}}{a_{21}}.\] Indeed \(\Psi=c+i\chi s\) gives \(\mu=a_{11}+i\chi a_{21}\). The formula is meromorphic where the denominator vanishes.}
\GSfield{Fate}{This is direct two-dimensional eigenvector algebra consistent with \citet[Ch.~II, \S1]{dubrovinmatveevnovikov1976nonlinear}. The required subtraction is a mathematical reconstruction, not a confidently identified hidden stroke.}
\GSfield{Present status}{As of November 2026, the stated reconstruction is a standard established result.}
\GSfield{Assessment}{literal statement: false as written; publication match: exact; verification: read the result; development: annotated; confidence: high}
\end{GSitem}
\begin{GSitem}{1974-340.18}{E}{The Wronskian normalization at the base point}
\GSfield{Kind; source}{claim; \GSpage{P000562}}
\GSfield{The notes}{The line reads \(\{c,\Psi\}|_{x_0}=i\Psi^{\prime}(x_0)=i\chi(x_0)\), with an overwritten middle symbol.}
\GSfield{Statement}{Using \(W(c,\Psi)=c\Psi^{\prime}-c^{\prime}\Psi\), \(c(x_0)=1,c^{\prime}(x_0)=0,\Psi(x_0)=1\), the correct identity is \(W(c,\Psi)=\Psi^{\prime}(x_0)=i\chi(x_0)\). The extra factor \(i\) before \(\Psi^{\prime}\) is inconsistent with these definitions.}
\GSfield{Reconstruction}{The overwritten glyph remains palaeographically uncertain; only the mathematical identity is fixed.}
\GSfield{Fate}{The check uses the preceding normalization and Wronskian constancy, the same standard construction as \citet[Ch.~II, \S1]{dubrovinmatveevnovikov1976nonlinear}.}
\GSfield{Present status}{As of November 2026, the stated reconstruction is a standard established result.}
\GSfield{Assessment}{literal statement: false as written; publication match: exact; verification: read the result; development: annotated; confidence: high}
\end{GSitem}
\begin{GSitem}{1974-340.19}{A}{Cancelling the closed-gap entire factor}
\GSfield{Kind; source}{construction; \GSpage{P000562}}
\GSfield{The notes}{The notes divide \(a_{21}\) by \(\sqrt{(1-\tau^2)/R}\), asserting that the quotient is entire.}
\GSfield{Statement}{In the finite-gap real periodic case, the zeros of \(1-\tau^2\) not included in \(R\) are precisely the closed-gap double zeros. Therefore \((1-\tau^2)/R\) has an entire square root \(H(E)\), up to sign, and \(a_{21}\) vanishes at its zeros. Division removes those simple zeros and gives an entire function \(P=a_{21}/H\).}
\GSfield{Reconstruction}{The multiplicities and removability are needed; a bare quotient of two entire functions need not be entire.}
\GSfield{Fate}{The finite-gap reduction is developed in \citet[Ch.~II, \S1]{dubrovinmatveevnovikov1976nonlinear} and announced in \citet[\S3]{dubrovinnovikov1974periodicity}.}
\GSfield{Present status}{As of November 2026, the stated reconstruction is a standard established result.}
\GSfield{Assessment}{literal statement: correct with missing hypotheses; publication match: exact; verification: read the result; development: annotated; confidence: high}
\end{GSitem}
\begin{GSitem}{1974-340.20}{A}{Why the reduced entire function is a polynomial}
\GSfield{Kind; source}{claim; \GSpage{P000562}}
\GSfield{The notes}{From \(\chi_R\sim k\) and the leading power of \(\sqrt R\), the notes conclude that the reduced function is a polynomial in \(E\).}
\GSfield{Statement}{The conclusion requires complex growth control in addition to removability. If an entire \(P\) satisfies a global bound \(|P(E)|\le C(1+|E|)^N\), Cauchy estimates give \(\deg P\le N\); the leading normalization then gives degree \(N\). Asymptotics on the positive real axis alone do not imply this: \(E^N+e^{-E}\) has the same leading behaviour there and is not a polynomial.}
\GSfield{Reconstruction}{The global polynomial bound is a sufficient editorial reconstruction, not a hypothesis visibly stated in the notes or a proof that the notes supplied it.}
\GSfield{Fate}{The established finite-gap polynomial conclusion is in \citet[Ch.~II, \S1]{dubrovinmatveevnovikov1976nonlinear}. The notebook omits the sectorial estimates and their global complex-analytic use; the elementary implication above audits that gap in the sketch.}
\GSfield{Present status}{As of November 2026, the stated reconstruction is a standard established result.}
\GSfield{Assessment}{literal statement: correct with missing hypotheses; publication match: partial; verification: read the result; development: annotated; confidence: high}
\end{GSitem}
\begin{GSitem}{1974-340.21}{A}{The polynomial representation of the Riccati variable}
\GSfield{Kind; source}{construction; \GSpage{P000562}}
\GSfield{The notes}{The final formulas are \(\chi_R=\sqrt R/P\), \(P=\prod(E-\gamma_i(x))\), and \(\chi=(\sqrt R-iP_x/2)/P\).}
\GSfield{Statement}{For the nonsingular finite-gap construction, the monic degree-\(N\) divisor polynomial \(P(E,x)\) gives \[\chi=\frac{\sqrt{R(E)}-\frac i2P_x(E,x)}{P(E,x)}.\] Compatibility with the Schrödinger equation is the identity \[R=(E-u)P^2+\tfrac12PP_{xx}-\tfrac14P_x^2,\] and its derivative yields \(P_{xxx}-4(u-E)P_x-2u_xP=0\). The notation \(\chi_R\) away from real bands means analytic continuation, not the nonholomorphic operation of taking a real part.}
\GSfield{Fate}{\citet[\S3]{dubrovinnovikov1974periodicity} and \citet[Ch.~II]{dubrovinmatveevnovikov1976nonlinear} supply the classical finite-gap formula. The companion\textquotesingle s separate compatibility audit expands the identity directly.}
\GSfield{Present status}{As of November 2026, the stated reconstruction is a standard established result.}
\GSfield{Assessment}{literal statement: correct with missing hypotheses; publication match: exact; verification: read the result; development: annotated; confidence: high}
\end{GSitem}
\begin{GSitem}{1974-340.22}{E}{The leading power of the spectral square root}
\GSfield{Kind; source}{claim; \GSpage{P000562}}
\GSfield{The notes}{The line gives \(\sqrt{R(E)}\sim E^N\), while the next numerator is \(kE^N\).}
\GSfield{Statement}{Since \(R\) is monic of degree \(2N+1\), one has \(\sqrt{R(E)}\sim E^{N+1/2}=kE^N\) on a chosen sectorial square-root branch. Division by a monic degree-\(N\) polynomial then gives \(\chi_R\sim k\).}
\GSfield{Fate}{The correction follows from the degree of the polynomial written in the notes and agrees with \citet[\S3]{dubrovinnovikov1974periodicity}.}
\GSfield{Present status}{As of November 2026, the stated reconstruction is a standard established result.}
\GSfield{Assessment}{literal statement: false as written; publication match: exact; verification: read the result; development: annotated; confidence: high}
\end{GSitem}
\begin{GSitem}{1974-340.23}{E}{Conserved Riccati integrals and the stationary variational equation}
\GSfield{Kind; source}{construction; \GSpage{P000562-P000563}}
\GSfield{The notes}{The integrals \(I_n[u]=\int_0^T\chi_{2n+3}(x)dx\) are introduced, followed by
\[\frac{\delta}{\delta u}\Bigl(\sum c_nI_n\Bigr)=d.\]}
\GSfield{Statement}{For smooth periodic \(u\), the odd Riccati densities integrate to local conserved functionals of the KdV hierarchy, up to normalization and total derivatives. A stationary linear combination of flows \(u_t=\partial_x\delta(\sum c_nI_n)/\delta u\) satisfies \(\delta(\sum c_nI_n)/\delta u=d\), where the sum is finite and \(d\) is a spatial constant.}
\GSfield{Fate}{\citet[\S2, equations (6)--(8), Theorem~2.1]{novikov1974periodic} publishes the stationary variational construction and its complete integrability in July 1974, before the September talk. \citet[\S2]{dubrovinnovikov1974periodicity} recalls and attributes this result; its new converse is treated separately in item 1974-340.7.}
\GSfield{Present status}{As of November 2026, the stated reconstruction is a standard established result.}
\GSfield{Assessment}{literal statement: correct with missing hypotheses; publication match: exact; verification: read the result; development: annotated; confidence: high}
\end{GSitem}
\begin{GSitem}{1974-340.24}{E}{The derivative order of the Riccati densities}
\GSfield{Kind; source}{claim; \GSpage{P000563}}
\GSfield{The notes}{The odd coefficient is written \(\chi_{2n+1}=Q(u,u^{\prime},\ldots,u^{(n-1)})\).}
\GSfield{Statement}{Before integrating by parts, \(\chi_{2n+1}\) generally contains derivatives through order \(2n\). Its leading linear term is \((-1)^{n+1}u^{(2n)}\); already \(\chi_3=u_{xx}-u^2\) contradicts the written order. Integrated densities may be replaced by lower-order representatives modulo total derivatives, but that is a different statement.}
\GSfield{Fate}{The recursion in item 1974-340.16 verifies the correction directly; compare the density construction of \citet[\S2]{dubrovinnovikov1974periodicity}. The doubtful exponent is retained in the source transcription and emended only mathematically.}
\GSfield{Present status}{As of November 2026, the stated reconstruction is a standard established result.}
\GSfield{Assessment}{literal statement: false as written; publication match: exact; verification: read the result; development: annotated; confidence: high}
\end{GSitem}

\GStalk{341}{1974-09-23}{Enrico Giusti; I. M. Gelfand (intervention)}{P000564-P000565}{Singular weak solutions and partial regularity of elliptic systems}
Giusti states the energy formulation, contrasts scalar and two-dimensional regularity with vectorial counterexamples, and gives an analytic coefficient tensor for the singular map \(x/|x|\). The first coefficient in the copied tensor is incompatible with that solution; an exact residual and a repaired uniformly elliptic tensor are proved in the audit. The discussion ends with a variational-existence question and Gelfand\textquotesingle s unspecified proposal for an analogue of Hironaka resolution. The dilogarithm material above Giusti\textquotesingle s heading belongs to record 342.

\begin{GSitem}{1974-341.1}{E}{The energy weak formulation of a quasilinear system}
\GSfield{Kind; source}{definition; \GSpage{P000564-P000565}}
\GSfield{The notes}{The weak equation sums \(a_{ij}^{hk}(x,u)D_i u^hD_j\varphi^k\); the following strong equation instead writes \(D_i u^k\) inside the sum over \(h\).}
\GSfield{Statement}{For \(u\in W^{1,2}_{\mathrm{loc}}(\Omega;\mathbb R^N)\), bounded coefficients and compactly supported smooth tests, the weak system is \[\sum_{i,j,h,k}\int_\Omega a_{ij}^{hk}(x,u)D_i u^hD_j\varphi^k\,dx=0.\] It is equivalent to \(\sum_{i,j,h}D_j(a_{ij}^{hk}(x,u)D_i u^h)=0\) in distributions for each \(k\). Here \(D_i=\partial/\partial x_i\), and the unknown argument lies in \(\mathbb R^N\).}
\GSfield{Fate}{\citet[\S3.2, equation (3.5)]{mingione2006regularity} uses the same energy-system formulation. The repeated component index in the strong equation is a correctable mismatch with the preceding weak form.}
\GSfield{Present status}{As of November 2026, the stated reconstruction is a standard established result.}
\GSfield{Assessment}{literal statement: false as written; publication match: exact; verification: read the result; development: annotated; confidence: high}
\end{GSitem}
\begin{GSitem}{1974-341.2}{E}{Strong ellipticity and the distinction from coefficient smoothness}
\GSfield{Kind; source}{definition; \GSpage{P000564-P000565}}
\GSfield{The notes}{The displayed quadratic form is bounded below by \(|\xi|^2\).}
\GSfield{Statement}{The relevant system assumption is \(\sum a_{ij}^{hk}(x,v)\xi_i^h\xi_j^k\ge\nu|\xi|^2\), for all real \(N\times n\) matrices \(\xi\), together with a uniform upper bound \(|a(x,v)|\le\Lambda\). Analyticity of \(a\) as a function of \(v\) does not imply continuity of the composite \(a(x,u(x))\) if \(u\) is discontinuous.}
\GSfield{Reconstruction}{The uniform upper bound, omitted in the displayed inequality, is part of the standard energy estimates.}
\GSfield{Fate}{\citet[\S3.2]{mingione2006regularity} explains exactly this mechanism in the Giusti--Miranda example. Ellipticity is an algebraic lower bound, not an automatic smoothness conclusion for systems.}
\GSfield{Present status}{As of November 2026, the stated reconstruction is a standard established result.}
\GSfield{Assessment}{literal statement: correct with missing hypotheses; publication match: exact; verification: read the result; development: annotated; confidence: high}
\end{GSitem}
\begin{GSitem}{1974-341.3}{E}{The Sobolev energy space and zero trace}
\GSfield{Kind; source}{definition; \GSpage{P000565}}
\GSfield{The notes}{The notes define \(W^{1,2}\) by square-integrable function and gradient and describe \(W^{1,2}_0\) as zero boundary values.}
\GSfield{Statement}{For an open set \(\Omega\), \(W^{1,2}(\Omega;\mathbb R^N)\) consists of functions with all first weak derivatives in \(L^2\). The space \(W_0^{1,2}\) is the closure of \(C_c^\infty\) in that norm. On a Lipschitz domain this is precisely the subspace with zero Sobolev trace.}
\GSfield{Fate}{This is the standard framework explicitly used in \citet[\S3]{mingione2006regularity}. The trace formulation needs a regular boundary, whereas the closure definition does not.}
\GSfield{Present status}{As of November 2026, the stated reconstruction is a standard established result.}
\GSfield{Assessment}{literal statement: correct with missing hypotheses; publication match: exact; verification: read the result; development: annotated; confidence: high}
\end{GSitem}
\begin{GSitem}{1974-341.4}{E}{Scalar or two-dimensional regularity in the energy class}
\GSfield{Kind; source}{theorem; \GSpage{P000565}}
\GSfield{The notes}{The cases \(N=1\) or \(n\le2\) are called regular, but the line writes \(W^{1,1}\); smooth coefficients are then said to give smooth solutions.}
\GSfield{Statement}{For the homogeneous second-order divergence systems under boundedness and uniform strong ellipticity, scalar energy weak solutions are locally Hölder continuous; in two dimensions the same Hölder conclusion holds for systems, using higher integrability and Sobolev embedding. The relevant class is \(W^{1,2}\). For smooth dependence on \((x,u)\), elliptic bootstrap applies on a region where the solution is already regular; it does not smooth every vectorial weak solution in higher dimensions.}
\GSfield{Fate}{\citet[\S3.1 and Theorem~4.1]{mingione2006regularity} explains the two-dimensional energy argument; \citet[pp.~28--29]{johnstara1990regularity} states the scalar and two-variable alternatives. These are classical pre-talk facts.}
\GSfield{Present status}{As of November 2026, the stated reconstruction is a standard established result.}
\GSfield{Assessment}{literal statement: false as written; publication match: exact; verification: read the result; development: annotated; confidence: high}
\end{GSitem}
\begin{GSitem}{1974-341.5}{E}{The Giusti coefficient tensor and the singular radial solution}
\GSfield{Kind; source}{construction; \GSpage{P000565}}
\GSfield{The notes}{For \(n=N\ge3\), the first factor has coefficient \(1/(n-2)\), the second \(4/(n-2)\), and \(u(x)=x/|x|\) is claimed to solve the homogeneous system.}
\GSfield{Statement}{For \(A_{ij}^{hk}(v)=\delta_{ij}\delta_{hk}+(\delta_{jk}+\alpha v_jv_k/(1+|v|^2))(\delta_{ih}+\beta v_iv_h/(1+|v|^2))\), substitution gives \[\sum_{i,j,h}\partial_j(A_{ij}^{hk}(u)\partial_i u^h)=(n-1)\left(\frac{\alpha(n-2)}2-2\right)\frac{x_k}{|x|^3}.\] Thus the written \(\alpha=1/(n-2)\) fails away from the origin. Taking \(\alpha=\beta=4/(n-2)\) gives analytic bounded uniformly elliptic coefficients and a discontinuous weak solution \(u\in W^{1,2}(B_1;\mathbb R^n)\), with trace \(u=x\).}
\GSfield{Reconstruction}{The replacement of the first 1 by 4 is a mathematically forced reconstruction within this tensor ansatz and agrees with the later published formula. It is not a claim to have read a hidden 4 in the scan.}
\GSfield{Fate}{\citet[equation (3.3)]{mingione2006regularity} prints both factors with \(4/(n-2)\), attributing the construction to Giusti--Miranda (1968). The original paper was not retrieved. The exact contraction, residual, ellipticity and weak extension at zero are proved in this companion\textquotesingle s audit.}
\GSfield{Present status}{As of November 2026, the corrected energy-class example is verified; the literal tensor is incompatible with the stated solution.}
\GSfield{Assessment}{literal statement: false as written; publication match: partial; verification: indirect; development: annotated; confidence: high}
\end{GSitem}
\begin{GSitem}{1974-341.6}{F}{A three-variable, two-component counterexample without a specified system}
\GSfield{Kind; source}{claim; \GSpage{P000565}}
\GSfield{The notes}{The notes assert an example with \(n=3,N=2\), giving neither its coefficient formula nor a reference.}
\GSfield{Statement}{The exact claim depends on whether coefficients are merely measurable in \(x\), analytic in \((x,u)\), or whether lower-order quadratic gradient terms are allowed. The preceding radial map has \(N=n\) and does not itself establish the two-component case.}
\GSfield{Fate}{\citet[Theorem~3.1]{frehse2007irregular} gives a later two-real-component example in dimension three through a complex scalar equation with measurable spatial coefficients (preprint July 2007). Its discussion of the 1973 example involves a different semilinear equation. These are partial contextual matches, not verification of the same analytic homogeneous class intended here.}
\GSfield{Present status}{As of November 2026, counterexamples in neighbouring classes are known, but the precise 1974 assertion and its coefficient class have not been identified.}
\GSfield{Assessment}{literal statement: not determinable; publication match: partial; verification: read the result; development: annotated; confidence: low}
\end{GSitem}
\begin{GSitem}{1974-341.7}{F}{The unverified uniqueness assertion for singular Dirichlet data}
\GSfield{Kind; source}{claim; \GSpage{P000565}}
\GSfield{The notes}{For \(|x|=1\), \(u_i=x_i\), the notes write \(\exists!\) solution and say that it is nonregular.}
\GSfield{Statement}{The corrected tensor supplies at least one singular \(W^{1,2}\) solution with trace \(x\). Uniqueness among all weak solutions does not follow from positive definiteness of \(A(u)\), because the coefficient depends on the unknown. No monotonicity of the full nonlinear operator, admissible subclass or uniqueness proof is recorded.}
\GSfield{Fate}{\citet[\S3.2]{mingione2006regularity} distinguishes a variational construction from the homogeneous system. The source matching and the companion\textquotesingle s residual proof establish existence only; no exact uniqueness theorem was verified.}
\GSfield{Present status}{As of November 2026, existence is verified for the reconstructed system; the uniqueness claim in the notebook remains unverified and is not used.}
\GSfield{Assessment}{literal statement: not determinable; publication match: none found; verification: unresolved lead; development: annotated; confidence: low}
\end{GSitem}
\begin{GSitem}{1974-341.8}{E}{A relatively closed singular set of dimension strictly below n minus two}
\GSfield{Kind; source}{theorem; \GSpage{P000565}}
\GSfield{The notes}{For continuous coefficients, the regular set \(\Omega_0\) is open, its complement \(K\) is relatively closed, \(\mathcal H^{n-2}(K)=0\), and \(\dim K<n-2\).}
\GSfield{Statement}{For \(n\ge3\), let \(u\in W^{1,2}_{\mathrm{loc}}(\Omega;\mathbb R^N)\) solve \(\operatorname{div}(A(x,u)Du)=0\), with bounded uniformly strongly elliptic coefficients uniformly continuous in \((x,v)\). There is an open regular set \(\Omega_u\) where \(u\) is locally Hölder of every exponent below one, and \(\dim_H(\Omega\setminus\Omega_u)<n-2\). In particular \(\mathcal H^{n-2}(K)=0\). The strict dimension statement uses more than the latter zero-measure statement alone.}
\GSfield{Reconstruction}{Uniform bounds, uniform continuity, the energy class and the meaning of regularity are restored. The original source merely says continuous coefficients.}
\GSfield{Fate}{\citet[\S1.1, equations (1.1)--(1.2)]{defilippismingionenowak2026partial} restates this precise classical local theorem and attributes its foundation to Giusti--Miranda and Morrey (1968). The 2025 preprint, revised July 2026, is used to verify the hypotheses and bound, not to date the original result to 2026.}
\GSfield{Present status}{As of November 2026, the stated reconstruction is a standard established result.}
\GSfield{Assessment}{literal statement: correct with missing hypotheses; publication match: exact; verification: indirect; development: annotated; confidence: high}
\end{GSitem}
\begin{GSitem}{1974-341.9}{F}{The compressed higher-order partial-regularity bound}
\GSfield{Kind; source}{claim; \GSpage{P000565}}
\GSfield{The notes}{For order \(m>2\), the notes write \(\dim K\le n-2\).}
\GSfield{Statement}{The notebook does not specify whether \(m\) is the differential order or half-order, which derivatives are required to be regular, or the coefficient dependence and growth. Therefore the displayed dimension bound is not a complete theorem.}
\GSfield{Fate}{\citet[Theorems 1--2]{giusti1969partial} treats arbitrary-order quasilinear systems before the talk (received 1968; published 1969), but the opened scan did not permit a reliable exact-hypothesis comparison with this one-line bound. No specific modern higher-order theorem is substituted.}
\GSfield{Present status}{As of November 2026, the source does not determine a sufficiently precise claim for an open-problem status.}
\GSfield{Assessment}{literal statement: not determinable; publication match: partial; verification: unresolved lead; development: annotated; confidence: low}
\end{GSitem}
\begin{GSitem}{1974-341.10}{E}{Variational problems with nonregular minimizers}
\GSfield{Kind; source}{question; \GSpage{P000565}}
\GSfield{The notes}{The notes ask whether there are variational problems with nonregular solutions.}
\GSfield{Statement}{In the broad class of vector-valued uniformly elliptic variational integrals \(\int F(x,u,Du)\), are there energy-class minimizing maps that are not continuous? Yes: analytic dependence of the quadratic coefficient tensor on \(u\) is compatible with a discontinuous local minimizer in sufficiently large dimensions. The question does not impose autonomous dependence only on \(Du\).}
\GSfield{Reconstruction}{“Solutions” is read as minimizers, the stronger variational interpretation; the distinction from merely stationary points is explicit.}
\GSfield{Fate}{\citet[\S3.2, equations (3.3)--(3.4)]{mingione2006regularity} explicitly records the Giusti--Miranda 1968 variational example. Thus the broad question was already answered before this talk. Stronger later autonomous examples are not treated as an answer to an unstated restriction.}
\GSfield{Present status}{As of November 2026, the stated reconstruction is a standard established result.}
\GSfield{Assessment}{literal statement: correct with missing hypotheses; publication match: exact; verification: indirect; development: annotated; confidence: high; problem \hyperref[prob:GP-1974-02]{GP-1974-02}}
\end{GSitem}
\begin{GSitem}{1974-341.11}{F}{Gelfand’s proposed analogue of resolution of singularities}
\GSfield{Kind; source}{question; \GSpage{P000565}; I. M. Gelfand}
\GSfield{The notes}{Gelfand asks for an analogue of Hironaka resolution of singularities for elliptic equations.}
\GSfield{Statement}{No class of equations, admissible modification, transformed unknown, or required regularity after modification is specified. Resolving analytic zero sets, desingularizing coefficients and resolving singular weak solutions are distinct possible problems. The analogy alone does not select a definite conjecture.}
\GSfield{Fate}{The partial-regularity and counterexample literature above concerns the size and occurrence of singular sets, not a universal resolution operation. Searches for an exact publication of this elliptic-resolution programme did not identify one; Hironaka\textquotesingle s algebraic theorem cannot be transferred without a specified analytic statement.}
\GSfield{Present status}{As of November 2026, the recorded programme remains insufficiently specified for a mathematical open-problem status or a claimed new solution.}
\GSfield{Assessment}{literal statement: not determinable; publication match: none found; verification: unresolved lead; development: none; confidence: low}
\end{GSitem}

\GStalk{342}{1974-09-23}{I. M. Gelfand}{P000563-P000567}{Combinatorial characteristic classes and two dilogarithm integrals}
The opening of P000564 announces the work of Gabrielov, Gelfand and Losik in combinatorial topology and an analytic approach. It gives a non-elementary logarithmic integral, a Basel evaluation with the wrong sign, and a second, correct, golden-ratio evaluation said to come from topology. No characteristic-class construction is actually recorded. The following Giusti lecture belongs to record 341; the separately headed recap on P000566 is assigned to record 343. These are editorial passage boundaries inside overlapping catalogue page ranges.

\begin{GSitem}{1974-342.1}{A}{Analytic expressions in combinatorial characteristic-class calculations}
\GSfield{Kind; source}{method; \GSpage{P000564, opening}}
\GSfield{The notes}{\ru{О раб. Габриэлова, И. Гельфанда и Лосика по комбинаторной топологии}; \ru{Интересно сделать аналитически}. The logarithmic integrals are said to occur in a finite combinatorial problem.}
\GSfield{Statement}{The identifiable programme is to express characteristic classes of triangulated smooth manifolds by local combinatorial data, with analytic functions entering the formula. The page does not specify a triangulation, cochain, normalization, or the particular class, and hence does not state a theorem that can be reconstructed uniquely.}
\GSfield{Fate}{\citet[pp.~103--115]{gabrielovgelfandlosik1975computation} and \citet[pp.~186--202]{gabrielovgelfandlosik1975calculus} are the directly attributable publications in 1975 (received 2 December 1974 and 31 March 1975). Their subject and authors match; the notebook does not support an exact theorem match.}
\GSfield{Present status}{As of November 2026, the publication family is identified. The particular finite combinatorial problem invoked by this paragraph remains unspecified; no new open problem is inferred.}
\GSfield{Assessment}{literal statement: not determinable; publication match: partial; verification: publisher metadata checked; development: none; confidence: medium}
\end{GSitem}
\begin{GSitem}{1974-342.2}{E}{The logarithmic integral and the Rogers dilogarithm}
\GSfield{Kind; source}{claim; \GSpage{P000564, first display}}
\GSfield{The notes}{The notes write \(\int_{1-x}^{x}\{\log|t|/(1-t)+\log|1-t|/t\}\,dt\), saying it cannot be evaluated in elementary functions.}
\GSfield{Statement}{For \(0<x<1\), set \(L(x)=\operatorname{Li}_2(x)+\tfrac12\log x\log(1-x)\). Then the displayed integral equals \(-2[L(x)-L(1-x)]\). This is a dilogarithmic antiderivative, not an assertion that every individual special value is non-elementary.}
\GSfield{Reconstruction}{The domain and the explicit identification with the real Rogers dilogarithm are editorial; the notebook gives only the integral.}
\GSfield{Fate}{This is classical dilogarithm calculus, documented with the derivative and reflection formula in \citet[Chapter II, \S1.A, pp.~23--24]{zagier2007dilogarithm}. The modern reference does not date the formula to 2007.}
\GSfield{Present status}{As of November 2026, this is an established result in the stated scope; the qualification in the statement is essential.}
\GSfield{Assessment}{literal statement: correct; publication match: exact; verification: read the result; development: none; confidence: high}
\end{GSitem}
\begin{GSitem}{1974-342.3}{E}{The sign of the Basel integral}
\GSfield{Kind; source}{claim; \GSpage{P000564, Basel display}}
\GSfield{The notes}{The scan explicitly has \(\int_0^1\log|1-t|\,dt/t=+\pi^2/6\), followed by a positive power-series sum.}
\GSfield{Statement}{The convergent improper integral is \(\int_0^1\log(1-t)\,dt/t=-\zeta(2)=-\pi^2/6\). Termwise integration of \(-\log(1-t)=\sum_{n\ge1}t^n/n\) is justified by monotone convergence.}
\GSfield{Fate}{The corrected identity is classical; see \citet[Eqs.~25.12.1--2]{nist2026polylogarithms}. The error is already in the handwriting, not introduced by the English translation.}
\GSfield{Present status}{As of November 2026, this is an established result in the stated scope; the qualification in the statement is essential.}
\GSfield{Assessment}{literal statement: false as written; publication match: exact; verification: read the result; development: annotated; confidence: high}
\end{GSitem}
\begin{GSitem}{1974-342.4}{E}{A correct golden-ratio dilogarithm evaluation}
\GSfield{Kind; source}{claim; \GSpage{P000564, second definite integral}}
\GSfield{The notes}{For \(a=\sqrt5-2\), the notes give \(\int_{-a}^a\{\log|t|/(1-t)+\log|1-t|/t\}\,dt=-\pi^2/6\), adding \ru{из топологических соображений}.}
\GSfield{Statement}{The improper integral through \(t=0\) exists and the stated value is correct. It follows from the reflection and five-term equations for the real Rogers dilogarithm; the notebook supplies no topological argument.}
\GSfield{Fate}{The analytic identity reduces to classical functional equations, as in \citet[Chapter II, \S1.A]{zagier2007dilogarithm}. A complete verification and two independent numerical checks are supplied in the verification appendix. This establishes the reading without attributing our proof to the speaker.}
\GSfield{Present status}{As of November 2026, this is an established result in the stated scope; the qualification in the statement is essential.}
\GSfield{Assessment}{literal statement: correct; publication match: exact; verification: read the result; development: annotated; confidence: high}
\end{GSitem}

\GStalk{343}{1974-09-23}{I. M. Gelfand}{P000565-P000570}{Finite-gap Schrödinger operators, Bloch functions, and an entire-function proposal}
P000566 is explicitly headed as Gelfand\textquotesingle s repetition of Dubrovin\textquotesingle s previous talk. It reconstructs periodic spectral bands using monodromy, then packages the finite-gap theory in a logarithmic derivative, quasimomentum, and a polynomial divisor. P000568 appends Gelfand\textquotesingle s question about entire functions of order one and a Dubrovin assertion on the zeros and poles of the Bloch function. The entire-function paragraph lacks the admissibility and normalization conditions necessary for a determinate inverse theorem.

\begin{GSitem}{1974-343.1}{E}{Periodic Hill operator and bounded Bloch solutions}
\GSfield{Kind; source}{definition; \GSpage{P000566, opening}}
\GSfield{The notes}{\(-y^{\prime\prime}+u(x)y=Ey\); the written periodicity is \(u(x+T)=u(y)\). The spectrum is described as the energies having bounded eigenfunctions.}
\GSfield{Statement}{For real \(T\)-periodic \(u\in C^\infty(\mathbb R)\), let \(L=-d^2/dx^2+u\) on \(L^2(\mathbb R)\), domain \(H^2\). Its spectrum consists precisely of real energies with a nonzero bounded solution of \(Ly=Ey\). These are generalized eigenfunctions, generally not \(L^2(\mathbb R)\) eigenvectors.}
\GSfield{Reconstruction}{Real smooth periodic potential and operator domain are stated explicitly; the displayed u(y) is emended to u(x).}
\GSfield{Fate}{Classical Floquet theory, recapitulated in \citet[Chapter 2, \S1]{dubrovinmatveevnovikov1976nonlinear}; compare earlier occurrence 1973-326.2 and 1973-326.8.}
\GSfield{Present status}{As of November 2026, this is an established result in the stated scope; the qualification in the statement is essential.}
\GSfield{Assessment}{literal statement: false as written; publication match: exact; verification: read the result; development: none; confidence: high}
\end{GSitem}
\begin{GSitem}{1974-343.2}{A}{Characterizing finite-gap potentials by stationary KdV}
\GSfield{Kind; source}{question; \GSpage{P000566, Novikov problem and theorem}}
\GSfield{The notes}{\ru{Проблема Новикова: когда есть конечн. число запрещ. зон?} The answer is written as a nonlinear ordinary differential equation \(K(u,u^{\prime},\ldots,u^{(2N+1)})=0\) from the KdV hierarchy.}
\GSfield{Statement}{A real smooth periodic potential has finitely many open bounded spectral gaps if and only if it is stationary for a nontrivial finite real linear combination of the KdV hierarchy flows (equivalently, its Schr\"odinger operator commutes with an odd-order scalar differential operator). The least nondegenerate order corresponds to the number of open gaps; a higher stationary equation need not have all its gaps open.}
\GSfield{Reconstruction}{The phrase KdV equations is interpreted as stationary finite linear combinations, not the evolutionary PDE \(u_t=6uu_x-u_{xxx}\). Publication within 1974 is clear; precise issue-day priority relative to 23 September is not asserted.}
\GSfield{Fate}{\citet[\S2]{novikov1974periodic}, received 18 February 1974 and published July 1974, proves stationary KdV implies finite-gap. The converse is explicit in \citet[\S3]{dubrovinnovikov1974periodicity}, received 3 June 1974 and published later in 1974. The full account is \citet[Chapter 2, \S2]{dubrovinmatveevnovikov1976nonlinear}. Class A refers to that converse and full characterization, not to the already published forward implication. Earlier occurrences include 1973-326.15 and 1974-340.7.}
\GSfield{Present status}{As of November 2026, this is an established result in the stated scope; the qualification in the statement is essential.}
\GSfield{Assessment}{literal statement: correct with missing hypotheses; publication match: exact; verification: read the result; development: none; confidence: high; problem \hyperref[prob:GP-1974-10]{GP-1974-10}}
\end{GSitem}
\begin{GSitem}{1974-343.3}{E}{Monodromy and its discriminant}
\GSfield{Kind; source}{construction; \GSpage{P000566, fundamental solutions}}
\GSfield{The notes}{The normalized \(c,s\) solutions, the translation matrix \(\widehat T=(\alpha_{ij})\), \(\det\widehat T=1\), and \(\lambda^2-2\tau\lambda+1=0\) with \(2\tau=\alpha_{11}+\alpha_{22}\) are written.}
\GSfield{Statement}{Normalize \(c(x_0)=s^{\prime}(x_0)=1\), \(c^{\prime}(x_0)=s(x_0)=0\). The period map on the two-dimensional solution space has determinant one by the constant Wronskian. Its eigenvalues are \(\tau(E)\pm\sqrt{\tau(E)^2-1}\); every matrix entry and \(\tau\) is entire in \(E\).}
\GSfield{Fate}{Classical monodromy construction; \citet[Chapter 2, \S1, Eqs.~(2.1.1)--(2.1.8)]{dubrovinmatveevnovikov1976nonlinear}. Earlier occurrence: 1973-326.7.}
\GSfield{Present status}{As of November 2026, this is an established result in the stated scope; the qualification in the statement is essential.}
\GSfield{Assessment}{literal statement: correct; publication match: exact; verification: read the result; development: none; confidence: high}
\end{GSitem}
\begin{GSitem}{1974-343.4}{E}{Allowed and forbidden intervals from the discriminant}
\GSfield{Kind; source}{theorem; \GSpage{P000566, band criterion and drawings}}
\GSfield{The notes}{An allowed energy has a Bloch function with \(|\lambda|=1\); \(\tau^2\le1\) is allowed and \(\tau^2>1\) forbidden. Two drawings correct the oscillation of the discriminant.}
\GSfield{Statement}{For the real periodic problem, \(\sigma(L)=\{E\in\mathbb R:|\tau(E)|\le1\}\). Periodic and antiperiodic eigenvalues give band endpoints. Critical points of the discriminant above the spectral bottom lie in the closures of spectral gaps, and extrema alternate at levels at least1 and at most−1 respectively. A closed gap means a tangency at \(+1\) or \(-1\), not that the entire discriminant becomes constant.}
\GSfield{Fate}{Classical Hill spectral theory; \citet[Chapter 2, \S1]{dubrovinmatveevnovikov1976nonlinear}. The notebook\textquotesingle s later \(\max\) and \(\min\) both equal 1 is a sign omission, visible after the corrected sketch.}
\GSfield{Present status}{As of November 2026, this is an established result in the stated scope; the qualification in the statement is essential.}
\GSfield{Assessment}{literal statement: correct with missing hypotheses; publication match: exact; verification: read the result; development: annotated; confidence: high}
\end{GSitem}
\begin{GSitem}{1974-343.5}{E}{Closed gaps have scalar monodromy}
\GSfield{Kind; source}{theorem; \GSpage{P000566, double roots}}
\GSfield{The notes}{If \(\tau^2(E)=1\) has a multiple root, the notes infer \(\alpha_{12}=\alpha_{21}=0\), distinguishing a Jordan band edge from a diagonal one.}
\GSfield{Statement}{For a real periodic Hill operator, a double periodic or antiperiodic eigenvalue has two independent periodic or antiperiodic solutions, so its monodromy is \(+I\) or \(-I\). Thus both off-diagonal entries vanish in every basis; a simple edge generally has a nontrivial Jordan block. This uses the differential equation and self-adjoint boundary problem, not the characteristic polynomial of an arbitrary matrix alone.}
\GSfield{Fate}{Classical closed-gap criterion; \citet[Chapter 2, \S1]{dubrovinmatveevnovikov1976nonlinear}; earlier 1973-326.12.}
\GSfield{Present status}{As of November 2026, this is an established result in the stated scope; the qualification in the statement is essential.}
\GSfield{Assessment}{literal statement: correct with missing hypotheses; publication match: exact; verification: read the result; development: none; confidence: high}
\end{GSitem}
\begin{GSitem}{1974-343.6}{A}{Removing double branch points from the spectral curve}
\GSfield{Kind; source}{construction; \GSpage{P000566-P000567, square root}}
\GSfield{The notes}{There are \(2N+1\) simple finite roots \(E_j\) of \(\tau^2-1\), the other roots are double; a radical involving \((\tau^2-1)/\prod(E-E_j)\) is called single-valued.}
\GSfield{Statement}{If the real periodic spectrum has exactly \(N\) open bounded gaps and distinct endpoints \(E_1<\cdots<E_{2N+1}\), then \((\tau^2(E)-1)/\prod_j(E-E_j)\) is entire with only even-order zeros, hence has an entire square root. The remaining two-sheeted compact curve is \(w^2=\prod_j(E-E_j)\), of genus \(N\).}
\GSfield{Reconstruction}{The quotient under one square root is the mathematical reading, not a claim that the scan is unambiguous.}
\GSfield{Fate}{The finite-gap reduction is in \citet[Chapter 2, \S\S1--2]{dubrovinmatveevnovikov1976nonlinear}. The radical must include the quotient; the handwritten overbar alone does not decide its full scope. The spectral-curve and divisor statements are explicitly announced in \citet[\S3]{dubrovinnovikov1974periodicity}, received 3 June 1974 and published later in 1974; this is a firmer contemporary match than the later survey alone.}
\GSfield{Present status}{As of November 2026, this is an established result in the stated scope; the qualification in the statement is essential.}
\GSfield{Assessment}{literal statement: correct with missing hypotheses; publication match: exact; verification: read the result; development: annotated; confidence: high}
\end{GSitem}
\begin{GSitem}{1974-343.7}{E}{The Riccati logarithmic derivative}
\GSfield{Kind; source}{definition; \GSpage{P000567, properties1--3}}
\GSfield{The notes}{\(\chi=-i\Psi^{\prime}/\Psi\); it is independent of the normalizing point \(x_0\), periodic in \(x\), and satisfies \(-i\chi^{\prime}+\chi^2+u-E=0\).}
\GSfield{Statement}{For a chosen nonzero Bloch solution, away from its zeros, the logarithmic derivative is independent of scalar normalization, is \(T\)-periodic, and obeys that Riccati equation. Branch choice in spectral parameter remains; independence of \(x_0\) means that the same Bloch line is selected.}
\GSfield{Fate}{Classical Riccati substitution, written in precisely this convention in \citet[Chapter 2, \S1, Eqs.~(2.1.10)--(2.1.12)]{dubrovinmatveevnovikov1976nonlinear}.}
\GSfield{Present status}{As of November 2026, this is an established result in the stated scope; the qualification in the statement is essential.}
\GSfield{Assessment}{literal statement: correct with missing hypotheses; publication match: exact; verification: read the result; development: none; confidence: high}
\end{GSitem}
\begin{GSitem}{1974-343.8}{E}{Amplitude and phase of a Bloch solution}
\GSfield{Kind; source}{theorem; \GSpage{P000567, properties4,7,8}}
\GSfield{The notes}{In an allowed band the notes write
\[\begin{aligned}
\chi_I&=\chi_R^{\prime}/(2\chi_R),\\
\Psi&=\sqrt{\chi_R(x_0,E)/\chi_R(x,E)}
       \exp\Bigl(i\int_{x_0}^x\chi_R\Bigr),\\
|\Psi|^2&=\chi_R(x_0,E)/\chi_R(x,E).
\end{aligned}\]}
\GSfield{Statement}{For real \(u,E\) with \(E\) inside an open spectral band, choose a nonreal Bloch line and its conjugate, with nonzero Wronskian. Then \(a=\operatorname{Re}\chi\) never vanishes, \(\operatorname{Im}\chi=a^{\prime}/(2a)\), and the displayed normalized amplitude-phase formula holds with a continuous square root equal to 1 at \(x=x_0\). It need not hold at a band endpoint without a limiting interpretation.}
\GSfield{Fate}{\citet[Chapter 2, \S1, Eqs.~(2.1.12)--(2.1.14)]{dubrovinmatveevnovikov1976nonlinear} gives these formulas.}
\GSfield{Present status}{As of November 2026, this is an established result in the stated scope; the qualification in the statement is essential.}
\GSfield{Assessment}{literal statement: correct with missing hypotheses; publication match: exact; verification: read the result; development: none; confidence: high}
\end{GSitem}
\begin{GSitem}{1974-343.9}{E}{Formal high-energy densities and parity}
\GSfield{Kind; source}{construction; \GSpage{P000567, properties5--6}}
\GSfield{The notes}{For \(E=k^2\), \(\chi\sim k+\sum_{n\ge1}\chi_n/(2k)^n\); the coefficients are differential polynomials in \(u\). The real part contains only odd inverse powers.}
\GSfield{Statement}{For smooth real \(u\), the formal Riccati solution with leading term \(k\) is obtained recursively. In this convention \(\chi_1=-u\), \(\chi_2=-iu^{\prime}\), \(\chi_3=u^{\prime\prime}-u^2\). Odd coefficients are real and even coefficients imaginary; even terms are total derivatives and integrate to zero over a period. A formal expansion is not a claim of convergence.}
\GSfield{Reconstruction}{The first coefficients and formal-versus-convergent distinction are supplied editorially. Signs follow the notebook\textquotesingle s chi=-i psi prime/psi convention.}
\GSfield{Fate}{\citet[Chapter 1, \S1 and Chapter 2, Eq.~(2.1.18)]{dubrovinmatveevnovikov1976nonlinear}; compare 1973-326.19.}
\GSfield{Present status}{As of November 2026, this is an established result in the stated scope; the qualification in the statement is essential.}
\GSfield{Assessment}{literal statement: correct with missing hypotheses; publication match: exact; verification: read the result; development: none; confidence: high}
\end{GSitem}
\begin{GSitem}{1974-343.10}{E}{Quasimomentum as a generating function}
\GSfield{Kind; source}{construction; \GSpage{P000567, property9}}
\GSfield{The notes}{\(\lambda=e^{ip(E)}\), \(p(E)=\int_{x_0}^{x_0+T}\chi_R(x,E)\,dx+2\pi n\); the Euler equations for its expansion are called KdV equations.}
\GSfield{Statement}{For a selected Bloch branch, \(p\) is defined modulo \(2\pi\). Its high-energy expansion yields integrals of local differential polynomials invariant under isospectral KdV flows. Hamiltonian evolution generated by these integrals gives the KdV hierarchy; critical points of finite linear combinations give stationary equations. The Euler--Lagrange equation alone is not the evolutionary KdV equation.}
\GSfield{Reconstruction}{The distinction between variational critical-point equations and Hamiltonian time evolution repairs an abbreviated phrase.}
\GSfield{Fate}{\citet[Chapter 2, Eq.~(2.1.15) and \S2]{dubrovinmatveevnovikov1976nonlinear}. The stationary hierarchy mechanism is already in \citet[\S2]{novikov1974periodic}. The cited Novikov article appeared inJuly 1974, before the23 September catalogue date; this occurrence is therefore expository.}
\GSfield{Present status}{As of November 2026, this is an established result in the stated scope; the qualification in the statement is essential.}
\GSfield{Assessment}{literal statement: correct with missing hypotheses; publication match: exact; verification: read the result; development: none; confidence: high}
\end{GSitem}
\begin{GSitem}{1974-343.11}{A}{Variational derivative of quasimomentum}
\GSfield{Kind; source}{theorem; \GSpage{P000567, property10}}
\GSfield{The notes}{\(\delta p(E)/\delta u(x)=-1/(2\chi_R(x,E))\).}
\GSfield{Statement}{For real smooth periodic \(u\), a spectral parameter in a band interior, and a fixed local branch of quasimomentum, its first variation under periodic \(v\) is \(D p(u)[v]=-\tfrac12\int_0^T v(x)/\chi_R(x,E)\,dx\). The local choice removes the additive \(2\pi\) ambiguity; analytic continuation requires consistent sheet choice.}
\GSfield{Fate}{\citet[Chapter 2, Eq.~(2.1.17)]{dubrovinmatveevnovikov1976nonlinear} gives the negative-sign formula. The identity is already announced in \citet[\S3]{dubrovinnovikov1974periodicity}, received 3 June 1974, but the opened English version prints a positive sign. With the present convention \(p(E;u_0)=T\sqrt{E-u_0}\) for constant potential, direct differentiation gives \(\delta p/\delta u(x)=-1/(2\sqrt{E-u_0})\), so the negative sign is forced. The notebook agrees with this check.}
\GSfield{Present status}{As of November 2026, this is an established result in the stated scope; the qualification in the statement is essential.}
\GSfield{Assessment}{literal statement: correct with missing hypotheses; publication match: exact; verification: read the result; development: annotated; confidence: high}
\end{GSitem}
\begin{GSitem}{1974-343.12}{E}{Expressing the logarithmic derivative through monodromy}
\GSfield{Kind; source}{construction; \GSpage{P000567, property11}}
\GSfield{The notes}{\(\chi=\sqrt{1-\tau^2}/\alpha_{21}+\tfrac i2(\alpha_{11}-\alpha_{22})/\alpha_{21}\).}
\GSfield{Statement}{Use the coefficient-basis convention \(T=\begin{pmatrix}c(T)&c^{\prime}(T)\\s(T)&s^{\prime}(T)\end{pmatrix}\), so \(\alpha_{21}=s(T)\). Solving the Bloch eigenvector equation gives the displayed expression, with the sign of the square root selecting the multiplier. Where \(s(T)=0\), interpret it meromorphically or use the other chart.}
\GSfield{Reconstruction}{The manuscript leaves its matrix-entry convention implicit. The transposed coefficient convention is selected so the written alpha21 formula is valid.}
\GSfield{Fate}{The two basis conventions are distinguished in \citet[Chapter 2, Eq.~(2.1.19)]{dubrovinmatveevnovikov1976nonlinear}. Changing to the usual fundamental-state matrix interchanges the off-diagonal indices.}
\GSfield{Present status}{As of November 2026, this is an established result in the stated scope; the qualification in the statement is essential.}
\GSfield{Assessment}{literal statement: correct with missing hypotheses; publication match: exact; verification: read the result; development: none; confidence: high}
\end{GSitem}
\begin{GSitem}{1974-343.13}{A}{The finite-gap polynomial divisor representation}
\GSfield{Kind; source}{construction; \GSpage{P000567-P000568, property12}}
\GSfield{The notes}{\(R(E)=\prod_{j=1}^{2N+1}(E-E_j)\), \(P(E,x)=\prod_{j=1}^{N}(E-\gamma_j(x))\), \(\chi=\sqrt R/P-\tfrac i2P_x/P\).}
\GSfield{Statement}{For a nonsingular real finite-gap Schr\"odinger potential with \(N\) open gaps, choose the branch of \(\sqrt R\) behaving as \(E^{N+1/2}\), and the monic degree-\(N\) divisor polynomial \(P\). Then the logarithmic derivative has the displayed meromorphic expression on \(w^2=R(E)\). Its zeros and poles must satisfy the Schr\"odinger compatibility identity, so an arbitrary pair of polynomials does not suffice.}
\GSfield{Fate}{\citet[Chapter 2, \S\S2--3]{dubrovinmatveevnovikov1976nonlinear}; the relation between stationary operators and the spectral polynomial appears in \citet[\S2]{novikov1974periodic}. The spectral-curve and divisor statements are explicitly announced in \citet[\S3]{dubrovinnovikov1974periodicity}, received 3 June 1974 and published later in 1974; this is a firmer contemporary match than the later survey alone.}
\GSfield{Present status}{As of November 2026, this is an established result in the stated scope; the qualification in the statement is essential.}
\GSfield{Assessment}{literal statement: correct with missing hypotheses; publication match: exact; verification: read the result; development: annotated; confidence: high}
\end{GSitem}
\begin{GSitem}{1974-343.14}{F}{Entire functions of order one in place of finite-gap polynomials}
\GSfield{Kind; source}{problem; \GSpage{P000568, Gelfand problem}}
\GSfield{The notes}{\ru{Пусть R и P --- целые ф-и 1-го пор. роста. Тогда им соотв. потенциалы u(x). Какие?} The next line says that periodicity is not required.}
\GSfield{Statement}{The strongest literal content is a request to identify potentials associated with an entire-function extension of the preceding \(R,P\) expression. The page supplies no normalization, admissibility, reality, zero-divisor conditions or precise meaning of the growth variable. Consequently there is no uniquely determined global inverse problem to classify as solved or open.}
\GSfield{Fate}{Finite-gap inverse data are developed in \citet[Chapter 2, \S3]{dubrovinmatveevnovikov1976nonlinear}. The verification appendix proves an explicit local compatibility criterion and a normalization obstruction for the notebook ansatz; this is an editorial audit, not a historical solution of an unspecified global question. \citet[Theorem 5.1, pp.~535--536]{marchenko1975characterization} supplies a precise real-periodic inverse characterization; \citet[Hypotheses3.1--3.2 and Eq.~(3.14)]{gesztesy1999infinite}, first circulated7 July 1999, treats an admissible infinite-genus setting. Both impose substantial extra data and hypotheses. They are thematic or partial correspondences, not exact solutions of this unnormalized entire-function question.}
\GSfield{Present status}{As of November 2026, the intended exact entire-function classification remains unidentifiable from this paragraph. The local compatibility issue can be checked, but no residual research conjecture is manufactured.}
\GSfield{Assessment}{literal statement: not determinable; publication match: thematic only; verification: unresolved lead; development: annotated; confidence: medium}
\end{GSitem}
\begin{GSitem}{1974-343.15}{A}{Bloch function on the hyperelliptic curve}
\GSfield{Kind; source}{theorem; \GSpage{P000568, Dubrovin theorem}; B. A. Dubrovin}
\GSfield{The notes}{Dubrovin states that \(\Psi(x,x_0,E)\) is single-valued on \(\Gamma_N\), meromorphic away from infinity, with \(N\) zeros above \(\gamma_i(x)\) and \(N\) poles above \(\gamma_i(x_0)\).}
\GSfield{Statement}{For nonsingular finite-gap spectral data with a nonspecial divisor of degree \(N\), the normalized Baker--Akhiezer function is single-valued and meromorphic on the curve minus its point at infinity; its polar divisor is bounded by the initial divisor and its zero divisor by the moving divisor, with \(N\) points counting multiplicity before cancellation. At \(x=x_0\) the function is 1, so coincident zero and pole divisors cancel. At infinity it has the prescribed exponential singularity.}
\GSfield{Reconstruction}{Multiplicity, divisor cancellation, nonspeciality and the exponential behavior at infinity are made explicit; the literal count fails at the normalization point.}
\GSfield{Fate}{\citet[Chapter 2, \S1, p.~77 and \S3]{dubrovinmatveevnovikov1976nonlinear} supplies the divisor description and theta-function inverse construction. The notebook proves circulation by 23 September 1974; the relevant announcements appeared 1974--1975. The spectral-curve and divisor statements are explicitly announced in \citet[\S3]{dubrovinnovikov1974periodicity}, received 3 June 1974 and published later in 1974; this is a firmer contemporary match than the later survey alone.}
\GSfield{Present status}{As of November 2026, this is an established result in the stated scope; the qualification in the statement is essential.}
\GSfield{Assessment}{literal statement: correct with missing hypotheses; publication match: exact; verification: read the result; development: none; confidence: high}
\end{GSitem}

\GStalk{344}{1974-09-30}{Graeme Segal}{P000568-P000570}{Configuration spaces, scanning, and positive and negative particles}
P000569--P000571 reports Dusa McDuff\textquotesingle s work, beginning with Segal\textquotesingle s published approximation of iterated loop spaces by configurations. It distinguishes stable homology from an actual homotopy equivalence, introduces cancellation of opposite charges, and sketches a local-to-global proof using particles that disappear at a boundary. The notation \(\mathrm{Map}\) must be read with its specified basepoints; omitting them changes the theorem. The final gluing paragraph continues beyond the catalogue\textquotesingle s last page.

\begin{GSitem}{1974-344.1}{E}{Unordered configurations and the electric-field map}
\GSfield{Kind; source}{construction; \GSpage{P000569, definition and electric-field construction}}
\GSfield{The notes}{\(C_k(\mathbb R^n)\) consists of unordered subsets of cardinality \(k\). Unit charges define \(f_S\), with each particle sent to infinity and infinity sent to zero, and the resulting sphere map has degree \(k\).}
\GSfield{Statement}{For \(n>1\), choose a decaying radial field with a pole of fixed local degree at each point of a configuration. Its one-point compactification gives a continuous map \(C_k(\mathbb R^n)\to\Omega^n_k S^n\), where the source sphere is based at infinity and the target at zero. Orient the target so that each positive pole contributes degree \(+1\); then the total degree is \(k\).}
\GSfield{Reconstruction}{The based mapping space, n>1 for the electrostatic description, and orientation convention are made explicit. Scanning gives the corresponding construction also in dimension 1.}
\GSfield{Fate}{The construction is already explicit in \citet[\S1, p.~213]{segal1973configuration}, published in 1973. Its description agrees unusually closely with the notebook.}
\GSfield{Present status}{As of November 2026, this is an established result in the stated scope; the qualification in the statement is essential.}
\GSfield{Assessment}{literal statement: correct with missing hypotheses; publication match: exact; verification: read the result; development: none; confidence: high}
\end{GSitem}
\begin{GSitem}{1974-344.2}{E}{Stable homology approximation and group completion}
\GSfield{Kind; source}{theorem; \GSpage{P000569, homology display; P000570, stable limit}}
\GSfield{The notes}{\(H_i(C_k(\mathbb R^n))\cong H_i(\mathrm{Map}_k(S^n,S^n))\) for \(i\ll k\). The right-hand side is homotopically independent of \(k\), while configurations stabilize by adding particles.}
\GSfield{Statement}{For each fixed \(i\), scanning induces an integral-homology isomorphism \(H_i(C_k(\mathbb R^n);\mathbb Z)\to H_i(\Omega^n_kS^n;\mathbb Z)\) for all sufficiently large \(k\). Adding a particle at infinity defines stabilization after choosing an end, and translation between components of the grouplike loop space identifies their homotopy types. This is a stable homology assertion, not a finite-\(k\) homotopy equivalence.}
\GSfield{Fate}{\citet[\S1 and main theorem]{segal1973configuration} proves the group-completion formulation and states this consequence. The notebook gives no numerical stable range, so none is attributed to the 1974 lecture.}
\GSfield{Present status}{As of November 2026, this is an established result in the stated scope; the qualification in the statement is essential.}
\GSfield{Assessment}{literal statement: correct with missing hypotheses; publication match: exact; verification: read the result; development: none; confidence: high}
\end{GSitem}
\begin{GSitem}{1974-344.3}{E}{Configuration spaces as classifying spaces}
\GSfield{Kind; source}{remark; \GSpage{P000570, Fuchs intervention}; D. B. Fuchs}
\GSfield{The notes}{Fuchs says that \(C_k(\mathbb R^n)\) gives different classifying spaces, without naming the groups or a limiting dimension.}
\GSfield{Statement}{Two standard instances make the remark precise: \(C_k(\mathbb R^2)\) is a classifying space for the braid group, and \(C_k(\mathbb R^\infty)\) is a classifying space for \(\Sigma_k\). For a fixed finite \(n\ge3\), the ordered configuration space is not generally contractible, so the unordered space cannot simply be identified with \(B\Sigma_k\).}
\GSfield{Reconstruction}{The two examples are editorial explanations of a general remark, not transcribed assertions.}
\GSfield{Fate}{The braid and stable symmetric-group interpretations form the background of \citet[\S1]{segal1973configuration}; the notebook itself specifies neither example. The statement is therefore an identified but only partially matched expository remark.}
\GSfield{Present status}{As of November 2026, this is an established result in the stated scope; the qualification in the statement is essential.}
\GSfield{Assessment}{literal statement: correct with missing hypotheses; publication match: partial; verification: indirect; development: none; confidence: medium}
\end{GSitem}
\begin{GSitem}{1974-344.4}{B}{The topology of signed configurations}
\GSfield{Kind; source}{definition; \GSpage{P000570, opposite charges}}
\GSfield{The notes}{Finite subsets carry charges of either sign; opposite charges disappear when they meet. Components are indexed by the total charge \(k\in\mathbb Z\).}
\GSfield{Statement}{Let \(C(M)=\coprod_{r\ge0}C_r(M)\). Form \(C^{\pm}(M)\) from \(C(M)\times C(M)\) by cancelling a point common to the positive and negative configurations, with the quotient topology. Particles of the same sign remain distinct. For connected Euclidean space, the invariant \(\#S_+-\#S_-\) labels the path components.}
\GSfield{Fate}{This is McDuff\textquotesingle s positive-and-negative-particle functor, published in 1975: \citet{mcduff1975configuration}. The precise quotient is also given in \citet[\S1; Example 11.5]{kallel1998particles}. The notebook records its seminar circulation by 30 September 1974.}
\GSfield{Present status}{As of November 2026, this is an established result in the stated scope; the qualification in the statement is essential.}
\GSfield{Assessment}{literal statement: correct; publication match: exact; verification: read the result; development: none; confidence: high}
\end{GSitem}
\begin{GSitem}{1974-344.5}{B}{The correct loop-space target for cancelling particles}
\GSfield{Kind; source}{conjecture; \GSpage{P000570, crossed comparison and corrected target}}
\GSfield{The notes}{A tentative \(C_k^{\pm}(\mathbb R^n)\simeq\mathrm{Map}_k(S^n,S^n)\) is labelled \ru{гипотеза}; it is immediately followed by a target \(Y_n\) of dimension \(2n\).}
\GSfield{Statement}{For the cancellation topology just defined, the identified target is \(Y_n=(S^n\times S^n)/\Delta(S^n)\), with the diagonal collapsed to the basepoint. Scanning gives \(C^{\pm}(\mathbb R^n)\simeq\Omega^nY_n\), componentwise with charge corresponding to degree. Replacing \(Y_n\) by \(S^n\) is not valid for this space of particles.}
\GSfield{Reconstruction}{The notebook leaves \(Y_n\) unnamed. Its quotient model is identified from the later text, not read into the handwriting.}
\GSfield{Fate}{\citet{mcduff1975configuration} is the publication announced in the talk; \citet[Example 11.5]{kallel1998particles} explicitly gives the quotient target and the equivalence, attributing it to McDuff. Thus the tentative sphere target is corrected by the very theory being reported, rather than retained as an open conjecture.}
\GSfield{Present status}{As of November 2026, this is an established result in the stated scope; the qualification in the statement is essential.}
\GSfield{Assessment}{literal statement: false as written; publication match: exact; verification: read the result; development: annotated; confidence: high; problem \hyperref[prob:GP-1974-11]{GP-1974-11}}
\end{GSitem}
\begin{GSitem}{1974-344.6}{B}{Scanning a manifold into a section space}
\GSfield{Kind; source}{construction; \GSpage{P000570, tangent-bundle compactification}}
\GSfield{The notes}{\(C_k(M)\to\Gamma_k^{\mathrm{cpt}}(T_M^{\mathrm{comp}})\), with infinity added to each tangent fibre. A finite-range electric field is said to make the construction rigorous.}
\GSfield{Statement}{For a smooth \(n\)-manifold, choose a Riemannian metric and sufficiently small scanning balls. A configuration determines a compactly supported section of the fibrewise one-point compactification \(\dot TM\); outside the scanning support it is the infinity section. The component of the section space is the one containing scanned \(k\)-particle configurations. Local charts are glued through the tangent-bundle transition maps.}
\GSfield{Fate}{This is the scanning construction of \citet{mcduff1975configuration}; the bundle and relative-section formulation are made explicit in \citet[Theorem 1.1]{kallel1998particles}. Its status in the notebook is an announcement of McDuff\textquotesingle s result, published the following year.}
\GSfield{Present status}{As of November 2026, this is an established result in the stated scope; the qualification in the statement is essential.}
\GSfield{Assessment}{literal statement: correct with missing hypotheses; publication match: exact; verification: read the result; development: none; confidence: high}
\end{GSitem}
\begin{GSitem}{1974-344.7}{B}{Stable homology of configurations on manifolds}
\GSfield{Kind; source}{theorem; \GSpage{P000570, stable manifold theorem}}
\GSfield{The notes}{The manifold scanning map is said to induce a homology isomorphism for \(i\ll k\), without specifying compactness, boundary conditions or a stable range.}
\GSfield{Statement}{For a connected smooth open manifold of dimension at least2 that is the interior of a compact manifold with nonempty boundary, choose an end for stabilization. In each fixed degree, scanning from \(C_k(M)\) to the corresponding compactly supported section component induces an integral-homology isomorphism for all sufficiently large \(k\). The stabilized space and section space are homology equivalent. Closed-manifold and relative variants require their own formulation.}
\GSfield{Reconstruction}{The stated finite-type open-manifold scope is a safe precise reconstruction. The notebook says only manifold; no assertion for arbitrary spaces or boundary conventions is inferred.}
\GSfield{Fate}{The stable scanning theorem originates in \citet{mcduff1975configuration}; \citet[Theorem 1.1]{kallel1998particles} proves a broader particle-space formulation. The restricted open-manifold statement is sufficient to make the unqualified notebook sentence rigorous.}
\GSfield{Present status}{As of November 2026, this is an established result in the stated scope; the qualification in the statement is essential.}
\GSfield{Assessment}{literal statement: correct with missing hypotheses; publication match: partial; verification: read the result; development: none; confidence: high}
\end{GSitem}
\begin{GSitem}{1974-344.8}{B}{Particles disappearing at a boundary}
\GSfield{Kind; source}{construction; \GSpage{P000570, failed restriction and boundary quotient}}
\GSfield{The notes}{The map \(C(M)\to C(U)\times C(V)\) by restriction is marked discontinuous. For manifolds with boundary, \(\overline C(U)\) identifies configurations agreeing on the interior, so boundary particles disappear.}
\GSfield{Statement}{Cardinality changes when a particle crosses the edge of an open piece, preventing continuity into the disjoint-union configuration topology. The relative space \(C(U,\partial U)\), formed by deleting particles on \(\partial U\), gives the required continuity at the exit locus. Collar choices provide compatible restriction maps for the gluing argument.}
\GSfield{Fate}{The relative configuration construction is part of McDuff\textquotesingle s scanning proof; see \citet{mcduff1975configuration} and the explicit relative particle construction in \citet[\S6]{kallel1998particles}.}
\GSfield{Present status}{As of November 2026, this is an established result in the stated scope; the qualification in the statement is essential.}
\GSfield{Assessment}{literal statement: correct with missing hypotheses; publication match: exact; verification: read the result; development: none; confidence: high}
\end{GSitem}
\begin{GSitem}{1974-344.9}{B}{Relative configurations and the local disk calculation}
\GSfield{Kind; source}{theorem; \GSpage{P000570, boundary theorem and disk proof}}
\GSfield{The notes}{For \(U\) with nonempty boundary, \(\overline C(U)\simeq\Gamma(\dot TU)\). For a disk, radially pushing points to the boundary leaves at most one particle, and both sides become \(S^n\).}
\GSfield{Statement}{In the relative scanning model for a compact connected smooth manifold with nonempty boundary, particles are allowed to disappear on all of the boundary and the target consists of sections with unrestricted boundary values. The local model is \(C(D^n,\partial D^n)\simeq D^n/\partial D^n\cong S^n\), while \(\Gamma(\dot TD^n)\simeq\mathrm{Map}(D^n,S^n)\simeq S^n\). For disconnected manifolds each component must meet the exit boundary.}
\GSfield{Reconstruction}{Connectedness, the exit boundary, and unrestricted section values at that boundary are supplied to distinguish this statement from compactly supported scanning on the interior.}
\GSfield{Fate}{The disk calculation follows from \citet[Proposition9.1 and Lemma 11.1]{kallel1998particles}. The notebook attributes the general relative theorem to McDuff; \citet{mcduff1975configuration} was verified bibliographically, but its full text was inaccessible in this run. The exact all-boundary generalization is therefore reported with a partial match, rather than claimed checked from that original paper. Boundary values fixed to infinity would produce a different target.}
\GSfield{Present status}{As of November 2026, this is an established result in the stated scope; the qualification in the statement is essential.}
\GSfield{Assessment}{literal statement: correct with missing hypotheses; publication match: partial; verification: indirect; development: none; confidence: high}
\end{GSitem}
\begin{GSitem}{1974-344.10}{B}{Local-to-global proof by homotopy pullback}
\GSfield{Kind; source}{method; \GSpage{P000571, two covering squares and homotopy lifting}}
\GSfield{The notes}{The two restriction squares for \(\overline C\) and sections are compared. Homotopies lift in the restriction maps; equivalences on \(U,V,U\cap V\) then give one on \(U\cup V\).}
\GSfield{Statement}{With collared manifold pieces and the appropriate relative configuration spaces, the restriction diagrams must model homotopy pullbacks (or the corresponding quasifibration sequences). An objectwise equivalence on the other three vertices then implies one on the union. A bare square of ordinary configuration spaces does not have this property; the exit-boundary construction and homotopy-lifting hypotheses do the work.}
\GSfield{Reconstruction}{Homotopy pullback is the modern precise interpretation of the drawn squares; no strict sheaf property is asserted.}
\GSfield{Fate}{The proof mechanism belongs to the scanning and quasifibration argument of \citet{mcduff1975configuration}; \citet[\S6 and Proposition9.1]{kallel1998particles} gives explicit relative quasifibrations and the handle induction. The notebook sketches the method but does not specify enough collars or maps for a complete proof. The exact original 1975 proof was not accessible, so the publication match is partial.}
\GSfield{Present status}{As of November 2026, this is an established result in the stated scope; the qualification in the statement is essential.}
\GSfield{Assessment}{literal statement: correct with missing hypotheses; publication match: partial; verification: read the result; development: none; confidence: high}
\end{GSitem}

\GStalk{345}{1974-09-30}{A. M. Polyakov}{P000568-P000573}{Lattice gauge fields, Wilson loops, and thermodynamic phases}
P000569 names Polyakov and a difference analogue of Yang--Mills theory under30 September. The substantive formulas occur onP000573--P000574, after the7 October heading and Arnold\textquotesingle s intervention. The catalogue date is preserved, while this manuscript displacement is recorded. Compact-group link variables and plaquette products lead to Gibbs integrals and loop observables. Sinai then specializes to an angular \(U(1)\) model. Its Bessel partition function must not be transferred unchanged to \(SU(2)\), and the final area/perimeter remark does not specify a complete phase-transition theorem.

\begin{GSitem}{1974-345.1}{E}{Compact-group link variables and plaquette holonomy}
\GSfield{Kind; source}{definition; \GSpage{P000573, Polyakov opening and plaquette sketch}}
\GSfield{The notes}{Oriented lattice edges carry \(B_{\tau,\alpha}\in SU(2)\), or another compact group. The oriented square has \(F_{\tau,\alpha\beta}=B_{\tau,\alpha}B_{\tau+a_\alpha,\beta}B_{\tau+a_\beta,\alpha}^{-1}B_{\tau,\beta}^{-1}\).}
\GSfield{Statement}{On a finite square lattice, assign \(U_{xy}\in G\) to each oriented link, with \(U_{yx}=U_{xy}^{-1}\). The ordered product around a plaquette is its holonomy. Under vertex gauge transformations \(U_{xy}\mapsto h_xU_{xy}h_y^{-1}\), that product is conjugated at its basepoint, so a character or a real class function of it is gauge invariant.}
\GSfield{Fate}{The construction is in \citet[\S III, especially the non-Abelian extension]{wilson1974confinement}. Wilson\textquotesingle s manuscript was received 12 June 1974 and the CLNS-262 preprint is catalogued February 1974; publication was15 October 1974. Thus it was already circulating before either possible notebook meeting date.}
\GSfield{Present status}{As of November 2026, this is an established result in the stated scope; the qualification in the statement is essential.}
\GSfield{Assessment}{literal statement: correct; publication match: exact; verification: read the result; development: none; confidence: high}
\end{GSitem}
\begin{GSitem}{1974-345.2}{E}{The finite-volume lattice Gibbs measure}
\GSfield{Kind; source}{definition; \GSpage{P000573, energy and partition integral; P000574, Gibbs density}}
\GSfield{The notes}{\(E[B]=\sum\chi(F)\), \(Z=\int\prod dB\,e^{-E/T}\), and \(p(B)=Z^{-1}e^{-\beta E(B)}\), with \(\chi\) described as trace.}
\GSfield{Statement}{For compact \(G\), normalized Haar measure on each unoriented link and a real bounded plaquette class function \(V\), define \(E(U)=\sum_pV(U_p)\), \(Z_\Lambda(\beta)=\int e^{-\beta E(U)}\prod_edU_e\), and the probability measure with this normalized density. For \(SU(2)\) the fundamental trace is real; for general compact groups one must use a real character combination such as \(\operatorname{Re}\chi\). The sign and normalization of \(V\) must be fixed before comparing coupling constants.}
\GSfield{Fate}{\citet[\S III.B]{wilson1974confinement} gives the finite-volume Euclidean lattice integral. The notebook\textquotesingle s positive-trace energy is a legitimate sign convention but differs from the common ferromagnetic action \(-\operatorname{Re}\operatorname{tr}U_p\).}
\GSfield{Present status}{As of November 2026, this is an established result in the stated scope; the qualification in the statement is essential.}
\GSfield{Assessment}{literal statement: correct with missing hypotheses; publication match: exact; verification: read the result; development: none; confidence: high}
\end{GSitem}
\begin{GSitem}{1974-345.3}{E}{The Wilson loop as a gauge-invariant correlation}
\GSfield{Kind; source}{definition; \GSpage{P000573, correlation display; P000574, contour average}}
\GSfield{The notes}{\(\Pi(C)=Z^{-1}\int\chi(\prod_C B)e^{-\beta E(B)}\prod dB\); \(C\) is a closed contour.}
\GSfield{Statement}{For an oriented closed lattice path \(C\), use the path-ordered product of link variables, replacing a reversed edge by its inverse. Its character is independent of the chosen starting vertex, and its Gibbs expectation is gauge invariant. A normalized trace divides by the representation dimension; the notebook uses the unnormalized trace.}
\GSfield{Fate}{\citet[Eq.~(3.13) and \S II]{wilson1974confinement} gives this loop observable and its physical interpretation. The recorded expression is an explicit lattice Wilson loop, not an arbitrary two-point correlation function.}
\GSfield{Present status}{As of November 2026, this is an established result in the stated scope; the qualification in the statement is essential.}
\GSfield{Assessment}{literal statement: correct with missing hypotheses; publication match: exact; verification: read the result; development: none; confidence: high}
\end{GSitem}
\begin{GSitem}{1974-345.4}{B}{Exact integration in two dimensions}
\GSfield{Kind; source}{method; \GSpage{P000573, planar exact-solvability statement}}
\GSfield{The notes}{\ru{Плоская зад. решается точно} follows the finite \(N\times N\) lattice definition.}
\GSfield{Statement}{For a simply connected finite planar lattice with open boundary and a central plaquette weight \(w:G\to\mathbb R_{>0}\), a spanning-tree gauge choice reduces the partition integral to independent plaquette integrals. If \(w=\sum_\rho a_\rho\chi_\rho\), the disk amplitude with boundary holonomy \(g\) is \(\sum_\rho d_\rho(a_\rho/d_\rho)^A\chi_\rho(g)\), where \(A\) is the number of plaquettes. Integrating the boundary gives \(Z=a_{\mathrm{triv}}^A\). Periodic boundary identifications impose additional global constraints.}
\GSfield{Reconstruction}{The region, boundary treatment, central weight and character normalization are supplied to make the phrase exactly solvable precise.}
\GSfield{Fate}{The later exact character-expansion formulation is explicitly derived in \citet[Appendix A.1, Eqs.~(A.6)--(A.11)]{hirasawa2020complex}. This verifies the statement in a specified planar setting; it is not evidence that the 2020 paper was the first solution, or that the notebook\textquotesingle s unspecified boundary conditions are immaterial.}
\GSfield{Present status}{As of November 2026, this is an established result in the stated scope; the qualification in the statement is essential.}
\GSfield{Assessment}{literal statement: correct with missing hypotheses; publication match: partial; verification: read the result; development: none; confidence: high}
\end{GSitem}
\begin{GSitem}{1974-345.5}{B}{Coupling group-valued spins to a lattice connection}
\GSfield{Kind; source}{construction; \GSpage{P000573, two spin-energy displays}}
\GSfield{The notes}{The spin energy \(\sum\chi(g_\tau^{-1}g_{\tau+a_\alpha})\) is followed by \(\sum\chi(g_\tau^{-1}g_{\tau+a_\alpha}B_{\tau,\alpha})+E[B]\).}
\GSfield{Statement}{Group-valued spins with a global symmetry can be coupled to link transport to make the interaction locally gauge invariant. For example, with \(g_x\mapsto g_xh_x^{-1}\) and a reverse transporter \(B_{xy}\mapsto h_yB_{xy}h_x^{-1}\), the factor \(\chi(g_x^{-1}g_yB_{xy})\) is invariant. Equivalently one can use forward links and the conventional \(\chi(g_x^{-1}U_{xy}g_y)\), with the corresponding left action on spins. Mixing these conventions destroys invariance.}
\GSfield{Reconstruction}{The transformation laws are an explicit editorial consistency check. The two energy displays are retained as one gauging construction rather than treated as two unrelated models. In particular the reverse-link convention used to rescue this ordering cannot simultaneously be inserted unchanged into the preceding forward-oriented plaquette formula; one must reverse those links too, or use the conventional forward-link interaction.}
\GSfield{Fate}{Gauge-invariant matter coupling is already present in \citet[\S III.A]{wilson1974confinement}; \citet{baliandrouffeitzykson1974gauge} develops the coupled group-spin model, received 3 June 1974 and published 15 November 1974. The notebook does not specify its transport convention, so the displayed ordering has only a partial exact match.}
\GSfield{Present status}{As of November 2026, this is an established result in the stated scope; the qualification in the statement is essential.}
\GSfield{Assessment}{literal statement: correct with missing hypotheses; publication match: partial; verification: read the result; development: none; confidence: high}
\end{GSitem}
\begin{GSitem}{1974-345.6}{E}{Euclidean continuation and transfer matrices}
\GSfield{Kind; source}{remark; \GSpage{P000573, field theory and statistical mechanics}}
\GSfield{The notes}{Field theory is identified with statistical mechanics in Euclidean \(\mathbb R^n\); the parenthesis says \ru{преобр. Фурье по времени}.}
\GSfield{Statement}{The relevant formal passage from Lorentzian quantum dynamics to Euclidean Gibbs weights is continuation to imaginary time. For a lattice theory, a suitable positive transfer operator relates Euclidean correlation functions to a quantum Hamiltonian. An ordinary Fourier transform in real time alone does not make the Minkowski action into a positive Euclidean Gibbs energy.}
\GSfield{Fate}{\citet[\S I and \S III.B]{wilson1974confinement} explicitly uses imaginary time and the transfer-matrix construction. The notebook\textquotesingle s Fourier-transform wording is an imprecise shorthand requiring correction; no general constructive-equivalence theorem is stated.}
\GSfield{Present status}{As of November 2026, this is an established result in the stated scope; the qualification in the statement is essential.}
\GSfield{Assessment}{literal statement: false as written; publication match: exact; verification: read the result; development: annotated; confidence: high}
\end{GSitem}
\begin{GSitem}{1974-345.7}{B}{When the Bessel partition function is valid}
\GSfield{Kind; source}{claim; \GSpage{P000573, Sinai partition function; P000574, angular specialization}; Ya. G. Sinai}
\GSfield{The notes}{Sinai writes \(\log Z(\beta,V)=V\log I_0(\beta)\). The next page has \(B=e^{iA}\) and \(E=\sum\cos(\operatorname{curl}A)\).}
\GSfield{Statement}{For the planar \(U(1)\) model with normalized link-angle measure, open boundary, and \(A\) plaquettes, \(Z=I_0(\beta)^A\), where \(I_m(\beta)=(2\pi)^{-1}\int_{-\pi}^{\pi}e^{\beta\cos\theta}\cos(m\theta)\,d\theta\). The notebook\textquotesingle s opposite action sign gives the same \(I_0\). This is not the \(SU(2)\) formula: for fundamental trace and weight \(e^{\beta\operatorname{tr}U}\), a normalized one-plaquette integral is \(I_1(2\beta)/\beta\). If \(V\) counts sites rather than plaquettes, the written equality is only asymptotic.}
\GSfield{Reconstruction}{U(1), normalized Haar measure, open boundary and plaquette count are required. The SU(2) comparison is editorial and not attributed to Sinai.}
\GSfield{Fate}{The general exact integration is verified in \citet[Appendix A.1 and A.3]{hirasawa2020complex}; the elementary Haar integrations are written out in the verification appendix. These distinguish an Abelian specialization from the preceding \(SU(2)\) example. No specific publication by Sinai carrying this exact formula was identified.}
\GSfield{Present status}{As of November 2026, this is an established result in the stated scope; the qualification in the statement is essential.}
\GSfield{Assessment}{literal statement: correct with missing hypotheses; publication match: partial; verification: read the result; development: annotated; confidence: high}
\end{GSitem}
\begin{GSitem}{1974-345.8}{E}{Boundary independence of the pressure}
\GSfield{Kind; source}{theorem; \GSpage{P000574, thermodynamic limit}; Ya. G. Sinai}
\GSfield{The notes}{\(\lim_{V\to\infty}V^{-1}\log Z(\beta,V)=f(\beta)\) is asserted to be independent of boundary conditions.}
\GSfield{Statement}{For translation-invariant bounded finite-range interactions on compact lattice variables, with normalized product a priori measure, pressure exists along van Hove boxes and is unchanged by boundary conditions whose change affects only a boundary layer of bounded energy per interaction. Indeed boundary choices change \(\log Z\) by at most a constant times the number of boundary interactions, which is \(o(V)\). This is a statement about pressure, not uniqueness of the infinite-volume Gibbs state.}
\GSfield{Reconstruction}{The shape of the limit, bounded finite-range interactions and ordinary boundary perturbations are supplied. Boundary independence alone does not rule out phase coexistence.}
\GSfield{Fate}{The thermodynamic-limit framework precedes the talk; \citet[Theorem 1]{ruelle1967variational} treats lattice pressure for summable interactions in the finite-state setting. The bounded compact-variable formulation here is its standard boundary-estimate extension, not an exact quotation of that theorem. The explicit planar model also has pressure \(\log I_0(\beta)\).}
\GSfield{Present status}{As of November 2026, this is an established result in the stated scope; the qualification in the statement is essential.}
\GSfield{Assessment}{literal statement: correct with missing hypotheses; publication match: partial; verification: read the result; development: none; confidence: high}
\end{GSitem}
\begin{GSitem}{1974-345.9}{E}{Nonanalyticity in the infinite-volume limit}
\GSfield{Kind; source}{definition; \GSpage{P000574, phase-transition sentence}; Ya. G. Sinai}
\GSfield{The notes}{Singularities of \(f(\beta)\) are called phase transitions, accompanied heuristically by a change in typical configurations.}
\GSfield{Statement}{A finite-volume partition integral with bounded energy is entire in complex \(\beta\) and positive for real \(\beta\), so its real logarithm is locally analytic. Thermodynamic singularities can arise only after the infinite-volume limit. Nonanalytic pressure is a thermodynamic phase-transition criterion, while a full classification of Gibbs phases also involves states and order parameters. In the written planar \(U(1)\) model, \(\log I_0(\beta)\) is analytic at every finite real \(\beta\).}
\GSfield{Fate}{The distinction between finite volume and thermodynamic phases is classical; compare \citet[\S I]{wilson1974confinement} and \citet{ruelle1967variational}. The analytic \(I_0\) formula prevents reading this paragraph as a proof of a finite-temperature transition in that particular two-dimensional example.}
\GSfield{Present status}{As of November 2026, this is an established result in the stated scope; the qualification in the statement is essential.}
\GSfield{Assessment}{literal statement: correct with missing hypotheses; publication match: partial; verification: read the result; development: annotated; confidence: high}
\end{GSitem}
\begin{GSitem}{1974-345.10}{F}{Area and perimeter laws require a specified model}
\GSfield{Kind; source}{remark; \GSpage{P000574, final area/length comparison}; Ya. G. Sinai}
\GSfield{The notes}{The last two lines say that in one phase the correlation is proportional to area and in another to contour length.}
\GSfield{Statement}{The intended diagnostic is \(-\log|\langle W(C)\rangle|\) growing proportionally to the area enclosed by a large loop or to its perimeter. The expectation itself decays exponentially. Dimension, gauge group, matter content, action, family of loops and limiting regime must be specified before the phrase becomes a theorem. In the planar pure \(U(1)\) model with positive \(\beta\), \(|\langle W(C)\rangle|=(I_1(\beta)/I_0(\beta))^{A(C)}\) for a simple loop, so it has an area law at every finite \(\beta>0\).}
\GSfield{Fate}{\citet[\S II and \S IV]{wilson1974confinement} supplies the area/perimeter criterion and a strong-coupling analysis, but it does not turn the notebook\textquotesingle s unspecified assertion into a universal two-phase theorem. The exact planar conclusion follows from independent plaquette integration; no modern Yang--Mills mass-gap problem is silently substituted.}
\GSfield{Present status}{As of November 2026, the physical diagnostic is identifiable, but no definite residual conjecture is specified here. The adjacent two-dimensional Abelian example is exactly solvable and has no finite-coupling area/perimeter phase transition.}
\GSfield{Assessment}{literal statement: not determinable; publication match: thematic only; verification: unresolved lead; development: annotated; confidence: medium}
\end{GSitem}

\GStalk{346}{1974-10-07}{A. G. Kushnirenko}{P000571-P000573}{Local Milnor numbers from Newton polyhedra}
The dated pageP000572 defines the multiplicity of an isolated critical point and asks how to recover it from Taylor coefficients. Kushnirenko gives the Newton lower bound and its generic equality in two and three dimensions, then onP000573 states face nondegeneracy through a finite-codimension monomial-ring quotient. The argument invokes Hochster\textquotesingle s Cohen--Macaulay theorem. The paper announcing these results was received on30 September 1974, a week before the recorded lecture, and appeared in 1975.

\begin{GSitem}{1974-346.1}{E}{Milnor number as local gradient multiplicity}
\GSfield{Kind; source}{definition; \GSpage{P000572, opening and margin}}
\GSfield{The notes}{For an analytic germ without constant or linear terms, \(\mu(f)\) is the multiplicity of the zero solution of \(\partial f/\partial z_i=0\), also described as the maximum number of solutions to \(\partial f/\partial z_i=\varepsilon_i\) for small \(\varepsilon\).}
\GSfield{Statement}{Let \(f\in\mathbb C\{z_1,\ldots,z_n\}\) have an isolated critical point at 0. Then \(\mu(f)=\dim_{\mathbb C}\mathbb C\{z\}/(f_{z_1},\ldots,f_{z_n})\). For a sufficiently small isolating neighborhood and a sufficiently small generic regular value \(\varepsilon\), the gradient equation has exactly \(\mu(f)\) simple solutions there. Nongeneric perturbations require multiplicities.}
\GSfield{Reconstruction}{Isolation, the local neighborhood, genericity and the distinction between distinct and multiple roots are made explicit.}
\GSfield{Fate}{This local-algebra invariant is part of the established singularity theory preceding the talk. \citet[\S1]{kushnirenko1975newton} fixes the same definition; its use here is expository, while the Newton formula is the new result.}
\GSfield{Present status}{As of November 2026, this is an established result in the stated scope; the qualification in the statement is essential.}
\GSfield{Assessment}{literal statement: correct with missing hypotheses; publication match: exact; verification: read the result; development: none; confidence: high}
\end{GSitem}
\begin{GSitem}{1974-346.2}{A}{The Newton boundary and its alternating volume number}
\GSfield{Kind; source}{construction; \GSpage{P000572, task and Newton diagrams}}
\GSfield{The notes}{\ru{Задача: вычислить} \(\mu(f)\) from Taylor coefficients. The lower Newton polygon gives \(\nu=2S-a-b+1\); in dimension 3 the formula is \(6V-2(S_{12}+S_{13}+S_{23})+(L_1+L_2+L_3)-1\).}
\GSfield{Statement}{For a convenient germ, meaning its support meets every coordinate axis, take the compact faces of \(\operatorname{conv}(\operatorname{supp}f+\mathbb R_{\ge0}^n)\) and the region \(\Gamma_-\) swept out between them and 0. If \(V_j\) is the sum of ordinary \(j\)-dimensional volumes of intersections of \(\Gamma_-\) with coordinate \(j\)-planes and \(V_0=1\), then \(\nu(f)=\sum_{j=0}^n(-1)^{n-j}j!V_j\). Its cases \(n=2,3\) are exactly the written formulas.}
\GSfield{Reconstruction}{Convenience and the meaning of each coordinate-section volume are supplied. The general sum packages the two formulas actually recorded; it is not attributed as an additional notebook display.}
\GSfield{Fate}{The definition and formulas are in \citet[\S\S2--3]{kushnirenko1975newton}; received 30 September 1974, published 1975. The arbitrary-dimensional expression is also the organizing definition in the full 1976 paper, \citet{kouchnirenko1976polyedres}.}
\GSfield{Present status}{As of November 2026, this is an established result in the stated scope; the qualification in the statement is essential.}
\GSfield{Assessment}{literal statement: correct with missing hypotheses; publication match: exact; verification: read the result; development: none; confidence: high; problem \hyperref[prob:GP-1974-12]{GP-1974-12}}
\end{GSitem}
\begin{GSitem}{1974-346.3}{A}{Kushnirenko lower bound and equality under nondegeneracy}
\GSfield{Kind; source}{theorem; \GSpage{P000572, two theorems}}
\GSfield{The notes}{Both dimension 2 and dimension 3 displays state \(\mu(f)\ge\nu(f)\), with equality under nondegeneracy; the condition holds almost always.}
\GSfield{Statement}{For a convenient holomorphic germ with an isolated critical point, \(\mu(f)\ge\nu(f)\). If for every compact face \(\Delta\) the face polynomial \(f_\Delta\) has no critical point in \((\mathbb C^*)^n\), then equality holds. At a fixed convenient Newton boundary, the exceptional face coefficients lie in a proper algebraic subset, so equality holds generically. Equality for every degenerate germ is not claimed.}
\GSfield{Fate}{Exact match with \citet[TheoremI(i)--(iii)]{kushnirenko1975newton}. The 1975 announcement gives the statement; \citet{kouchnirenko1976polyedres}, published 1976, supplies the extended proof. The recorded talk establishes circulation by 7 October 1974 and the received date establishes the announcement existed a week earlier.}
\GSfield{Present status}{As of November 2026, this is an established result in the stated scope; the qualification in the statement is essential.}
\GSfield{Assessment}{literal statement: correct with missing hypotheses; publication match: exact; verification: read the result; development: none; confidence: high; problem \hyperref[prob:GP-1974-12]{GP-1974-12}}
\end{GSitem}
\begin{GSitem}{1974-346.4}{A}{Nondegeneracy of an edge polynomial}
\GSfield{Kind; source}{claim; \GSpage{P000572, edge discriminant}}
\GSfield{The notes}{The illustration gives \(\beta^2\ne4\alpha\gamma\), and the text calls for the discriminant of each edge to be nonzero, because an edge is essentially a one-variable polynomial.}
\GSfield{Statement}{After removing a monomial and using a primitive lattice coordinate along an edge, its face polynomial becomes a one-variable polynomial. Nonzero roots must be simple. For a quadratic edge polynomial with nonzero endpoint coefficients \(\alpha,\gamma\), this is \(\beta^2-4\alpha\gamma\ne0\). In two variables this test must be applied to every compact edge; in higher dimensions it does not replace the conditions on higher-dimensional faces.}
\GSfield{Reconstruction}{The lattice coordinate, removal of monomial factors and nonzero endpoint coefficients are explicit safeguards against an overbroad discriminant test.}
\GSfield{Fate}{This is the planar face test for the nondegeneracy condition in \citet[\S8]{kushnirenko1975newton}. The precise quadratic discriminant follows by elementary differentiation; it is an explanatory specialization of the new general theorem.}
\GSfield{Present status}{As of November 2026, this is an established result in the stated scope; the qualification in the statement is essential.}
\GSfield{Assessment}{literal statement: correct with missing hypotheses; publication match: exact; verification: read the result; development: none; confidence: high}
\end{GSitem}
\begin{GSitem}{1974-346.5}{A}{Face nondegeneracy as a finite-codimension ideal}
\GSfield{Kind; source}{definition; \GSpage{P000573, first formula}}
\GSfield{The notes}{For the cone \(P(\Delta)\) over a face, \(\dim_{\mathbb C}\mathbb C[P(\Delta)]/(z_i\partial f_\Delta/\partial z_i)<\infty\) is written.}
\GSfield{Statement}{For a compact face \(\Delta\) disjoint from 0, let \(P(\Delta)=\operatorname{cone}(\Delta)\cap\mathbb Z_{\ge0}^n\). In the semigroup ring \(\mathbb C[P(\Delta)]\), require the ideal generated by all logarithmic derivatives \(z_i\partial f_\Delta/\partial z_i\) to have finite vector-space codimension. Requiring this for every face is the algebraic nondegeneracy formulation used in the announcement.}
\GSfield{Fate}{The notation and condition agree exactly with \citet[\S8]{kushnirenko1975newton}. They connect the geometric volume formula to a finite-dimensional quotient; no global polynomial gradient quotient is being taken at this stage.}
\GSfield{Present status}{As of November 2026, this is an established result in the stated scope; the qualification in the statement is essential.}
\GSfield{Assessment}{literal statement: correct with missing hypotheses; publication match: exact; verification: read the result; development: none; confidence: high}
\end{GSitem}
\begin{GSitem}{1974-346.6}{E}{Cohen--Macaulay semigroup rings in the proof}
\GSfield{Kind; source}{method; \GSpage{P000573, proof sentence}}
\GSfield{The notes}{The proof is said to use commutative algebra and Hochster\textquotesingle s assertion that \(\mathbb C[P(\Delta)]\) is Cohen--Macaulay.}
\GSfield{Statement}{For a finitely generated saturated affine semigroup \(P\), its semigroup algebra over a field is Cohen--Macaulay. The lattice points of the rational polyhedral cone over \(\Delta\) satisfy these hypotheses. This property supplies regular-sequence and Hilbert-series control when the face-gradient ideal has finite codimension.}
\GSfield{Reconstruction}{Saturation and finite generation are essential. An arbitrary ring generated by an arbitrary selection of monomials need not be Cohen--Macaulay.}
\GSfield{Fate}{Hochster\textquotesingle s theorem was published in September 1972, before the lecture: \citet[main theorem, p.~318]{hochster1972invariants}. Its use in this Newton-polyhedron argument is reported by the notebook and developed in \citet{kouchnirenko1976polyedres}.}
\GSfield{Present status}{As of November 2026, this is an established result in the stated scope; the qualification in the statement is essential.}
\GSfield{Assessment}{literal statement: correct with missing hypotheses; publication match: exact; verification: read the result; development: none; confidence: high}
\end{GSitem}

\GStalk{347}{1974-10-07}{V. I. Arnold}{P000571-P000576}{A polynomial residue-sum identity}
Arnold\textquotesingle s contribution is the single paragraph near the top ofP000573, between Kushnirenko\textquotesingle s proof remark and Polyakov\textquotesingle s lattice-field discussion. It gives a sum of polynomial values divided by the Jacobian over \(n^2\) intersection points. The written degree bound is \(n-4\), which is a valid weaker bound under the standard hypotheses; it is not silently changed to the sharp classical bound.

\begin{GSitem}{1974-347.1}{E}{Euler--Jacobi vanishing of the total residue}
\GSfield{Kind; source}{theorem; \GSpage{P000573, Arnold paragraph}}
\GSfield{The notes}{For two degree-\(n\) polynomials \(P_n,Q_n\), the notes give
\[\sum_{P_n=Q_n=0}\frac{R_{n-4}(x,y)}
 {\det\bigl(\partial(P_n,Q_n)/\partial(x,y)\bigr)}=0,\]
summing over all \(n^2\) zeros.}
\GSfield{Statement}{Let \(P,Q\in\mathbb C[x,y]\) have degree \(n\ge4\), with leading homogeneous parts having no common nonzero zero, and suppose their \(n^2\) affine intersection points are simple. Then for every polynomial \(R\) of degree at most \(2n-3\), \(\sum_{P(a)=Q(a)=0}R(a)/(P_xQ_y-P_yQ_x)(a)=0\). In particular the notebook\textquotesingle s degree at most \(n-4\) satisfies the theorem.}
\GSfield{Reconstruction}{Transversality, absence of intersections at infinity and \(n\ge4\) are explicit. The stronger \(2n-3\) bound is a literature comparison; \(n-4\) is the legible original reading, not an OCR error.}
\GSfield{Fate}{This is the classical Euler--Jacobi formula, not a newly posed residue problem. \citet[Corollary, p.~237]{khovanskii1978eulerjacobi} states the general degree bound and explicitly treats it as classical, while proving a Newton-polyhedron generalization suggested by Arnold. The 1978 result is a later development of the theme, not evidence that this 1974 formula was new.}
\GSfield{Present status}{As of November 2026, this is an established result in the stated scope; the qualification in the statement is essential.}
\GSfield{Assessment}{literal statement: correct with missing hypotheses; publication match: exact; verification: read the result; development: annotated; confidence: high}
\end{GSitem}

\GStalk{348}{1974-10-14}{Graeme Segal}{P000574-P000580}{Classifying spaces, spherical buildings, and the Solomon--Tits theorem}
The actual Segal text is P000575--P000578, headed 14 October 1974. It constructs nerves, compares covers with spaces, proves the building's wedge-of-spheres homotopy type by a hyperplane induction, and constructs a group classifying space's universal cover. Bundle classification ends mid-construction. The catalogue's wider page envelope includes adjoining talks; P000574, P000579 and P000580 are not assigned to Segal's argument. This exposition continues 1973-308 and 1973-309, with a building connection to 1973-321.5. The base case exposes the missing reduced-homology symbol in the Steinberg-module formula.

\begin{GSitem}{1974-348.1}{E}{The nerve and realization of a category}
\GSfield{Kind; source}{construction; \GSpage{P000575}}
\GSfield{The notes}{Objects give vertices, morphisms edges, composable pairs filled triangles, and composable triples tetrahedra. The graph alone does not retain composition.}
\GSfield{Statement}{For a small category \(\mathcal C\), set \(N_k\mathcal C=\operatorname{Fun}([k],\mathcal C)\), with faces given by omission/composition and degeneracies by identities. Its geometric realization is \(B\mathcal C\).}
\GSfield{Reconstruction}{The face and degeneracy maps make the pictured construction precise. A smallness convention is supplied.}
\GSfield{Fate}{\citet[\S2]{segal1968classifying}, received 14 April 1967 and published 1968. Earlier occurrences: 1971-255.2 and 1973-308.11.}
\GSfield{Present status}{As of November 2026: established under the stated hypotheses; this occurrence supplies no unresolved problem.}
\GSfield{Assessment}{literal statement: correct; publication match: exact; verification: read the result; development: none; confidence: high}
\end{GSitem}
\begin{GSitem}{1974-348.2}{E}{Partially ordered sets as thin categories}
\GSfield{Kind; source}{definition; \GSpage{P000575}}
\GSfield{The notes}{The notes identify partially ordered sets with categories having at most one morphism between two objects.}
\GSfield{Statement}{A poset \(P\) defines a category with one arrow \(x\to y\) exactly when \(x\leq y\). Conversely a thin category determines a preorder; it determines a partial order on its objects if it is skeletal. Its nondegenerate nerve consists of strict finite chains.}
\GSfield{Reconstruction}{Antisymmetry, or equivalently passage to a skeleton, is missing from the stated converse. Two distinct isomorphic objects in a thin category contradict that converse literally.}
\GSfield{Fate}{\citet[\S2]{segal1968classifying}; \citet[\S1]{quillen1973higher}. Both precede the talk.}
\GSfield{Present status}{As of November 2026: established under the stated hypotheses; this occurrence supplies no unresolved problem.}
\GSfield{Assessment}{literal statement: correct with missing hypotheses; publication match: exact; verification: read the result; development: annotated; confidence: high}
\end{GSitem}
\begin{GSitem}{1974-348.3}{E}{A group as a one-object category}
\GSfield{Kind; source}{construction; \GSpage{P000575}}
\GSfield{The notes}{A group gives the opposite extreme: one object, group elements as morphisms, and group multiplication as composition.}
\GSfield{Statement}{For a discrete group \(G\), the one-object category with automorphism group \(G\) has \(N_kG=G^k\). Its realization is a model of \(BG\).}
\GSfield{Reconstruction}{The discrete case is used here, consistently with the later universal-cover assertion; a topological group requires a topological nerve and different homotopy-group statements.}
\GSfield{Fate}{\citet[\S3]{segal1968classifying}. Compare 1973-308.12.}
\GSfield{Present status}{As of November 2026: established under the stated hypotheses; this occurrence supplies no unresolved problem.}
\GSfield{Assessment}{literal statement: correct with missing hypotheses; publication match: exact; verification: read the result; development: none; confidence: high}
\end{GSitem}
\begin{GSitem}{1974-348.4}{E}{A contractible-cover model for a space}
\GSfield{Kind; source}{claim; \GSpage{P000575}}
\GSfield{The notes}{For an open cover \(\{U_\alpha\}\), intersections are unions of cover members; order is inclusion. If the sets are contractible, the notes assert \(|\mathcal C|\simeq X\), followed by \ru{Добавл. пересечения}.}
\GSfield{Statement}{A sufficient precise version is the nerve theorem for a numerable good open cover of a paracompact space: every nonempty finite intersection is contractible, and \(X\) is homotopy equivalent to the nerve. The category of nonempty finite indexed intersections, ordered by reverse inclusion of index sets, realizes its barycentric subdivision.}
\GSfield{Reconstruction}{This is a sufficient good-cover reconstruction, not a proof of the more general basis-like formulation in the notes. Numerability, finite-intersection hypotheses, and retention of intersection labels are explicit additions.}
\GSfield{Fate}{\citet[\S4, Proposition 4.1]{segal1968classifying} supplies the numerable-cover comparison (1968).}
\GSfield{Present status}{As of November 2026: established under the stated hypotheses; this occurrence supplies no unresolved problem.}
\GSfield{Assessment}{literal statement: correct with missing hypotheses; publication match: partial; verification: read the result; development: annotated; confidence: medium}
\end{GSitem}
\begin{GSitem}{1974-348.5}{E}{The sphere and the octahedral order complex}
\GSfield{Kind; source}{construction; \GSpage{P000576}}
\GSfield{The notes}{Two hemispherical sets have an annular intersection \(w\), so the three-object diagram gives only an interval. Splitting \(w=W_1\cup W_2\) and its overlap into two components produces the drawn octahedral surface.}
\GSfield{Statement}{Take six contractible open sets in three levels \(\{U,V\}\), \(\{W_1,W_2\}\), \(\{X_1,X_2\}\), each lower-level set included in each higher-level set and no other strict inclusions. Their order complex is \(S^0*S^0*S^0\cong S^2\), the octahedral boundary.}
\GSfield{Reconstruction}{The written \(x_i\in(W_1\cap W_2)_{\rm comp}\) are interpreted as component labels; actual open components, rather than selected points, are required for this cover argument. The schematic diagram does not specify a unique cover.}
\GSfield{Fate}{The cover method is in \citet[\S4]{segal1968classifying}; the six-element model is directly reconstructed from the page.}
\GSfield{Present status}{As of November 2026: established under the stated hypotheses; this occurrence supplies no unresolved problem.}
\GSfield{Assessment}{literal statement: correct with missing hypotheses; publication match: partial; verification: read the result; development: annotated; confidence: high}
\end{GSitem}
\begin{GSitem}{1974-348.6}{E}{The spherical building of a vector space}
\GSfield{Kind; source}{definition; \GSpage{P000576}}
\GSfield{The notes}{For a finite-dimensional vector space over any field, the category of proper subspaces and inclusions is called its building.}
\GSfield{Statement}{For a field \(K\) and \(n=\dim_KV\geq2\), let \(\mathcal B(V)=\{W:0<W<V\}\), ordered by inclusion. Its order complex is the spherical building of type \(A_{n-1}\). Both zero and \(V\) are excluded.}
\GSfield{Reconstruction}{The exclusion of zero is supplied by the subsequent definition \(0<W\leq H\) on the same page. Including zero would make the nerve contractible.}
\GSfield{Fate}{\citet[Example 8.2 and end of \S8]{quillen1978poset} checks the model; \citet{tits1974buildings} expanded the 1968 seminar notes in 1974. Compare 1973-321.5.}
\GSfield{Present status}{As of November 2026: established under the stated hypotheses; this occurrence supplies no unresolved problem.}
\GSfield{Assessment}{literal statement: correct with missing hypotheses; publication match: exact; verification: read the result; development: annotated; confidence: high}
\end{GSitem}
\begin{GSitem}{1974-348.7}{E}{The Solomon--Tits homotopy type}
\GSfield{Kind; source}{theorem; \GSpage{P000576}}
\GSfield{The notes}{The building is homotopy equivalent to a wedge of spheres of dimension \(n-2\). The number depends on the field and is infinite for an infinite field.}
\GSfield{Statement}{For \(n\geq2\), \(|\mathcal B(V)|\simeq\bigvee_I S^{n-2}\). For \(K=\mathbb F_q\), \(|I|=q^{\binom n2}\); for an infinite field the index set is infinite. Wedges carry their CW topology, and \(S^0\) is a pointed two-point set.}
\GSfield{Reconstruction}{\ours{} The finite-field count makes the notebook's field-dependence explicit: hyperplanes not containing a fixed line number \(q^{n-1}\), so the induction gives \(r_n=q^{n-1}r_{n-1}\), \(r_2=q\). The same induction gives infinitely many summands for an infinite field.}
\GSfield{Fate}{\citet[Example 8.2 and pp.117--118]{quillen1978poset} gives a later proof of the established theorem. The notebook attributes it to Solomon--Tits, not to new work of Segal.}
\GSfield{Present status}{As of November 2026: established under the stated hypotheses; this occurrence supplies no unresolved problem.}
\GSfield{Assessment}{literal statement: correct with missing hypotheses; publication match: exact; verification: read the result; development: annotated; confidence: high}
\end{GSitem}
\begin{GSitem}{1974-348.8}{E}{A maximum contracts a poset nerve}
\GSfield{Kind; source}{claim; \GSpage{P000576}}
\GSfield{The notes}{If an ordered set has a maximum, its realization is contractible; a cone is drawn.}
\GSfield{Statement}{If a nonempty poset \(P\) has a maximum \(m\), each simplex extends to a simplex with vertex \(m\); hence \(|P|\) is a cone. More generally, an initial or terminal object contracts a category's nerve.}
\GSfield{Reconstruction}{The general categorical version is contextual; the page needs only the maximum case.}
\GSfield{Fate}{\citet[\S1, Proposition 2, Corollary 2]{quillen1973higher}. Earlier occurrence: 1973-309.5.}
\GSfield{Present status}{As of November 2026: established under the stated hypotheses; this occurrence supplies no unresolved problem.}
\GSfield{Assessment}{literal statement: correct; publication match: exact; verification: read the result; development: none; confidence: high}
\end{GSitem}
\begin{GSitem}{1974-348.9}{E}{The cover by hyperplane cones}
\GSfield{Kind; source}{construction; \GSpage{P000576}}
\GSfield{The notes}{The maximal elements are hyperplanes; \(|\mathcal C|=\bigcup_H|\mathcal C_H|\), where \(\mathcal C_H=\{W:0<W\leq H\}\); each member is contractible.}
\GSfield{Statement}{Each \(|\mathcal C_H|\) is a cone on \(H\). Every finite proper-subspace flag has its largest member contained in a hyperplane, so these simplicial subcomplexes cover \(|\mathcal B(V)|\).}
\GSfield{Reconstruction}{The pieces are subcomplexes, not asserted to be open subsets.}
\GSfield{Fate}{The contraction is checked by \citet[\S1, Proposition 2]{quillen1973higher}; compare the dual line-cover proof in \citet[end of \S8]{quillen1978poset}.}
\GSfield{Present status}{As of November 2026: established under the stated hypotheses; this occurrence supplies no unresolved problem.}
\GSfield{Assessment}{literal statement: correct; publication match: partial; verification: read the result; development: none; confidence: high}
\end{GSitem}
\begin{GSitem}{1974-348.10}{E}{Contracting the fixed-line part of the building}
\GSfield{Kind; source}{claim; \GSpage{P000576}}
\GSfield{The notes}{Fix a line \(L\), and put \(Y=\bigcup_{H\supset L}|\mathcal C_H|\). The notes contract \(Y\) using intersections of the cone cover.}
\GSfield{Statement}{Let \(P_L=\{0<W<V:W+L<V\}\). Then \(Y=|P_L|\). The order-preserving map \(W\mapsto W+L\) satisfies \(W\leq W+L\geq L\); realization is homotopic both to the identity and to the constant map at \(L\), contracting \(Y\).}
\GSfield{Reconstruction}{\ours{} This direct order-homotopy argument supplies the finite-intersection step for arbitrary fields and arbitrary numbers of hyperplanes.}
\GSfield{Fate}{\citet[\S1, Proposition 2]{quillen1973higher} supplies the homotopy principle. The particular contraction is reconstructed from the notebook.}
\GSfield{Present status}{As of November 2026: established under the stated hypotheses; this occurrence supplies no unresolved problem.}
\GSfield{Assessment}{literal statement: correct; publication match: partial; verification: read the result; development: annotated; confidence: high}
\end{GSitem}
\begin{GSitem}{1974-348.11}{E}{The nonempty finite-intersection criterion}
\GSfield{Kind; source}{claim; \GSpage{P000577}}
\GSfield{The notes}{If all finite intersections of the covering pieces are contractible and nonempty, their union is contractible. The words excluding an empty intersection are partly overwritten but visible.}
\GSfield{Statement}{For a numerable open good cover, or a cover of a simplicial complex by subcomplexes, suppose the intersection of every nonempty finite subfamily is nonempty and contractible. The nerve is then a full simplex and the union is contractible.}
\GSfield{Reconstruction}{The requirement concerns every finite subfamily, not merely those whose intersection happens to be nonempty. The cover topology is supplied. The subcomplex version applies to the building.}
\GSfield{Fate}{The nerve argument is used in \citet[end of \S8]{quillen1978poset}; the open-cover comparison is in \citet[\S4]{segal1968classifying}.}
\GSfield{Present status}{As of November 2026: established under the stated hypotheses; this occurrence supplies no unresolved problem.}
\GSfield{Assessment}{literal statement: correct with missing hypotheses; publication match: partial; verification: read the result; development: annotated; confidence: high}
\end{GSitem}
\begin{GSitem}{1974-348.12}{E}{Collapsing the common part gives a wedge}
\GSfield{Kind; source}{claim; \GSpage{P000577}}
\GSfield{The notes}{Lemma 1: if \(X=Y\cup Y_1\cup\cdots\cup Y_k\), all pieces are contractible and \(Y_i\cap Y_j\subset Y\) for distinct indices, then \(X\simeq\bigvee_iY_i/(Y_i\cap Y)\).}
\GSfield{Statement}{Let \(X\) be a CW complex covered by subcomplexes \(Y,Y_i\), with \(Y\) nonempty contractible, \(Y_i\cap Y\ne\varnothing\), and \(Y_i\cap Y_j\subset Y\) for \(i\ne j\). Then \(X\to X/Y\) is a homotopy equivalence and \(X/Y\cong\bigvee_iY_i/(Y_i\cap Y)\). Contractibility of the \(Y_i\) is used in the next step.}
\GSfield{Reconstruction}{CW/cofibration and nonempty-intersection hypotheses are added. The result extends to arbitrary families of subcomplexes with CW weak topology, as required over an infinite field.}
\GSfield{Fate}{This elementary quotient fact is used in the 1974 talk. A modern check is \citet[Proposition 0.17]{hatcher2002algebraic}.}
\GSfield{Present status}{As of November 2026: established under the stated hypotheses; this occurrence supplies no unresolved problem.}
\GSfield{Assessment}{literal statement: correct with missing hypotheses; publication match: partial; verification: read the result; development: annotated; confidence: high}
\end{GSitem}
\begin{GSitem}{1974-348.13}{E}{A cofiber in a contractible space is a suspension}
\GSfield{Kind; source}{claim; \GSpage{P000577}}
\GSfield{The notes}{Lemma 2 states \(Z/Z_0\simeq SZ_0\) when \(Z\) is contractible and \(Z_0\subset Z\); cone and suspension diagrams explain the replacement.}
\GSfield{Statement}{For a CW pair \((Z,Z_0)\) with \(Z\) contractible and \(Z_0\ne\varnothing\), the quotient is homotopy equivalent to the unreduced suspension \(SZ_0\). Use the mapping cone of the inclusion, which is nullhomotopic.}
\GSfield{Reconstruction}{The cofibration condition and nonemptiness qualify the unqualified subset notation.}
\GSfield{Fate}{A later precise reference for this established cofiber fact is \citet[Examples 0.13--0.14]{hatcher2002algebraic}.}
\GSfield{Present status}{As of November 2026: established under the stated hypotheses; this occurrence supplies no unresolved problem.}
\GSfield{Assessment}{literal statement: correct with missing hypotheses; publication match: exact; verification: read the result; development: annotated; confidence: high}
\end{GSitem}
\begin{GSitem}{1974-348.14}{E}{The hyperplane intersection and dimension induction}
\GSfield{Kind; source}{method; \GSpage{P000577}}
\GSfield{The notes}{Lemma 3 identifies \(|\mathcal C_H|\cap Y\) with the building of \(H\), yielding reduction by dimension.}
\GSfield{Statement}{If \(H\) does not contain \(L\), then \(\mathcal C_H\cap P_L=\{W:0<W<H\}\). Indeed \(H+L=V\), while \(W<H\) implies \(\dim(W+L)\leq n-1\). For \(n\geq3\), the preceding decomposition therefore gives \(|\mathcal B(V)|\simeq\bigvee_{H\not\supset L} S|\mathcal B(H)|\).}
\GSfield{Reconstruction}{\ours{} The condition \(H\not\supset L\) is restored from the preceding line. For distinct such hyperplanes, \(\dim(H\cap H')=n-2\), so their cone intersection lies in \(Y\), as required. The base case is separate.}
\GSfield{Fate}{This is the notebook's proof of the theorem checked in \citet[\S8]{quillen1978poset}; that source uses a different, dual cover proof.}
\GSfield{Present status}{As of November 2026: established under the stated hypotheses; this occurrence supplies no unresolved problem.}
\GSfield{Assessment}{literal statement: correct with missing hypotheses; publication match: partial; verification: read the result; development: annotated; confidence: high}
\end{GSitem}
\begin{GSitem}{1974-348.15}{E}{The dimension-two base case}
\GSfield{Kind; source}{claim; \GSpage{P000577}}
\GSfield{The notes}{The first induction step is \(\dim V=2\), where the building is a collection of points (lines); the phrase about a wedge is overwritten.}
\GSfield{Statement}{For \(\dim_KV=2\), \(|\mathcal B(V)|=\mathbb P^1(K)\) with the discrete topology. A chosen base line identifies it with a wedge of copies of \(S^0\), one for every remaining line. Over \(\mathbb F_q\) there are \(q+1\) points and \(q\) copies of \(S^0\).}
\GSfield{Reconstruction}{\ours{} A wedge of one-point spaces would be a point; the intended spheres are pointed two-point spaces. This also fixes the homology convention below.}
\GSfield{Fate}{The base case agrees with \citet[Example 8.2]{quillen1978poset} and is directly visible in the row of points.}
\GSfield{Present status}{As of November 2026: established under the stated hypotheses; this occurrence supplies no unresolved problem.}
\GSfield{Assessment}{literal statement: not determinable; publication match: partial; verification: read the result; development: annotated; confidence: high}
\end{GSitem}
\begin{GSitem}{1974-348.16}{E}{Suspending a wedge of spheres}
\GSfield{Kind; source}{claim; \GSpage{P000577}}
\GSfield{The notes}{The suspension of a wedge of spheres is again a wedge of spheres.}
\GSfield{Statement}{For a based CW complex \(A=\bigvee_I S^d\), \(d\geq0\), the reduced suspension is \(\Sigma A\cong\bigvee_I S^{d+1}\). The unreduced suspension \(SA\) is homotopy equivalent to \(\Sigma A\), by collapsing the basepoint interval.}
\GSfield{Reconstruction}{The based CW convention includes the disconnected \(d=0\) case in the induction.}
\GSfield{Fate}{The standard argument is checked in \citet[Example 0.10]{hatcher2002algebraic}; the notebook uses it in 1974.}
\GSfield{Present status}{As of November 2026: established under the stated hypotheses; this occurrence supplies no unresolved problem.}
\GSfield{Assessment}{literal statement: correct with missing hypotheses; publication match: exact; verification: read the result; development: none; confidence: high}
\end{GSitem}
\begin{GSitem}{1974-348.17}{E}{The Steinberg module and reduced homology}
\GSfield{Kind; source}{definition; \GSpage{P000577}}
\GSfield{The notes}{The displayed line identifies \(H_{n-2}(|\mathcal C|;\mathbb Z)\) with the Steinberg module of \(V\).}
\GSfield{Statement}{The integral Steinberg module is \(\operatorname{St}(V)=\widetilde H_{n-2}(|\mathcal B(V)|;\mathbb Z)\), with its natural \(\operatorname{GL}(V)\)-action. For \(n\geq3\), ordinary and reduced homology agree in this degree; for \(n=2\), reduced homology is necessary.}
\GSfield{Reconstruction}{\ours{} Over \(\mathbb F_q\) in dimension two, ordinary \(H_0\) has rank \(q+1\), while the reduced group has rank \(q\). The literal formula fails at the base case retained in this talk.}
\GSfield{Fate}{The building interpretation belongs to the older Solomon--Tits theory; \citet[\S\S8--9]{quillen1978poset} checks its reduced-homology dimension.}
\GSfield{Present status}{As of November 2026: established under the stated hypotheses; this occurrence supplies no unresolved problem.}
\GSfield{Assessment}{literal statement: false as written; publication match: partial; verification: read the result; development: annotated; confidence: high}
\end{GSitem}
\begin{GSitem}{1974-348.18}{E}{The classifying space of a discrete group is aspherical}
\GSfield{Kind; source}{theorem; \GSpage{P000577}}
\GSfield{The notes}{For the category of a group, \(|\mathcal C|=BG\), \(\pi_1(|\mathcal C|)=G\), and the universal covering is contractible.}
\GSfield{Statement}{For a discrete group \(G\), the nerve realization \(BG\) is a connected CW complex with \(\pi_1(BG)\cong G\) and \(\pi_k(BG)=0\) for \(k>1\).}
\GSfield{Reconstruction}{Discreteness is essential: for a well-behaved topological group the usual formula is \(\pi_k(BG)\cong\pi_{k-1}(G)\), not asphericity in general.}
\GSfield{Fate}{\citet[\S3]{segal1968classifying}, published six years before this exposition.}
\GSfield{Present status}{As of November 2026: established under the stated hypotheses; this occurrence supplies no unresolved problem.}
\GSfield{Assessment}{literal statement: correct with missing hypotheses; publication match: exact; verification: read the result; development: none; confidence: high}
\end{GSitem}
\begin{GSitem}{1974-348.19}{E}{The contractible chaotic category}
\GSfield{Kind; source}{construction; \GSpage{P000577-P000578}}
\GSfield{The notes}{\(\widetilde{\mathcal C}\) has objects the elements of \(G\) and one morphism between every ordered pair. Its realization is called contractible because every element is linked to \(1\).}
\GSfield{Statement}{The unique morphisms \(g\to1\) form a natural transformation from the identity functor to the constant functor at \(1\). Realization gives a contraction. Simplices are arbitrary finite strings of elements of \(G\).}
\GSfield{Reconstruction}{Connectedness of a graph alone is insufficient; coherent higher simplices and the natural transformation give the actual argument.}
\GSfield{Fate}{\citet[\S3 and Proposition 2.1]{segal1968classifying}.}
\GSfield{Present status}{As of November 2026: established under the stated hypotheses; this occurrence supplies no unresolved problem.}
\GSfield{Assessment}{literal statement: correct; publication match: exact; verification: read the result; development: annotated; confidence: high}
\end{GSitem}
\begin{GSitem}{1974-348.20}{E}{The universal-cover functor and multiplication convention}
\GSfield{Kind; source}{construction; \GSpage{P000578}}
\GSfield{The notes}{The functor sends each object to the unique object and \(g\to g'\) to \(g^{-1}g'\). Its realization is the universal covering.}
\GSfield{Statement}{Adopt the convention that a path labelled first \(a\), then \(b\), has composite label \(ab\). Then \((g^{-1}g')(g'^{-1}g'')=g^{-1}g''\), so this is a functor. Left multiplication on strings by discrete \(G\) gives a principal covering \(EG\to BG\), with \(EG\) the contractible realization of the chaotic category.}
\GSfield{Reconstruction}{An opposite composition convention requires changing the formula consistently. The scan is not emended merely because conventions differ.}
\GSfield{Fate}{\citet[\S3]{segal1968classifying} gives the construction and the extra local-triviality issue for general topological groups.}
\GSfield{Present status}{As of November 2026: established under the stated hypotheses; this occurrence supplies no unresolved problem.}
\GSfield{Assessment}{literal statement: correct with missing hypotheses; publication match: exact; verification: read the result; development: none; confidence: high}
\end{GSitem}
\begin{GSitem}{1974-348.21}{E}{Transition functions and a bundle's classifying map}
\GSfield{Kind; source}{method; \GSpage{P000578}}
\GSfield{The notes}{A bundle \(E\to X\), a trivializing cover, intersections, and \(X\to BG\) are drawn. The last sentence about the intersection category stops before completing the construction.}
\GSfield{Statement}{For a numerable principal \(G\)-bundle over a paracompact base, continuous transition functions define a functor from the topological \v{C}ech category, retaining points in intersections, to the one-object topological category \(G\). Realization and the cover equivalence give the classifying map up to homotopy, using a standard numerable universal-bundle model.}
\GSfield{Reconstruction}{The completion is editorial. The notes do not specify discrete versus topological \(G\), a principal bundle, or how point-dependent transitions enter the functor. A category of unlabelled open sets alone cannot encode arbitrary transition functions.}
\GSfield{Fate}{The completion was published in \citet[\S4, Proposition 4.1]{segal1968classifying}; the surviving final sentence is a partial match.}
\GSfield{Present status}{As of November 2026: established under the stated hypotheses; this occurrence supplies no unresolved problem.}
\GSfield{Assessment}{literal statement: not determinable; publication match: partial; verification: read the result; development: annotated; confidence: high}
\end{GSitem}

\GStalk{349}{1974-10-21}{A. G. Kushnirenko}{P000578-P000580}{Global critical-point counts and Newton polyhedra}
The upper portion ofP000579 treats all isolated critical points of a polynomial on \(\mathbb C^n\). Its polytope is the convex hull of 0 and the support, rather than the local lower Newton boundary. The total Milnor number is bounded above by the associated Newton number, with equality under nondegeneracy at infinity. The normalized-gradient degree explains the total multiplicity; Hochster\textquotesingle s theorem and Mumford\textquotesingle s toric viewpoint are named as tools.

\begin{GSitem}{1974-349.1}{A}{The total Milnor number of a polynomial}
\GSfield{Kind; source}{question; \GSpage{P000579, polynomial and counting question}}
\GSfield{The notes}{For \(P:\mathbb C^n\to\mathbb C\), all critical points are assumed isolated. The question is how to count them with multiplicities.}
\GSfield{Statement}{When the gradient ideal has finite colength, the finite sum of local critical-point multiplicities is
\[\begin{aligned}
\mu_{\mathrm{tot}}(P)
 &=\dim_{\mathbb C}\mathbb C[z_1,\ldots,z_n]/(P_{z_1},\ldots,P_{z_n})\\
 &=\sum_{a\in\operatorname{Crit}(P)}\mu_a(P).
\end{aligned}\] The question is to bound and generically compute this number from the support of \(P\).}
\GSfield{Fate}{The equality with the sum of local algebras is \citet[\S\S4--5]{kushnirenko1975newton}; the ensuing Newton bound is its TheoremII. The paper was received 30 September 1974 and published 1975; the notebook records the global discussion on21 October 1974.}
\GSfield{Present status}{As of November 2026, this is an established result in the stated scope; the qualification in the statement is essential.}
\GSfield{Assessment}{literal statement: correct; publication match: exact; verification: read the result; development: none; confidence: high; problem \hyperref[prob:GP-1974-13]{GP-1974-13}}
\end{GSitem}
\begin{GSitem}{1974-349.2}{A}{The Newton upper bound and equality at infinity}
\GSfield{Kind; source}{theorem; \GSpage{P000579, convex hull and two inequalities}}
\GSfield{The notes}{Take the convex hull of 0 and the exponents with nonzero coefficients. In two dimensions the count is at most \(2S-a-b+1\), with equality under nondegeneracy.}
\GSfield{Statement}{Let \(P\) have a finite critical locus and be convenient, so its support meets every coordinate axis. For \(\Delta_\infty=\operatorname{conv}(\{0\}\cup\operatorname{supp}P)\), define \(\nu_\infty\) by the same alternating sum of coordinate-section volumes as in record 346. Then \(\sum\mu_a(P)\le\nu_\infty\). Equality holds if every face at infinity is nondegenerate in the logarithmic-derivative sense; in that case the critical locus is finite. In dimension 2 this is \(2S-a-b+1\).}
\GSfield{Reconstruction}{Convenience and nondegeneracy on every relevant face at infinity are supplied; arbitrary support or a single tested face does not justify the formula.}
\GSfield{Fate}{Exact match with \citet[TheoremII(i)--(iii)]{kushnirenko1975newton}, including genericity at fixed support; \citet{kouchnirenko1976polyedres} gives the full treatment. This upper bound and the local lower bound concern different Newton regions and are compatible.}
\GSfield{Present status}{As of November 2026, this is an established result in the stated scope; the qualification in the statement is essential.}
\GSfield{Assessment}{literal statement: correct with missing hypotheses; publication match: exact; verification: read the result; development: none; confidence: high; problem \hyperref[prob:GP-1974-13]{GP-1974-13}}
\end{GSitem}
\begin{GSitem}{1974-349.3}{E}{Global multiplicity as the degree of the normalized gradient}
\GSfield{Kind; source}{method; \GSpage{P000579, degree-of-map argument}}
\GSfield{The notes}{The map on \(|z|=R\) sends \(z\) to \((P_{z_1},\ldots,P_{z_n})/|\operatorname{grad}P|\); the sum of Milnor numbers equals its degree.}
\GSfield{Statement}{Choose \(R\) so every isolated critical point lies strictly inside the ball and the gradient is nonzero on its boundary. With complex orientations, \(\deg(\nabla P/\|\nabla P\|:S_R^{2n-1}\to S^{2n-1})=\sum_a\mu_a(P)\). This is additivity of the degree of a holomorphic map with isolated zeros. It does not require the gradient to be a proper map of all of \(\mathbb C^n\).}
\GSfield{Fate}{The local and global degree interpretation is stated explicitly in \citet[\S1.1, Eqs.~(1)--(3), p.~2]{kouchnirenko1976polyedres}, which credits the established singularity theory preceding the talk. Compare \citet[\S7]{milnor1968singular} and the support-dependent count in \citet[\S\S4--5]{kushnirenko1975newton}. The notebook uses the degree principle as a proof idea, not as a new independent theorem.}
\GSfield{Present status}{As of November 2026, this is an established result in the stated scope; the qualification in the statement is essential.}
\GSfield{Assessment}{literal statement: correct with missing hypotheses; publication match: exact; verification: read the result; development: none; confidence: high}
\end{GSitem}
\begin{GSitem}{1974-349.4}{E}{Normal cone semigroup algebras are Cohen--Macaulay}
\GSfield{Kind; source}{theorem; \GSpage{P000579, Hochster sentence}}
\GSfield{The notes}{The ring of monomials whose exponents lie in a cone is said to be Cohen--Macaulay; Melvin Hochster is named at the head of the page.}
\GSfield{Statement}{For a rational polyhedral cone \(\sigma\subset\mathbb R_{\ge0}^n\), the saturated finitely generated semigroup \(\sigma\cap\mathbb Z^n\) defines a normal affine semigroup ring \(\mathbb C[\sigma\cap\mathbb Z^n]\), and this ring is Cohen--Macaulay. Saturation is a hypothesis, not a consequence for an arbitrary selected set of monomial generators.}
\GSfield{Fate}{\citet[main theorem, p.~318]{hochster1972invariants}, published September 1972. This is the same proof ingredient already mentioned in 1974-346.6; the second occurrence is retained because the global counting argument invokes it explicitly.}
\GSfield{Present status}{As of November 2026, this is an established result in the stated scope; the qualification in the statement is essential.}
\GSfield{Assessment}{literal statement: correct with missing hypotheses; publication match: exact; verification: read the result; development: none; confidence: high}
\end{GSitem}
\begin{GSitem}{1974-349.5}{F}{Toric compactifications as the geometric interpretation}
\GSfield{Kind; source}{remark; \GSpage{P000579, Mumford line}}
\GSfield{The notes}{The final line before Polyakov says only \ru{Мамфорд --- торич. вложения}.}
\GSfield{Statement}{The identifiable method is to treat the cone monomial rings as affine charts of toric varieties and use a toric compactification adapted to the Newton polytope. No particular fan, compactification theorem or intersection formula is specified on this page.}
\GSfield{Fate}{The connection is consistent with the toric proof language of the Newton-polyhedron literature, including \citet[proof of the theorem]{khovanskii1978eulerjacobi}. That is a thematic comparison only; this isolated name does not identify a distinct theorem of Mumford or justify a new open-problem entry.}
\GSfield{Present status}{As of November 2026, the toric-geometric method is identifiable, but the exact Mumford result intended by this isolated line is not. No mathematical question is stated.}
\GSfield{Assessment}{literal statement: not determinable; publication match: thematic only; verification: unresolved lead; development: none; confidence: medium}
\end{GSitem}

\GStalk{350}{1974-10-21}{A. M. Polyakov}{P000578-P000582}{Yang--Mills connections, curvature, and characteristic forms}
The lower portion ofP000579 returns from the lattice model to the continuum action. Polyakov identifies the gauge potential with connection coefficients, the field strength with curvature, and a quadratic curvature expression with a Pontryagin form. This is an expository geometric dictionary. The handwritten action has misplaced indices and parentheses, and its square must encompass the entire curvature.

\begin{GSitem}{1974-350.1}{E}{The Yang--Mills curvature-squared action}
\GSfield{Kind; source}{construction; \GSpage{P000579, action display}}
\GSfield{The notes}{The displayed action contains \(\partial_\alpha b_{\alpha\beta}-\partial_\beta b_\alpha-[b_\alpha,b_\beta]^2\) inside an incomplete trace parenthesis.}
\GSfield{Statement}{Fix a compact Lie group, an invariant positive inner product on its Lie algebra, and a connection \(A=\sum_\alpha A_\alpha dx^\alpha\) on Euclidean four-space. In the convention \(F_A=dA+A\wedge A\), \(F_{\alpha\beta}=\partial_\alpha A_\beta-\partial_\beta A_\alpha+[A_\alpha,A_\beta]\), and \(S(A)=(2g^2)^{-1}\int\langle F_A,F_A\rangle\,d^4x\), for fields for which the integral exists. A change of Lie-algebra convention may change the commutator sign, but the whole curvature is squared.}
\GSfield{Reconstruction}{The invariant inner product, Euclidean convention and full outer square are editorial repairs. No existence theorem for the quantized theory is present.}
\GSfield{Fate}{\citet[Eq.~(4) and the field-equations section]{yangmills1954conservation} gives the curvature and invariant Lagrangian. The theory predates this lecture by twenty years. The manuscript display cannot literally be a curvature-squared action and requires an explicit emendation.}
\GSfield{Present status}{As of November 2026, this is an established result in the stated scope; the qualification in the statement is essential.}
\GSfield{Assessment}{literal statement: false as written; publication match: exact; verification: read the result; development: annotated; confidence: high}
\end{GSitem}
\begin{GSitem}{1974-350.2}{E}{Gauge potentials as local connection forms}
\GSfield{Kind; source}{definition; \GSpage{P000579, connection and Riemann-curvature comparison}}
\GSfield{The notes}{An element of a Lie algebra at every point is identified with a connection; taking the coefficients to be Christoffel symbols gives the curvature tensor. The fibre is called a Lie algebra, with SU(2) named as example.}
\GSfield{Statement}{A connection is locally a Lie-algebra-valued one-form, not merely one Lie-algebra element at each point. On overlaps its local forms obey the affine gauge-transformation law. Its curvature is a two-form valued in the adjoint bundle; the analogous construction for a linear connection on the tangent bundle yields the Riemann curvature tensor. The Lie algebra for the named group is \(\mathfrak{su}(2)\).}
\GSfield{Reconstruction}{The one-form index, overlap law and distinction between SU(2) and its Lie algebra correct the notebook\textquotesingle s compressed dictionary.}
\GSfield{Fate}{The standard connection-curvature dictionary is reviewed in \citet[\S2, Eq.~(2.10)]{chernsimons1974characteristic}; Yang--Mills gauge covariance is already in \citet[Eq.~(3)]{yangmills1954conservation}. Both are earlier than 21 October 1974.}
\GSfield{Present status}{As of November 2026, this is an established result in the stated scope; the qualification in the statement is essential.}
\GSfield{Assessment}{literal statement: correct with missing hypotheses; publication match: exact; verification: read the result; development: none; confidence: high}
\end{GSitem}
\begin{GSitem}{1974-350.3}{E}{The quadratic curvature form representing a Pontryagin class}
\GSfield{Kind; source}{construction; \GSpage{P000579, final Pontryagin expression}}
\GSfield{The notes}{The contracted expression \(R^b_{a,\delta_1\delta_2}R^a_{b,\delta_3\delta_4}\) is labelled a Pontryagin class.}
\GSfield{Statement}{For a real vector bundle with a metric connection and skew-symmetric curvature matrix \(\Omega\), the closed four-form \(-\tfrac1{8\pi^2}\operatorname{tr}(\Omega\wedge\Omega)\) represents the real image of \(p_1\). The written component contraction supplies the coefficient pattern only after alternating the four differential indices. A trace in a different group representation has a different normalization; for the fundamental complex rank-two \(SU(2)\) bundle one normally speaks of the second Chern class.}
\GSfield{Reconstruction}{The exterior product, normalization, real bundle and de Rham representative are necessary to turn the component mnemonic into a precise statement.}
\GSfield{Fate}{This is the Chern--Weil theorem, reviewed in \citet[Theorem 2.12 and Eq.~(5.7)]{chernsimons1974characteristic}, published January 1974 and explicitly presented there as established connection theory. No new characteristic-class construction is asserted on the notebook page.}
\GSfield{Present status}{As of November 2026, this is an established result in the stated scope; the qualification in the statement is essential.}
\GSfield{Assessment}{literal statement: correct with missing hypotheses; publication match: exact; verification: read the result; development: none; confidence: high}
\end{GSitem}

\GStalk{351}{1974-10-28}{M. I. Graev}{P000580-P000585}{Irreducible representations of current groups and continuous tensor products}
P000580--P000583 describes the Vershik--Gelfand--Graev construction for continuous maps from an interval to \(SL_2(\mathbb R)\). Complementary-series representations are embedded compatibly under successive subdivisions; a second realization uses formal coherent vectors and a logarithmic kernel. The detailed construction was already published in 1973, and is therefore expository in this 1974 lecture. The reference to Functional Analysis, volume 2, number 1, points instead to the earlier 1968 Gelfand--Graev paper. Gelfand\textquotesingle s intervening question about Kolmogorov dimension and the drawing of negative squares lack enough normalization data for a unique identification.

\begin{GSitem}{1974-351.1}{E}{The interval current group}
\GSfield{Kind; source}{definition; \GSpage{P000580, matrix and programme}}
\GSfield{The notes}{The matrix entries of \(g(x)\) are continuous real functions on \([0,1]\), with \(\det g(x)=1\). The announced method of constructing irreducible representations is compared with integration.}
\GSfield{Statement}{Let \(G^X=C([0,1],SL_2(\mathbb R))\) with pointwise multiplication and the uniform topology. A unitary representation is understood to be strongly continuous. The construction recorded here factors through the pointwise quotient \(PSL_2(\mathbb R)\), so it is not required to detect central sign-valued currents.}
\GSfield{Reconstruction}{Uniform topology, strong continuity and the central quotient are explicit. The 1968 precursor was checked bibliographically; the 1973 construction was read in full relevant sections.}
\GSfield{Fate}{This is the precise setup of \citet[Introduction and \S2]{vershikgelfandgraev1973functions}, received 15 June 1973 and published 1973. The notebook\textquotesingle s volume 2, number 1 citation identifies \citet{gelfandgraev1968quaternion}, received 16 October 1967; it is a precursor, not the publication containing every formula on these pages.}
\GSfield{Present status}{As of November 2026, this is an established result in the stated scope; the qualification in the statement is essential.}
\GSfield{Assessment}{literal statement: correct with missing hypotheses; publication match: exact; verification: read the result; development: none; confidence: high}
\end{GSitem}
\begin{GSitem}{1974-351.2}{E}{Finite evaluation representations and tensor products}
\GSfield{Kind; source}{construction; \GSpage{P000581, two-point picture}; I. M. Gelfand}
\GSfield{The notes}{Gelfand replaces the interval by two points and lets \(g(x_1),g(x_2)\) act separately in \(H_1\otimes H_2\).}
\GSfield{Statement}{At distinct points \(x_1,\ldots,x_r\), choose irreducible unitary representations \(\pi_j\) of \(SL_2(\mathbb R)\). The formula \(\pi(g)=\bigotimes_j\pi_j(g(x_j))\) defines an irreducible unitary current-group representation. Evaluation onto the finite product is surjective because this Lie group is connected and the interval permits interpolation. Conversely, irreducible representations of the finite product decompose as outer tensor products because \(SL_2(\mathbb R)\) is type I.}
\GSfield{Reconstruction}{Distinct evaluation points, irreducibility and the distinction between an outer tensor product and restriction to the diagonal subgroup are supplied.}
\GSfield{Fate}{The finite-product representation statement appears explicitly in \citet[\S1.3, Eq.~(14)]{vershikgelfandgraev1973functions}. It is background for the continuous tensor product, and recurs from 1969-180.4.}
\GSfield{Present status}{As of November 2026, this is an established result in the stated scope; the qualification in the statement is essential.}
\GSfield{Assessment}{literal statement: correct with missing hypotheses; publication match: exact; verification: read the result; development: none; confidence: high}
\end{GSitem}
\begin{GSitem}{1974-351.3}{E}{Characters and distributions, with the circle-valued qualification}
\GSfield{Kind; source}{construction; \GSpage{P000581, Abelian example}; I. M. Gelfand}
\GSfield{The notes}{For \(g(x)=e^{i\alpha(x)}\), an irreducible representation is written \(e^{i\ell(\alpha)}\), with \(\ell\) a linear functional; smooth test functions give distributions.}
\GSfield{Statement}{For the additive group of real smooth functions, continuous unitary characters are \(\alpha\mapsto e^{i\ell(\alpha)}\), with \(\ell\) a real continuous linear functional. Such functionals are distributions. To descend this formula to \(C^\infty([0,1],S^1)\), require \(\ell(1)\in\mathbb Z\), since two smooth lifts differ by a constant integral multiple of \(2\pi\). For continuous test functions the appropriate continuous dual consists of signed measures, rather than arbitrary distributions.}
\GSfield{Reconstruction}{The distinction between the additive group and its circle-valued quotient, and the dependence on the test-function topology, are editorial qualifications.}
\GSfield{Fate}{\citet[Introduction, pp.~87--88]{vershikgelfandgraev1973functions} uses the additive real current group, for which the formula has no period obstruction. The notebook\textquotesingle s circle notation requires the extra integrality condition. Related earlier occurrence:1969-180.9.}
\GSfield{Present status}{As of November 2026, this is an established result in the stated scope; the qualification in the statement is essential.}
\GSfield{Assessment}{literal statement: correct with missing hypotheses; publication match: exact; verification: read the result; development: annotated; confidence: high}
\end{GSitem}
\begin{GSitem}{1974-351.4}{E}{A spatially uniform irreducible current-group representation}
\GSfield{Kind; source}{question; \GSpage{P000581, uniformity question}; I. M. Gelfand}
\GSfield{The notes}{Gelfand asks whether there is an irreducible representation \ru{которое всюду устроено одинаково}.}
\GSfield{Statement}{For \(X=[0,1]\), \(G=PSL_2(\mathbb R)\), and a positive multiple \(m=\lambda\,dx\) of Lebesgue measure, construct a nontrivial irreducible strongly continuous representation of \(C(X,G)\) whose equivalence class is invariant under measure-preserving transformations of \(X\) that act on the chosen current group by pullback. In the continuous model these transformations may be taken to be measure-preserving homeomorphisms; the step-function construction accommodates the corresponding measurable formulation. Require also that the restriction to currents supported in any nonempty open subinterval is nontrivial. This full-support condition excludes finite sets of evaluation points and follows from the positive measure on every such interval in the construction.}
\GSfield{Reconstruction}{Uniformity is interpreted through the measure-invariance property explicitly stated in the matching paper. This is an existence theorem for the specified group and topology, not a classification of all current-group representations. Full support is the explicit reconstruction of everywhere in the question; measure-preserving homeomorphisms of an interval alone would be too small a symmetry group to exclude evaluation examples.}
\GSfield{Fate}{This existence question is answered by \citet[Theorems2.2, 3.3 and 3.4]{vershikgelfandgraev1973functions}, already published in 1973. The result is related to 1969-180.2 and GP-1969-02, but only the earlier title is available: neither identity of their hypotheses nor resolution of the broader earlier problem is asserted.}
\GSfield{Present status}{As of November 2026, solved in the stated PSL2(R), uniform-topology setting by the 1973 multiplicative-integral construction. The classifying problem GP-1969-01 and the broader interval-support question GP-1969-02 are not thereby declared solved.}
\GSfield{Assessment}{literal statement: correct with missing hypotheses; publication match: exact; verification: read the result; development: none; confidence: high; problem \hyperref[prob:GP-1974-14]{GP-1974-14}}
\end{GSitem}
\begin{GSitem}{1974-351.5}{E}{The complementary series in the real-line model}
\GSfield{Kind; source}{construction; \GSpage{P000581, complementary-series norm and final action}}
\GSfield{The notes}{For \(0<\lambda<1\), \(\|f\|_\lambda^2=\iint |t-s|^{-\lambda}f(t)\overline{f(s)}\,dt\,ds\), and \(T_\lambda(g)f(t)=f((\alpha t+\gamma)/(\beta t+\delta))|\beta t+\delta|^{\lambda-2}\).}
\GSfield{Statement}{Complete a suitable dense test-function space in the positive kernel norm. The displayed operators extend to irreducible unitary representations for \(0<\lambda<1\), and factor through \(PSL_2(\mathbb R)\). Positivity can be seen from the Fourier multiplier \(c_\lambda|\xi|^{\lambda-1}\), with \(c_\lambda>0\). The completed space can contain generalized functions and is not simply a pointwise class of ordinary functions with an absolutely convergent integral.}
\GSfield{Reconstruction}{Completion, the dense test space and the row-vector convention for the fractional-linear transformation are made explicit.}
\GSfield{Fate}{Exact real-line formulas in \citet[\S1.1, Eqs.~(3)--(4)]{vershikgelfandgraev1973functions}; the authors describe this as the established supplementary series. The parameter interval is the one used for the integral kernel, and does not exhaust equivalent realizations obtained by continuation.}
\GSfield{Present status}{As of November 2026, this is an established result in the stated scope; the qualification in the statement is essential.}
\GSfield{Assessment}{literal statement: correct with missing hypotheses; publication match: exact; verification: read the result; development: none; confidence: high}
\end{GSitem}
\begin{GSitem}{1974-351.6}{E}{The trivial representation is an accumulation point}
\GSfield{Kind; source}{claim; \GSpage{P000581, representations close to one}; I. M. Gelfand}
\GSfield{The notes}{Gelfand stresses the existence of infinite-dimensional representations close to 1; a second parameter is written \(\rho\approx1\), beside \(\rho=1\) labelled the identity representation.}
\GSfield{Statement}{The spherical matrix coefficients of the complementary series tend to 1 as \(\lambda\downarrow0\), uniformly on compact subsets of the group, after normalizing the spherical vector. Thus the trivial representation is not isolated in the unitary dual. An endpoint pullback formula with exponent0 is not itself trivial on every function: the trivial representation is obtained through the appropriate invariant subspace or quotient.}
\GSfield{Reconstruction}{Closeness is stated in terms of matrix coefficients, not operator-norm convergence of operators on unidentified varying Hilbert spaces.}
\GSfield{Fate}{\citet[Introduction and Lemma 1.1]{vershikgelfandgraev1973functions} gives precisely the limiting spherical function. The two notebook parameters describe endpoint realizations and must not be equated without a convention change.}
\GSfield{Present status}{As of November 2026, this is an established result in the stated scope; the qualification in the statement is essential.}
\GSfield{Assessment}{literal statement: correct with missing hypotheses; publication match: exact; verification: read the result; development: none; confidence: high}
\end{GSitem}
\begin{GSitem}{1974-351.7}{F}{An unidentified notion of representation dimension}
\GSfield{Kind; source}{question; \GSpage{P000581, Kolmogorov-dimension intervention}; I. M. Gelfand}
\GSfield{The notes}{The page labels a \ru{Задача} about dimension in Kolmogorov\textquotesingle s sense, refers to \ru{один японец}, and tentatively equates this dimension with a multiplier of Plancherel measure after choosing a kernel.}
\GSfield{Statement}{The notes do not define the dimension, the kernel, its normalization, the reference measure in the Plancherel decomposition, or the named author. Metric entropy dimension, formal degree and spectral multiplicity are different notions. No precise equality or research problem can therefore be recovered uniquely.}
\GSfield{Fate}{The matching current-group paper \citet{vershikgelfandgraev1973functions} does not supply this intervening definition. Searches for the Russian and English combinations of Kolmogorov dimension, Plancherel measure and representations produced no verified exact identification. Araki\textquotesingle s appearance elsewhere in current-group literature is insufficient to identify the unnamed Japanese mathematician.}
\GSfield{Present status}{As of November 2026, the intended dimension and reference remain unidentified. This is retained as an unresolved historical reading, with no fabricated formal-degree theorem and no problem-index entry.}
\GSfield{Assessment}{literal statement: not determinable; publication match: none found; verification: unresolved lead; development: none; confidence: low}
\end{GSitem}
\begin{GSitem}{1974-351.8}{E}{Finite-dimensional symmetric-power representations}
\GSfield{Kind; source}{construction; \GSpage{P000581, homogeneous-polynomial comparison}; I. M. Gelfand}
\GSfield{The notes}{The substitution \(f(x_1,x_2)\mapsto f(\alpha x_1+\beta x_2,\gamma x_1+\delta x_2)\) is applied to homogeneous polynomials of degree \(n\), with weight \(n\) said to approximate dimension.}
\GSfield{Statement}{The space of homogeneous polynomials of degree \(n\ge0\) in two variables has basis \(x_1^n,x_1^{n-1}x_2,\ldots,x_2^n\), and dimension \(n+1\). Linear substitution, using inverse or right-action convention consistently, gives the symmetric-power representation of \(SL_2\). Over \(\mathbb C\) it is irreducible; for the real noncompact group and \(n>0\) it is not a finite-dimensional unitary representation with a positive invariant norm.}
\GSfield{Reconstruction}{The exact dimension and the nonunitary character of the finite-dimensional nontrivial examples are supplied; weight and dimension are not identified as equal.}
\GSfield{Fate}{This is the elementary algebraic comparison preceding the unitary construction. The notebook does not cite a distinct theorem. Its role is contrasted with the infinite-dimensional unitary representations in \citet[\S1.1]{vershikgelfandgraev1973functions}; the dimension \(n+1\) follows directly by counting the displayed monomial basis.}
\GSfield{Present status}{As of November 2026, this is an established result in the stated scope; the qualification in the statement is essential.}
\GSfield{Assessment}{literal statement: correct with missing hypotheses; publication match: partial; verification: indirect; development: none; confidence: high}
\end{GSitem}
\begin{GSitem}{1974-351.9}{E}{The endpoint invariant functional and its trivial quotient}
\GSfield{Kind; source}{construction; \GSpage{P000581, lambda-zero norm and invariant integral}}
\GSfield{The notes}{At \(\lambda=0\), the norm degenerates to \(|\int f(t)\,dt|^2\); the last display says that integration is invariant under the exponent−2 fractional-linear action.}
\GSfield{Statement}{On a suitable test space the limiting form is \(B_0(f,h)=(\int f)\overline{(\int h)}\). Its kernel is the codimension-one invariant subspace of zero-integral functions. The quotient is the trivial one-dimensional representation, since change of variables gives \(\int T_0(g)f=\int f\). The zero-integral subspace obtains a different nondegenerate invariant form from the first derivative of the kernel at 0; it is not given a positive norm by the zero form itself.}
\GSfield{Reconstruction}{The invariant subspace and the one-dimensional quotient are distinguished; this removes the apparent claim that the action on the entire test space becomes trivial.}
\GSfield{Fate}{\citet[\S6.3, p.~128]{vershikgelfandgraev1973functions} explicitly uses this invariant functional and its kernel to build the logarithmic model. The elementary endpoint calculation links the two realizations on the notebook page.}
\GSfield{Present status}{As of November 2026, this is an established result in the stated scope; the qualification in the statement is essential.}
\GSfield{Assessment}{literal statement: correct with missing hypotheses; publication match: exact; verification: read the result; development: none; confidence: high}
\end{GSitem}
\begin{GSitem}{1974-351.10}{F}{Negative squares in an incompletely normalized continuation}
\GSfield{Kind; source}{remark; \GSpage{P000581, analytic-continuation diagram}}
\GSfield{The notes}{A line marked0,1,2,3,4 has a crossed annotation between1 and 2, then one negative square between2 and 3 and two between3 and 4.}
\GSfield{Statement}{A distributional analytic continuation of a kernel form requires its scalar normalization, domain, treatment at poles, and any moment conditions or quotient. Those data are absent here. The drawing therefore does not determine a signature theorem. In the standard normalized circle model the supplementary series also has a positive realization for \(1<\lambda<2\), equivalent to the parameter \(2-\lambda\).}
\GSfield{Fate}{The latter positive continuation is stated in \citet[\S1.1, p.~90, footnote1]{vershikgelfandgraev1973functions}. It explains why the crossed1--2 label should not be restored. No verified source was found for the exact subsequent negative-square counts under the unspecified notebook normalization.}
\GSfield{Present status}{As of November 2026, the positive complementary-series continuation is known, but the exact indefinite form behind the sketch is unidentifiable. The numerical signature labels are preserved as source evidence without being promoted to a theorem.}
\GSfield{Assessment}{literal statement: not determinable; publication match: partial; verification: unresolved lead; development: none; confidence: low}
\end{GSitem}
\begin{GSitem}{1974-351.11}{E}{Dyadic refinement of the step-current group}
\GSfield{Kind; source}{construction; \GSpage{P000582, successive interval diagrams}}
\GSfield{The notes}{The constant subgroup is \(G\), the two-half subgroup is \(G^2\), then \(G^4\), and so on, with representations on tensor products of \(H_{\lambda/2^n}\).}
\GSfield{Statement}{For the dyadic partition \(\mathcal P_n\) of the interval and \(m=\lambda dx\), \(0<\lambda<1\), each cell has measure \(\lambda/2^n\). Diagonal duplication embeds \(G^{2^n}\) into \(G^{2^{n+1}}\). The associated outer tensor product acts in \(\mathcal H_n=\bigotimes_{j=1}^{2^n}H_{\lambda/2^n}\). Dyadic step functions approximate every continuous current uniformly, but an arbitrary step function is not itself a continuous current.}
\GSfield{Reconstruction}{The measure scaling and diagonal embeddings are supplied; the inductive union is not identified set-theoretically with the continuous-current group.}
\GSfield{Fate}{This is the dyadic specialization of the directed Borel-partition construction in \citet[\S2.1--2]{vershikgelfandgraev1973functions}. The distinction between the step group and the continuous group is explicit in that paper.}
\GSfield{Present status}{As of November 2026, this is an established result in the stated scope; the qualification in the statement is essential.}
\GSfield{Assessment}{literal statement: correct with missing hypotheses; publication match: exact; verification: read the result; development: none; confidence: high}
\end{GSitem}
\begin{GSitem}{1974-351.12}{E}{Multiplicity one under diagonal restriction}
\GSfield{Kind; source}{theorem; \GSpage{P000582, supplementary summand once}}
\GSfield{The notes}{The notes state \(H_{\lambda/2}\otimes H_{\lambda/2}=H_\lambda\oplus\text{rest}\), with \(H_\lambda\) occurring once.}
\GSfield{Statement}{If \(a,b>0\) and \(a+b<1\), the restriction of \(T_a\otimes T_b\) from \(G\times G\) to the diagonal \(G\) contains \(T_{a+b}\) with multiplicity one. The remaining decomposition contains only principal continuous and discrete series. Consequently the required intertwining isometry is unique up to a scalar of absolute value1.}
\GSfield{Reconstruction}{Restriction to the diagonal subgroup is essential; the tensor product remains irreducible as an outer representation of the full product group.}
\GSfield{Fate}{Exact statement in \citet[Theorem 1.4]{vershikgelfandgraev1973functions}, which attributes the tensor-product decomposition to Pukanszky\textquotesingle s 1961 work. The 1973 paper is the directly read primary source for this use; the original 1961 proof was not checked here.}
\GSfield{Present status}{As of November 2026, this is an established result in the stated scope; the qualification in the statement is essential.}
\GSfield{Assessment}{literal statement: correct with missing hypotheses; publication match: exact; verification: read the result; development: none; confidence: high}
\end{GSitem}
\begin{GSitem}{1974-351.13}{E}{The isometric diagonal delta embedding}
\GSfield{Kind; source}{construction; \GSpage{P000582, delta-function embedding}}
\GSfield{The notes}{The intertwiner is written \(f(t)\mapsto f(t)\delta(t-s)\), labelled a unitary embedding.}
\GSfield{Statement}{For \(a,b>0\), \(a+b<1\), the distribution \(f(t)\delta(t-s)\) belongs to the completed kernel tensor space and defines an isometry \(H_{a+b}\hookrightarrow H_a\otimes H_b\). Pairing two such diagonal distributions multiplies the kernels to give \(|t-s|^{-(a+b)}\). This is an isometric injection, not a surjective unitary isomorphism and not a claim that delta is an ordinary \(L^2\) function.}
\GSfield{Fate}{\citet[Lemmas1.7--1.8 and Theorems1.6--1.7]{vershikgelfandgraev1973functions} gives the precise distributional interpretation and coherent multiple embeddings. The explicit delta formula removes the otherwise arbitrary phase choice under refinement.}
\GSfield{Present status}{As of November 2026, this is an established result in the stated scope; the qualification in the statement is essential.}
\GSfield{Assessment}{literal statement: correct with missing hypotheses; publication match: exact; verification: read the result; development: none; confidence: high}
\end{GSitem}
\begin{GSitem}{1974-351.14}{E}{An irreducible continuous tensor-product representation}
\GSfield{Kind; source}{theorem; \GSpage{P000582, completed limit and extension}}
\GSfield{The notes}{The iterated tensor spaces are nested, their inductive limit is completed, and the step-function representation is said to extend to continuous functions.}
\GSfield{Statement}{Use compatible isometric intertwiners satisfying composition under refinement, then complete \(\bigcup_n\mathcal H_n\). This produces a unitary step-current representation. Matrix-coefficient continuity in the uniform topology permits a unique strongly continuous extension to continuous currents. With the finite positive measure used here, the multiplicity-one construction yields irreducibility; irreducibility does not follow merely from the existence of an arbitrary inductive limit.}
\GSfield{Reconstruction}{Compatibility, the topology in which continuity is checked, and the extra irreducibility argument are made explicit.}
\GSfield{Fate}{\citet[Theorems2.1--2.2, Lemma 2.2 and Theorem 3.3]{vershikgelfandgraev1973functions} proves compatibility, independence of phase choices, irreducibility and extension. These are already 1973 results, not new 1974 announcements.}
\GSfield{Present status}{As of November 2026, this is an established result in the stated scope; the qualification in the statement is essential.}
\GSfield{Assessment}{literal statement: correct with missing hypotheses; publication match: exact; verification: read the result; development: none; confidence: high; problem \hyperref[prob:GP-1974-14]{GP-1974-14}}
\end{GSitem}
\begin{GSitem}{1974-351.15}{E}{The centered logarithmic Hilbert space}
\GSfield{Kind; source}{construction; \GSpage{P000582-P000583, tilde and logarithmic form}}
\GSfield{The notes}{The exponential scalar product uses \(-\log|t_1-t_2|\). The following page writes \(\widetilde f(x,t)=f(x,t)-1/[\pi(1+x^2)]\), although \(\int f(x,t)\,dt=1\).}
\GSfield{Statement}{Set \(\rho(t)=1/[\pi(1+t^2)]\) and \(h_x(t)=f(x,t)-\rho(t)\), so \(\int h_x=0\). On real zero-integral test functions define \(B(h,k)=-\iint\log|t-s|h(t)k(s)\,dt\,ds\). Its completion is a real Hilbert space; for complex functions insert conjugation in the second argument. The form arises by differentiating \(|t-s|^{-\lambda}\) at 0 after the constant term has vanished.}
\GSfield{Reconstruction}{The real-label version matches the paper; sesquilinear complexification is stated separately. Smoothness and decay may first be imposed on a dense test space before Hilbert completion.}
\GSfield{Fate}{Exact line-model construction in \citet[\S6.3, p.~128]{vershikgelfandgraev1973functions}. The variable in the Cauchy density is \(t\), not the interval coordinate \(x\). With the written \(x\), the subtracted function is constant in \(t\) and is not integrable, so the centering cannot work.}
\GSfield{Present status}{As of November 2026, this is an established result in the stated scope; the qualification in the statement is essential.}
\GSfield{Assessment}{literal statement: false as written; publication match: exact; verification: read the result; development: annotated; confidence: high}
\end{GSitem}
\begin{GSitem}{1974-351.16}{E}{Formal coherent vectors rather than ordinary functions}
\GSfield{Kind; source}{construction; \GSpage{P000582, formal combinations and exponential product}}
\GSfield{The notes}{Finite formal linear combinations of the labels \(f(x,t)\) are introduced; the notes emphasize that the symbol labelled by zero need not be the zero vector.}
\GSfield{Statement}{Let \(\mathcal K=L^2(X,m;H_0)\), where \(H_0\) is the centered logarithmic Hilbert space. For real \(h\in\mathcal K\), introduce a formal vector \(e(h)\) and define \(\langle e(h),e(k)\rangle=\exp\langle h,k\rangle_{\mathcal K}\). Extend sesquilinearly, quotient any null vectors and complete. The exponential kernel is positive by its tensor-power series. In particular \(e(0)\) has norm 1, although the label 0 is the zero vector of \(\mathcal K\).}
\GSfield{Reconstruction}{The label map is nonlinear. Distinct variables \(t_1,t_2\) are required in the two kernel factors; the measure \(m\) records the parameter of the representation.}
\GSfield{Fate}{\citet[\S6.1, Lemma 6.2, and \S6.3]{vershikgelfandgraev1973functions} gives this construction and its logarithmic integral form. It also recurs from the exponential-Hilbert-space discussion indexed as 1969-180.6; no identity with every earlier model is assumed.}
\GSfield{Present status}{As of November 2026, this is an established result in the stated scope; the qualification in the statement is essential.}
\GSfield{Assessment}{literal statement: correct with missing hypotheses; publication match: exact; verification: read the result; development: none; confidence: high}
\end{GSitem}
\begin{GSitem}{1974-351.17}{E}{The scalar cocycle restoring unitarity}
\GSfield{Kind; source}{construction; \GSpage{P000583, representation display}}
\GSfield{The notes}{The endpoint fractional-linear action is multiplied by an unspecified \(\lambda(g,f)\).}
\GSfield{Statement}{In the real-label model put \(b_g=T_0(g)\rho-\rho\), \(A_gh=T_0(g)h\), and let \(Q\) be the inner product after integration over \(X\). Then an explicit normalized action is \[U_ge(h)=\exp\{-Q(A_gh,b_g)-\tfrac12Q(b_g,b_g)\}\,e(A_gh+b_g).\] The cocycle identity \(b_{g_1g_2}=b_{g_1}+A_{g_1}b_{g_2}\) makes this an ordinary unitary representation for the real model. Omitting the scalar generally fails to preserve the exponential kernel.}
\GSfield{Reconstruction}{The explicit multiplier and the ordinary-versus-projective distinction are supplied from the matching construction; the notebook records only the unnamed scalar factor.}
\GSfield{Fate}{\citet[\S6.1, Lemmas6.1--6.3; \S6.3 and Theorem 6.2]{vershikgelfandgraev1973functions} provides the multiplier and proves equivalence with the preceding irreducible construction. Complex holomorphic restrictions can instead produce projective representations; they are not what this unqualified real-line display specifies.}
\GSfield{Present status}{As of November 2026, this is an established result in the stated scope; the qualification in the statement is essential.}
\GSfield{Assessment}{literal statement: correct with missing hypotheses; publication match: exact; verification: read the result; development: none; confidence: high}
\end{GSitem}

\GStalk{352}{}{—}{P000583-P000585}{Year 1975: catalogue heading}
This MCCME record is a year heading, not a separate mathematical talk. Its catalogue range overlaps the end of Graev\textquotesingle s record 351 and the beginning of the dated 17 February session. No mathematical occurrence is assigned merely to the heading.


\GStalk{353}{1975-02-17}{V. N. Tulovskii}{P000583-P000586}{Wavefront sets, products, and Schwartz kernels}
The dated heading begins on P000584 and explicitly says \ru{ИОФ --- I и II}. The notes review microlocal regularity, multiplication of distributions, and propagation through kernels; the composition formula continues onto P000585, before Sinai\textquotesingle s heading. This revisits 1971-234 and especially 1973-329. Gelfand intervenes with slicing and a proposed extension from linear to polynomial integration. The mathematical backbone is already in \citet[\S2.5]{hormander1971fourier}.

\begin{GSitem}{1974-353.1}{E}{The smooth wavefront set}
\GSfield{Kind; source}{definition; \GSpage{P000584}}
\GSfield{The notes}{\((x_0,\xi_0)\notin\mathrm{WF}(u)\) when a cutoff \(\varphi(x_0)\ne0\) makes \(\widetilde{u\varphi}\) rapidly decreasing in an angular neighbourhood of \(\xi_0\).}
\GSfield{Statement}{We use \(\widehat u(\xi)=\int e^{-ix\cdot\xi}u(x)\,dx\) throughout the microlocal reconstruction. For an open \(\Omega\subset\mathbb R^n\) and \(u\in\mathcal D'(\Omega)\), a nonzero covector \((x_0,\xi_0)\) is outside \(\mathrm{WF}(u)\) iff some \(\varphi\in C_c^\infty(\Omega)\), equal to one near \(x_0\), and an open cone \(V\ni\xi_0\) satisfy \(\sup_{\xi\in V}(1+|\xi|)^N|\widehat{\varphi u}(\xi)|<\infty\) for every integer \(N\ge0\).}
\GSfield{Fate}{Published in \citet[Definition 2.5.2 and pp.~121--122]{hormander1971fourier} in 1971; the equivalent Fourier description is explicit on p.~122. Earlier seminar occurrences are 1971-234.1 and 1973-329.1.}
\GSfield{Present status}{As of November 2026: established in the stated scope; the notebook is an expository occurrence, not a new open problem.}
\GSfield{Assessment}{literal statement: correct; publication match: exact; verification: read the result; development: none; confidence: high}
\end{GSitem}
\begin{GSitem}{1974-353.2}{E}{Gelfand’s hyperplane slicing interpretation}
\GSfield{Kind; source}{remark; \GSpage{P000584}; I. M. Gelfand}
\GSfield{The notes}{Gelfand relates the smoothness of a one-variable integral over parallel hyperplanes to membership of the normal covector in the wavefront.}
\GSfield{Statement}{For a compactly supported localized distribution \(v=\chi u\), the pushforward under \(x\mapsto x\cdot\theta\) has one-dimensional Fourier transform \(\widehat v(t\theta)\). Absence from \(\mathrm{WF}(u)\) requires such rapid decay uniformly for \(\theta\) in an angular neighbourhood and with localization near \(x_0\). A single unlocalized slice integral cannot detect every covector, because cancellations are possible.}
\GSfield{Reconstruction}{The cutoff, angular uniformity, and warning about cancellation make precise the one-line analogy; they are editorial qualifications.}
\GSfield{Fate}{The precise directional statement follows from the Fourier characterization in \citet[p.~122]{hormander1971fourier}; this repeats the explanatory intervention 1973-329.3. It is an interpretation of an old definition, not an asserted new inversion theorem.}
\GSfield{Present status}{As of November 2026: established in the stated scope; the notebook is an expository occurrence, not a new open problem.}
\GSfield{Assessment}{literal statement: correct with missing hypotheses; publication match: partial; verification: read the result; development: none; confidence: high}
\end{GSitem}
\begin{GSitem}{1974-353.3}{E}{One-sided analytic continuation and directional regularity}
\GSfield{Kind; source}{claim; \GSpage{P000584}}
\GSfield{The notes}{For \(\xi_0>0\), the notes give \(u=\varphi+u_1\) locally, with \(\varphi\in C_c^\infty\) and \(u_1\) continuing analytically into the upper half-plane.}
\GSfield{Statement}{A holomorphic function in an upper half-disc with at most polynomial growth towards the boundary has a distributional boundary value. With \(\widehat u(\xi)=\int e^{-ix\xi}u(x)\,dx\), its smooth wavefront has only positive covectors on the boundary interval, and therefore excludes negative covectors; adding a smooth function does not change this. Reversing the Fourier convention reverses the sign. The notebook specifies neither that convention nor the boundary-value growth condition, so its positive-half-ray wording is not by itself a convention-independent contradiction.}
\GSfield{Reconstruction}{We give a sufficient distributional-boundary-value statement, with polynomial growth and the same negative-exponent Fourier transform as the rest of the chapter. The notebook's upper-half-plane and positive-excluded-ray wording can agree only with the opposite transform convention; no such convention is recorded.}
\GSfield{Fate}{The item is an expository boundary-value interpretation of the wavefront framework in \citet[Section 2.5]{hormander1971fourier}, not an exact theorem quotation established from that paper. An independent explicit sign check is \citet[Appendix A, equation (A.4) and following calculation]{bohrlefeuvrepaternain2024invariant}: \(\mathrm{WF}((x+i0)^{-1})=\{(0,\xi):\xi>0\}\). The modern check does not establish a new historical priority or verify an unqualified converse in the notebook.}
\GSfield{Present status}{As of November 2026: established in the stated scope; the notebook is an expository occurrence, not a new open problem.}
\GSfield{Assessment}{literal statement: not determinable; publication match: partial; verification: read the result; development: none; confidence: medium}
\end{GSitem}
\begin{GSitem}{1974-353.4}{E}{Projection onto singular support}
\GSfield{Kind; source}{theorem; \GSpage{P000584}}
\GSfield{The notes}{\(\pi(\mathrm{WF}(u))=\operatorname{sing\,supp}u\), with \(\pi:T^*\Omega\to\Omega\).}
\GSfield{Statement}{For every distribution on a smooth manifold, the base projection of its smooth wavefront set is exactly its singular support.}
\GSfield{Fate}{\citet[Theorem 2.5.3]{hormander1971fourier}, published 1971, is an exact match. The earlier occurrence is 1973-329.4.}
\GSfield{Present status}{As of November 2026: established in the stated scope; the notebook is an expository occurrence, not a new open problem.}
\GSfield{Assessment}{literal statement: correct; publication match: exact; verification: read the result; development: none; confidence: high}
\end{GSitem}
\begin{GSitem}{1974-353.5}{E}{Frequency projection and rapid Fourier decay}
\GSfield{Kind; source}{theorem; \GSpage{P000584}}
\GSfield{The notes}{\(\pi'\mathrm{WF}(u)\) is called the smallest closed cone outside which the Fourier transform decreases rapidly.}
\GSfield{Statement}{For \(u\in\mathcal E'(\mathbb R^n)\), let \(K=\{\xi\ne0:\exists x,\ (x,\xi)\in\mathrm{WF}(u)\}\). Then \(K\) is closed in \(\mathbb R^n\setminus0\), and \(\widehat u\) decreases faster than every power in each closed cone disjoint from \(K\). Conversely, this property determines the smallest such closed cone. For a general local distribution apply the statement to compactly supported localizations.}
\GSfield{Reconstruction}{Compact support is supplied from the cited theorem.}
\GSfield{Fate}{Exact compact-support version: \citet[Proposition 2.5.5]{hormander1971fourier}. Compact support is missing in the notebook; without a suitable global hypothesis even an ordinary Fourier transform need not be defined.}
\GSfield{Present status}{As of November 2026: established in the stated scope; the notebook is an expository occurrence, not a new open problem.}
\GSfield{Assessment}{literal statement: correct with missing hypotheses; publication match: exact; verification: read the result; development: none; confidence: high}
\end{GSitem}
\begin{GSitem}{1974-353.6}{E}{Multiplication under the no-opposite-covectors condition}
\GSfield{Kind; source}{construction; \GSpage{P000584}}
\GSfield{The notes}{For \(\mathcal D'_\Gamma=\{u:\mathrm{WF}(u)\subset\Gamma\}\), and \(\Gamma_1+\Gamma_2\) avoiding zero, the notes give a bilinear map into \(\mathcal D'_{(\Gamma_1+\Gamma_2)\cup\Gamma_1\cup\Gamma_2}\).}
\GSfield{Statement}{If \(u,v\in\mathcal D'(\Omega)\) and there is no \((x,\xi)\in\mathrm{WF}(u)\) with \((x,-\xi)\in\mathrm{WF}(v)\), their product is canonically defined by pullback of \(u\otimes v\) to the diagonal. Its wavefront is contained in \(\mathrm{WF}(u)\cup\mathrm{WF}(v)\cup\{(x,\xi+\eta):(x,\xi)\in\mathrm{WF}(u),(x,\eta)\in\mathrm{WF}(v)\}\).}
\GSfield{Reconstruction}{No continuity claim for an unspecified topology on the open-cone spaces is added.}
\GSfield{Fate}{\citet[Theorem 2.5.10]{hormander1971fourier} proves the closed-cone version, predating the talk. The occurrence 1973-329.11 has the same criterion. The notebook calls its containing cones open; the pointwise noncancellation criterion remains the decisive condition.}
\GSfield{Present status}{As of November 2026: established in the stated scope; the notebook is an expository occurrence, not a new open problem.}
\GSfield{Assessment}{literal statement: correct with missing hypotheses; publication match: exact; verification: read the result; development: none; confidence: high}
\end{GSitem}
\begin{GSitem}{1974-353.7}{E}{A distribution kernel defines an operator}
\GSfield{Kind; source}{construction; \GSpage{P000584}}
\GSfield{The notes}{\(A\in\mathcal D'(\Omega_x\times\Omega_y)\) is paired with \(u(y)v(x)\), defining \(A:C_c^\infty(\Omega_y)\to\mathcal D'(\Omega_x)\).}
\GSfield{Statement}{A distribution \(K\) on \(X\times Y\), with densities understood, defines a continuous linear operator by \(\langle Au,v\rangle=\langle K,v\otimes u\rangle\). Conversely every continuous operator \(C_c^\infty(Y)\to\mathcal D'(X)\) has a unique Schwartz kernel.}
\GSfield{Reconstruction}{Continuity and the density convention are stated explicitly.}
\GSfield{Fate}{This is the established Schwartz kernel formalism used in \citet[\S2.5, before Theorem 2.5.12]{hormander1971fourier}; earlier occurrence 1973-329.5.}
\GSfield{Present status}{As of November 2026: established in the stated scope; the notebook is an expository occurrence, not a new open problem.}
\GSfield{Assessment}{literal statement: correct; publication match: exact; verification: read the result; development: none; confidence: high}
\end{GSitem}
\begin{GSitem}{1974-353.8}{E}{Kernel extensions with both axis wavefronts absent}
\GSfield{Kind; source}{theorem; \GSpage{P000584-P000585}}
\GSfield{The notes}{If \(\mathrm{WF}(A)\) has no points with \(\xi=0\) or \(\eta=0\), then \(A:C_c^\infty(Y)\to C^\infty(X)\) and \(A:\mathcal E'(Y)\to\mathcal D'(X)\).}
\GSfield{Statement}{For \(K\in\mathcal D'(X\times Y)\), assume \(\mathrm{WF}(K)\) avoids both sets \(\{(x,\xi;y,0):\xi\ne0\}\) and \(\{(x,0;y,\eta):\eta\ne0\}\). Kernel action then maps tests to smooth functions and extends canonically and continuously to compactly supported distributions. Proper support is needed only for corresponding noncompact-input or compact-output assertions.}
\GSfield{Fate}{\citet[Theorems 2.5.12--2.5.14]{hormander1971fourier} establish these operations. This repeats 1973-329.6.}
\GSfield{Present status}{As of November 2026: established in the stated scope; the notebook is an expository occurrence, not a new open problem.}
\GSfield{Assessment}{literal statement: correct; publication match: exact; verification: read the result; development: none; confidence: high}
\end{GSitem}
\begin{GSitem}{1974-353.9}{E}{The primed relation and propagation through a kernel}
\GSfield{Kind; source}{theorem; \GSpage{P000585}}
\GSfield{The notes}{\(\mathrm{WF}'(A)=\{(x,\xi;y,-\eta):(x,\xi;y,\eta)\in\mathrm{WF}(A)\}\), and \(\mathrm{WF}(Au)\subset\mathrm{WF}'(A)\circ\mathrm{WF}(u)\).}
\GSfield{Statement}{Under the no-axis assumptions of the preceding item, for every \(u\in\mathcal E'(Y)\), \(\mathrm{WF}(Au)\subset\{(x,\xi):\exists(y,\eta)\in\mathrm{WF}(u),\ (x,\xi;y,\eta)\in\mathrm{WF}'(K)\}\). The minus sign belongs to the input covector of the kernel.}
\GSfield{Fate}{Exact statement in \citet[Theorem 2.5.14]{hormander1971fourier}; compare 1973-329.7--8. The preceding assumptions must remain attached to this shorthand formula.}
\GSfield{Present status}{As of November 2026: established in the stated scope; the notebook is an expository occurrence, not a new open problem.}
\GSfield{Assessment}{literal statement: correct; publication match: exact; verification: read the result; development: none; confidence: high}
\end{GSitem}
\begin{GSitem}{1974-353.10}{F}{Gelfand proposes polynomial as well as linear integration}
\GSfield{Kind; source}{remark; \GSpage{P000585}; I. M. Gelfand}
\GSfield{The notes}{\ru{Идея Г.} is followed by integrating linear functions and polynomials; a curved outline is cut by parallel segments.}
\GSfield{Statement}{The evidence supports a proposal to replace linear level sets by polynomial ones in an integral-geometric construction. It does not identify the family of polynomials, the integration measure, a target space, or an inversion or range question. Thus it determines no unique theorem or problem.}
\GSfield{Fate}{Generalized Radon transforms are thematically related to the canonical-relation calculus of \citet[\S4.1]{hormander1971fourier}, but that relation does not identify this particular proposal. The scan does not justify reading the small sketch as a specified polynomial transform.}
\GSfield{Present status}{As of November 2026: the precise intended construction is not recoverable from this entry; no claim of an undeveloped idea is made.}
\GSfield{Assessment}{literal statement: not determinable; publication match: thematic only; verification: unresolved lead; development: none; confidence: low}
\end{GSitem}
\begin{GSitem}{1974-353.11}{E}{Composition of wavefront relations}
\GSfield{Kind; source}{theorem; \GSpage{P000585}}
\GSfield{The notes}{For \(A:\mathcal E'(Y)\to\mathcal D'(X)\), \(B:\mathcal E'(Z)\to\mathcal E'(Y)\), the notes write \(\mathrm{WF}'(A\circ B)\subset\mathrm{WF}'(A)\circ\mathrm{WF}'(B)\).}
\GSfield{Statement}{For properly supported distribution kernels whose wavefronts avoid both coordinate axes, composition is defined and its primed wavefront lies in the composite canonical relation. More generally one needs the wavefront transversality condition at the intermediate variable and must include the horizontal and vertical error sets when these are nonempty.}
\GSfield{Reconstruction}{The earlier no-axis hypothesis is retained and support conditions are made explicit.}
\GSfield{Fate}{\citet[Theorem 2.5.15]{hormander1971fourier} gives the general inclusion, with the simpler formula after removal of the two axis terms; compare 1973-329.9.}
\GSfield{Present status}{As of November 2026: established in the stated scope; the notebook is an expository occurrence, not a new open problem.}
\GSfield{Assessment}{literal statement: correct with missing hypotheses; publication match: exact; verification: read the result; development: none; confidence: high}
\end{GSitem}

\GStalk{354}{1975-02-17}{Ya. G. Sinai}{P000584-P000587}{Configuration measures and finite-range Gibbs interactions}
Sinai's contribution starts after Tulovskii's composition formula on P000585; P000586 contains only the product-measure continuation. Gelfand participates in the configuration-space formulation. The following dated session, on P000587, belongs to record 355. This is preparatory exposition for the later phase-diagram theorem, connected to the earlier Gibbs discussions 1970-195 and 1970-200.

\begin{GSitem}{1974-354.1}{E}{Locally finite configurations and their measurable sets}
\GSfield{Kind; source}{definition; \GSpage{P000585}; I. M. Gelfand}
\GSfield{The notes}{Configurations are countable subsets \(X\subset\mathbb R^n\) without accumulation points; count events \(\mathrm C_{K,j}\) generate the measurable sets. The displayed set mistakenly binds \(K\), not \(X\).}
\GSfield{Statement}{Let \(\Gamma\) be the locally finite simple counting measures on \(\mathbb R^n\), identified with their point sets. Equip \(\Gamma\) with the sigma-algebra generated by \(N_B(X)=|X\cap B|\) for bounded Borel \(B\). For a fixed compact \(K\), \(\mathrm C_{K,j}=\{X\in\Gamma:N_K(X)=j\}\).}
\GSfield{Fate}{The standard measurable configuration formalism appears in \citet[Section 5]{ruelle1970superstable}; the notation in this talk is expository. Simplicity must be changed to integer-valued counting measures when the intensity has atoms.}
\GSfield{Present status}{As of November 2026: established in the stated scope; the notebook is an expository occurrence, not a new open problem.}
\GSfield{Assessment}{literal statement: false as written; publication match: exact; verification: read the result; development: none; confidence: high}
\end{GSitem}
\begin{GSitem}{1974-354.2}{E}{The Poisson reference measure}
\GSfield{Kind; source}{construction; \GSpage{P000585}}
\GSfield{The notes}{\(\rho(\mathrm C_{K,j})=m(K)^j e^{-m(K)}/j!\), for a locally compact measured space.}
\GSfield{Statement}{For a locally finite diffuse Radon measure \(m\) on a second-countable locally compact space, the Poisson process has independent counts on disjoint bounded Borel sets, and \(N_B\) is Poisson with mean \(m(B)\). On a region \(K\), its symmetric \(j\)-particle measure is \(e^{-m(K)}m^{\otimes j}/j!\), pushed to unordered configurations.}
\GSfield{Reconstruction}{Local finiteness, diffuseness, a standard measurable space, and joint count independence are explicit additions.}
\GSfield{Fate}{This is the usual Poisson reference in \citet[Section 5]{ruelle1970superstable}. Independence of disjoint counts completes the displayed one-set marginal formula; no new probability measure is proposed here.}
\GSfield{Present status}{As of November 2026: established in the stated scope; the notebook is an expository occurrence, not a new open problem.}
\GSfield{Assessment}{literal statement: correct with missing hypotheses; publication match: exact; verification: read the result; development: none; confidence: high}
\end{GSitem}
\begin{GSitem}{1974-354.3}{E}{Finite interaction range gives conditional locality}
\GSfield{Kind; source}{remark; \GSpage{P000585}}
\GSfield{The notes}{\ru{Взаимод. конечно}: \ru{независимость}; the sketch separates an inner region from the exterior by a collar of width one.}
\GSfield{Statement}{For a finite-range Gibbs specification, the conditional law inside \(O\), given its entire exterior configuration, depends only on exterior particles within the interaction range of \(O\). Consequently the interior and the farther exterior are conditionally independent given the collar. Unconditional independence does not follow.}
\GSfield{Fate}{The precise conditional-distribution framework was already established by \citet[Section 1]{dobrushin1968conditional}. The next talk supplies the missing conditional-law formula explicitly (1974-355.3--4).}
\GSfield{Present status}{As of November 2026: established in the stated scope; the notebook is an expository occurrence, not a new open problem.}
\GSfield{Assessment}{literal statement: correct with missing hypotheses; publication match: partial; verification: read the result; development: none; confidence: high}
\end{GSitem}
\begin{GSitem}{1974-354.4}{E}{Finite-range many-particle interaction potentials}
\GSfield{Kind; source}{definition; \GSpage{P000585}}
\GSfield{The notes}{\(U(x_1,\ldots,x_j)=0\) if the diameter is at least one. A proposed infinite hard-core value is crossed out.}
\GSfield{Statement}{A symmetric family \(U_j\) is of range one if every \(U_j(x_1,\ldots,x_j)\) vanishes when the set of its arguments has diameter at least one. For a finite configuration, its energy sums \(U_j\) over unordered \(j\)-element subsets. Finite range and finite interaction order are separate conditions.}
\GSfield{Fate}{Finite-range interactions are standard in the conditional specification approach of \citet[Section 1]{dobrushin1968conditional}. The deleted line cannot be used as an assumption; the more explicit pair-potential model appears on P000587.}
\GSfield{Present status}{As of November 2026: established in the stated scope; the notebook is an expository occurrence, not a new open problem.}
\GSfield{Assessment}{literal statement: correct with missing hypotheses; publication match: partial; verification: read the result; development: none; confidence: high}
\end{GSitem}
\begin{GSitem}{1974-354.5}{F}{The unexplained termination of the potential sequence}
\GSfield{Kind; source}{claim; \GSpage{P000585}}
\GSfield{The notes}{\ru{Посл. обрыв., т.к. становится \(\equiv+\infty\)} follows the partially cancelled hard-core clause.}
\GSfield{Statement}{The fragment does not identify which sequence terminates, or whether it means a maximum admissible particle number, a finite list of potentials, or an energy tending to infinity. A hard core bounds the number of particles in a bounded region, but finite interaction range alone does not bound interaction order.}
\GSfield{Fate}{Neither the crossed-out condition nor the following density resolves the intended assertion. Searches in the Dobrushin/Ruelle specification literature did not identify an exact theorem corresponding to this phrase.}
\GSfield{Present status}{As of November 2026: not identifiable; no open problem is extracted from the cancelled shorthand.}
\GSfield{Assessment}{literal statement: not determinable; publication match: none found; verification: unresolved lead; development: none; confidence: low}
\end{GSitem}
\begin{GSitem}{1974-354.6}{E}{Finite-volume Gibbs density and particle-number sectors}
\GSfield{Kind; source}{construction; \GSpage{P000585-P000586}}
\GSfield{The notes}{The notes write \(d(X)=e^{-\sum U_j}\), call it \(d(\mathrm{Gibbs})/d\mu\), and then use \(m^{\otimes n}\) on \(K^n\).}
\GSfield{Statement}{For a fixed exterior condition \(Y\), define
\[dG_K^Y=Z_K(Y)^{-1}e^{-H_K(X\mid Y)}d\Pi_{m,K}(X),\]
assuming \(0<Z_K(Y)<\infty\). Conditional on a specified particle number \(n\), one may instead normalize \(e^{-H_K}\) against the symmetric pushforward of \(m^{\otimes n}\). The grand-canonical version includes all \(n\), with \(1/n!\) factors.}
\GSfield{Reconstruction}{The two ensemble interpretations are separated rather than choosing one silently.}
\GSfield{Fate}{\citet[Section 5]{ruelle1970superstable} treats the normalized finite-volume measures and their limits. The notebook omits normalization and passes between labelled coordinates and unordered configurations without specifying the ensemble; P000587 restores the partition function.}
\GSfield{Present status}{As of November 2026: established in the stated scope; the notebook is an expository occurrence, not a new open problem.}
\GSfield{Assessment}{literal statement: correct with missing hypotheses; publication match: exact; verification: read the result; development: none; confidence: high}
\end{GSitem}

\GStalk{355}{1975-02-24}{Ya. G. Sinai}{P000585-P000591}{Gibbs specifications, uniqueness, and low-temperature phase diagrams}
The notes begin on P000587 (the overwritten day is tentatively read 24) and end on P000589. Continuum Gibbs kernels lead to existence and uniqueness, then to finite-spin lattice systems and the low-temperature theorem soon published by Pirogov and Sinai. The first part is older Dobrushin--Ruelle theory; the phase-diagram part is a prepublication occurrence, preceding the papers received on 21 April 1975. Gelfand asks about a \(p\)-adic analogue without specifying the underlying model.

\begin{GSitem}{1974-355.1}{E}{Poisson configurations on an infinite measured metric space}
\GSfield{Kind; source}{definition; \GSpage{P000587}}
\GSfield{The notes}{The Poisson count formula is repeated on \(\mathbb R^2\), then on \((M,m)\); \(m(M)=\infty\) is the interesting case, with unspecified homogeneity.}
\GSfield{Statement}{Under the standard local-finiteness and diffuseness assumptions of 1974-354.2, infinite total intensity is compatible with a locally finite process. The count in each relatively compact \(O\) has mean \(m(O)\); disjoint counts are independent. Homogeneity is an extra symmetry assumption and is unnecessary for the definition.}
\GSfield{Fate}{This repeats the standard reference process in \citet[Section 5]{ruelle1970superstable}; it does not assert a new existence theorem for arbitrary measurable metric spaces.}
\GSfield{Present status}{As of November 2026: established in the stated scope; the notebook is an expository occurrence, not a new open problem.}
\GSfield{Assessment}{literal statement: correct with missing hypotheses; publication match: partial; verification: read the result; development: none; confidence: high}
\end{GSitem}
\begin{GSitem}{1974-355.2}{E}{Hard-core finite-range pair energy}
\GSfield{Kind; source}{definition; \GSpage{P000587}}
\GSfield{The notes}{\(U_1\) is bounded, \(U_2\) symmetric, infinite at distances \(<\varepsilon\) and zero at distances \(\ge r_0\). The intermediate range is unspecified.}
\GSfield{Statement}{For a finite inside configuration \(X\) and exterior \(Y\), use \(H_O(X\mid Y)=\sum_{x\in X}U_1(x)+\sum_{\{x,x'\}\subset X}U_2(x,x')+\sum_{x\in X,y\in Y}U_2(x,y)\). Each inside pair is counted once. A lower bound and appropriate regularity on the annulus \(\varepsilon\le d(x,y)<r_0\) are needed for subsequent existence claims.}
\GSfield{Fate}{\citet[Sections 1 and 5]{ruelle1970superstable} gives stability and regularity conditions for such continuum energies. Hard core and finite range alone are not a complete theorem hypothesis on an arbitrary metric measure space.}
\GSfield{Present status}{As of November 2026: established in the stated scope; the notebook is an expository occurrence, not a new open problem.}
\GSfield{Assessment}{literal statement: correct with missing hypotheses; publication match: partial; verification: read the result; development: none; confidence: high}
\end{GSitem}
\begin{GSitem}{1974-355.3}{E}{The DLR conditional probability equation}
\GSfield{Kind; source}{definition; \GSpage{P000587}}
\GSfield{The notes}{The conditional density in \(O\) is the exponential of the one-body, inside-pair and inside--outside energies, divided by \(\Xi(Y;O)\).}
\GSfield{Statement}{A probability \(\mu\) is Gibbs for a specified energy if, for every bounded region \(O\), its regular conditional law given the outside configuration equals \(\gamma_O(dX\mid Y)=\Xi(Y;O)^{-1}e^{-H_O(X\mid Y)}\Pi_{m,O}(dX)\), for \(\mu\)-almost every \(Y\). Here \(\Xi\) is the integral of the numerator, finite and positive on admissible boundary conditions.}
\GSfield{Fate}{This is the continuum conditional Gibbs formalism already treated in \citet[Section 5]{ruelle1970superstable}; compare the finite-state formulation of \citet[Section 1]{dobrushin1968conditional}. Earlier seminar context: 1970-195.4 and 1970-200.4.}
\GSfield{Present status}{As of November 2026: established in the stated scope; the notebook is an expository occurrence, not a new open problem.}
\GSfield{Assessment}{literal statement: correct; publication match: exact; verification: read the result; development: none; confidence: high}
\end{GSitem}
\begin{GSitem}{1974-355.4}{E}{The boundary collar is sufficient exterior information}
\GSfield{Kind; source}{theorem; \GSpage{P000587}}
\GSfield{The notes}{Everything depends only on the points of \(Y\) in the \(r_0\)-neighbourhood of \(O\).}
\GSfield{Statement}{For the finite-range pair energy just defined, \(\gamma_O(\,\cdot\mid Y)\) depends only on \(Y\cap\{y\notin O:d(y,O)<r_0\}\). This is the spatial Markov property of the specification, not independence of the unconditioned field.}
\GSfield{Fate}{The assertion follows directly from the displayed energy and is part of the established specification formalism in \citet[Section 1]{dobrushin1968conditional}. It resolves the shorthand in 1974-354.3.}
\GSfield{Present status}{As of November 2026: established in the stated scope; the notebook is an expository occurrence, not a new open problem.}
\GSfield{Assessment}{literal statement: correct; publication match: partial; verification: read the result; development: none; confidence: high}
\end{GSitem}
\begin{GSitem}{1974-355.5}{E}{Gelfand interprets conditioning as disintegration}
\GSfield{Kind; source}{remark; \GSpage{P000587}; I. M. Gelfand}
\GSfield{The notes}{Gelfand invokes a measurable partition and a general Fubini theorem; the measure on exterior configurations is induced by inverse images.}
\GSfield{Statement}{On standard Borel configuration spaces, the restriction map \(r:X\mapsto X\cap O^c\) has a pushforward probability \(r_*\mu\) and a conditional probability kernel \(\mu_Y\), unique \(r_*\mu\)-almost everywhere, satisfying \(\int F\,d\mu=\int(\int F\,d\mu_Y)\,d(r_*\mu)(Y)\). The kernel is supported on the corresponding fibres almost everywhere.}
\GSfield{Fate}{\citet[Section 3.1, pp.~125--129]{rokhlin1949measure}, received 15 July 1947 and published in 1949, proves the canonical conditional-measure construction and essential uniqueness for measurable partitions of Lebesgue spaces. The cited section was read directly. Gelfand's analogy is an expository one; Dobrushin returns to the precise measure-theoretic formulation in record 358.}
\GSfield{Present status}{As of November 2026: established in the stated scope; the notebook is an expository occurrence, not a new open problem.}
\GSfield{Assessment}{literal statement: correct with missing hypotheses; publication match: partial; verification: read the result; development: none; confidence: high}
\end{GSitem}
\begin{GSitem}{1974-355.6}{E}{Consistency for nested regions}
\GSfield{Kind; source}{theorem; \GSpage{P000587-P000588}}
\GSfield{The notes}{A diagram with \(O_2\subset O_1\) accompanies integration over all ways to fill in the smaller missing configuration.}
\GSfield{Statement}{For a Gibbs specification and bounded \(O_2\subset O_1\), \(\gamma_{O_1}\gamma_{O_2}=\gamma_{O_1}\), where kernels act successively on bounded measurable observables. This expresses the tower property and is the compatibility condition required before seeking a global measure.}
\GSfield{Reconstruction}{The kernel identity is a precise reconstruction of the nested-region explanation.}
\GSfield{Fate}{\citet[Section 1]{dobrushin1968conditional} formulates compatible conditional distributions before its existence theorem. The diagram is interpreted using the normalized kernel on the preceding page.}
\GSfield{Present status}{As of November 2026: established in the stated scope; the notebook is an expository occurrence, not a new open problem.}
\GSfield{Assessment}{literal statement: correct with missing hypotheses; publication match: partial; verification: read the result; development: none; confidence: high}
\end{GSitem}
\begin{GSitem}{1974-355.7}{E}{Existence from finite-volume limits}
\GSfield{Kind; source}{theorem; \GSpage{P000588}}
\GSfield{The notes}{A Gibbs measure exists; exhaust \(M\) by compact \(O_k\), construct measures there and choose a convergent subsequence. The notes call each local configuration space compact.}
\GSfield{Statement}{A verified sufficient continuum theorem is this: for a translation-invariant superstable, lower-regular pair potential with \(\int_{\mathbb R^d}|1-e^{-\beta U(x)}|\,dx<\infty\), the finite-volume free-boundary measures have subsequential local limits giving tempered Gibbs states. In the hard-core model, local particle numbers are bounded on compact closures, but passage through boundaries and discontinuous potentials still needs justification.}
\GSfield{Reconstruction}{Ruelle’s explicit hypotheses replace the unspecified general metric-space claim.}
\GSfield{Fate}{\citet[Theorem 5.5]{ruelle1970superstable} proves the stated version in 1970. The unrestricted configuration space is not compact: arbitrarily many particles or colliding particles obstruct the naive argument. Thus the notebook sketch is not a proof at its stated generality.}
\GSfield{Present status}{As of November 2026: established under standard stability and regularity assumptions; the notebook’s arbitrary-space compactness argument remains insufficient, not an open existence problem.}
\GSfield{Assessment}{literal statement: false as written; publication match: partial; verification: read the result; development: none; confidence: high}
\end{GSitem}
\begin{GSitem}{1974-355.8}{E}{One-dimensional uniqueness}
\GSfield{Kind; source}{theorem; \GSpage{P000588}}
\GSfield{The notes}{For \(M=\mathbb R\) with Lebesgue measure, the Gibbs measure is said to be unique.}
\GSfield{Statement}{The intended assertion is uniqueness for the usual one-dimensional finite-range gas at finite positive temperature and activity, with a hard core and bounded finite interaction away from the core. It is not a uniqueness theorem for every measurable pair potential satisfying only the two displayed endpoint values, nor for arbitrary infinite-range interactions.}
\GSfield{Fate}{One-dimensional uniqueness belongs to the earlier Dobrushin--Ruelle theory. \citet[Section 2]{dobrushin1968conditional} contains the finite-state one-dimensional result; the continuum version in precisely the notebook’s scope was not recovered in the checked sources. Compare the older seminar occurrence 1970-195.5.}
\GSfield{Present status}{As of November 2026: standard one-dimensional finite-range uniqueness is established; the literal missing regularity and interaction assumptions prevent an exact source match here.}
\GSfield{Assessment}{literal statement: correct with missing hypotheses; publication match: partial; verification: indirect; development: none; confidence: medium}
\end{GSitem}
\begin{GSitem}{1974-355.9}{E}{Uniqueness at sufficiently low activity}
\GSfield{Kind; source}{theorem; \GSpage{P000588}}
\GSfield{The notes}{For \(U_1=0\) and \(m\le c\) times Lebesgue measure, sufficiently small \(c\) gives uniqueness.}
\GSfield{Statement}{For the translation-invariant superstable, lower-regular continuum pair potentials of 1974-355.7, the tempered Gibbs state is unique at sufficiently small positive activity \(z\). The smallness threshold depends on temperature and the potential. An inhomogeneous intensity dominated by \(c\,dx\) requires a corresponding uniform contraction estimate.}
\GSfield{Fate}{\citet[Theorem 5.7]{ruelle1970superstable} is an exact low-activity theorem for the homogeneous subclass. The notebook’s more general dominated measure is not separately verified here.}
\GSfield{Present status}{As of November 2026: established in the stated scope; the notebook is an expository occurrence, not a new open problem.}
\GSfield{Assessment}{literal statement: correct with missing hypotheses; publication match: partial; verification: read the result; development: none; confidence: medium}
\end{GSitem}
\begin{GSitem}{1974-355.10}{E}{Lattice gas with a Bernoulli reference measure}
\GSfield{Kind; source}{construction; \GSpage{P000588}}
\GSfield{The notes}{Replace \(\mathbb R^d\) by \(\mathbb Z^d\), \(d\ge2\), and use the symmetric Bernoulli measure on occupation variables.}
\GSfield{Statement}{A simple lattice gas is a spin field \(\eta\in\{0,1\}^{\mathbb Z^d}\). On a finite region \(\Lambda\), use the product Bernoulli-\(1/2\) reference and density proportional to \(e^{-\beta H_\Lambda(\eta\mid\xi)}\). A one-site field can absorb a different positive Bernoulli parameter. Finite-range energies define consistent finite-volume kernels.}
\GSfield{Fate}{This is a special case of \citet[Section 1]{dobrushin1968conditional}, published 1968, and the lattice setting of \citet[Section 1]{pirogovsinai1975phase}. The move to \(d\ge2\) prepares phase coexistence.}
\GSfield{Present status}{As of November 2026: established in the stated scope; the notebook is an expository occurrence, not a new open problem.}
\GSfield{Assessment}{literal statement: correct; publication match: exact; verification: read the result; development: none; confidence: high}
\end{GSitem}
\begin{GSitem}{1974-355.11}{E}{Weak-interaction or high-temperature uniqueness}
\GSfield{Kind; source}{theorem; \GSpage{P000588}}
\GSfield{The notes}{For sufficiently small \(U_1,U_2\), the Gibbs measure is unique.}
\GSfield{Statement}{For a finite-alphabet specification on \(\mathbb Z^d\), let \(c_{ij}\) be the supremum total-variation change in the conditional distribution at site \(i\) when only the exterior spin at \(j\) is changed. If \(\sup_i\sum_jc_{ij}<1\), there is at most one Gibbs measure; compactness and quasilocality supply existence. Bounded finite-range interactions satisfy this criterion at sufficiently high temperature.}
\GSfield{Reconstruction}{The dimensionless contraction condition replaces an unspecified norm-smallness assertion.}
\GSfield{Fate}{\citet[Theorem 2]{dobrushin1968conditional} gives the contraction criterion; \citet{dobrushin1968uniqueness} develops uniqueness and phase transitions. Smallness of the one-site field alone is not essential: the dependence on other sites is what is controlled.}
\GSfield{Present status}{As of November 2026: established in the stated scope; the notebook is an expository occurrence, not a new open problem.}
\GSfield{Assessment}{literal statement: correct with missing hypotheses; publication match: partial; verification: read the result; development: none; confidence: high; problem \hyperref[prob:GP-1974-20]{GP-1974-20}}
\end{GSitem}
\begin{GSitem}{1974-355.12}{E}{Periodic ground configurations and local energy differences}
\GSfield{Kind; source}{definition; \GSpage{P000588-P000589}}
\GSfield{The notes}{Translation-invariant \(U_1=\alpha\), even \(U_2(x-y)\), periodic configurations, and minimization against finite changes define the ground states.}
\GSfield{Statement}{For a finite-range translation-invariant lattice Hamiltonian, a configuration \(g\) is a ground state if every finite modification \(\eta\) satisfies \(H(\eta)-H(g)\ge0\), where only terms meeting the modified region enter the finite difference. Periodicity is an additional property, not part of the general ground-state definition.}
\GSfield{Fate}{The local-energy definition and periodic reference states are the framework of \citet[Section 1]{pirogovsinai1975phase}. The 1975 talk predates that paper but uses an older standard notion of ground configuration.}
\GSfield{Present status}{As of November 2026: established in the stated scope; the notebook is an expository occurrence, not a new open problem.}
\GSfield{Assessment}{literal statement: correct with missing hypotheses; publication match: partial; verification: read the result; development: none; confidence: high}
\end{GSitem}
\begin{GSitem}{1974-355.13}{E}{A finite set of periodic reference ground states}
\GSfield{Kind; source}{claim; \GSpage{P000589}}
\GSfield{The notes}{\(G(U)\) is taken to be finite, with \(|G(U)|=k\).}
\GSfield{Statement}{The phase-diagram theorem assumes that the model has exactly a specified finite family \(g_1,\ldots,g_k\) of periodic ground configurations, modulo the stated periodicity convention. Finite-range translation invariance does not imply this hypothesis; highly degenerate models can have infinitely many ground configurations.}
\GSfield{Fate}{This is a hypothesis of \citet[Section 1 and Theorem 1]{pirogovsinai1975phase}, not a universal theorem. The exact low-temperature conclusion depends also on Peierls stability and independent perturbations.}
\GSfield{Present status}{As of November 2026: established in the stated scope; the notebook is an expository occurrence, not a new open problem.}
\GSfield{Assessment}{literal statement: correct with missing hypotheses; publication match: exact; verification: read the result; development: none; confidence: high}
\end{GSitem}
\begin{GSitem}{1974-355.14}{A}{Peierls energy cost for defects}
\GSfield{Kind; source}{definition; \GSpage{P000589}}
\GSfield{The notes}{A site lies on the boundary if its neighbourhood agrees with no ground state; \(H(X')-H(X)\ge g\,\Gamma(X')\).}
\GSfield{Statement}{Choose a fixed coarse-graining scale large enough to recognize the periodic ground patterns and the interaction range. A bad cube is one not locally matching any reference pattern. The Peierls condition requires the energy excess of a finite perturbation of a ground state to be at least \(\tau>0\) times the number of bad cubes, with \(\tau\) uniform.}
\GSfield{Reconstruction}{The unspecified neighbourhood and boundary cardinality are replaced by the finite-scale definition in the paper.}
\GSfield{Fate}{\citet[Definitions 1.2--1.3]{pirogovsinai1975phase} gives the matching definition. This particular contour stability hypothesis is part of the Pirogov--Sinai theorem announced here and published later in 1975.}
\GSfield{Present status}{As of November 2026: established in the stated scope; the notebook is an expository occurrence, not a new open problem.}
\GSfield{Assessment}{literal statement: correct with missing hypotheses; publication match: exact; verification: read the result; development: none; confidence: high}
\end{GSitem}
\begin{GSitem}{1974-355.15}{A}{Independent parameters unfold the ground-state energies}
\GSfield{Kind; source}{construction; \GSpage{P000589}}
\GSfield{The notes}{\(U_\mu=U+\sum_{j=1}^{k-1}\mu_jU^{(j)}\), near \(\mu=0\), and a condition written \(G(U_\mu)\le G(U)\).}
\GSfield{Statement}{Choose \(k-1\) perturbing interactions whose energy-density differences on \(g_1,\ldots,g_k\), relative to one reference state, form a nonsingular \((k-1)\times(k-1)\) matrix. The ensuing theorem is local in this parameter neighbourhood, with the listed periodic ground states as its reference phases. A set of ground states is not numerically ordered.}
\GSfield{Reconstruction}{The nonsingularity condition is recovered from the matched theorem; the ambiguous set inequality is not silently converted into a cardinality bound.}
\GSfield{Fate}{\citet[Definition 1.4 and Theorem 1]{pirogovsinai1975phase} supplies the independence condition. The notebook’s inequality is plausibly shorthand for inclusion/no new reference states, but its intended meaning is not certain.}
\GSfield{Present status}{As of November 2026: established in the stated scope; the notebook is an expository occurrence, not a new open problem.}
\GSfield{Assessment}{literal statement: correct with missing hypotheses; publication match: partial; verification: read the result; development: none; confidence: medium}
\end{GSitem}
\begin{GSitem}{1974-355.16}{A}{Low-temperature coexistence and phase-diagram strata}
\GSfield{Kind; source}{theorem; \GSpage{P000589}}
\GSfield{The notes}{For large \(\beta\), a point \(\mu(\beta)\) has \(k\) Gibbs distributions, \(k\) emanating curves have \(k-1\), and \(\binom{k}{2}\) surfaces have \(k-2\).}
\GSfield{Statement}{Under the finite-state, finite-range, finitely many periodic ground states, Peierls, and independent-perturbation hypotheses, the low-temperature phase diagram near the reference point has a coexistence locus where \(k\) stable periodic pure phases coexist. Its adjacent local strata correspond to subsets of the stable phases, with the expected codimensions. Counts refer to stable pure phases, not all Gibbs probability measures; convex mixtures already give infinitely many measures when two phases coexist.}
\GSfield{Fate}{This is the central match to \citet[Theorem 1]{pirogovsinai1975phase}, with proof continued in \citet{pirogovsinai1976phase}. Both papers were received on 21 April 1975, after this 24 February notebook occurrence, and appeared in 1975 and 1976 respectively. The notebook’s unqualified uniqueness and total-measure count are stronger than the safely matched statement.}
\GSfield{Present status}{As of November 2026: the Pirogov--Sinai local phase-diagram theory is established. This entry is evidence of an early seminar presentation, not a still-open conjecture; no completeness claim about all nonperiodic or interface Gibbs states is inferred.}
\GSfield{Assessment}{literal statement: false as written; publication match: partial; verification: read the result; development: none; confidence: high}
\end{GSitem}
\begin{GSitem}{1974-355.17}{F}{Gelfand asks for a p-adic Gibbs analogue}
\GSfield{Kind; source}{question; \GSpage{P000589}; I. M. Gelfand}
\GSfield{The notes}{\ru{Вопрос Г. Гиббсовские распределения в \(p\)-адическом случае}.}
\GSfield{Statement}{The note asks for a \(p\)-adic analogue but does not say whether the sites, spin space, interaction values, or probability values are \(p\)-adic, or whether existence, uniqueness, phase transitions, or a particular construction is sought. These interpretations lead to different mathematical problems.}
\GSfield{Fate}{No exact specification is recoverable from this question. A concrete later thematic direction is \citet{rozikovkhakimov2013padic}, received 16 October 2012 and published in 2013, whose publisher abstract discusses \(p\)-adic probability-valued Gibbs measures and Markov fields on countable graphs. That is distinct from real-valued Gibbs measures on a \(p\)-adic spatial domain, and the notebook does not choose between such possibilities. No solution of Gelfand's precise intention is asserted.}
\GSfield{Present status}{As of November 2026: no definite historical problem is recoverable; the broad prompt is not classified as an undeveloped research gap.}
\GSfield{Assessment}{literal statement: not determinable; publication match: thematic only; verification: publisher metadata checked; development: none; confidence: low}
\end{GSitem}

\GStalk{356}{1975-03-03}{V. N. Tulovskii}{P000589-P000594}{Oscillatory integrals and Lagrangian distributions}
The dated session occupies P000590--P000592, ending before 10 March. It constructs oscillatory distributions, their conic Lagrangians, and Fourier integral kernels. Kirillov inserts the affine geometry of Lagrangian leaves. Except for that intervention, the explicit reference remains Hörmander's already published FIO theory. The corollary about possible leaf topologies is too strong.

\begin{GSitem}{1974-356.1}{E}{Testing wavefront regularity with nonlinear phases}
\GSfield{Kind; source}{theorem; \GSpage{P000590}}
\GSfield{The notes}{The linear Fourier phase is replaced by \(\Phi\) with \(\Phi'(x_0)=\xi_0\), and the pairing with \(\varphi e^{i\lambda\Phi}\) is rapidly decreasing.}
\GSfield{Statement}{Fix \(\widehat u(\xi)=\langle u,e^{-ix\cdot\xi}\rangle\). Microlocal smoothness at \((x_0,\xi_0)\) is equivalent to rapid decay of \(\langle u,\chi e^{-it\psi}\rangle\) as \(t\to+\infty\), for cutoffs supported sufficiently near \(x_0\) and smooth real phases whose differentials stay in a sufficiently small cone around \(\xi_0\). The estimates are uniform over suitable bounded families of such phases and cutoffs.}
\GSfield{Reconstruction}{A single negative-sign Fourier convention is used here and below; the original convention is not specified.}
\GSfield{Fate}{This is the nonlinear phase form of the Fourier criterion and nonstationary-phase estimates in \citet[Sections 2.3 and 2.5]{hormander1971fourier}. The notebook's plus sign is consistent with the opposite Fourier convention; its missing local and uniform quantifiers require repair.}
\GSfield{Present status}{As of November 2026: established in the stated scope; the notebook is an expository occurrence, not a new open problem.}
\GSfield{Assessment}{literal statement: correct with missing hypotheses; publication match: partial; verification: read the result; development: none; confidence: high}
\end{GSitem}
\begin{GSitem}{1974-356.2}{E}{Defining an oscillatory distribution by cutoff removal}
\GSfield{Kind; source}{construction; \GSpage{P000590-P000591}}
\GSfield{The notes}{For \(a\in S^m_{1,0}\), \(\langle A,u\rangle=\lim_{\varepsilon\to0}\int a(x,\theta)e^{i\Phi(x,\theta)}u(x)\chi(\varepsilon\theta)\,dx\,d\theta\).}
\GSfield{Statement}{Let \(\Phi\) be smooth real and positively homogeneous of degree one for \(\theta\ne0\), with \(d_{x,\theta}\Phi\ne0\), and \(a\in S^m_{1,0}\), with harmless smooth extension at bounded \(\theta\). For a smooth compactly supported cutoff \(\chi=1\) near zero, integration by parts defines the indicated distributional limit independently of \(\chi\). Tests are compactly supported in the \(x\)-variables.}
\GSfield{Reconstruction}{The meaning of phase function is supplied explicitly.}
\GSfield{Fate}{\citet[Section 1.2 and Proposition 2.5.7]{hormander1971fourier} contains this construction; compare 1973-329.12--13. Homogeneity and nonvanishing phase differential are part of the phase-function definition, not optional hypotheses.}
\GSfield{Present status}{As of November 2026: established in the stated scope; the notebook is an expository occurrence, not a new open problem.}
\GSfield{Assessment}{literal statement: correct with missing hypotheses; publication match: exact; verification: read the result; development: none; confidence: high}
\end{GSitem}
\begin{GSitem}{1974-356.3}{E}{The linear-phase delta example}
\GSfield{Kind; source}{remark; \GSpage{P000591}}
\GSfield{The notes}{For \(a=1\), \(\Phi=\langle x,\theta\rangle\), the notebook gives \(\langle A,u\rangle=u(0)\).}
\GSfield{Statement}{With ordinary Lebesgue \(d\theta\) and the Fourier convention fixed above, \(\int_{\mathbb R^n}e^{ix\cdot\theta}\,d\theta=(2\pi)^n\delta_0\). Thus the stated value \(u(0)\) uses the normalized measure \((2\pi)^{-n}d\theta\). The distribution is concentrated on the critical set \(x=0\).}
\GSfield{Fate}{This is Fourier inversion, the basic example behind \citet[Section 1.2]{hormander1971fourier}. The notebook does not fix its measure normalization, so this is a convention clarification rather than a certain mathematical error.}
\GSfield{Present status}{As of November 2026: established in the stated scope; the notebook is an expository occurrence, not a new open problem.}
\GSfield{Assessment}{literal statement: correct with missing hypotheses; publication match: partial; verification: read the result; development: none; confidence: high}
\end{GSitem}
\begin{GSitem}{1974-356.4}{E}{Removing the nonstationary part of an amplitude}
\GSfield{Kind; source}{method; \GSpage{P000591}}
\GSfield{The notes}{Split \(a=a_1+a_2\) near a point where \(\Phi'_\theta\ne0\); the corresponding \(A_2\) is smooth.}
\GSfield{Statement}{On a closed conic support where \(|\partial_\theta\Phi|\) is bounded away from zero after angular localization, choose a differential operator \(L\) in \(\theta\) with \(Le^{i\Phi}=e^{i\Phi}\). Iterating its transpose improves amplitude integrability and proves that the resulting oscillatory integral is smooth in \(x\).}
\GSfield{Fate}{The nonstationary-phase integration-by-parts method is established in \citet[Section 1.2]{hormander1971fourier}. The conic cutoff in the sketch records localization, not an additional singularity.}
\GSfield{Present status}{As of November 2026: established in the stated scope; the notebook is an expository occurrence, not a new open problem.}
\GSfield{Assessment}{literal statement: correct; publication match: exact; verification: read the result; development: none; confidence: high}
\end{GSitem}
\begin{GSitem}{1974-356.5}{E}{The critical-set bound for the wavefront}
\GSfield{Kind; source}{theorem; \GSpage{P000591}}
\GSfield{The notes}{\(C_\Phi=\{\Phi'_\theta=0\}\), \(\lambda_\Phi(x,\theta)=(x,\Phi'_x)\), and \(\mathrm{WF}(A)\subset\Lambda_\Phi=\lambda_\Phi(C_\Phi)\).}
\GSfield{Statement}{For a phase and symbol as above, the wavefront of its oscillatory integral is contained in the closed conic set of \((x,\partial_x\Phi(x,\theta))\) with \(\partial_\theta\Phi=0\) and \((x,\theta)\) in the conic support of the amplitude. One takes the relevant closure when supports or charts are not proper.}
\GSfield{Fate}{\citet[Proposition 2.5.7]{hormander1971fourier} is the exact analytic source. The formula is a containment; amplitudes may vanish or cancel, so equality is not automatic.}
\GSfield{Present status}{As of November 2026: established in the stated scope; the notebook is an expository occurrence, not a new open problem.}
\GSfield{Assessment}{literal statement: correct; publication match: exact; verification: read the result; development: none; confidence: high}
\end{GSitem}
\begin{GSitem}{1974-356.6}{E}{Nondegenerate phases parametrize Lagrangians}
\GSfield{Kind; source}{definition; \GSpage{P000591}}
\GSfield{The notes}{The rank of \((\Phi''_{\theta x},\Phi''_{\theta\theta})\) is \(N\) on \(C_\Phi\); then \(\lambda_\Phi\) is locally a diffeomorphism and \(\Lambda_\Phi\) is conic Lagrangian.}
\GSfield{Statement}{If the \(N\) differentials \(d(\partial_{\theta_j}\Phi)\) are independent on \(C_\Phi\), this critical set is smooth of dimension \(n\). The map \(\lambda_\Phi\) is a conic Lagrangian immersion and locally identifies \(C_\Phi\) with its image. It is not a local diffeomorphism onto the \(2n\)-dimensional cotangent bundle.}
\GSfield{Fate}{This is the standard nondegenerate phase construction in \citet[Section 3.1]{hormander1971fourier}. The notebook's local-diffeomorphism phrase is correct only with the image as target.}
\GSfield{Present status}{As of November 2026: established in the stated scope; the notebook is an expository occurrence, not a new open problem.}
\GSfield{Assessment}{literal statement: correct with missing hypotheses; publication match: exact; verification: read the result; development: none; confidence: high}
\end{GSitem}
\begin{GSitem}{1974-356.7}{E}{Local phase functions for every conic Lagrangian}
\GSfield{Kind; source}{theorem; \GSpage{P000591}}
\GSfield{The notes}{Conversely every smooth conic Lagrangian can locally be obtained from a nondegenerate phase, and one may take \(N=n\).}
\GSfield{Statement}{Every germ of a smooth conic Lagrangian in \(T^*X\setminus0\), \(\dim X=n\), has a local positively homogeneous nondegenerate phase representation with at most \(n\) phase variables. Additional variables may be introduced by the usual stabilization on a conic chart, with the accompanying amplitude-order shift.}
\GSfield{Fate}{\citet[Section 3.1, especially Theorems 3.1.3 and 3.1.6]{hormander1971fourier} establishes local parametrization and phase equivalence. The statement is local in the Lagrangian, not a claim of one global generating function.}
\GSfield{Present status}{As of November 2026: established in the stated scope; the notebook is an expository occurrence, not a new open problem.}
\GSfield{Assessment}{literal statement: correct; publication match: exact; verification: read the result; development: none; confidence: high}
\end{GSitem}
\begin{GSitem}{1974-356.8}{E}{Euler homogeneity proves the Lagrangian condition}
\GSfield{Kind; source}{method; \GSpage{P000591}}
\GSfield{The notes}{On \(C_\Phi\), \(\Phi'_\theta=0\) implies \(\Phi=0\), then \(\Phi'_x\,dx=0\); hence \(\xi\,dx=0\) on the image.}
\GSfield{Statement}{Euler's identity gives \(\Phi=\theta\cdot\partial_\theta\Phi=0\) on \(C_\Phi\). Restricting \(d\Phi\) gives \(\partial_x\Phi\,dx=0\), so the pullback of the Liouville form vanishes. Its exterior derivative therefore vanishes too; together with nondegeneracy and dimension \(n\), the image is Lagrangian.}
\GSfield{Fate}{This is the elementary proof in the phase geometry of \citet[Section 3.1]{hormander1971fourier}. Homogeneity is essential to the stronger vanishing of the Liouville form.}
\GSfield{Present status}{As of November 2026: established in the stated scope; the notebook is an expository occurrence, not a new open problem.}
\GSfield{Assessment}{literal statement: correct; publication match: exact; verification: read the result; development: none; confidence: high}
\end{GSitem}
\begin{GSitem}{1974-356.9}{E}{Equivalent phases define the same distribution class}
\GSfield{Kind; source}{theorem; \GSpage{P000591}}
\GSfield{The notes}{Two phases parametrizing the same \(\Lambda\) produce the same class of \(A\).}
\GSfield{Statement}{The microlocal class \(I^m(X,\Lambda)\) is independent of the chosen nondegenerate parametrization. In Hörmander's normalization a phase with \(N\) variables uses amplitude order \(m+(\dim X-2N)/4\). Changing \(N\) changes that amplitude order; the principal symbol also includes the stationary-phase signature/Maslov and density factors.}
\GSfield{Reconstruction}{The order shift suppressed in the notebook is restored.}
\GSfield{Fate}{\citet[Theorem 3.2.1 and Definition 3.2.2]{hormander1971fourier} supplies the exact phase-invariance theorem. Equal raw amplitude orders for phases with different \(N\) would be incorrect.}
\GSfield{Present status}{As of November 2026: established in the stated scope; the notebook is an expository occurrence, not a new open problem.}
\GSfield{Assessment}{literal statement: correct with missing hypotheses; publication match: exact; verification: read the result; development: none; confidence: high}
\end{GSitem}
\begin{GSitem}{1974-356.10}{E}{Fourier integral kernels and gluing over phase charts}
\GSfield{Kind; source}{construction; \GSpage{P000591-P000592}}
\GSfield{The notes}{For \(X,Y\), a conic Lagrangian in \(T^*(X\times Y)\setminus0\) gives a kernel by an oscillatory integral; in general use a sum of phase charts.}
\GSfield{Statement}{A Fourier integral kernel is a Lagrangian distribution on \(X\times Y\), locally represented by nondegenerate oscillatory integrals with compatible order. Partitions of unity and locally finite sums produce the global distribution. To extend the resulting operator to compactly supported distribution inputs, exclude the input axis wavefront; proper support controls noncompact inputs.}
\GSfield{Fate}{This is the FIO construction of \citet[Sections 3.2 and 4.1]{hormander1971fourier}. The associated operator relation reverses the sign of the \(Y\)-covector, as already recorded in 1974-353.9.}
\GSfield{Present status}{As of November 2026: established in the stated scope; the notebook is an expository occurrence, not a new open problem.}
\GSfield{Assessment}{literal statement: correct with missing hypotheses; publication match: exact; verification: read the result; development: none; confidence: high}
\end{GSitem}
\begin{GSitem}{1974-356.11}{E}{Kirillov: affine structures on Lagrangian leaves}
\GSfield{Kind; source}{theorem; \GSpage{P000592}; A. A. Kirillov}
\GSfield{The notes}{If a symplectic manifold has a Lagrangian foliation, its leaves carry an invariantly defined affine structure.}
\GSfield{Statement}{For a Lagrangian foliation with tangent bundle \(L\), the symplectic form identifies \(L\) with its annihilator in \(T^*M\). Transporting the Bott partial connection on that annihilator gives a flat torsion-free connection along each leaf. This is an intrinsic affine structure, without a completeness or trivial-holonomy conclusion.}
\GSfield{Fate}{A checked modern formulation and construction are \citet[Proposition 1.9 and Theorem 1.10]{loustausanders2017bilagrangian}, whose introduction attributes the result to Weinstein (1971). The 1971 original was not accessible for direct theorem verification; the modern result was read. The notebook is therefore an expository occurrence, not evidence of priority for Kirillov.}
\GSfield{Present status}{As of November 2026: established in the stated scope; the notebook is an expository occurrence, not a new open problem.}
\GSfield{Assessment}{literal statement: correct with missing hypotheses; publication match: exact; verification: read the result; development: none; confidence: medium}
\end{GSitem}
\begin{GSitem}{1974-356.12}{E}{Basic Hamiltonians give parallel commuting fields}
\GSfield{Kind; source}{construction; \GSpage{P000592}; A. A. Kirillov}
\GSfield{The notes}{For functions constant along the leaves, the covectors \(d\varphi\) correspond to parallel vector fields; these commute.}
\GSfield{Statement}{If \(f\) and \(g\) are locally constant along a Lagrangian foliation, their Hamiltonian fields are tangent to its leaves and \(\{f,g\}=0\); hence they commute. Locally independent basic differentials give a parallel frame for the leafwise affine connection. Such a frame need not be globally single-valued or complete.}
\GSfield{Fate}{This is the Hamiltonian description accompanying \citet[Theorem 1.10]{loustausanders2017bilagrangian}. The notebook calls the vectors gradients, but the construction uses the symplectic isomorphism, not a Riemannian metric.}
\GSfield{Present status}{As of November 2026: established in the stated scope; the notebook is an expository occurrence, not a new open problem.}
\GSfield{Assessment}{literal statement: correct; publication match: partial; verification: read the result; development: none; confidence: high}
\end{GSitem}
\begin{GSitem}{1974-356.13}{E}{The asserted Euclidean-or-torus leaf dichotomy}
\GSfield{Kind; source}{claim; \GSpage{P000592}; A. A. Kirillov}
\GSfield{The notes}{\ru{Следствие: слои --- \(\mathbb R^n\) или торы}. An uncertain abbreviation follows as an example of a Lagrangian foliation.}
\GSfield{Statement}{The dichotomy is false. For \(T^*(S^1\times\mathbb R)\) with the standard form, the leaves with both momenta constant are Lagrangian cylinders \(S^1\times\mathbb R\). An affine structure alone gives no such topological classification. A connected leaf with \(n\) globally independent complete commuting fields is instead \(\mathbb R^n/\Gamma\), with discrete \(\Gamma\), hence may have any cylinder rank.}
\GSfield{Fate}{The leafwise result verified in \citet[Theorem 1.10]{loustausanders2017bilagrangian} does not assert the notebook corollary. The editorial counterexample is checked in the mathematical audit. The following abbreviation likely denotes a cotangent bundle (\ru{кокасательное расслоение}); it should not be turned into a person named Kokas.}
\GSfield{Present status}{As of November 2026: the literal corollary is false; the affine-connection theorem is established, and no open classification problem is inferred.}
\GSfield{Assessment}{literal statement: false as written; publication match: partial; verification: read the result; development: annotated; confidence: high}
\end{GSitem}
\begin{GSitem}{1974-356.14}{E}{Recognizing FIOs by rapidly oscillating tests}
\GSfield{Kind; source}{claim; \GSpage{P000592}}
\GSfield{The notes}{A kernel whose wavefront lies in \(\Lambda\) and which acts well on transverse oscillating exponentials is an FIO; an asymptotic series is sketched.}
\GSfield{Statement}{For a classical Lagrangian distribution of order \(m\) on a \(d\)-manifold, pairings with \(\chi e^{-it\psi}\), where \(\operatorname{graph}(d\psi)\) meets \(\Lambda\) transversely, have stationary-phase expansions with leading scale \(t^{m-d/4}\) and phase values at the intersections. A converse requires uniform symbol estimates in test-phase parameters, not merely a separate pointwise asymptotic for each \(\psi\).}
\GSfield{Fate}{\citet[Theorems 3.2.4--3.2.5 and the following characterization paragraph]{hormander1971fourier} establishes the forward statement and describes the converse. Record 357 repeats the question explicitly. The notebook suppresses the symbol order and essential uniformity.}
\GSfield{Present status}{As of November 2026: established in the stated scope; the notebook is an expository occurrence, not a new open problem.}
\GSfield{Assessment}{literal statement: correct with missing hypotheses; publication match: partial; verification: read the result; development: none; confidence: high}
\end{GSitem}

\GStalk{357}{1975-03-10}{V. N. Tulovskii}{P000592-P000596}{FIO normal forms, propagation, and characteristic dimension}
The session starts on P000593 and ends with Bernstein's short intervention on P000594. Tulovskii restates FIO kernels, a testing characterization, the normal form for a real principal symbol, and propagation of singularities. Gelfand asks about a \(p\)-adic version. Bernstein switches from the analytic wavefront of a distribution to the algebraic characteristic variety of a differential-operator ideal; that change of meaning is essential.

\begin{GSitem}{1974-357.1}{E}{FIO operator action on compactly supported distributions}
\GSfield{Kind; source}{construction; \GSpage{P000593}}
\GSfield{The notes}{The oscillatory kernel formula is repeated, with \(A:\mathcal E'(Y)\to\mathcal D'(X)\).}
\GSfield{Statement}{A distribution kernel represented by nondegenerate phase charts defines this extension when its wavefront avoids covectors \((x,0;y,\eta)\), \(\eta\ne0\). The other axis condition gives smooth output on smooth test inputs. Merely being a Lagrangian distribution on \(X\times Y\) does not ensure both operator properties.}
\GSfield{Fate}{This is the combination of \citet[Theorems 2.5.12--2.5.14 and Section 4.1]{hormander1971fourier}, already discussed in 1974-353.8 and 1974-356.10.}
\GSfield{Present status}{As of November 2026: established in the stated scope; the notebook is an expository occurrence, not a new open problem.}
\GSfield{Assessment}{literal statement: correct with missing hypotheses; publication match: exact; verification: read the result; development: none; confidence: high}
\end{GSitem}
\begin{GSitem}{1974-357.2}{E}{Pseudodifferential operators as diagonal FIOs}
\GSfield{Kind; source}{remark; \GSpage{P000593}}
\GSfield{The notes}{In the PDO case \(\Lambda=\{(x,\xi;x,-\xi)\}\).}
\GSfield{Statement}{The unprimed kernel relation of a pseudodifferential operator is contained in the nonzero conormal bundle of the diagonal. After reversing the input covector, it becomes the identity relation on \(T^*X\setminus0\). Classical symbols produce the familiar oscillatory kernel with phase \((x-y)\cdot\xi\).}
\GSfield{Fate}{\citet[Section 4.1]{hormander1971fourier} places pseudodifferential operators in the FIO class. Earlier occurrences include 1971-234.3 and 1973-329.5--7.}
\GSfield{Present status}{As of November 2026: established in the stated scope; the notebook is an expository occurrence, not a new open problem.}
\GSfield{Assessment}{literal statement: correct; publication match: exact; verification: read the result; development: none; confidence: high}
\end{GSitem}
\begin{GSitem}{1974-357.3}{F}{Gelfand asks for a p-adic FIO construction}
\GSfield{Kind; source}{question; \GSpage{P000593}; I. M. Gelfand}
\GSfield{The notes}{\ru{Вопрос Г.: хорошо бы в \(p\)-адическом случае построить}. The intended object is the just-discussed calculus.}
\GSfield{Statement}{The notes request a non-Archimedean analogue without specifying a field, character, symbol class, canonical relations, or closure operations. There is no unique exact theorem to match. A later concrete analogue uses distributional wavefront sets over a \(p\)-adic field and oscillatory integrals with a fixed additive character.}
\GSfield{Fate}{\citet[Sections 1--2]{heifetz1985padic}, published in 1985, constructs a \(p\)-adic wavefront theory and stationary-phase methods, including pullbacks, products and kernel operations. This is a substantial thematic continuation, but it does not establish that a particular full FIO calculus requested here was completed. No exact solution date is assigned.}
\GSfield{Present status}{As of November 2026: relevant non-Archimedean microlocal theories exist, but the notebook prompt is too underspecified to classify as a solved or unsolved definite problem.}
\GSfield{Assessment}{literal statement: not determinable; publication match: thematic only; verification: read the result; development: none; confidence: medium}
\end{GSitem}
\begin{GSitem}{1974-357.4}{E}{When is a distribution kernel a Fourier integral kernel?}
\GSfield{Kind; source}{question; \GSpage{P000593}}
\GSfield{The notes}{Tulovskii answers: wavefront inside a Lagrangian, and transverse exponential tests have an asymptotic expansion determined by that Lagrangian.}
\GSfield{Statement}{A precise recognition problem is to characterize classical Lagrangian distributions \(I^m_{\mathrm{cl}}(X\times Y,\Lambda)\) by their pairings with localized transverse oscillatory tests, uniformly in smooth phase families and with the order \(m-(\dim X+\dim Y)/4\) prescribed. Wavefront containment alone does not characterize this class; singular amplitudes of arbitrary order can have the same wavefront set.}
\GSfield{Fate}{\citet[Theorems 3.2.4--3.2.5 and the following paragraph]{hormander1971fourier} already gives the asymptotic test framework and explicitly indicates its converse characterization. The question repeats 1974-356.14 and is expository, with missing uniform symbol bounds rather than an identified research gap.}
\GSfield{Present status}{As of November 2026: established in the stated scope; the notebook is an expository occurrence, not a new open problem.}
\GSfield{Assessment}{literal statement: correct with missing hypotheses; publication match: partial; verification: read the result; development: none; confidence: high}
\end{GSitem}
\begin{GSitem}{1974-357.5}{E}{Microlocal straightening of a real principal symbol}
\GSfield{Kind; source}{theorem; \GSpage{P000593}}
\GSfield{The notes}{For an order-one PDO with real \(p\), \(p(x_0,\xi_0)=0\), and \(dp(x_0,\xi_0)\ne0\), there are FIOs \(A,B\) with \(P\equiv B\partial_{X_n}A\) and \(BA\equiv I\) microlocally.}
\GSfield{Statement}{The standard homogeneous FIO normal form additionally requires that \(H_p\) not be tangent to the radial cone direction at the point. For a classical order-one operator of real principal type under that nonradial condition, elliptic local FIOs quantizing inverse homogeneous canonical transformations conjugate \(P\), modulo smoothing terms, to \(D_{X_n}=-i\partial_{X_n}\), with harmless elliptic factors according to normalization.}
\GSfield{Fate}{\citet[Proposition 6.1.4 and its application]{duistermaathormander1972fourierii} explicitly includes the nonradial hypothesis. The notebook omits it. The symbol \(p=x\xi\) at \((0,1)\) has \(dp\ne0\) but radial Hamiltonian field, and cannot be straightened to a nonradial coordinate symbol by a homogeneous canonical transformation; see the audit.}
\GSfield{Present status}{As of November 2026: the qualified normal-form theorem is established; the dp-only homogeneous formulation is false.}
\GSfield{Assessment}{literal statement: false as written; publication match: partial; verification: read the result; development: annotated; confidence: high}
\end{GSitem}
\begin{GSitem}{1974-357.6}{E}{Propagation of singularities for real principal type}
\GSfield{Kind; source}{theorem; \GSpage{P000593}}
\GSfield{The notes}{For \(Pu=f\), \(\mathrm{WF}(u)\setminus\mathrm{WF}(f)\subset\{p=0\}\) and it is invariant under \(H_p=\sum(p_{x_j}\partial_{\xi_j}-p_{\xi_j}\partial_{x_j})\).}
\GSfield{Statement}{For a scalar classical PDO with real homogeneous principal symbol of principal type, elliptic regularity gives \(\mathrm{WF}(u)\setminus\mathrm{WF}(Pu)\subset\operatorname{Char}(P)\). In the open complement of \(\mathrm{WF}(Pu)\), this set is a union of bicharacteristic segments of the principal Hamilton field. Invariance is understood only while the segment stays in that open region; it does not continue through singular forcing without qualification.}
\GSfield{Reconstruction}{The principal-type scope and restriction to nonsingular forcing are made explicit.}
\GSfield{Fate}{\citet[Theorem 6.1.1]{duistermaathormander1972fourierii} is the published 1972 result. The notebook's Hamilton field is the negative of the common convention, which reverses orientation but not the unparameterized bicharacteristics.}
\GSfield{Present status}{As of November 2026: established in the stated scope; the notebook is an expository occurrence, not a new open problem.}
\GSfield{Assessment}{literal statement: correct with missing hypotheses; publication match: exact; verification: read the result; development: none; confidence: high}
\end{GSitem}
\begin{GSitem}{1974-357.7}{E}{Bernstein: characteristic variety of an operator ideal}
\GSfield{Kind; source}{definition; \GSpage{P000594}; I. N. Bernstein}
\GSfield{The notes}{The \(\mathrm{WF}\) of a left ideal of polynomial-coefficient differential operators is the common zero set of its principal symbols.}
\GSfield{Statement}{Over a characteristic-zero field, let \(D\) be the Weyl algebra and \(I\subsetneq D\) a left ideal. With the order filtration and its induced good quotient filtration, \(\operatorname{Char}(D/I)\subset T^*\mathbb A^n\) is the zero set of \(\operatorname{gr}I\). Symbols of all elements of \(I\), not merely an arbitrary generating list, must be used. The zero section is retained.}
\GSfield{Fate}{\citet[Section 1, Definition 1.6]{bernstein1971modules} develops characteristic carriers for differential-operator modules. The word \(\mathrm{WF}\) here is an algebraic analogy, not the smooth wavefront of one distribution.}
\GSfield{Present status}{As of November 2026: established in the stated scope; the notebook is an expository occurrence, not a new open problem.}
\GSfield{Assessment}{literal statement: correct with missing hypotheses; publication match: partial; verification: read the result; development: none; confidence: high}
\end{GSitem}
\begin{GSitem}{1974-357.8}{E}{Bernstein’s lower bound on characteristic dimension}
\GSfield{Kind; source}{theorem; \GSpage{P000594}; I. N. Bernstein}
\GSfield{The notes}{\(\dim\mathrm{WF}\ge\dim X\).}
\GSfield{Statement}{For a nonzero finitely generated Weyl-algebra module over a field of characteristic zero, its dimension is at least \(n\). For the order characteristic variety of \(D/I\), this gives \(\dim\operatorname{Char}(D/I)\ge n\). The hypothesis \(I\ne D\) is essential; no corresponding assertion holds for the empty characteristic variety of the zero module.}
\GSfield{Fate}{\citet[Theorem 1.3, proved in Section 5]{bernstein1972continuation} proves the dimension inequality; the paper was submitted on 21 June 1972 and published in 1972. The 1971 paper supplies the filtered-module framework, not the exact inequality locator. A modern full order-characteristic statement is \citet[Theorem 3.38]{mustata2023dmodules}. Thus the 1975 intervention is expository.}
\GSfield{Present status}{As of November 2026: established in the stated scope; the notebook is an expository occurrence, not a new open problem.}
\GSfield{Assessment}{literal statement: correct with missing hypotheses; publication match: exact; verification: read the result; development: none; confidence: high}
\end{GSitem}

\GStalk{358}{1975-03-17}{R. L. Dobrushin}{P000594-P000598}{Random fields specified by conditional laws}
The actual talk occupies P000595--P000596. It distinguishes spin, particle, and generalized random fields, formulates the inverse problem of prescribed conditional distributions, then explains regular conditional probabilities and measurable-map disintegration. This returns to the framework needed in Sinai's two talks. The final sentence, about conditioning on countably many real functions, completes the available record.

\begin{GSitem}{1974-358.1}{E}{Finite-spin random fields}
\GSfield{Kind; source}{definition; \GSpage{P000595}}
\GSfield{The notes}{\(X\) is finite, \(T\) countable, and a random field is a probability measure on \(X^T\); the physical interpretation is a spin system.}
\GSfield{Statement}{Equip the finite alphabet \(X\) with its discrete sigma-algebra and \(X^T\) with the product sigma-algebra. A probability law there defines the coordinate random variables \(\eta_t\); conversely their consistent finite-dimensional laws determine that probability uniquely. No Gibbs property is part of this definition.}
\GSfield{Fate}{This is the starting framework of \citet[Section 1]{dobrushin1968conditional}; its use in the notebook is exposition of established probability theory.}
\GSfield{Present status}{As of November 2026: established in the stated scope; the notebook is an expository occurrence, not a new open problem.}
\GSfield{Assessment}{literal statement: correct; publication match: exact; verification: read the result; development: none; confidence: high}
\end{GSitem}
\begin{GSitem}{1974-358.2}{E}{Point random fields}
\GSfield{Kind; source}{definition; \GSpage{P000595}}
\GSfield{The notes}{Realizations are subsets of \(\mathbb R^\nu\) meeting every bounded region in finitely many points; a measure on realizations is a point random field.}
\GSfield{Statement}{A simple locally finite point field is a probability measure on the configuration space with the count-generated sigma-algebra from 1974-354.1. Local finiteness applies to relatively compact Borel regions; integer-valued counting measures allow multiplicity when desired.}
\GSfield{Fate}{The continuum configuration framework is explicitly used in \citet[Section 5]{ruelle1970superstable}. This repeats Sinai's model while distinguishing it from a field with a fixed site set.}
\GSfield{Present status}{As of November 2026: established in the stated scope; the notebook is an expository occurrence, not a new open problem.}
\GSfield{Assessment}{literal statement: correct; publication match: exact; verification: read the result; development: none; confidence: high}
\end{GSitem}
\begin{GSitem}{1974-358.3}{E}{Generalized random fields valued in distributions}
\GSfield{Kind; source}{definition; \GSpage{P000595}}
\GSfield{The notes}{The third object is a probability measure on \(\mathcal S^*(\mathbb R^\nu)\), linked to quantum field theory.}
\GSfield{Statement}{A tempered generalized random field may be represented by a probability measure on \(\mathcal S'(\mathbb R^\nu)\), with evaluation maps \(u\mapsto\langle u,\varphi\rangle\) measurable for every Schwartz test function. Such a field need not admit pointwise values. The mere definition does not construct a measure for an interacting quantum field model.}
\GSfield{Fate}{\citet[Section 1]{minlos1959generalized} develops generalized random processes and their extension to measures. Its main results were reported on 11 April 1958 and published in 1959, well before this session. The notes offer only the object definition and motivation.}
\GSfield{Present status}{As of November 2026: established in the stated scope; the notebook is an expository occurrence, not a new open problem.}
\GSfield{Assessment}{literal statement: correct; publication match: exact; verification: read the result; development: none; confidence: high}
\end{GSitem}
\begin{GSitem}{1974-358.4}{E}{The Markov-chain comparison}
\GSfield{Kind; source}{definition; \GSpage{P000596}}
\GSfield{The notes}{The notes write \(\Pr(x_n\mid x_k,\ k\le n)=\Pr(x_n\mid x_{n-1})\).}
\GSfield{Statement}{The Markov property is \(\Pr(X_n\in A\mid\sigma(X_k:k<n))=\Pr(X_n\in A\mid X_{n-1})\) almost surely. The conditioning must exclude \(X_n\) itself; conditioning on \(\sigma(X_k:k\le n)\) gives the indicator of \(X_n\in A\), and does not express the usual Markov property.}
\GSfield{Fate}{The chain analogy in \citet[Section 1]{dobrushin1968conditional} motivates finite-range spatial conditional distributions. The scan clearly has \(k\le n\), so the inequality is a faithful notebook error, not an uncertain transcription.}
\GSfield{Present status}{As of November 2026: established in the stated scope; the notebook is an expository occurrence, not a new open problem.}
\GSfield{Assessment}{literal statement: false as written; publication match: partial; verification: read the result; development: none; confidence: high}
\end{GSitem}
\begin{GSitem}{1974-358.5}{E}{Construct all fields with prescribed exterior conditional laws}
\GSfield{Kind; source}{problem; \GSpage{P000596}}
\GSfield{The notes}{The problem is to prescribe, for each finite \(V\subset T\), the conditional distribution in \(V\) given the configuration in \(T\setminus V\), then construct the corresponding field measures.}
\GSfield{Statement}{Given a proper consistent family of probability kernels \((\gamma_V)_{V\Subset T}\) on \(X^T\), determine the set \(\mathcal G(\gamma)=\{\mu:\mu\gamma_V=\mu\text{ for every finite }V\}\), including whether it is nonempty and when it has one element. For finite \(X\) and countable \(T\), quasilocality of a specification suffices for nonemptiness. A contraction condition such as 1974-355.11 gives uniqueness; coexistence shows it is not automatic.}
\GSfield{Fate}{\citet[Section 1 and Theorems 1--2]{dobrushin1968conditional}, received 20 April 1967 and published in 1968, formulates this exact inverse problem and proves these sufficient existence/uniqueness results. The 1975 statement is an expository programme; it is not a claim that every compatible measurable specification has a universal simple classification.}
\GSfield{Present status}{As of November 2026: the finite-alphabet quasilocal existence and weak-dependence uniqueness statements are established. The general classification remains a model-dependent programme, not a single unresolved conjecture extracted from this entry.}
\GSfield{Assessment}{literal statement: correct with missing hypotheses; publication match: exact; verification: read the result; development: none; confidence: high; problem \hyperref[prob:GP-1974-20]{GP-1974-20}}
\end{GSitem}
\begin{GSitem}{1974-358.6}{E}{Regular conditional probability as a measurable kernel}
\GSfield{Kind; source}{definition; \GSpage{P000596}}
\GSfield{The notes}{For a sub-sigma-algebra, conditional probabilities are probability measures depending measurably on the conditioning data and satisfying an integral identity.}
\GSfield{Statement}{A regular conditional probability of \(\mu\) given \(\mathcal B\subset\mathcal A\) is a probability kernel \(K(\omega,A)\) such that \(K(\cdot,A)\) is \(\mathcal B\)-measurable and \(\int_BK(\omega,A)\,\mu(d\omega)=\mu(A\cap B)\) for every \(A\in\mathcal A,B\in\mathcal B\). For each fixed \(A\), it is a version of \(\mathbb E_\mu[1_A\mid\mathcal B]\).}
\GSfield{Fate}{This is the kernel form of the canonical systems in \citet[Section 3.1, pp.~125--129]{rokhlin1949measure}. The defining identities are meaningful independently of whether such a probability-valued version exists on an arbitrary space.}
\GSfield{Present status}{As of November 2026: established in the stated scope; the notebook is an expository occurrence, not a new open problem.}
\GSfield{Assessment}{literal statement: correct; publication match: partial; verification: read the result; development: none; confidence: high}
\end{GSitem}
\begin{GSitem}{1974-358.7}{E}{Existence and almost-everywhere uniqueness of conditional kernels}
\GSfield{Kind; source}{theorem; \GSpage{P000596}}
\GSfield{The notes}{In a broad class of cases, regular conditional probabilities exist and are unique.}
\GSfield{Statement}{For a probability on a standard Borel space and conditioning by a measurable map to another standard Borel space, a disintegration exists and is unique outside a set of conditioning values of pushforward measure zero. In equivalent countably generated conditioning formulations the same almost-everywhere qualification applies. Pointwise uniqueness on null fibres is false.}
\GSfield{Fate}{\citet[Section 3.1, pp.~125--129]{rokhlin1949measure} proves existence for measurable partitions of Lebesgue spaces and essential uniqueness. Received 15 July 1947, it appeared in Russian in 1949. The notebook's broad-class phrase is correct but supplies neither the space hypotheses nor the null-set qualification.}
\GSfield{Present status}{As of November 2026: established in the stated scope; the notebook is an expository occurrence, not a new open problem.}
\GSfield{Assessment}{literal statement: correct with missing hypotheses; publication match: partial; verification: read the result; development: none; confidence: high}
\end{GSitem}
\begin{GSitem}{1974-358.8}{E}{Fubini’s theorem as conditional measure on a rectangle}
\GSfield{Kind; source}{remark; \GSpage{P000596}}
\GSfield{The notes}{For the square with planar Lebesgue measure, conditioning on a vertical coordinate gives Lebesgue measure along the corresponding segment.}
\GSfield{Statement}{On \([0,1]^2\) with product Lebesgue probability and projection \((x,y)\mapsto x\), a disintegration is \(K_x=\delta_x\otimes dy\). Hence \(\int F(x,y)\,dx\,dy=\int(\int F(x,y)\,dy)\,dx\) for integrable \(F\). The explicit kernel can be chosen for every \(x\), although uniqueness is only almost everywhere.}
\GSfield{Fate}{This is the elementary product example of \citet[Section 3.1]{rokhlin1949measure}; it explains the general-Fubini remark already made by Gelfand in 1974-355.5.}
\GSfield{Present status}{As of November 2026: established in the stated scope; the notebook is an expository occurrence, not a new open problem.}
\GSfield{Assessment}{literal statement: correct; publication match: partial; verification: read the result; development: none; confidence: high}
\end{GSitem}
\begin{GSitem}{1974-358.9}{E}{Conditioning on a map or a countable family of observables}
\GSfield{Kind; source}{construction; \GSpage{P000596}}
\GSfield{The notes}{A measurable \(\xi:X\to F\) defines the conditioning partition; a countable collection of real functions can be treated similarly.}
\GSfield{Statement}{For standard Borel \(X,F\) and measurable \(\xi\), define \(\nu=\xi_*\mu\). Conditional kernels \(K_y\) satisfy \(\mu(A)=\int_FK_y(A)\,\nu(dy)\) and are concentrated on \(\xi^{-1}(y)\) for \(\nu\)-almost every \(y\). A countable sequence \((\xi_j)\) is a single map into the standard Borel product \(\mathbb R^{\mathbb N}\).}
\GSfield{Fate}{The measurable partitions generated by one or countably many real functions are explicit examples in \citet[Section 3.1, p.~126]{rokhlin1949measure}. This closes the foundational loop needed for Gibbs exterior conditioning.}
\GSfield{Present status}{As of November 2026: established in the stated scope; the notebook is an expository occurrence, not a new open problem.}
\GSfield{Assessment}{literal statement: correct with missing hypotheses; publication match: exact; verification: read the result; development: none; confidence: high}
\end{GSitem}

\GStalk{359}{1975-03-31}{A. M. Polyakov}{P000596-P000599}{Spontaneously broken symmetry and particles as field configurations}
Only P000597 belongs to the actual session. It contains a title, a particle metaphor, and names: Skyrme, Nielsen, Olesen (written Olsen), Faddeev, Polyakov and 't Hooft. There are no equations or hypotheses from which a specific soliton theorem can be recovered. The final sentence announces Nevanlinna theory following Griffiths for the next meeting.

\begin{GSitem}{1974-359.1}{F}{Particle-like structures in theories with broken symmetry}
\GSfield{Kind; source}{remark; \GSpage{P000597}}
\GSfield{The notes}{Polyakov discusses spontaneously broken symmetry and describes particles as bubbles in the ether of quantum field theory, followed by a list of researchers.}
\GSfield{Statement}{The recoverable mathematical theme is that nonlinear field configurations may carry particle-like localized energy and nontrivial topology. The note gives no field space, gauge group, energy functional, boundary conditions, stability criterion, or topological charge, and therefore specifies no particular existence or classification result.}
\GSfield{Fate}{Polyakov's \citet{polyakov1974particle} is an obvious contemporary primary lead: its 1974 publication predates this talk. Nielsen--Olesen vortices are another distinct lead. The names and metaphor do not license importing either complete model into the notebook. The institutional publication record verifies Polyakov's title and date, but an exact notebook-to-theorem match cannot be established.}
\GSfield{Present status}{As of November 2026: the broad historical theme is identifiable; the specific mathematical statement and its fate are not. No open problem is created from this three-line record.}
\GSfield{Assessment}{literal statement: not determinable; publication match: thematic only; verification: publisher metadata checked; development: none; confidence: low}
\end{GSitem}

\GStalk{360}{1975-04-07}{James King (attribution marked ? in the notebook)}{P000598-P000602}{Residue currents, divisors, and the first main theorem}
The actual session occupies P000598--P000601; P000602 begins the following dated continuation. The notebook itself places a question mark after King, although its closing bibliographical pointer in the continuation identifies Griffiths--King's Acta paper. The talk passes from logarithmic residues and integration currents to the divisor--Chern-class identity, Bézout, growth of holomorphic maps, and the first main theorem. Its mathematical content is already in \citet[Sections 1 and 5]{griffithsking1973nevanlinna}, published in 1973. The raw graph-volume integrals and the integrated Nevanlinna characteristic must be distinguished.

\begin{GSitem}{1974-360.1}{E}{Two programmes for holomorphic mapping theory}
\GSfield{Kind; source}{remark; \GSpage{P000598}}
\GSfield{The notes}{The opening lists hyperbolic analysis, in which negative curvature restricts holomorphic maps, and counting volumes of divisors by residue calculus.}
\GSfield{Statement}{These are two announced methods, not two quantified conjectures: differential inequalities for maps into targets with suitably negative curvature, and current identities that compare divisor counting with growth. The first is made precise in record 361; the second is developed in the present talk.}
\GSfield{Fate}{\citet[Sections 1, 5, and Appendix]{griffithsking1973nevanlinna} contains both parts of this programme. The note does not assert that every notion of negative curvature yields the same mapping restriction.}
\GSfield{Present status}{As of November 2026: established in the stated scope; the notebook is an expository occurrence, not a new open problem.}
\GSfield{Assessment}{literal statement: correct; publication match: partial; verification: read the result; development: none; confidence: high}
\end{GSitem}
\begin{GSitem}{1974-360.2}{E}{Logarithmic residue as integration around a normal circle}
\GSfield{Kind; source}{construction; \GSpage{P000598}}
\GSfield{The notes}{The notes write \(\eta=(dg/2\pi ig)\wedge\zeta+\xi\), then \(\lim_{\varepsilon\to0}\int_{C(\varepsilon)}\eta=\zeta|_D\), calling \(\eta\) an arbitrary form off \(D=\{g=0\}\).}
\GSfield{Statement}{If \(D\) is a smooth reduced hypersurface, \(dg\ne0\) along \(D\), and \(\zeta,\xi\) are smooth forms of degrees \(r-1,r\), the displayed logarithmic form has residue \(\zeta|_D\). The limit is integration along the positively oriented normal circle fibres of a shrinking tube. Over an \((r-1)\)-cycle \(c\subset D\), it reads \(\lim_{\varepsilon\to0}\int_{\operatorname{tube}_\varepsilon(c)}\eta=\int_c\zeta\). An arbitrary form on the complement need not admit this logarithmic decomposition.}
\GSfield{Reconstruction}{Fibre integration distinguishes the form-valued residue from a scalar integral over a whole cycle.}
\GSfield{Fate}{The local residue/current method is the one used in \citet[Section 1(a), especially Proposition 1.1]{griffithsking1973nevanlinna}. The tube description is a standard local interpretation of the logarithmic model; the notebook suppresses the smooth-divisor and simple-pole assumptions.}
\GSfield{Present status}{As of November 2026: established in the stated scope; the notebook is an expository occurrence, not a new open problem.}
\GSfield{Assessment}{literal statement: correct with missing hypotheses; publication match: partial; verification: read the result; development: none; confidence: high}
\end{GSitem}
\begin{GSitem}{1974-360.3}{E}{Currents and their distributional exterior derivative}
\GSfield{Kind; source}{definition; \GSpage{P000598}}
\GSfield{The notes}{A current is a functional on compactly supported forms; an \(L^1_{\mathrm{loc}}\) form \(\lambda\) gives \([\lambda](\varphi)=\int\lambda\wedge\varphi\), and \(dT(\varphi)=\pm T(d\varphi)\).}
\GSfield{Statement}{On an oriented real \(2n\)-manifold, a degree-\(r\) current is a continuous linear functional on compactly supported smooth \((2n-r)\)-forms. For a locally integrable \(r\)-form, \([\lambda]\) is defined by integration and \(dT(\varphi)=(-1)^{r+1}T(d\varphi)\). A smooth oriented submanifold defines an integration current of degree its real codimension.}
\GSfield{Fate}{These conventions are introduced in \citet[Section 1(a)]{griffithsking1973nevanlinna}. The missing continuity and degree convention are supplied; the notebook explicitly leaves the sign open rather than giving a wrong fixed sign.}
\GSfield{Present status}{As of November 2026: established in the stated scope; the notebook is an expository occurrence, not a new open problem.}
\GSfield{Assessment}{literal statement: correct with missing hypotheses; publication match: exact; verification: read the result; development: none; confidence: high}
\end{GSitem}
\begin{GSitem}{1974-360.4}{E}{The residue current measures the singular part of differentiation}
\GSfield{Kind; source}{definition; \GSpage{P000598}}
\GSfield{The notes}{For a locally integrable form smooth off the exceptional set, \(\operatorname{Res}\eta=d[\eta]-[d\eta]=[\zeta|_D]\).}
\GSfield{Statement}{Suppose \(\eta\in L^1_{\mathrm{loc}}\) is smooth off \(D\), and its ordinary derivative there also has a locally integrable extension \(d_{\mathrm{reg}}\eta\). Then \(d[\eta]-[d_{\mathrm{reg}}\eta]\) is supported on \(D\). For the logarithmic model of 1974-360.2 it equals the pushforward \(i_*[\zeta|_D]\), where \(i:D\hookrightarrow V\). The second integrability requirement is necessary to make the difference meaningful in this form.}
\GSfield{Reconstruction}{The notation distinguishes the current derivative from the ordinary derivative on the complement.}
\GSfield{Fate}{The Poincaré--Lelong current calculation of \citet[Proposition 1.1 and its proof]{griffithsking1973nevanlinna} gives the model. No assertion that every locally integrable form has a locally integrable ordinary derivative is valid.}
\GSfield{Present status}{As of November 2026: established in the stated scope; the notebook is an expository occurrence, not a new open problem.}
\GSfield{Assessment}{literal statement: correct with missing hypotheses; publication match: partial; verification: read the result; development: none; confidence: high}
\end{GSitem}
\begin{GSitem}{1974-360.5}{E}{The one-variable residue normalization}
\GSfield{Kind; source}{method; \GSpage{P000598}}
\GSfield{The notes}{For \(\eta=dz/z\), the notes give \(d[\eta](\varphi)-[\eta](d\varphi)=\lim\int\varphi\,dz=\varphi(0)\).}
\GSfield{Statement}{With the current convention in 1974-360.3 and counterclockwise circles, \(d[dz/z]=2\pi i\,\delta_0\), while the ordinary derivative of \(dz/z\) is zero off the origin. Equivalently \(d[dz/(2\pi iz)]=\delta_0\). The contour integrand is \(\varphi\,dz/z\). For a degree-one current, \(d[\eta](\varphi)=[\eta](d\varphi)\), so the notebook's written left-hand difference is identically zero.}
\GSfield{Fate}{This is the elementary one-variable calculation underlying \citet[Proposition 1.1]{griffithsking1973nevanlinna}. The scan confirms the missing denominator/normalization and the wrong second term; these are faithful notebook errors.}
\GSfield{Present status}{As of November 2026: established in the stated scope; the notebook is an expository occurrence, not a new open problem.}
\GSfield{Assessment}{literal statement: false as written; publication match: partial; verification: read the result; development: none; confidence: high}
\end{GSitem}
\begin{GSitem}{1974-360.6}{E}{The divisor residue includes vanishing multiplicities}
\GSfield{Kind; source}{theorem; \GSpage{P000598}}
\GSfield{The notes}{For holomorphic \(g\), \(\operatorname{Res}[dg/(2\pi ig)]=[Z]\), with \(Z=\sum m_iZ_i\) and \(m_i=\operatorname{ord}_{Z_i}g\).}
\GSfield{Statement}{If \(g\) is a nonzero holomorphic function on a complex manifold, its divisor is \(\operatorname{div}(g)=\sum_i m_iZ_i\). As currents, \(dd^c\log|g|^2=\sum_i m_i[Z_i]\), with \(d^c=(i/4\pi)(\bar\partial-\partial)\). Equivalently the logarithmic differential has this residue current. The weights are part of the divisor; the reduced zero set alone loses them.}
\GSfield{Fate}{\citet[Proposition 1.1, equation (1.2)]{griffithsking1973nevanlinna} is an exact Poincaré--Lelong match. The statement includes singular zero sets, via their integration currents.}
\GSfield{Present status}{As of November 2026: established in the stated scope; the notebook is an expository occurrence, not a new open problem.}
\GSfield{Assessment}{literal statement: correct; publication match: exact; verification: read the result; development: none; confidence: high}
\end{GSitem}
\begin{GSitem}{1974-360.7}{E}{A divisor determines a line bundle and canonical section}
\GSfield{Kind; source}{construction; \GSpage{P000599}}
\GSfield{The notes}{Local equations \(g_i\) for a codimension-one \(X\subset V\) give ratios \(f_{ij}=g_i/g_j\), a bundle \(L_X\), and a section \(\sigma=\{g_i\}\) vanishing on \(X\).}
\GSfield{Statement}{For an effective Cartier divisor \(D\) on a complex manifold, local equations with invertible holomorphic ratios define \(\mathcal O(D)\) and its canonical holomorphic section \(\sigma\), whose divisor is \(D\). A bare subset of codimension one is insufficient: it must be an analytic divisor with its multiplicities. Frame-transition conventions determine whether one writes the ratio or its reciprocal.}
\GSfield{Fate}{The divisor and metric setup used by \citet[Sections 0(a) and 1(c)]{griffithsking1973nevanlinna} is exactly this construction. The notebook suppresses the analytic/Cartier and transition-convention qualifications.}
\GSfield{Present status}{As of November 2026: established in the stated scope; the notebook is an expository occurrence, not a new open problem.}
\GSfield{Assessment}{literal statement: correct with missing hypotheses; publication match: exact; verification: read the result; development: none; confidence: high}
\end{GSitem}
\begin{GSitem}{1974-360.8}{E}{The Hermitian Poincaré--Lelong identity and its sign}
\GSfield{Kind; source}{theorem; \GSpage{P000599}}
\GSfield{The notes}{For \(\|\sigma\|^2=h_i|g_i|^2\), the notes put \(\eta=(2\pi i)^{-1}\partial\log\|\sigma\|^2\), identify \([d\eta]\) with \(c(L_D)\), and later write \([D]=d[\eta]+[c(L_D)]\).}
\GSfield{Statement}{Fix \(d^c=(i/4\pi)(\bar\partial-\partial)\), so \(dd^c=(i/2\pi)\partial\bar\partial\). If \(h=\|e\|^2\) in a holomorphic frame, \(c_1(L,h)=-dd^c\log h\). For a nonzero holomorphic section, \(dd^c\log\|\sigma\|^2=[D]-c_1(L,h)\). Thus the ordinary \(d\eta\) off \(D\), for the notebook's \(\eta\), is \(-c_1(L,h)\), whereas the current derivative has the additional \([D]\) term.}
\GSfield{Fate}{\citet[Proposition 1.14]{griffithsking1973nevanlinna} gives the global current identity. The notebook's two sign assignments are inconsistent under a fixed Chern-form convention; the later \([D]=d[\eta]+c_1\) has the standard sign.}
\GSfield{Present status}{As of November 2026: established in the stated scope; the notebook is an expository occurrence, not a new open problem.}
\GSfield{Assessment}{literal statement: false as written; publication match: partial; verification: read the result; development: none; confidence: high}
\end{GSitem}
\begin{GSitem}{1974-360.9}{E}{The divisor class is Poincaré dual to the first Chern class}
\GSfield{Kind; source}{theorem; \GSpage{P000599}}
\GSfield{The notes}{The current de Rham theorem identifies the zero-set current with the Chern class; for closed \(\psi\), \(\int_D\psi=\int_Vc(L_D)\wedge\psi\).}
\GSfield{Statement}{For a divisor \(D\) on a compact complex \(n\)-manifold, its current cohomology class equals \(c_1(\mathcal O(D))\). Consequently, for a closed smooth \((2n-2)\)-form \(\psi\), \(\int_D\psi=\int_Vc_1(\mathcal O(D),h)\wedge\psi\), counting divisor multiplicities. On a noncompact manifold, take compactly supported test forms or otherwise justify all integrals and boundary terms.}
\GSfield{Fate}{This is the cohomological interpretation immediately after \citet[Proposition 1.14, pp.~162--163]{griffithsking1973nevanlinna}. It is already established before the talk.}
\GSfield{Present status}{As of November 2026: established in the stated scope; the notebook is an expository occurrence, not a new open problem.}
\GSfield{Assessment}{literal statement: correct with missing hypotheses; publication match: exact; verification: read the result; development: none; confidence: high}
\end{GSitem}
\begin{GSitem}{1974-360.10}{E}{Bézout by products of divisor classes}
\GSfield{Kind; source}{theorem; \GSpage{P000599}}
\GSfield{The notes}{For two curves in \(\mathbb P^2\), the notes identify the number of intersection points with the product of Chern classes and obtain \(d_Xd_Y\).}
\GSfield{Statement}{For plane projective curves \(X,Y\) over \(\mathbb C\) with no common irreducible component, the sum of local intersection multiplicities is \(\deg X\,\deg Y\). With a hyperplane form \(\omega\) normalized by \(\int_{\mathbb P^2}\omega^2=1\), the cohomological product is \((\deg X)(\deg Y)\int\omega^2\). It counts distinct points only when every intersection is transverse. The product of arbitrary singular currents is not automatically defined; here one uses proper algebraic intersection or an appropriate regularization.}
\GSfield{Fate}{The divisor-current and proper-intersection formalism of \citet[Section 1(c), Propositions 1.14--1.15]{griffithsking1973nevanlinna} supplies the argument. The notebook omits properness and multiplicities in the phrase 'number of points'. This is classical exposition, not a new proof claiming priority.}
\GSfield{Present status}{As of November 2026: established in the stated scope; the notebook is an expository occurrence, not a new open problem.}
\GSfield{Assessment}{literal statement: correct with missing hypotheses; publication match: partial; verification: read the result; development: none; confidence: high}
\end{GSitem}
\begin{GSitem}{1974-360.11}{E}{Integration currents on singular analytic sets}
\GSfield{Kind; source}{theorem; \GSpage{P000600}}
\GSfield{The notes}{The notes say the theorem remains valid for \(X\) with singularities, followed by 'Lelong, 1953'. }
\GSfield{Statement}{For a pure-dimensional analytic set \(X\), integration over its regular locus has locally finite mass and extends to a closed positive current \([X]\) on the ambient manifold. Analytic cycles are handled by summing these currents with their multiplicities. This justifies the divisor-current formulas even when the divisor is singular; it does not define arbitrary products of such currents.}
\GSfield{Fate}{\citet[Theorems 7--8, pp.~257--258]{lelong1957integration} proves the integration-current extension and additivity; the paper is from 1957. Its bibliography includes a 1953 Lelong contribution, but the notebook's precise 1953 attribution was not verified. The checked 1957 publication is sufficient to establish that the result predates the seminar.}
\GSfield{Present status}{As of November 2026: established in the stated scope; the notebook is an expository occurrence, not a new open problem.}
\GSfield{Assessment}{literal statement: correct with missing hypotheses; publication match: partial; verification: read the result; development: none; confidence: medium}
\end{GSitem}
\begin{GSitem}{1974-360.12}{E}{Pullback volume counts the multiplicity of a holomorphic map}
\GSfield{Kind; source}{construction; \GSpage{P000600}}
\GSfield{The notes}{For \(f:\mathbb C^n\to V\) with compact Kähler target, \(\operatorname{vol}(f(B(r)))=\int_{B(r)}f^*(\omega^n)\).}
\GSfield{Statement}{The integral \(n!^{-1}\int_{B(r)}f^*\omega^n\) is the parametrized \(2n\)-dimensional volume of \(f\), counting local covering multiplicity. It equals the ordinary image volume only under suitable almost-everywhere injectivity; it vanishes where the complex rank is below \(n\). The factor \(n!\) is omitted if volume is defined using \(\omega^n\) itself.}
\GSfield{Fate}{\citet[Section 5(a)--(b)]{griffithsking1973nevanlinna} uses these pullback growth integrals. The notebook's equality needs the multiplicity interpretation; it is not an identity for unweighted images of arbitrary maps.}
\GSfield{Present status}{As of November 2026: established in the stated scope; the notebook is an expository occurrence, not a new open problem.}
\GSfield{Assessment}{literal statement: correct with missing hypotheses; publication match: partial; verification: read the result; development: none; confidence: high}
\end{GSitem}
\begin{GSitem}{1974-360.13}{E}{Graph volume and the mixed raw growth integrals}
\GSfield{Kind; source}{construction; \GSpage{P000600}}
\GSfield{The notes}{For the graph above \(B(r)\), the notes expand \(\int_{\Gamma(r)}(\varphi+\omega)^n\) and define \(T_k(f,r)=\int_{B(r)}\varphi^{n-k}\wedge f^*(\omega^k)\).}
\GSfield{Statement}{Let \(\omega_0\) be the Euclidean Kähler form and put \(t_k(f,r)=\int_{B(r)}\omega_0^{n-k}\wedge f^*\omega^k\), \(0\le k\le n\). Then \(n!\operatorname{Vol}(\Gamma_f|_{B(r)})=\sum_{k=0}^n\binom nk t_k(f,r)\) for the product metric. These are raw ball integrals; a Nevanlinna characteristic is obtained by an additional radial integration and normalization.}
\GSfield{Reconstruction}{The lower-case notation is editorial, explicitly separating two quantities denoted T in the notebook.}
\GSfield{Fate}{The graph-growth method and higher characteristic forms are discussed in \citet[Section 5, pp.~180--183]{griffithsking1973nevanlinna}. Changing the notebook's raw \(T_k\) to \(t_k\) here prevents confusion with the integrated characteristic on the next page.}
\GSfield{Present status}{As of November 2026: established in the stated scope; the notebook is an expository occurrence, not a new open problem.}
\GSfield{Assessment}{literal statement: correct; publication match: partial; verification: read the result; development: none; confidence: high}
\end{GSitem}
\begin{GSitem}{1974-360.14}{F}{The unspecified comparison of higher growth with first growth}
\GSfield{Kind; source}{claim; \GSpage{P000600}}
\GSfield{The notes}{The notes say 'in many cases' \(T_k=O(T_1)\), adding 'of one order'. }
\GSfield{Statement}{No family of maps, target hypothesis, range of \(k\), or choice between raw and integrated characteristics is specified. Therefore the phrase gives no definite universal estimate. The graph-volume identity by itself does not imply domination of every higher mixed integral by the first.}
\GSfield{Fate}{\citet[Section 5(b), pp.~182--183]{griffithsking1973nevanlinna} discusses the different higher growth quantities and the difficulty of reducing their information to first growth. This is thematic support for the surrounding subject, not verification of an unqualified \(O\)-estimate. No exact source or modern open question is assigned to 'many cases'. }
\GSfield{Present status}{As of November 2026: the intended restricted comparison is not identifiable; it is not classified as an undeveloped general theorem.}
\GSfield{Assessment}{literal statement: not determinable; publication match: thematic only; verification: read the result; development: none; confidence: low}
\end{GSitem}
\begin{GSitem}{1974-360.15}{E}{Pulling a divisor identity back by a holomorphic map}
\GSfield{Kind; source}{construction; \GSpage{P000600}}
\GSfield{The notes}{For a positive line bundle with \(c(L)=\omega\) and \(D=\{\sigma=0\}\), the notes write \(d[\eta]-[\omega]=[D]\) and the same formula after pullback by \(f\).}
\GSfield{Statement}{Let \(L\to V\) have a smooth Hermitian metric with positive curvature form \(\omega=c_1(L,h)\), and let \(f:\mathbb C^n\to V\) have image not contained in the zero divisor \(D\) of \(\sigma\). The pullback section \(f^*\sigma\) defines \(D_f=f^*D\), with multiplicity, and \(dd^c\log\|f^*\sigma\|^2=[D_f]-f^*\omega\). Thus \([D_f]=dd^c\log\|f^*\sigma\|^2+f^*\omega\). This is justified by applying Poincaré--Lelong to the pullback section, not by assuming that every current has a pullback under every map.}
\GSfield{Fate}{\citet[Proposition 1.14 and Section 5(b)]{griffithsking1973nevanlinna} gives precisely the identity used in the first main theorem. The minus sign in the notebook's rearranged formula conflicts with its correct formula on P000599.}
\GSfield{Present status}{As of November 2026: established in the stated scope; the notebook is an expository occurrence, not a new open problem.}
\GSfield{Assessment}{literal statement: false as written; publication match: partial; verification: read the result; development: none; confidence: high}
\end{GSitem}
\begin{GSitem}{1974-360.16}{E}{The missing boundary term in the counting calculation}
\GSfield{Kind; source}{method; \GSpage{P000600}}
\GSfield{The notes}{After multiplication by \(\chi_{B(r)}\varphi^{n-1}\), a displayed integral over apparently \(S(r)\) of \(\varphi^{n-1}\wedge f^*\omega\) is equated to divisor volume.}
\GSfield{Statement}{For almost every radius, the current identity gives \(\int_{D_f\cap B(r)}\omega_0^{n-1}=\int_{B(r)}f^*\omega\wedge\omega_0^{n-1}+\int_{S(r)}d^c\log\|f^*\sigma\|^2\wedge\omega_0^{n-1}\), interpreted by smooth approximation if needed. The first integral is over the ball, while the boundary integral has degree \(2n-1\) and contains the logarithmic derivative. Integrating the \(2n\)-form \(f^*\omega\wedge\omega_0^{n-1}\) over a sphere is dimensionally inappropriate.}
\GSfield{Reconstruction}{The sign follows the fixed ddc convention; the stated identity supplies the missing Stokes term.}
\GSfield{Fate}{This is the Stokes step in the twice-integrated Poincaré--Lelong argument of \citet[Proposition 5.11 and its reference to Section 3]{griffithsking1973nevanlinna}. The notebook has suppressed the essential boundary term; the faint domain letter does not repair that omission.}
\GSfield{Present status}{As of November 2026: established in the stated scope; the notebook is an expository occurrence, not a new open problem.}
\GSfield{Assessment}{literal statement: false as written; publication match: partial; verification: read the result; development: none; confidence: high}
\end{GSitem}
\begin{GSitem}{1974-360.17}{E}{Integrated divisor counting and the radius exponent}
\GSfield{Kind; source}{definition; \GSpage{P000600-P000601}}
\GSfield{The notes}{The notes pass from raw divisor volume to
\[N(D_f,r)=\int_0^r\Bigl(\int_{D_f\cap B(\rho)}\varphi^{n-1}\Bigr)
 \,\frac{d\rho}{\rho^{n+1}},\]
while the preceding line uses the classical \(d\rho/\rho\) notation.}
\GSfield{Statement}{Use ordinary Euclidean radius, \(\omega_0=(i/2)\sum_jdz_j\wedge d\bar z_j\), and \(dd^c=(i/2\pi)\partial\bar\partial\). Set \(n_D(t)=\int_{D_f\cap B(t)}\omega_0^{n-1}\). For fixed \(r_0>0\) and \(r\ge r_0\), define \(N_f(D,r)=\pi^{1-n}\int_{r_0}^r n_D(t)t^{-(2n-1)}\,dt\) and \(T_f(r)=\pi^{1-n}\int_{r_0}^r(\int_{B(t)}f^*\omega\wedge\omega_0^{n-1})t^{-(2n-1)}\,dt\). For \(n=1\), the weight is \(dt/t\). Equivalently one normalizes the raw masses by \(t^{2n-2}\) and then integrates with \(dt/t\).}
\GSfield{Reconstruction}{The Euclidean form, ordinary radius, base radius, and factor \(\pi^{1-n}\) are fixed editorially. They follow from \(\int_{B(t)}dd^c u\wedge\omega_0^{n-1}=(\pi^{n-1}/2)t^{2n-1}(d/dt)\operatorname{Avg}_{S(t)}u\); this independently checks the dimension-one reduction.}
\GSfield{Fate}{The counting and characteristic definitions used in \citet[Sections 3(a) and 5(b)]{griffithsking1973nevanlinna} have the Jensen radial normalization. The notebook's \(\rho^{n+1}\) is not the Euclidean-radius exponent \(2n-1\) (they happen to agree at \(n=2\)); it already fails the one-dimensional check.}
\GSfield{Present status}{As of November 2026: established in the stated scope; the notebook is an expository occurrence, not a new open problem.}
\GSfield{Assessment}{literal statement: false as written; publication match: partial; verification: read the result; development: none; confidence: high}
\end{GSitem}
\begin{GSitem}{1974-360.18}{E}{The first main theorem bounds divisor counting by growth}
\GSfield{Kind; source}{theorem; \GSpage{P000601}}
\GSfield{The notes}{The final inequality is \(N(D_f,r)\le T(f,r)+O(1)\), called the first Nevanlinna inequality.}
\GSfield{Statement}{For the compact Hermitian target and nonzero pullback section above, scale \(\sigma\) so that \(\|\sigma\|\le1\). With the definitions of 1974-360.17, put \(m_f(D,r)=\operatorname{Avg}_{S(r)}\log(1/\|f^*\sigma\|)\ge0\), where the average uses Euclidean surface probability measure. The integrated Poincaré--Lelong identity is \(N_f(D,r)+m_f(D,r)=T_f(r)+m_f(D,r_0)\). Therefore \(N_f(D,r)\le T_f(r)+O(1)\), with a constant depending on the section, metric, and fixed \(r_0\), but not on \(r\).}
\GSfield{Fate}{\citet[Proposition 5.11 and equation (5.12), p.~184]{griffithsking1973nevanlinna} is the exact source, published in 1973, two years before this talk. This is an established first main theorem; no second-main-theorem or defect assertion is silently added. It connects thematically with the earlier Picard/Nevanlinna discussion 1965-71.1.}
\GSfield{Present status}{As of November 2026: established in the stated scope; the notebook is an expository occurrence, not a new open problem.}
\GSfield{Assessment}{literal statement: correct; publication match: exact; verification: read the result; development: none; confidence: high}
\end{GSitem}

\GStalk{361}{1975-04-14}{James King}{P000602}{Nevanlinna theory continued: the volume Schwarz lemma}
The dated continuation occupies P000602 only. It recalls divisor counting, then defines the curvature of a volume form and states an Ahlfors-type volume-decreasing theorem for maps from the ball or polydisc. The closing references name Griffiths and King, with an additional Griffiths--Cowen note and an uncertain eigenfunction-related remark. The precise comparison argument predates the talk in \citet[Section 2, Proposition 2.7]{griffiths1971canonical}; the final fragment does not specify a spectral problem.

\begin{GSitem}{1974-361.1}{E}{The first main inequality recalled with its constant suppressed}
\GSfield{Kind; source}{theorem; \GSpage{P000602}}
\GSfield{The notes}{The continuation repeats \(N(D_f,r)\le T(f,r)\) and labels an iterated integral the order of growth; its inner ball is written \(B(r)\) although the outer variable is \(t\).}
\GSfield{Statement}{Under the hypotheses and compatible counting definitions of 1974-360.17--18, the invariant statement is \(N_f(D,r)\le T_f(r)+O(1)\). A chosen additive normalization of \(T_f\) may absorb the bounded term for a fixed divisor, but the term cannot simply be dropped for arbitrary normalizations. The inner ball in a radial growth integral is \(B(t)\), and raw ball masses require the dimension-dependent normalization described there.}
\GSfield{Fate}{\citet[Proposition 5.11 and equation (5.12)]{griffithsking1973nevanlinna} proves the recalled result. This is a second occurrence of the same established first main theorem, not a new strengthened inequality.}
\GSfield{Present status}{As of November 2026: established in the stated scope; the notebook is an expository occurrence, not a new open problem.}
\GSfield{Assessment}{literal statement: correct with missing hypotheses; publication match: partial; verification: read the result; development: none; confidence: high}
\end{GSitem}
\begin{GSitem}{1974-361.2}{E}{The sign and normalization of the Ricci form of a volume}
\GSfield{Kind; source}{definition; \GSpage{P000602}}
\GSfield{The notes}{First \(\operatorname{Ric}\Omega=i\partial\bar\partial\log g\) for \(\Omega=(i/2)g\,dz\wedge d\bar z\); then the general formula has \(i/(2\pi)\). The text says its scalar coefficient is minus Gaussian curvature.}
\GSfield{Statement}{For a positive smooth volume form locally \(\Omega=a\,dV_{\mathrm{hol}}\), define \(R_\Omega=dd^c\log a=(i/2\pi)\partial\bar\partial\log a\). The form is invariant under holomorphic coordinate changes, because the logarithm of the Jacobian modulus squared is pluriharmonic. This is the negative of the usual normalized Ricci form. In one complex dimension, with metric area \(\Omega=(i/2)g\,dz\wedge d\bar z\) and Gaussian curvature \(K_g\), \(i\partial\bar\partial\log g=-K_g\Omega\), while \(R_\Omega=-K_g\Omega/(2\pi)\).}
\GSfield{Reconstruction}{The notation R is editorial and makes explicit that the notebook uses the negative Ricci convention.}
\GSfield{Fate}{The volume-form curvature convention and coordinate invariance are used in \citet[Section 2]{griffiths1971canonical}. The notebook has the correct negative-Gaussian sign, but shifts the normalization by \(2\pi\) between two displays; those displays must not be combined without that factor.}
\GSfield{Present status}{As of November 2026: established in the stated scope; the notebook is an expository occurrence, not a new open problem.}
\GSfield{Assessment}{literal statement: correct with missing hypotheses; publication match: partial; verification: read the result; development: none; confidence: high}
\end{GSitem}
\begin{GSitem}{1974-361.3}{E}{The volume-curvature inequalities needed for comparison}
\GSfield{Kind; source}{definition; \GSpage{P000602}}
\GSfield{The notes}{The two explicit assumptions are \(\operatorname{Ric}\Omega>0\) and \((\operatorname{Ric}\Omega)^n\ge\Omega\).}
\GSfield{Statement}{On a complex \(n\)-manifold with a positive \(C^2\) volume form \(\Omega\), assume \(R_\Omega\) is a positive definite real \((1,1)\)-form and \(R_\Omega^n\ge\Omega\) as positive top-degree forms, in the same \(dd^c\) convention. This determinant inequality fixes the necessary scale. Multiplication of \(\Omega\) by a positive constant leaves \(R_\Omega\) unchanged but changes the comparison inequality.}
\GSfield{Fate}{\citet[Lemma 2.3 and the proof of Proposition 2.7]{griffiths1971canonical} supplies volume forms with such normalized curvature control in the canonical algebraic setting and the corresponding maximum-principle argument. The notebook's explicit analytic inequalities are more informative than its later shorthand 'negative curvature'. }
\GSfield{Present status}{As of November 2026: established in the stated scope; the notebook is an expository occurrence, not a new open problem.}
\GSfield{Assessment}{literal statement: correct; publication match: partial; verification: read the result; development: none; confidence: high}
\end{GSitem}
\begin{GSitem}{1974-361.4}{E}{Ahlfors-type volume decrease from the ball or polydisc}
\GSfield{Kind; source}{theorem; \GSpage{P000602}}
\GSfield{The notes}{With the Bergman volume on \(B^n\) or \(\Delta^n\), every holomorphic map to \(V^n\) with negatively curved volume satisfies \(f^*\Omega_V\le\Phi_B\).}
\GSfield{Statement}{Normalize the complete invariant volume \(\Phi_U\) on the unit ball or polydisc \(U\) so that \(R_{\Phi_U}>0\) and \(R_{\Phi_U}^n=\Phi_U\). Let \(V\) have complex dimension \(n\). If a positive \(C^2\) volume form \(\Omega_V\) satisfies \(R_{\Omega_V}>0\) and \(R_{\Omega_V}^n\ge\Omega_V\), every holomorphic \(f:U\to V\) satisfies \(f^*\Omega_V\le\Phi_U\). At rank-deficient points the left side vanishes. The comparison follows by applying the maximum principle to the ratio \(f^*\Omega_V/\Phi_{U_s}\) on smaller balls or polydiscs \(U_s\) and letting \(s\) increase to one.}
\GSfield{Reconstruction}{The analytic curvature inequalities explicitly stated in the notes are used as hypotheses; the checked publication has the canonical-algebraic target as its stated setting. This is a faithful proof reconstruction, not a claim of a new generalization.}
\GSfield{Fate}{\citet[Proposition 2.7, pp.~444--445]{griffiths1971canonical}, received 27 March 1970 and published in May 1971, proves volume decrease for maps to complete canonical algebraic manifolds. Its proof gives the displayed analytic sufficient-condition version. The notebook's blanket curvature-\(-1\) description must be replaced by compatible source and target volume normalizations; the polydisc does not have constant holomorphic sectional curvature.}
\GSfield{Present status}{As of November 2026: established in the stated scope; the notebook is an expository occurrence, not a new open problem.}
\GSfield{Assessment}{literal statement: correct with missing hypotheses; publication match: partial; verification: read the result; development: none; confidence: high}
\end{GSitem}
\begin{GSitem}{1974-361.5}{F}{The unidentifiable eigenfunction-related closing fragment}
\GSfield{Kind; source}{remark; \GSpage{P000602}}
\GSfield{The notes}{After the Acta reference, the page has uncertain names resembling 'Herman \(\Rightarrow\) Joahim Weil', then \ru{Нетрив. формулир. нужна в теории соб. ф-й.} and 'Griffits, Cowen'. }
\GSfield{Statement}{The fragment contains no operator, eigenvalue problem, space of functions, proposed inequality, or quantified assertion. It cannot support a unique spectral-theoretic reconstruction, nor identify what a nontrivial formulation would say.}
\GSfield{Reconstruction}{The Weyl identification is a contextual inference from a checked primary bibliography and the re-examined scan. It does not silently replace the uncertain Russian phrase about eigenfunctions with the theory of entire functions.}
\GSfield{Fate}{The Acta pointer identifies \citet{griffithsking1973nevanlinna}. A strong bibliographical lead for the next line is \citet[Introduction, pp.~93--95]{cowengriffiths1976curves}, published in 1976, which discusses H. and J. Weyl's work on holomorphic curves. Thus the names likely mean Hermann and Joachim Weyl, and the final Cowen reference has a concrete thematic counterpart. Nevertheless the notebook gives no definite eigenfunction assertion, and the 1976 paper is not claimed to solve one.}
\GSfield{Present status}{As of November 2026: the bibliographical context is substantially identified, but the intended mathematical assertion remains unresolved; no definite problem is extracted.}
\GSfield{Assessment}{literal statement: not determinable; publication match: thematic only; verification: read the result; development: none; confidence: medium}
\end{GSitem}

\section{Developments}\label{dev:1974}
\noindent No class C or D item met the selection criterion for a new research development. The following subsections are editorial verification audits: they prove corrections, state local reconstruction limits, and check conventions. The markers \ours{} and \oursComputed{} identify work performed for this companion, without asserting novelty. The final subsection audits adjacent pages already indexed in the preceding year; its 26 rows are not additional current-year occurrences.

\subsection{Verification of the singular elliptic-system example}

This is an editorial verification of the formulas on P000565, not a claim of a new singular-solution construction. The two identical coefficients in the reconstruction below agree with the formula reproduced by \citet[\S3.2, equation (3.3)]{mingione2006regularity}. That review attributes the construction to Giusti and Miranda (1968); the original 1968 article was not retrieved. Thus the first handwritten coefficient is emended by mathematical reconstruction, supported indirectly by the historical source, rather than by claiming that the scan contains an invisible numeral.

\begin{proposition}[Residual of the tensor as copied]\leavevmode\par\noindent\ours\quad
Let \(n=N\ge3\), \(B=B_1(0)\subset\mathbb R^n\), and \(u(x)=x/|x|\) for \(x\ne0\), with any value assigned at the origin. For constants \(\alpha,\beta\), set
\[
 A_{ij}^{hk}(v)=\delta_{ij}\delta_{hk}
 +\left(\delta_{jk}+\alpha\frac{v_jv_k}{1+|v|^2}\right)
  \left(\delta_{ih}+\beta\frac{v_iv_h}{1+|v|^2}\right),
 \qquad
 F_j^k=\sum_{i,h}A_{ij}^{hk}(u)\partial_i u^h.
\]
Then \(u\in W^{1,2}(B;\mathbb R^n)\), and, away from the origin,
\[
 F_j^k=\frac1r\left[n\delta_{jk}
       +\left(-1+\frac{\alpha(n-1)}2\right)u_ju_k\right],
 \qquad
 \sum_j\partial_jF_j^k
 =(n-1)\left(\frac{\alpha(n-2)}2-2\right)\frac{x_k}{r^3},
 \quad r=|x|.
\]
In particular, the copied value \(\alpha=1/(n-2)\), with the second coefficient \(\beta=4/(n-2)\), gives the nonzero residual
\[
 -\frac32(n-1)\frac{x_k}{r^3}.
\]
It therefore does not make \(x/|x|\) solve the displayed homogeneous system.
\end{proposition}
\begin{proof}
Direct differentiation gives
\[
 \partial_i u^h=\frac{\delta_{ih}-u_i u_h}{r},
 \qquad |Du|^2=\frac{n-1}{r^2}.
\]
Because \(u\) is bounded, the inner boundary term used to identify each distributional first derivative is \(O(\varepsilon^{n-1})\) and tends to zero. Thus these punctured classical derivatives are indeed its weak derivatives. The radial integral of \(|Du|^2\) is a constant times \(\int_0^1r^{n-3}\,dr<\infty\). Moreover,
\[
 \sum_{i,h}\left(\delta_{ih}+\frac\beta2u_i u_h\right)
                    \partial_i u^h=\frac{n-1}{r},
\]
because \(\sum_i u_i\partial_i u^h=0\). Substituting this contraction proves the flux formula, including its independence of \(\beta\). Finally
\[
 \sum_j\partial_j\!\left(\frac{\delta_{jk}}r\right)
 =-\frac{x_k}{r^3},\qquad
 \sum_j\partial_j\!\left(\frac{x_jx_k}{r^3}\right)
 =(n-2)\frac{x_k}{r^3}.
\]
Combining these identities gives the residual. A nonzero residual on an open subset of the punctured ball already excludes a distributional solution on the ball.
\end{proof}

\begin{proposition}[A uniformly elliptic reconstruction and its weak solution]\leavevmode\par\noindent\ours\quad
Take both coefficients equal to \(a=4/(n-2)\). Then \(A(v)\) is real analytic and uniformly bounded for \(v\in\mathbb R^n\), and for every real matrix \(\xi=(\xi_i^h)\),
\[
 |\xi|^2\le
 \sum_{i,j,h,k}A_{ij}^{hk}(v)\xi_i^h\xi_j^k
 \le (1+n+2a+a^2)|\xi|^2.
\]
The map \(u=x/|x|\) is an energy-class weak solution on \(B\), has boundary trace \(u(x)=x\) on \(\partial B\), and has no continuous representative at the origin.
\end{proposition}
\begin{proof}
Put \(C(v)=I+a\,vv^T/(1+|v|^2)\). Since the two factors coincide, the quadratic form is exactly
\[
 |\xi|^2+\left(\sum_{i,h}C_{ih}(v)\xi_i^h\right)^2.
\]
Writing \(s=|v|^2/(1+|v|^2)\in[0,1)\), one has
\(\|C(v)\|_{\mathrm F}^2=n+2as+a^2s^2\le n+2a+a^2\).
Cauchy--Schwarz proves the upper bound; the square proves the lower bound. The denominator \(1+|v|^2\) never vanishes on real space, so the coefficients are real analytic; the same representation proves uniform boundedness.

The preceding proposition makes the divergence vanish on \(B\setminus\{0\}\). For \(\varphi\in C_c^\infty(B;\mathbb R^n)\), integrate by parts on \(B\setminus\overline{B_\varepsilon}\). The inner boundary term has magnitude at most \(C\|\varphi\|_\infty\varepsilon^{n-1}\varepsilon^{-1}=O(\varepsilon^{n-2})\), because \(|F|\le C/r\). It tends to zero for \(n\ge3\). The integral of \(F:D\varphi\) over the removed ball also tends to zero, since \(r^{-1}\) is locally integrable. Consequently
\[
 \sum_{i,j,h,k}\int_B A_{ij}^{hk}(u)\partial_i u^h\partial_j\varphi^k\,dx=0.
\]
The map is smooth near \(\partial B\), with the claimed trace. Its different radial limits at the origin preclude continuity there. The result proves existence, and asserts neither uniqueness for the nonlinear Dirichlet problem nor that this particular homogeneous system is the Euler equation of a minimizing functional.
\end{proof}

\begin{remark}[Independent exact checks]\leavevmode\par\noindent\oursComputed\quad
The script \texttt{code/exact\_checks.py} differentiates the unsimplified tensor flux using first-order jets with rational arithmetic. At \((3,4,0,\ldots,0)\), in each dimension \(3,4,5,6\), it reproduces the residual above for the copied coefficients and obtains zero for the reconstruction. These finite checks supplement the dimension-independent proof; they do not replace it. The point may lie outside \(B\), because the algebraic flux identity being checked is valid on all of \(\mathbb R^n\setminus\{0\}\).
\end{remark}

\subsection{The genus-two Hamiltonian and spectral normalization}

The algebra below checks P000555 in the canonical convention \(\{q_i,p_j\}=\delta_{ij}\). The Hamiltonian system and its spectral construction were already published by \citet[\S2, pp.~64--65]{novikov1974periodic} before the September lecture. The coefficient \(5/8\) also follows by expanding the formula of \citet[\S4]{dubrovinnovikov1974periodicity}. This is a verification of the recorded formulas, without a novelty claim.

\begin{proposition}[Exact Poisson-bracket diagnostic]\leavevmode\par\noindent\oursComputed\quad
Let \(c,d,a\) be constants and put
\begin{align*}
 J_1(a)&=p_1p_2-\frac12q_2^2-\frac52q_1^2q_2-aq_1^4+cq_1^2-dq_1,\\
 J_2&=p_1^2-2q_1p_1p_2+2(q_2-c)p_2^2+D,\\
 D&=q_1^5+2cq_1^3+dq_1^2-4q_1q_2^2+4cq_1q_2-2dq_2.
\end{align*}
Then
\[
 \{J_1(a),J_2\}=(8a-5)q_1^3(q_1p_2-p_1).
\]
Thus the copied quartic coefficient \(a=5/2\) gives a generically nonzero bracket. The replacement \(a=5/8\) gives commuting Hamiltonians.
\end{proposition}
\begin{proof}
For the convention specified above,
\(\{F,G\}=\sum_i(F_{q_i}G_{p_i}-F_{p_i}G_{q_i})\).
Here
\begin{align*}
 (J_1)_{q_1}&=-5q_1q_2-4aq_1^3+2cq_1-d,&
 (J_1)_{q_2}&=-q_2-\tfrac52q_1^2,\\
 (J_1)_{p_1}&=p_2,&(J_1)_{p_2}&=p_1,\\
 (J_2)_{p_1}&=2p_1-2q_1p_2,&
 (J_2)_{p_2}&=-2q_1p_1+4(q_2-c)p_2,\\
 (J_2)_{q_1}&=-2p_1p_2+5q_1^4+6cq_1^2+2dq_1-4q_2^2+4cq_2,\\
 (J_2)_{q_2}&=2p_2^2-8q_1q_2+4cq_1-2d.
\end{align*}
Substitution and collecting coefficients gives the displayed factorization. The accompanying script verifies the polynomial identity symbolically over the rational numbers, with \(q_1,q_2,p_1,p_2,c,d,a\) independent.
\end{proof}

\begin{proposition}[Conditional spectral polynomial in an explicit normalization]\leavevmode\par\noindent\oursComputed\quad
Set \(J_1=J_1(5/8)\), let the spatial derivative be the Hamiltonian vector field \(D_xF=\{F,J_1\}\), and identify the Schrödinger potential as \(u=q_1\). Define
\[
 P(E)=E^2+\frac{q_1}2E+\frac{q_1^2-2q_2+2c}{16}.
\]
Then the squared-eigenfunction identity gives
\begin{align*}
 R(E)&=(E-u)P^2+\frac12P D_x^2P-\frac14(D_xP)^2\\
 &=E^5+\frac c4E^3-\frac d{16}E^2
   +\left(\frac{J_1}{32}+\frac{c^2}{64}\right)E
   -\frac{J_2}{256}-\frac{cd}{128}.
\end{align*}
Moreover \(D_xR=0\) and
\(D_x^3P-4(u-E)D_xP-2(D_xu)P=0\).
\end{proposition}
\begin{proof}
Hamilton's equations give
\[
 D_xq_1=p_2,\quad D_xq_2=p_1,\quad
 D_xp_1=5q_1q_2+\tfrac52q_1^3-2cq_1+d,\quad
 D_xp_2=q_2+\tfrac52q_1^2.
\]
In particular
\[
 D_xP=\frac{p_2}2E+\frac{q_1p_2-p_1}{8},\qquad
 D_x^2P=\frac{q_2+\frac52q_1^2}{2}E
          +\frac{p_2^2-4q_1q_2+2cq_1-d}{8}.
\]
Substitution and collection by powers of \(E\) prove the formula for \(R\). The two Hamiltonians are constant along this vector field by the preceding bracket identity, so \(D_xR=0\). Direct differentiation, or differentiating the squared-eigenfunction identity and canceling the nonzero polynomial \(P\), gives the third-order identity. All polynomial equalities have also been checked exactly by the supplied script.
\end{proof}

This last proposition is deliberately conditional on the stated choices of \(u\), spatial flow and \(P\). The notebook prints \(c^2/4\) in the coefficient of \(E\) and a positive constant term. The author-hosted Dubrovin--Novikov text yields \(c^2/64\) after translating its parameters, but retains a positive constant term and has other apparent typographic defects. Although \citet[Ch.~II, \S4, p.~105]{dubrovinmatveevnovikov1976nonlinear} supplies the same \(q_1=u\), \(q_2=u''-5u^2/2\) interpretation, a complete agreement of all source sign conventions has not been established. Accordingly the data record an unresolved normalization discrepancy; the conditional identity above is not presented as a recovered literal reading of the notebook.


\subsection{Checking the two logarithmic integrals on P000564}
This is a verification of recorded identities, not the development of a newly open problem. The scan itself supplies the positive sign in the first identity and the negative sign in the second. The two expressions therefore need separate checks.

\begin{proposition}[The Basel sign]\leavevmode\par\noindent\ours\quad
The improper integral satisfies
\[
 \int_0^1\frac{\log(1-t)}{t}\,dt=-\frac{\pi^2}{6}.
\]
\end{proposition}
\begin{proof}
For \(0\le t<1\), \(-\log(1-t)=\sum_{n\ge1}t^n/n\). Since all terms are nonnegative, monotone convergence permits integration after division by \(t\), giving \(\sum_{n\ge1}1/n^2=\pi^2/6\). This is the classical dilogarithm identity, recorded in \citet[Eqs.~25.12.1--2]{nist2026polylogarithms}. The notebook's sign is impossible even before evaluation, since its integrand is negative.
\end{proof}

For the next calculation let
\[
 L(x)=\operatorname{Li}_2(x)+\tfrac12\log x\log(1-x),\qquad0<x<1,
\]
and extend to \(x<0\) by \(L(x)=-L(x/(x-1))\), with \(L(0)=0\). We use the classical reflection and five-term identities, in the normalization of \citet[Chapter II, \S1.A, pp.~23--24]{zagier2007dilogarithm}:
\[
 L(x)+L(1-x)=\frac{\pi^2}{6},
\]
\[
 L(x)+L(y)=L(xy)+L\!\left(\frac{x(1-y)}{1-xy}\right)
                    +L\!\left(\frac{y(1-x)}{1-xy}\right),\quad0<x,y<1.
\]
These real equalities have no ambiguity modulo a period. They follow as well by differentiating both sides and evaluating a boundary value; the logarithms remain real in the specified domain.

\begin{proposition}[The symmetric golden-ratio integral]\leavevmode\par\noindent\ours\quad
Set \(a=\sqrt5-2\). Then
\[
 \int_{-a}^{a}\left(\frac{\log|s|}{1-s}
                   +\frac{\log|1-s|}{s}\right)\,ds=-\frac{\pi^2}{6}.
\]
\end{proposition}
\begin{proof}
On \((-1,0)\cup(0,1)\), direct differentiation of the real dilogarithm gives
\[
 -2L'(s)=\frac{\log|s|}{1-s}+\frac{\log|1-s|}{s}.
\]
The logarithmic singularity at zero is integrable, and \(L\) is continuous there. Let \(t=(\sqrt5-1)/2\), so \(t^2=1-t\), \(a=t^3=2t-1\), \(1-a=2t^2\), and \(1+a=2t\). The five-term identity with \(x=t,y=t^2\) reduces to
\[
 L(t)+L(t^2)=L(a)+L(1/2)+L(t^2/2).
\]
Reflection gives the left side \(\pi^2/6\), and \(L(1/2)=\pi^2/12\). Thus \(L(a)+L(t^2/2)=\pi^2/12\). The negative-argument rule gives
\(L(-a)=-L(a/(1+a))=-L(t^2/2)\). Integrating the derivative now yields
\(-2[L(a)-L(-a)]=-\pi^2/6\), as asserted.
\end{proof}

\begin{remark}[Independent computational check]\leavevmode\par\noindent\oursComputed\quad
The included standard-library script \texttt{code/check\_dilogarithm.py} first evaluates the absolutely convergent series
\[
 \log a\log\frac{1+a}{1-a}
       -4\sum_{j\ge0}\frac{a^{2j+1}}{(2j+1)^2},
\]
obtained by integrating the even part of the integrand. At 90-decimal precision it gives
\[
 -1.6449340668482264364724151666460251892189499012067984\ldots.
\]
After \(N\) terms the series tail is bounded by
\(4a^{2N+1}/((1-a^2)(2N+1)^2)\); at \(N=180\) this is below \(1.51\cdot10^{-231}\). This bound controls truncation, not decimal rounding.
An independent composite-Simpson calculation after \(s=a e^{-v}\), \(0\le v\le80\), gives errors relative to the double-precision value of \(-\pi^2/6\) of \(-7.37\cdot10^{-13}\) and \(-4.66\cdot10^{-14}\) for 20000 and 40000 subintervals respectively. The numerical agreement checks the signs and endpoints; the proof above establishes the exact identity. A second algebraic check verifies separately each golden-ratio substitution in the five-term equation. Neither check found an error in the second displayed integral.
\end{remark}

\subsection{Haar normalization in the planar lattice calculation}
The following calculation audits the hypotheses of the formulas on P000573--P000574. It is a classical finite-dimensional integration, not an advance on an open gauge-theory problem. The character-expansion version and the difference between a disk and a torus are given in \citet[Appendix A]{hirasawa2020complex}; the lattice variables and loop observable already occur in \citet[\S\S II--III]{wilson1974confinement}.

\begin{proposition}[Free planar boundary and one-plaquette factors]\leavevmode\par\noindent\ours\quad
Consider a rectangular planar square lattice with \(A\) plaquettes, normalized Haar measure on every link, and no fixed boundary holonomy. Let \(\beta\in\mathbb R\). For \(U(1)\) and plaquette weight \(e^{\beta\cos\theta}\),
\[
 Z_A(\beta)=I_0(\beta)^A,
 \qquad
 I_m(\beta)=\frac1{2\pi}\int_{-\pi}^{\pi}
       e^{\beta\cos\theta}\cos(m\theta)\,d\theta.
\]
For a simple oriented Wilson loop enclosing \(a\) plaquettes,
\[
 \langle W(C)\rangle=
    \left(\frac{I_1(\beta)}{I_0(\beta)}\right)^a.
\]
For \(SU(2)\) with the fundamental trace and one-plaquette weight
\(e^{\beta\operatorname{tr}U}\), the corresponding one-plaquette factor is
\[
 z_{SU(2)}(\beta)=\frac{I_1(2\beta)}{\beta},
 \qquad z_{SU(2)}(0)=1.
\]
\end{proposition}
\begin{proof}
A spanning-tree change of link variables uses left and right invariance of normalized Haar measure. On a planar rectangular disk the remaining variables may be chosen as independent plaquette holonomies. This can also be seen by peeling successive boundary plaquettes: integrate their free outer link, and Haar invariance makes the resulting factor independent of the other links. All remaining tree integrations equal 1. Thus the partition integral is the \(A\)-th power of its one-plaquette integral, with no extra gauge-volume factor.

For \(U(1)\), normalized Haar measure is \(d\theta/(2\pi)\), giving \(I_0\). The product of oriented plaquette holonomies inside a simple loop equals its boundary holonomy. The independent expectations inside the loop each equal \(I_1/I_0\), proving the loop formula.

For \(SU(2)\), a conjugacy class has eigenvalues \(e^{\pm i\theta}\), \(0\le\theta\le\pi\). Its fundamental trace is \(2\cos\theta\), and normalized Haar integration of a class function uses \((2/\pi)\sin^2\theta\,d\theta\). Hence
\[
 z_{SU(2)}(\beta)
 =\frac2\pi\int_0^\pi e^{2\beta\cos\theta}\sin^2\theta\,d\theta
 =I_0(2\beta)-I_2(2\beta)
 =\frac{I_1(2\beta)}{\beta}.
\]
For the last identity, integrate the derivative of
\(\sin\theta\,e^{2\beta\cos\theta}\) from 0 to \(\pi\); its boundary values vanish. The limit at 0 follows directly from the original normalized Haar integral.
\end{proof}

\begin{remark}[What the notebook normalization can mean]\leavevmode\par\noindent\ours\quad
The sign convention \(E=+\sum\cos\theta\) replaces \(\beta\) by \(-\beta\). Since \(I_0\) is even, the partition function is unchanged, while the loop expectation acquires \((-1)^a\). Thus its magnitude still satisfies the displayed area law. If \(V\) counts the sites of an \(N\)-by-\(N\) array, then \(A=(N-1)^2\), and \(\log Z=V\log I_0\) is only an asymptotic relation. Identifying opposite boundary edges changes the independent-plaquette argument and requires a global constraint. Finally, the different small-coupling expansions
\[
 I_0(\beta)=1+\frac{\beta^2}{4}+O(\beta^4),\qquad
 \frac{I_1(2\beta)}{\beta}=1+\frac{\beta^2}{2}+O(\beta^4)
\]
already distinguish the Abelian formula from the stated fundamental-trace \(SU(2)\) convention. The exact \(U(1)\) pressure \(\log I_0(\beta)\) is analytic for every finite real \(\beta\); the neighboring general phase-transition discussion cannot be a proof of a finite-coupling transition in that example.
\end{remark}


\subsection{What entire-function growth does and does not specify}
\label{audit:entire-data}

\paragraph{The historical question and its limits.}
After the polynomial formula for \(\chi\), P000568 asks
\ru{Пусть \(R\) и \(P\) --- целые ф-и 1-го пор. роста. Тогда им соотв. потенциалы \(u(x)\). Какие?}
and adds \ru{не обяз. периодических}. The scan confirms the first-order-growth
wording. The strongest supported historical question is to identify the potentials
represented by an entire-function replacement of the finite-gap formula. The
notebook supplies no normalization, admissible zeros, realness assumption, or
global regularity requirement. It therefore does not determine a unique global
inverse problem. This subsection is a verification annotation to that Class F
item, not a declaration that an unformulated open problem has been solved.

\paragraph{Precisely relevant later results.}
\citet[Thm.~5.1, pp.~535--536]{marchenko1975characterization} characterize
periodic and antiperiodic spectra for a real periodic Hill operator in specified
Sobolev classes using a normalized slit-domain conformal map and weighted
conditions on its slit heights. Their paper was received on 3 February 1975.
That theorem has a specified period, reality, normalization and spectral
asymptotics, absent from the seminar question. \citet[Hypotheses~3.1--3.2,
Eq.~(3.14), Lemmas~3.3--3.5]{gesztesy1999infinite}, first circulated on
7 July 1999, develop an infinite-genus construction with convergent spectral
products, an assumed Green-function identity and controlled high-energy
asymptotics. These are concrete realizations of the direction, not an exact
match to arbitrary entire functions of order one. No evidence found in these
sources establishes why Gelfand did not specify or pursue this formulation.

\paragraph{A bounded reconstruction that can be settled.}
Let \(I\subset\mathbb R\) be a nonempty open interval and \(U\subset\mathbb C\)
a connected open set. Suppose \(P(E,x)\ne0\) on \(U\times I\), is twice
continuously differentiable in \(x\), and is holomorphic in \(E\), with its first
two \(x\)-derivatives continuous. Let \(R\) be holomorphic on \(U\), and suppose
there is a holomorphic \(y\) on \(U\) such that \(y^2=R\). For example one can
choose a simply connected region avoiding the zeros of \(R\). All conclusions
below are local in this specified region; entire functions are a special case.

\begin{proposition}[Local compatibility criterion \ours]
\label{prop:entire-compatibility}
Define
\[
 Q(E,x)=E-\frac{R(E)}{P(E,x)^2}
       +\frac{P_{xx}(E,x)}{2P(E,x)}
       -\frac{P_x(E,x)^2}{4P(E,x)^2}.
\]
There is a continuous potential \(u:I\to\mathbb C\), independent of \(E\),
for which
\[
 \chi(E,x)=\frac{y(E)}{P(E,x)}-\frac{iP_x(E,x)}{2P(E,x)}
 \quad\hbox{satisfies}\quad
 -i\chi_x+\chi^2+u-E=0
\]
for every \((E,x)\in U\times I\), if and only if \(Q(E,x)\) is independent
of \(E\). In that case \(u(x)=Q(E,x)\), equivalently
\begin{equation}
 R(E)=(E-u(x))P(E,x)^2+\tfrac12P(E,x)P_{xx}(E,x)
                    -\tfrac14P_x(E,x)^2.
 \label{eq:entire-identity}
\end{equation}
For \(x_0\in I\), the nonzero function
\[
 \psi(E,x)=\exp\!\left(\int_{x_0}^{x}
       \left(\frac{i y(E)}{P(E,s)}+\frac{P_x(E,s)}{2P(E,s)}\right)\,ds\right)
\]
satisfies \(-\psi_{xx}+u\psi=E\psi\) and \(\psi(E,x_0)=1\).
\end{proposition}
\begin{proof}
Put \(a=y/P\) and \(b=-P_x/(2P)\), so that \(\chi=a+ib\).
Because \(y\) is independent of \(x\), \(a_x=2ab\). Also
\(b_x=-P_{xx}/(2P)+P_x^2/(2P^2)\). Thus
\[
 -i\chi_x+\chi^2
   =a^2-b^2+b_x+i(2ab-a_x)
   =\frac{R}{P^2}-\frac{P_{xx}}{2P}+\frac{P_x^2}{4P^2}.
\]
The residual is exactly \(u-Q\). This proves both implications and
\eqref{eq:entire-identity}. The exponential is defined without selecting a
square root of \(P\). Its logarithmic derivative is \(i\chi\), so
\(\psi_{xx}/\psi=i\chi_x-\chi^2\); the differential equation follows.
\end{proof}

Equation \eqref{eq:entire-identity} is the classical squared-eigenfunction
identity; compare \citet[Eq.~(2.10)]{gesztesy1999infinite}. The proof above is
included to fix the signs in the notebook's convention, not to claim novelty.

\begin{proposition}[Invariance under a common factor \ours]
Suppose the compatibility criterion holds. For a holomorphic function
\(h(E)\ne0\) independent of \(x\), set
\(\widetilde P=hP\), \(\widetilde R=h^2R\), and \(\widetilde y=hy\).
Then \(\widetilde\chi=\chi\), \(\widetilde Q=Q\), and the potential is unchanged.
\end{proposition}
\begin{proof}
In the quotients \(\widetilde y/\widetilde P\),
\(\widetilde P_x/\widetilde P\), \(\widetilde P_{xx}/\widetilde P\), and
\(\widetilde R/\widetilde P^2\), all factors \(h\) cancel. This proves each
assertion. The nonvanishing hypothesis prevents spurious domain exclusions;
common zeros would require an additional removability discussion.
\end{proof}

\begin{example}[Two entire pairs of order one \ours]
The pair \(P(E,x)=R(E)=e^E\) consists of nonvanishing entire functions of
order exactly one, constant in \(x\). It gives
\(Q(E,x)=E-e^{-E}\), which is not independent of \(E\); therefore it defines
no potential through the notebook's displayed formula. In contrast,
\(P(E,x)=e^E\), \(R(E)=Ee^{2E}\) gives \(Q=0\), the free potential. Here
\(y=e^E\sqrt E\) can be chosen on any simply connected region avoiding zero,
and \(\psi=e^{i\sqrt E(x-x_0)}\). Both entire pairs meet the stated growth
condition. In the second example the exponential factors cancel: no infinite
spectral genus is created. More generally multiplying finite polynomial data
by \(h=e^{E^m}\), \(m\ge1\), changes growth without changing the potential.
Thus growth order alone neither guarantees compatibility nor specifies a
new class of potentials.
\end{example}

\paragraph{Independent second check.}
For any two solutions \(f,g\) of \(v''=(u-E)v\), put \(P=fg\) and
\(W=fg'-f'g\). Direct differentiation, independent of the Riccati substitution,
gives
\[
 (E-u)P^2+\tfrac12PP''-\tfrac14(P')^2=-\tfrac14W^2,
 \qquad W'=0.
\]
This rederives the same identity with \(R=-W^2/4\). The exact-polynomial
program \texttt{code/check\_spectral\_compatibility.py} verifies both derivations
and common-factor invariance over the rational numbers. A sign-sensitive
nonconstant check uses \(P=E+x^{-2}\), \(R=E^3\), \(u=2x^{-2}\) on \(x>0\)
away from zeros of \(P\); symbolic expansion and sixteen rational substitutions
give the same potential. \oursComputed{} The complete output is included. All
checks passed; the proof, rather than the finite substitutions, establishes the
general assertions.

\paragraph{Remaining scope.}
The local criterion does not characterize global real, nonsingular,
self-adjoint or almost-periodic realizations. A definite further problem must
choose the spectral variable (\(E\) or \(\sqrt E\)), remove common factors,
specify zero/divisor and reality conditions, and control high-energy behaviour
and continuation through zeros of \(P\). A further bounded research run could
test a chosen normalized family against those conditions. The present evidence
does not select such a family, so no new open conjecture is entered in the
problem index and no claim is made that the historical programme remains open
in every precise interpretation.


\subsection{Two missing-hypothesis checks in the late microlocal talks}

The following elementary arguments audit statements on P000592--P000593.
They are editorial proofs, not claims of new research results. They do not
turn an imprecise notebook phrase into a class C or D problem. The
established comparison results are the affine-leaf connection described by
\citet[Proposition 1.9 and Theorem 1.10]{loustausanders2017bilagrangian}
and the homogeneous normal-form statement of
\citet[Proposition 6.1.4]{duistermaathormander1972fourierii}.

\begin{proposition}[A complete Lagrangian cylinder leaf]\leavevmode\par\noindent\ours\quad
The corollary on P000592 that Lagrangian leaves are either
\(\mathbb R^n\) or tori is false, even for leaves admitting a global
frame of complete commuting Hamiltonian vector fields.
\end{proposition}

\begin{proof}
Write \(S^1=\mathbb R/\mathbb Z\), and use coordinates
\[
 M=T^*(S^1\times\mathbb R)
   =\{(q_1,q_2,p_1,p_2)\},\qquad
 \omega=dq_1\wedge dp_1+dq_2\wedge dp_2 .
\]
Here \(dq_1\) is the globally defined invariant one-form on \(S^1\).
For each \((a,b)\in\mathbb R^2\), set
\[
 L_{a,b}=\{p_1=a,\ p_2=b\}.
\]
The map \(M\to\mathbb R^2\), \((q,p)\mapsto(p_1,p_2)\), is a
submersion, so these sets form a smooth foliation. Each leaf has real
dimension two, and \(\omega|_{L_{a,b}}=0\) because both \(dp_j\) vanish
on its tangent bundle. It is therefore Lagrangian in the real
four-dimensional symplectic manifold.

The leaf is diffeomorphic to \(S^1\times\mathbb R\). It cannot be
\(\mathbb R^2\), because its fundamental group is \(\mathbb Z\), and
it cannot be a two-torus, because it is noncompact. With the convention
\(\iota_{H_f}\omega=df\), the Hamiltonian fields of the global functions
\(p_1,p_2\) are \(\partial_{q_1},\partial_{q_2}\). They commute, are
independent everywhere, and are complete. Thus adding completeness to
the notebook sentence does not remove this counterexample.
\end{proof}

\begin{remark}
Three independent checks are available here: the dimension and
restriction of \(\omega\) establish the Lagrangian property; the
fundamental group excludes the plane; and noncompactness excludes the
torus. More generally, a connected leaf with \(n\) globally independent
complete commuting fields is a transitive \(\mathbb R^n\)-orbit with
discrete stabilizer, so mixed cylinder factors are allowed. An
arbitrary complete affine structure need not supply such a global frame:
the global-action conclusion must not be inferred from the
leafwise affine connection alone.
\end{remark}

\begin{proposition}[A radial point obstructs the stated homogeneous normal form]\leavevmode\par\noindent\ours\quad
The conditions \(p(\rho)=0\) and \(dp(\rho)\ne0\) do not suffice to
straighten a real homogeneous principal symbol to a coordinate covector
by a homogeneous canonical transformation, even if a nonvanishing
scalar factor is allowed.
\end{proposition}

\begin{proof}
On \(T^*\mathbb R^2\setminus0\), take
\[
 p(x,\xi)=x_1\xi_1,\qquad
 \rho=(0,0;1,0),\qquad
 R=\xi_1\partial_{\xi_1}+\xi_2\partial_{\xi_2}.
\]
The symbol is real and positively homogeneous of degree one.
At \(\rho\), it satisfies \(p=0\) and \(dp=dx_1\ne0\).
With
\(H_p=\sum_j(p_{\xi_j}\partial_{x_j}-p_{x_j}\partial_{\xi_j})\),
\[
 H_p=x_1\partial_{x_1}-\xi_1\partial_{\xi_1},
 \qquad H_p(\rho)=-R(\rho).
\]
In particular, the Hamilton field is a nonzero radial vector.

Let \(\kappa\) be a homogeneous symplectic local diffeomorphism. Since it
commutes with positive fibre dilations, differentiation of that
commutation relation gives \(\kappa_*R=R\). Symplectic covariance gives
\[
 \kappa_*H_p=H_{p\circ\kappa^{-1}}.
\]
Thus the transformed Hamilton field is still radial and nonzero at
\(\kappa(\rho)\). If the transformed defining symbol is further
multiplied by a smooth nonvanishing function \(a\), then, on its
characteristic hypersurface,
\[
 H_{a(p\circ\kappa^{-1})}
   =a\,H_{p\circ\kappa^{-1}},
\]
so radiality persists.

In contrast, for the coordinate symbol \(q(y,\eta)=\eta_2\),
\(H_q=\partial_{y_2}\). This is horizontal and nonzero, hence cannot
be proportional to the nonzero vertical radial vector at any
point of \(T^*\mathbb R^2\setminus0\), including points where
\(\eta_2=0\). The proposed homogeneous straightening is impossible.
\end{proof}

\begin{remark}
The two-dimensional example is the stabilization of the one-dimensional
symbol \(x\xi\). It avoids the incidental fact that the one-dimensional
coordinate characteristic set \(\eta=0\) misses the punctured cotangent
bundle. The obstruction uses three separately checked facts:
\(dp(\rho)\ne0\), homogeneity preserves the radial field, and a coordinate
symbol has a nonradial Hamilton field. It audits the canonical
transformation required by the usual elliptic FIO conjugation theorem;
it does not purport to rule out arbitrary unspecified kernel
factorizations. The missing nonradial condition is explicitly present
in the cited proposition of Duistermaat--Hörmander.
\end{remark}

\paragraph{Fourier sign used throughout.}
The reconstruction uses
\(\widehat u(\xi)=\int e^{-ix\xi}u(x)\,dx\).
Consequently \((x+i0)^{-1}\), the upper-half-plane boundary value of
\(1/z\), has positive, rather than negative, wavefront covectors.
The explicit positive-frequency representation and wavefront
calculation in
\citet[Appendix A, equation (A.4)]{bohrlefeuvrepaternain2024invariant}
independently checks this convention. P000584 does not specify its
Fourier sign, so its contrary excluded-half-ray wording is logged as
an ambiguity, not as a convention-independent mathematical error.


\subsection{Adjacent pages already indexed in the preceding year}

The supplied overlapping page ranges include P000550 and P000552--P000553, but these are not additional mathematical content of records 331--332. P000550 continues the preceding KdV discussion; P000552 begins a separately dated meeting, \(1/\mathrm{IV}\text{-}74\), headed by Yura Koblenz, with Gelfand and Shnirelman interventions. The 26 entries below audit all identifiable mathematical content on those pages against the supplied prior index. They retain prior occurrence identifiers and do not create a current-year record333. Status is assessed as of November 2026. The two current records on P000551 separately provide five compatibility aliases for existing prior-year source occurrences.

\paragraph{1973-330.6 (P000550): Polynomial-exponential solutions and Bargmann potentials.}
The notes assert that arbitrary interpolation parameters produce a fundamental pair \(P(x,\zeta)e^{i\zeta x}\), \(Q(x,\zeta)e^{-i\zeta x}\) for a Sturm--Liouville equation with a Bargmann potential. A valid reflectionless construction requires admissible scattering data and a nonsingular interpolation system; an arbitrary complex choice need not give a real nonsingular potential, nor a fundamental pair at the Wronskian zeros.

The construction belongs to established reflectionless inverse scattering. \citet[\S3, Theorem~1]{shabat2014interpolation} gives a precise interpolation theorem and explicitly recalls the 1974 Gelfand discussion on p.~305. Its general equation may have spectral-parameter-dependent potential; the real Bargmann specialization needs further conditions. This is a partial match, not validation of “arbitrary”.

\emph{Audit classification: E; verification: read the result.}

\paragraph{1973-330.7 (P000550): The interpolation problem and its missing normalization.}
The recorded conditions are \(P/Q=c_j e^{2i\zeta_jx}\) at finitely many nodes and \(P/Q\to1\) at infinity. The node count is overwritten as \(2n\,?\,1\); the degrees, admissibility, and rank condition are absent. Therefore the literal system does not determine a general existence or uniqueness theorem.

\citet[\S3, equations (44)--(49), Theorem~1]{shabat2014interpolation} specifies degrees \(m,n\), \(m+n+1\) nodes, normalization of the product of leading coefficients, and nondegeneracy. Its conventions do not identify the overwritten count or repair the notebook automatically. No open problem is extracted from this incomplete formulation.

\emph{Audit classification: F; verification: read the result.}

\paragraph{1973-330.8 (P000550): The polynomial Wronskian and the absent leading factor.}
The notes write \(P_xQ-Q_xP-2i\zeta PQ=\prod_j(\zeta-\zeta_j)\). Differentiating the interpolation equations does give zeros of the left side at the nodes. But if \(P,Q\) have degrees \(m,n\) and leading product \(pq\), its leading coefficient is \(-2ipq\), not automatically one. Its negative exponential sign is consistent with \(Pe^{-i\zeta x},Qe^{i\zeta x}\), opposite to the earlier displayed pair.

\citet[Theorem~1, equation (50)]{shabat2014interpolation} proves the corresponding Wronskian factorization with an explicit leading constant \(-2\) in its real spectral convention. The degree, normalization and exponential signs must be chosen consistently. This is a mathematical normalization audit, not a confident reconstruction of the overwritten spectral letters.

\emph{Audit classification: E; verification: read the result.}

\paragraph{1973-330.9 (P000550): Cubic time evolution of the interpolation data.}
The source evolves \(c_j(t)=c_j(0)e^{2i\zeta_j^3t}\). The structural assertion is that the scattering or interpolation coefficients evolve exponentially with a cubic spectral exponent under a KdV flow. The numerical factor and sign depend on the scaling of time, the Schrödinger operator and the chosen scattering functions.

\citet[equations (13), (18)]{shabat1973kdv} obtains linear cubic evolution of the scattering kernel; the paper was received 13 December 1972, presented 18 December 1972, and published in 1973. The cubic law is known before the source discussion; the notebook coefficient is not claimed to match another convention without conversion.

\emph{Audit classification: E; verification: read the result.}

\paragraph{1973-330.10 (P000550): The incomplete third-order auxiliary equation.}
The line reads \((4\partial_t+\partial_x^3-3q\partial_x)P=(\ )Q=0\). The second parenthesis is empty, and acting on polynomial factors after removing exponentials shifts every \(\partial_x\) to \(\partial_x\pm i\zeta\). A full Lax auxiliary equation also depends on its zeroth-order term and the normalization of the eigenfunctions.

The cubic auxiliary-operator mechanism is established in \citet{shabat1973kdv} and in \citet[\S2, equation (33)]{shabat2014interpolation}. Neither supplies an exact missing notebook operator under the unresolved conventions. This remains an incomplete formula, not an open mathematical problem.

\emph{Audit classification: F; verification: read the result.}

\paragraph{1973-330.11 (P000550): Logarithmic derivatives, the spectral shift and the Miura sign.}
For \(-y_{xx}+Vy=\zeta^2y\), the logarithmic derivative \(f=y_x/y\) satisfies \(f_x+f^2=V-\zeta^2\), away from zeros. The notebook omits this shift and pairs \(4f_t+f_{xxx}+3(f^2+\zeta^2)f_x=0\) with a plus-sign KdV equation for \(f_x+f^2\). These coefficients do not form the standard Miura identity. One consistent convention is \[4V_t+V_{xxx}-6VV_x=(\partial_x+2f)\bigl(4f_t+f_{xxx}-6(f^2+\zeta^2)f_x\bigr),\quad V=f_x+f^2+\zeta^2.\]

The Riccati and factorization mechanism is classical; compare \citet[equations (2), (18)]{shabat2014interpolation}. The displayed identity is an editorial direct algebraic check \oursComputed, offered as a consistent convention rather than a palaeographic emendation of every sign on the page.

\emph{Audit classification: E; verification: read the result.}

\paragraph{1973-333.1 (P000552): The sign of the Laplace--Beltrami operator.}
On a closed smooth Riemannian manifold, the coordinate expression \(|g|^{-1/2}\partial_i(|g|^{1/2}g^{ij}\partial_j)\) is nonpositive on functions. To retain the written nonnegative eigenvalues \(\lambda_j\) and heat trace \(\sum e^{-t\lambda_j}\), define \(L=-\operatorname{div}\nabla\ge0\). The scan has heavy cancellation before the coordinate expression; an intended minus cannot be read confidently.

The sign convention and spectral expansion are standard. \citet[\S1]{mckeansinger1967curvature} uses a nonpositive Laplacian and the opposite sign in the heat exponential; the mathematical correction is convention consistency, not an invented visible glyph.

\emph{Audit classification: E; verification: read the result.}

\paragraph{1973-333.2 (P000552): Heat-kernel expansion and trace.}
For a real orthonormal eigenbasis of \(L\ge0\), \(K(t,x,y)=\sum_j e^{-t\lambda_j}\varphi_j(x)\varphi_j(y)\) and \(\operatorname{Tr}(e^{-tL})=\int_MK(t,x,x)dV\), for \(t>0\). A complex orthonormal basis requires conjugation in the second factor.

This is established spectral theory for a smooth compact manifold; \citet[\S1 and heat-kernel construction]{mckeansinger1967curvature} and \citet[introduction]{gilkey1979recursion} state the same framework. No novelty or extra assertion about termwise convergence at \(t=0\) is intended.

\emph{Audit classification: E; verification: read the result.}

\paragraph{1973-333.3 (P000552): The small-time heat-trace expansion.}
For a closed smooth \(n\)-manifold, \(\operatorname{Tr}(e^{-tL})\sim(4\pi t)^{-n/2}\sum_{j\ge0}a_jt^j\) as \(t\downarrow0\), with \(a_0=\operatorname{vol}M\). This is an asymptotic expansion to each finite order, not an asserted convergent series. Complex-time versions require sectors inside \(\operatorname{Re}t>0\).

The notes attribute the expansion to Minakshisundaram (1949). The established result and its local-invariant interpretation are checked through \citet[introduction]{gilkey1979recursion}; \citet{mckeansinger1967curvature} predates the talk. These later explanatory sources are not used to change the historical priority.

\emph{Audit classification: E; verification: read the result.}

\paragraph{1973-333.4 (P000552): Curvature contractions.}
With a fixed curvature-sign convention, \(\operatorname{Ric}\) is a metric contraction of \(\operatorname{Rm}\), and \(\operatorname{Scal}=g^{ij}\operatorname{Ric}_{ij}\). Norms mean full metric contractions; \(|\operatorname{Rm}|^2=\sum R_{ijkl}^2\) holds in an orthonormal frame at the point. Several indices in the scan are overwritten and cannot be recovered uniquely.

These are standard tensor definitions. \citet[\S7]{mckeansinger1967curvature} and \citet[\S1]{gilkey1979recursion} provide the invariant contraction framework. The coordinate-free reconstruction avoids inventing unreadable indices.

\emph{Audit classification: E; verification: read the result.}

\paragraph{1973-333.5 (P000552): The integrated scalar-curvature coefficient.}
For the scalar nonnegative Laplacian on a smooth closed manifold, with the usual geometric scalar-curvature convention, \(a_1=\frac16\int_M\operatorname{Scal}\,dV\). It is an integrated coefficient, not the pointwise heat coefficient itself.

The formula is established before the talk; \citet[\S1, equation (5a), and \S4]{mckeansinger1967curvature} computes the scalar-curvature term. Sign conventions must be translated consistently with the Laplacian and curvature tensor.

\emph{Audit classification: E; verification: read the result.}

\paragraph{1973-333.6 (P000552): The quadratic integrated heat coefficient.}
Under the same closed-manifold hypotheses, \[a_2=\frac1{360}\int_M\bigl(2|\operatorname{Rm}|^2-2|\operatorname{Ric}|^2+5\operatorname{Scal}^2\bigr)dV.\] A local formula also contains a divergence term whose integral vanishes on a closed manifold; boundary problems require additional coefficients.

The three curvature terms match the established result of \citet[\S7]{mckeansinger1967curvature}; \citet[introduction and \S4]{gilkey1979recursion} supplies an independent recurrence derivation. The source integral needs its volume element, but no change to the written numerical coefficients.

\emph{Audit classification: E; verification: read the result.}

\paragraph{1973-333.7 (P000552): The missing off-diagonal Gaussian.}
Near the diagonal, in a normal neighborhood and as \(t\downarrow0\), the heat parametrix has the form \[K(t,x,y)\sim(4\pi t)^{-n/2}e^{-d_g(x,y)^2/(4t)}\sum_{j\ge0}u_j(x,y)t^j.\] The notebook displays the power series without the Gaussian. That abbreviated formula is valid on the diagonal, but is false as a literal off-diagonal expansion.

The local parametrix is classical; see the heat-kernel construction in \citet{mckeansinger1967curvature}. \citet[introduction, p.~202]{gilkey1979recursion} explicitly states off-diagonal rapid decay. The omitted exponential is an emendation of mathematical content.

\emph{Audit classification: E; verification: read the result.}

\paragraph{1973-333.8 (P000552): Sakai and the fourth heat coefficient.}
The source calls \(a_3\) the fourth coefficient, counting \(a_0\) first, and attributes a computation to Sakai. No formula for \(a_3\) is recorded, so the item is an attribution, not an occurrence of every term in the general cubic-curvature expression.

\citet{sakai1971eigenvalues}, Tohoku Math. J.23(4),589--603, appeared in1971. Publisher issue metadata were opened; the attribution is explicitly confirmed by \citet[introduction]{gilkey1979recursion}. The full Sakai formula was not independently read in this audit.

\emph{Audit classification: E; verification: indirect.}

\paragraph{1973-333.9 (P000552): Koblenz’s claimed independent fourth-coefficient computation.}
The speaker, named Yura Koblenz from Leningrad, is said to have computed \(a_3\) independently of Sakai. There is no formula, manuscript title, date of independent discovery or method in the notebook.

Targeted searches for Koblenz, heat coefficients and Laplace asymptotics did not identify an exact primary publication. Sakai’s known computation does not prove or disprove the claim of independence. The attribution remains unverified, rather than classified as a new open mathematical question.

\emph{Audit classification: F; verification: unresolved lead.}

\paragraph{1973-333.10 (P000552): The lattice-point and Riemann-zero analogy.}
Gelfand writes a circle lattice-point main term \(\pi r^2\) plus a schematic oscillatory sum and identifies the frequencies with zeros of the Riemann zeta function, “modulo the Riemann hypothesis”. The summation variables, coefficients and scaling are overwritten or undefined. The literal expression is not a verified lattice-point formula.

No precise claim is recovered from this analogy. \citet[introduction]{colindeverdiere1973geodesicsii} distinguishes exact Poisson formulas on flat tori from general geodesic trace asymptotics, but this does not identify the notebook’s asserted zeta-zero formula. It remains F, with no transfer of an unrelated modern theorem.

\emph{Audit classification: F; verification: read the result.}

\paragraph{1973-333.11 (P000552-P000553): Closed billiard paths and oscillatory spectral density.}
Shnirelman identifies oscillation lengths with closed billiard trajectories in a ball; the continuation displays \(f(\lambda)=f_0(\lambda)+\sum c_k\lambda^{\alpha_k}\sin(\ell_k\lambda+\omega_k)\). The variable must be a frequency for \(\ell_k\lambda\) to be the usual phase. Amplitudes, smoothing, boundary conditions, repeated trajectories and degeneracies are not specified.

The named Balian--Bloch1972 paper is a strong historical match, but its publisher page returned403 and the exact formula was not read. \citet[introduction]{colindeverdiere1973geodesicsii} explicitly cites their bounded-domain Green-function method and gives a related closed-manifold theorem under nondegeneracy conditions. The match to the literal billiard formula remains partial.

\emph{Audit classification: E; verification: indirect.}

\paragraph{1973-333.12 (P000552): Cauchy smoothing and convergence.}
The notes convolve the spectral measure with \(\varepsilon/(\lambda^2+\varepsilon^2)\). A normalized Cauchy kernel has an extra factor \(1/\pi\). One must also specify a finite spectral cutoff or a convergent regularization: for frequencies of \(\sqrt L\) in dimension \(n\ge2\), the positive tail \(\sum_j\varepsilon/\lambda_j^2\) diverges by Weyl growth. If \(\lambda_j\) denotes eigenvalues of \(L\), the bare sum instead converges in dimensions below4. A rapidly decreasing test function always defines a suitable smoothed spectral measure.

Smoothing is an established spectral method, not a theorem guaranteeing convergence of this unspecific formula. The Weyl-law framework is in \citet[\S1]{mckeansinger1967curvature}; the oscillatory context is in \citet{colindeverdiere1973geodesicsii}. The convergence distinction is an editorial check of the displayed kernel.

\emph{Audit classification: E; verification: read the result.}

\paragraph{1973-333.13 (P000553): The heat trace is not the displayed oscillatory trace.}
The scan literally equates \(\sum e^{i\lambda_jz}\) with \(\operatorname{Tr}(e^{-z\Delta})\). For eigenvalues \(\mu_j\) of \(L\ge0\), the heat trace is \(\sum e^{-z\mu_j}\), holomorphic for \(\operatorname{Re}z>0\). If \(\lambda_j=\sqrt{\mu_j}\), the half-wave sum is \(\operatorname{Tr}(e^{iz\sqrt L})\), holomorphic for \(\operatorname{Im}z>0\) and distributionally defined on the real axis. These are different operators and domains.

The later half-wave equation on the page supports a frequency interpretation, but does not repair the earlier equality literally. Compare the complex-time heat analysis of \citet{colindeverdiere1973geodesicsii} and the later spectral-trace paper \citet{duistermaatguillemin1975spectrum}; the latter was received20October1974 and publishedFebruary1975.

\emph{Audit classification: E; verification: read the result.}

\paragraph{1973-333.14 (P000553): The unspecified general spectral-gap question.}
The “twin problem” asks for \(\lambda_{k+1}-\lambda_k>c\lambda_k^{-\nu}\) for “general operators”. No operator class, multiplicity convention, generic topology, constants or frequency convention is fixed. Read universally with multiplicities, it is false even for the circle.

The adjacent circle and torus examples show why a qualified formulation is essential. Neither the wave-trace theorem nor the near-flat torus theorem proves a lower bound for an unspecified generic class. The item remains F; no modern universal gap conjecture is silently substituted.

\emph{Audit classification: F; verification: read the result.}

\paragraph{1973-333.15 (P000553): Exact spectral doublets on the circle.}
On a circle of length \(\ell\), the scalar Laplacian has eigenvalues \((2\pi m/\ell)^2\), each of multiplicity2 for \(m\ge1\). The half-wave eigenvalues are \(2\pi m/\ell\) with the same multiplicity. Every smooth metric on a circle is isometric to one specified by its total length, so metric perturbation preserves this doublet structure.

This elementary Fourier computation proves the particular counterexample to the unqualified lower-gap assertion. The notebook’s wider phrase about stable asymptotic double multiplicity should not be extended to an arbitrary operator perturbation without specifying its class. The nearby torus analogy is treated independently by \citet{shnirelman1975multiplicity}.

\emph{Audit classification: E; verification: read the result.}

\paragraph{1973-333.16 (P000553): Shnirelman’s all-index close-neighbor theorem.}
Let a smooth Riemannian metric on \(T^2\) be sufficiently close to a flat metric, and let \(\lambda_k\) be the ordered eigenvalues of \(\sqrt{-\Delta_g}\), with multiplicities. Then, for every \(N\), there is \(C_N>0\) such that for every \(k>1\), \[\min(\lambda_{k+1}-\lambda_k,\lambda_k-\lambda_{k-1})<C_N k^{-N}.\] The closeness is enough for the invariant-torus construction used in the proof; no explicit norm threshold is supplied in the two-page announcement.

\citet[Theorem and Lemmas1--3, pp.~265--266]{shnirelman1975multiplicity} states exactly this all-index result, explicitly strengthening Lazutkin’s positive-density subsequence theorem. It was received by the Moscow Mathematical Society7May1974 and published1975, after the1April1974 notebook meeting. Both primary pages were read. This is B relative to Koblenz’s main talk, with Shnirelman’s intervention explicitly preserved.

\emph{Audit classification: B; verification: read the result.}

\paragraph{1973-333.17 (P000553): The tentative three-torus comparison.}
The source says that in \(T^3\) “apparently” the preceding behavior no longer holds. It specifies neither a metric nor the quantifier of the alleged negation. Flat tori still have symmetry-related multiplicities, so the phrase cannot be read as an assertion about every three-torus metric.

\citet{shnirelman1975multiplicity} proves a two-dimensional theorem and does not supply the proposed three-dimensional negation. No exact primary theorem matching this tentative sentence was verified. It remains F rather than an eligible open problem.

\emph{Audit classification: F; verification: read the result.}

\paragraph{1973-333.18 (P000553): Half-wave propagation and its trace.}
With \(L=-\Delta_g\ge0\), the Cauchy problem \(u_t=i\sqrt L\,u\), \(u(0)=\delta_x\), has distribution kernel \(U(t,y,x)\) of \(e^{it\sqrt L}\). Its distributional trace is \(\int_M U(t,x,x)dV=\sum_k e^{it\sqrt{\mu_k}}\). The sum is not generally pointwise convergent for real \(t\).

The spectral construction is classical before the meeting. Its microlocal analysis uses the real-principal propagation framework of \citet[\S6.1]{duistermaathormander1972fourierii}; the later \citet{duistermaatguillemin1975spectrum} formalizes the trace relation. The source uses the same letter for energy eigenvalues and their square roots, requiring this distinction.

\emph{Audit classification: E; verification: read the result.}

\paragraph{1973-333.19 (P000553): Unit-speed singularity propagation and caustics.}
The page places kernel singularities at \(|x-y|=t\), draws a returning loop, and mentions bifurcations. On a Riemannian manifold the precise statement is propagation along the geodesic Hamilton flow of \(|\xi|_g\), at unit speed in the base. The equality \(d_g(x,y)=|t|\) describes the short-time normal-neighborhood front, not all returning geodesics or caustics globally.

\citet[Theorem6.1.1]{duistermaathormander1972fourierii} establishes propagation of wave-front singularities along real-principal bicharacteristics before the talk. The sketch does not define an additional classification of bifurcations. Its loop is accounted for as a propagation/return diagram, not a separate unknown theorem.

\emph{Audit classification: E; verification: read the result.}

\paragraph{1973-333.20 (P000553): Closed-geodesic contributions to a wave trace.}
Restricting the propagator to the diagonal and integrating selects fixed points of geodesic flow, hence closed geodesics. The singular support of the half-wave trace is contained in zero together with signed lengths of closed geodesics. Identifying a nonzero singular coefficient or a full expansion requires clean/nondegenerate fixed-point hypotheses and control of cancellations.

\citet[introduction and main theorem]{colindeverdiere1973geodesicsii} gives a pre-talk closed-manifold oscillatory spectral/geodesic relation under explicit nondegeneracy hypotheses. \citet{duistermaatguillemin1975spectrum}, receivedOctober1974 and publishedFebruary1975, is a later general trace-theoretic continuation; only its publisher metadata were checked here. These are related to, but not an exact theorem for, the notebook’s unspecified billiard amplitudes.

\emph{Audit classification: E; verification: read the result.}

\section{Problem index}\label{problems:1974}
Earlier-year occurrence identifiers are retained as text; current-year identifiers link to their entries.

\subsection{GP-1973-50: Monotonicity in a Hill spectral band}\label{prob:GP-1973-50}
For a real smooth periodic one-dimensional Schrödinger operator, is the Floquet phase strictly monotone inside a band when the Bloch branch is fixed by a positive-current orientation?
\par\noindent\textit{Occurrences:} 1973-326.9, \hyperref[item:1974-340.10]{1974-340.10}.
\par\noindent\textit{Status (2026-11):} As of November 2026, yes. With \(j=\operatorname{Im}(\overline\psi\psi_x)>0\), differentiation gives \(p^{\prime}(E)=(2j)^{-1}\int_{x_0}^{x_0+T}|\psi|^2dx>0\). The reciprocal branch has opposite orientation. This is classical Floquet theory; compare \citet[Ch.~II, \S1]{dubrovinmatveevnovikov1976nonlinear}.

\subsection{GP-1973-51: Supplementary data for the periodic inverse spectral problem}\label{prob:GP-1973-51}
What information, in addition to periodic spectral band data, determines a periodic Schrödinger potential?
\par\noindent\textit{Occurrences:} 1973-326.13, \hyperref[item:1974-336.15]{1974-336.15}, \hyperref[item:1974-339.3]{1974-339.3}, \hyperref[item:1974-339.4]{1974-339.4}, \hyperref[item:1974-340.8]{1974-340.8}.
\par\noindent\textit{Status (2026-11):} As of November 2026, in the smooth real nonsingular finite-gap class discussed here, a suitable degree-g divisor, including the sheet choice in each closed gap interval and a fixed base point, supplies the missing phase data and parametrizes the isospectral torus. \citet[\S3]{dubrovinnovikov1974periodicity} and \citet[Ch.~II, \S3]{dubrovinmatveevnovikov1976nonlinear} verify this answer. The present occurrences do not supply a theorem for arbitrary periodic potentials with infinitely many open gaps.

\subsection{GP-1974-01: Continuous functions on the product-sequence ray space}\label{prob:GP-1974-01}
For \(S=(\mathbb R^{\mathbb N}\setminus\{0\})/\mathbb R_{>0}\) with the quotient of the product topology, is every continuous real function constant?
\par\noindent\textit{Occurrences:} 1967-132.7, \hyperref[item:1974-337.6]{1974-337.6}.
\par\noindent\textit{Status (2026-11):} As of November 2026, yes; the exact example was published in 1967 by \citet[introduction]{gelfandfuchs1967topology}. The earlier occurrence had no separate exact problem identifier. This is a settled explanatory example, not a new open problem.

\subsection{GP-1974-02: Discontinuous minimizers of elliptic vectorial variational integrals}\label{prob:GP-1974-02}
Do uniformly elliptic vectorial variational problems with smooth or analytic coefficient dependence admit a discontinuous energy-class minimizer?
\par\noindent\textit{Occurrences:} \hyperref[item:1974-341.10]{1974-341.10}.
\par\noindent\textit{Status (2026-11):} As of November 2026, yes, in the broad class with dependence on the unknown itself; the Giusti--Miranda example dates from 1968, as verified through \citet[\S3.2]{mingione2006regularity}. This does not assert that the repaired homogeneous PDE in the audit is its Euler--Lagrange equation or that every dimension supports the same minimizing example.

\subsection{GP-1974-10: Finite-gap spectra and stationary KdV characterization}\label{prob:GP-1974-10}
For real smooth periodic potentials, characterize when the whole-line Schr\"odinger spectrum has only finitely many open bounded gaps, and relate this to stationary equations of the KdV hierarchy.
\par\noindent\textit{Occurrences:} 1973-326.15, \hyperref[item:1974-340.7]{1974-340.7}, \hyperref[item:1974-343.2]{1974-343.2}.
\par\noindent\textit{Status (2026-11):} As of November 2026, solved: stationary implies finite-gap in \citet[\S2]{novikov1974periodic}, published July 1974; the converse appears in \citet[\S3]{dubrovinnovikov1974periodicity}, received 3 June 1974 and published later that year. See also \citet[Chapter 2, \S2]{dubrovinmatveevnovikov1976nonlinear}. This is distinct from the supplementary-data inverse problem GP-1973-51 and the density problem GP-1973-52.

\subsection{GP-1974-11: The loop-space target for positive and negative particles}\label{prob:GP-1974-11}
Identify the homotopy type of finite signed configurations in \(\mathbb R^n\) in which opposite particles cancel, and decide whether the target is \(\Omega^nS^n\).
\par\noindent\textit{Occurrences:} \hyperref[item:1974-344.5]{1974-344.5}.
\par\noindent\textit{Status (2026-11):} As of November 2026, resolved by the McDuff model \(\Omega^n((S^n\times S^n)/\Delta S^n)\); the sphere-only target is incorrect for the stated cancellation topology. The 1974 page already replaces it by an unspecified2n-dimensional target; the publication is 1975. See \citet{mcduff1975configuration} and \citet[Example 11.5]{kallel1998particles}.

\subsection{GP-1974-12: Computing local Milnor numbers from Newton data}\label{prob:GP-1974-12}
For a holomorphic germ with an isolated critical point and a fixed convenient Newton boundary, find a lower bound for the Milnor number in terms of that boundary, and state coefficient conditions under which the bound is attained.
\par\noindent\textit{Occurrences:} \hyperref[item:1974-346.2]{1974-346.2}, \hyperref[item:1974-346.3]{1974-346.3}.
\par\noindent\textit{Status (2026-11):} As of November 2026, resolved by the Newton-number inequality and equality for face-nondegenerate germs. This does not say that the Newton boundary alone determines the Milnor number of every degenerate germ. \citet[TheoremI]{kushnirenko1975newton}, received 30 September 1974 and published 1975; full treatment in \citet{kouchnirenko1976polyedres}.

\subsection{GP-1974-13: Counting all isolated critical points from the Newton polytope}\label{prob:GP-1974-13}
For a convenient complex polynomial, bound the sum of Milnor numbers of its isolated critical points by the convex hull of 0 and its support, and characterize a generic class attaining the bound.
\par\noindent\textit{Occurrences:} \hyperref[item:1974-349.1]{1974-349.1}, \hyperref[item:1974-349.2]{1974-349.2}.
\par\noindent\textit{Status (2026-11):} As of November 2026, resolved by the global Newton-number theorem: the count is at most the alternating coordinate-volume number and equals it under nondegeneracy at infinity. See \citet[TheoremII]{kushnirenko1975newton}. The general mixed-volume root-counting theorem is related but is not substituted for this gradient-specific assertion.

\subsection{GP-1974-14: A measure-invariant irreducible representation of the PSL2 current group}\label{prob:GP-1974-14}
For \(C([0,1],PSL_2(\mathbb R))\) in the uniform topology, construct a nontrivial irreducible strongly continuous unitary representation whose equivalence class is unchanged by pullback along measure-preserving transformations acting on the current group, for a fixed positive Lebesgue measure. Require full interval support: currents supported in every nonempty open subinterval act nontrivially, so a finite evaluation representation is excluded.
\par\noindent\textit{Occurrences:} \hyperref[item:1974-351.4]{1974-351.4}, \hyperref[item:1974-351.14]{1974-351.14}.
\par\noindent\textit{Status (2026-11):} As of November 2026, solved in this precise scope by the 1973 multiplicative integral of complementary-series representations; see \citet[Theorems2.2, 3.3--3.4]{vershikgelfandgraev1973functions}. The earlier titles 1969-180.2/GP-1969-02 and 1969-180.1/GP-1969-01 concern related existence and classification questions, but their full formulations are unavailable here, so they are not merged or declared solved.

\subsection{GP-1974-20: Random fields with prescribed conditional distributions}\label{prob:GP-1974-20}
For a proper consistent specification \((\gamma_V)_{V\Subset T}\) on a finite-alphabet configuration space \(X^T\), determine all probability measures satisfying \(\mu\gamma_V=\mu\) for every finite \(V\), with particular attention to existence and uniqueness.
\par\noindent\textit{Occurrences:} \hyperref[item:1974-355.11]{1974-355.11}, \hyperref[item:1974-358.5]{1974-358.5}.
\par\noindent\textit{Status (2026-11):} As of November 2026: existence for quasilocal finite-alphabet specifications and uniqueness under Dobrushin contraction are established, as are model-dependent examples of phase coexistence. This is an expository recurrence of \citet[Section 1 and Theorems 1--2]{dobrushin1968conditional}, not a blanket claim that all specifications have been classified.

\section{Corrections to the transcription and translation}\label{corr:1974}
The evidence distinguishes a transcription correction from a mathematical emendation of wording faithfully preserved in the notebook.
\subsection*{P000550-P000553 (both; provenance): catalogue page ranges for records 331--332}
\noindent\textit{Transcription reads.} Entire overlapping ranges assigned to the two records\par
\noindent\textit{Critical reading.} Actual shared record content is on P000551\par
\noindent\textit{Evidence.} P000550 continues the preceding KdV discussion; P000552 is dated 1/IV-74 and begins the next spectral talk, continued on P000553.\par
\subsection*{P000550 (both; ambiguity): Bargmann interpolation hypotheses}
\noindent\textit{Transcription reads.} Any parameters produce a fundamental system\par
\noindent\textit{Critical reading.} Require nondegenerate interpolation and admissible real nonsingular data; exclude Wronskian zeros when claiming a fundamental pair\par
\noindent\textit{Evidence.} Shabat2014 Theorem1 states the rank and normalization assumptions; arbitrary complex data need not define a real regular reflectionless potential.\par
\subsection*{P000550 (both; faithful error): polynomial Wronskian normalization}
\noindent\textit{Transcription reads.} \(P_xQ-Q_xP-2i\zeta PQ=\prod(\zeta-\zeta_j)\)\par
\noindent\textit{Critical reading.} Supply leading factor \(-2ipq\), and reconcile exponential signs\par
\noindent\textit{Evidence.} For leading coefficients \(p,q\), the left side has leading coefficient \(-2ipq\). The earlier pair \(Pe^{+i\zeta x},Qe^{-i\zeta x}\) has the opposite sign in this Wronskian convention.\par
\subsection*{P000550 (both; ambiguity): Riccati spectral shift and Miura coefficients}
\noindent\textit{Transcription reads.} \(f_x+f^2=q\), plus-sign equations with coefficient 3\par
\noindent\textit{Critical reading.} For \(-y_{xx}+Vy=\zeta^2y\): \(f_x+f^2=V-\zeta^2\); use a consistently normalized Miura identity\par
\noindent\textit{Evidence.} The exact conditional identity is given in boundary row 1973-330.11. The overwritten sign and time scaling do not permit a unique literal replacement.\par
\subsection*{P000551 (both; provenance): date of records 331--332}
\noindent\textit{Transcription reads.} Catalogue: 25-Mar-75\par
\noindent\textit{Critical reading.} Manuscript: 25/III-74; catalogue date retained in metadata\par
\noindent\textit{Evidence.} The scan heading unambiguously ends in74; the same five source items occur in the previous-year index. Parent verification of the MCCME all-links page confirms its1975 date, so these IDs are compatibility aliases, not new occurrences. The individual catalogue record pages could not be retrieved.\par
\subsection*{P000551 (both; faithful error): principal-bundle extension}
\noindent\textit{Transcription reads.} \(\operatorname{Aut}(P)\to\operatorname{Diff}M\), kernel \(H^M\)\par
\noindent\textit{Critical reading.} Image \(\operatorname{Diff}(M)_P\); kernel \(\Gamma(\operatorname{Ad}P)\)\par
\noindent\textit{Evidence.} The shorter formula is valid for a trivial bundle. In general a diffeomorphism lifts exactly when it preserves the bundle isomorphism class.\par
\subsection*{P000551 (both; ambiguity): Fock-space normalization}
\noindent\textit{Transcription reads.} Coordinates divided by n! and norm weighted by 1/n!\par
\noindent\textit{Critical reading.} Specify the sector norm; standard coordinates use 1/sqrt(n!)\par
\noindent\textit{Evidence.} The handwritten factorials are retained; the ambiguity concerns conventions, not a demonstrable misread glyph.\par
\subsection*{P000552 (both; ambiguity): Laplacian sign}
\noindent\textit{Transcription reads.} Positive spectrum and \(\exp(-t\Delta)\), but coordinate \(\operatorname{div}\operatorname{grad}\)\par
\noindent\textit{Critical reading.} Use \(L=-\operatorname{div}\operatorname{grad}\) for positive eigenvalues, or reverse spectrum and heat signs\par
\noindent\textit{Evidence.} The scan is cancelled before the coordinate expression; sign consistency is clear mathematically but the original cancelled sign is not recoverable.\par
\subsection*{P000552 (both; faithful error): off-diagonal heat expansion}
\noindent\textit{Transcription reads.} \((4\pi t)^{-n/2}\) times power series near the diagonal\par
\noindent\textit{Critical reading.} Include \(\exp(-d_g(x,y)^2/(4t))\) off the diagonal\par
\noindent\textit{Evidence.} The omitted Gaussian is necessary for the classical local heat parametrix. The unmodified expression is only the diagonal form.\par
\subsection*{P000552 (both; ambiguity): Cauchy spectral smoothing}
\noindent\textit{Transcription reads.} Unqualified convolution with \(\varepsilon/(\lambda^2+\varepsilon^2)\)\par
\noindent\textit{Critical reading.} Specify cutoff or a convergent regularization; divide by \(\pi\) for a normalized Cauchy kernel\par
\noindent\textit{Evidence.} Weyl growth makes the positive frequency sum divergent in dimension at least2; for Laplacian energy eigenvalues it converges only below dimension4.\par
\subsection*{P000553 (both; faithful error): oscillatory trace versus heat trace}
\noindent\textit{Transcription reads.} \(\sum e^{i\lambda_kz}=\operatorname{Tr}(e^{-z\Delta})\)\par
\noindent\textit{Critical reading.} For \(\lambda_k=\sqrt{\mu_k}\), use \(\operatorname{Tr}(\exp(i z\sqrt L))\); heat trace uses \(\exp(-z\mu_k)\)\par
\noindent\textit{Evidence.} The scan confirms the two different exponents. The later half-wave equation supports the frequency interpretation.\par
\subsection*{P000553 (both; ambiguity): global propagation front}
\noindent\textit{Transcription reads.} \(|x-y|=t\)\par
\noindent\textit{Critical reading.} Use geodesic Hamilton flow; distance \(=t\) describes only a short-time normal-neighborhood front\par
\noindent\textit{Evidence.} Returning geodesics and caustics shown in the sketch need not realize the minimizing distance.\par
\subsection*{P000554 (both; ambiguity): Nash--De Giorgi--Moser parenthesis}
\noindent\textit{Transcription reads.} Hölder coefficients\par
\noindent\textit{Critical reading.} Bounded measurable, uniformly elliptic coefficients suffice for Hölder continuity of energy weak solutions\par
\noindent\textit{Evidence.} This separates the Hölder regularity of the solution from an unnecessary hypothesis on the coefficients.\par
\subsection*{P000554 (both; faithful error): Jacobian quotient}
\noindent\textit{Transcription reads.} \(H^1(R,\mathbb C)/H^1(R,\mathbb Z)\)\par
\noindent\textit{Critical reading.} \(H^1(R,\mathbb C)/(F^1H^1+H^1(R,\mathbb Z))\)\par
\noindent\textit{Evidence.} The scan genuinely gives the shorter quotient. Its complex dimension and noncompactness are incompatible with the Jacobian.\par
\subsection*{P000555 (both; faithful error): quartic term in V}
\noindent\textit{Transcription reads.} \(-5q_1^4/2\)\par
\noindent\textit{Critical reading.} \(-5q_1^4/8\)\par
\noindent\textit{Evidence.} The literal bracket equals \(15q_1^3(q_1p_2-p_1)\). The corrected coefficient gives zero and agrees with Dubrovin--Novikov (1974), section 4.\par
\subsection*{P000555 (both; ambiguity): parameters in Hamiltonian pair}
\noindent\textit{Transcription reads.} c,d and c1,d1\par
\noindent\textit{Critical reading.} Identify c=c1 and d=d1 for the displayed commuting-pair reconstruction\par
\noindent\textit{Evidence.} The notebook changes parameter labels without declaring the identification.\par
\subsection*{P000555 (both; ambiguity): coefficient of E in spectral quintic}
\noindent\textit{Transcription reads.} \(J_1/32+c_1^2/4\)\par
\noindent\textit{Critical reading.} Published normalization suggests \(J_1/32+c_1^2/64\)\par
\noindent\textit{Evidence.} Under c=8c2,d=c1 in the 1974 published formula, all the other coefficients match. The remaining discrepancy is retained as a normalization reconstruction.\par
\subsection*{P000555 (both; faithful error): Abel theorem}
\noindent\textit{Transcription reads.} isomorphism\par
\noindent\textit{Critical reading.} birational map; isomorphism on the nonspecial locus\par
\noindent\textit{Evidence.} The fibre over a divisor class is its complete linear system; special effective divisors give positive-dimensional fibres.\par
\subsection*{P000555-P000556 (both; faithful error): torus and Jacobian identification}
\noindent\textit{Transcription reads.} All common affine levels are abelian varieties; all isospectral sets are Jacobians\par
\noindent\textit{Critical reading.} Generic affine levels are open parts of compact Jacobians; the finite-gap real isospectral family is a real torus\par
\noindent\textit{Evidence.} Affine levels are not compact; the finite-dimensional theorem requires a nonsingular finite-gap curve.\par
\subsection*{P000555 (both; ambiguity): spectral constant term}
\noindent\textit{Transcription reads.} \(+J_2/256+cd/128\)\par
\noindent\textit{Critical reading.} In the explicitly fixed audit normalization: \(-J_2/256-cd/128\); literal emendation unresolved\par
\noindent\textit{Evidence.} Exact polynomial expansion with \(u=q_1\) and \(P=E^2+q_1E/2+(q_1^2-2q_2+2c)/16\) forces the displayed negative sign. Available published formulas use signs whose complete normalization correspondence has not been established.\par
\subsection*{P000556 (both; faithful error): Hilbert-cube pseudoboundary}
\noindent\textit{Transcription reads.} \(Q\setminus P\) is a Z-set\par
\noindent\textit{Critical reading.} \(Q\setminus P\) is a sigma-Z-set\par
\noindent\textit{Evidence.} It is a dense nonclosed union of the closed coordinate faces. The source definition explicitly requires a Z-set to be closed.\par
\subsection*{P000556 (both; ambiguity): maximal negligible set}
\noindent\textit{Transcription reads.} A maximal \(Z_\sigma\)-set exists\par
\noindent\textit{Critical reading.} No inclusion-maximal sigma-Z-subset of Q exists; intended universal meaning unspecified\par
\noindent\textit{Evidence.} Baire category gives a missing point, and adjoining that singleton preserves the sigma-Z property.\par
\subsection*{P000556 (both; faithful error): Fréchet-space classification}
\noindent\textit{Transcription reads.} All separable Fréchet spaces are homeomorphic\par
\noindent\textit{Critical reading.} All infinite-dimensional separable Fréchet spaces are homeomorphic\par
\noindent\textit{Evidence.} Finite-dimensional real spaces have distinct topological dimensions. Anderson (1967), p. 31, states the infinite-dimensional qualification.\par
\subsection*{P000556-P000557 (both; faithful error): mean-curvature sign convention}
\noindent\textit{Transcription reads.} \(\operatorname{div}(Du/W)=-H\), but \(F=\int W+\lambda\), \(\lambda_t=H\)\par
\noindent\textit{Critical reading.} The displayed functional yields div(Du/W)=H; change one sign consistently\par
\noindent\textit{Evidence.} Direct differentiation of the functional fixes the Euler--Lagrange sign.\par
\subsection*{P000557 (both; faithful error): changed one-dimensional left endpoint}
\noindent\textit{Transcription reads.} \(u(-1)=3,\ u(1)=3+A\)\par
\noindent\textit{Critical reading.} \(u(-1)=-3,\ u(1)=3+A\) for the subsequent argument\par
\noindent\textit{Evidence.} The handwritten plus sign is faithfully transcribed, but the next integral starts from -3 and compares with 3 cube-root(x).\par
\subsection*{P000557 (both; faithful error): integrated prescribed flux}
\noindent\textit{Transcription reads.} \(u^{\prime}/\sqrt{1+(u^{\prime})^2}=(1+x^{4/3})^{-1/2}\)\par
\noindent\textit{Critical reading.} Add an integration constant C before comparing boundary data\par
\noindent\textit{Evidence.} Different solutions of the same prescribed derivative of the flux differ by a constant in the flux. The bound at zero then forces C<=0.\par
\subsection*{P000558 (both; faithful error): BV compactness}
\noindent\textit{Transcription reads.} Bounded BV sets are relatively compact in L1\par
\noindent\textit{Critical reading.} Use the full BV norm on a bounded Lipschitz domain, or state local compactness\par
\noindent\textit{Evidence.} Unbounded domains admit escaping translates; total variation alone does not bound constants.\par
\subsection*{P000558 (both; faithful error): number of Bloch poles and zeros}
\noindent\textit{Transcription reads.} Has N poles and N zeros\par
\noindent\textit{Critical reading.} Degree-N moving and reference divisors, with common points cancelled\par
\noindent\textit{Evidence.} At x=x0 the normalized Bloch solution is identically one. A literal uncancelled count cannot hold there.\par
\subsection*{P000560 (both; ambiguity): Bloch equation uses potential symbol}
\noindent\textit{Transcription reads.} \(u(x+T)=\nu u(x)\)\par
\noindent\textit{Critical reading.} Use \(\psi\) for the Bloch eigenfunction\par
\noindent\textit{Evidence.} The same page has already reserved u for a periodic potential; the multiplier equation is for an eigenfunction.\par
\subsection*{P000560-P000561 (both; faithful error): discriminant extrema}
\noindent\textit{Transcription reads.} max and min >= 1\par
\noindent\textit{Critical reading.} maxima >= 1; minima <= -1, or absolute values >= 1\par
\noindent\textit{Evidence.} The drawing puts the lower extrema below -1. The real scalar Floquet theorem has the stated opposite signs.\par
\subsection*{P000561 (both; ambiguity): monodromy numerator}
\noindent\textit{Transcription reads.} \(a_{11}a_{22}\)\par
\noindent\textit{Critical reading.} \(a_{11}-a_{22}\)\par
\noindent\textit{Evidence.} No intervening operator is distinguishable in the scan. The normalized eigenvector identity mu=a11+i chi a21 requires the difference.\par
\subsection*{P000561 (both; ambiguity): phase integer}
\noindent\textit{Transcription reads.} \(p=\int\chi_R+\pi n\), overwritten coefficient\par
\noindent\textit{Critical reading.} \(p=\int\chi_R\pmod{2\pi}\)\par
\noindent\textit{Evidence.} The multiplier is exp(i p); its logarithm is defined modulo 2 pi. The preceding amplitude is positive and periodic on a real band.\par
\subsection*{P000562 (both; ambiguity): Wronskian factor}
\noindent\textit{Transcription reads.} \(W(c,\Psi)=i\Psi^{\prime}(x_0)=i\chi(x_0)\)\par
\noindent\textit{Critical reading.} \(W(c,\Psi)=\Psi^{\prime}(x_0)=i\chi(x_0)\)\par
\noindent\textit{Evidence.} The middle handwritten symbol is overwritten. The formula follows from c=1,cprime=0 and chi=-i Psiprime/Psi.\par
\subsection*{P000562 (both; faithful error): growth at infinity}
\noindent\textit{Transcription reads.} \(\sqrt R\sim E^N\)\par
\noindent\textit{Critical reading.} \(\sqrt R\sim E^{N+1/2}=kE^N\)\par
\noindent\textit{Evidence.} \(R\) has degree \(2N+1\), and the next source line already contains \(kE^N\).\par
\subsection*{P000562 (both; ambiguity): polynomiality implication}
\noindent\textit{Transcription reads.} Real-axis asymptotic implies entire function is polynomial\par
\noindent\textit{Critical reading.} Supply global polynomial growth, obtained through the appropriate complex estimates\par
\noindent\textit{Evidence.} An entire \(E^N+\exp(-E)\) has the same leading asymptotic on the positive real axis but is not a polynomial.\par
\subsection*{P000563 (both; ambiguity): Riccati coefficient derivative order}
\noindent\textit{Transcription reads.} \(\chi_{2n+1}=Q(u,\ldots,u^{(n-1)})\)\par
\noindent\textit{Critical reading.} Before integration by parts the maximum derivative order is \(2n\)\par
\noindent\textit{Evidence.} The recurrence gives \(\chi_3=u_{xx}-u^2\). The source exponent \(n-1\) is itself palaeographically provisional.\par
\subsection*{P000564 (both; faithful error): strong system component}
\noindent\textit{Transcription reads.} \(\sum_hD_j(a_{ij}^{hk}D_i u^k)=0\)\par
\noindent\textit{Critical reading.} \(\sum_hD_j(a_{ij}^{hk}D_i u^h)=0\)\par
\noindent\textit{Evidence.} The preceding weak form sums the h component of Du and leaves k as the equation index.\par
\subsection*{P000564 (both; faithful error): derivative index and coefficient argument}
\noindent\textit{Transcription reads.} \(D_i=\partial/\partial x\), \(t\in\mathbb R^n\)\par
\noindent\textit{Critical reading.} \(D_i=\partial/\partial x_i\), \(t\in\mathbb R^N\)\par
\noindent\textit{Evidence.} The system has n spatial variables and N unknown components. The literal dimensions are visible in the scan.\par
\subsection*{P000564 (both; faithful error): Basel integral and power-series explanation}
\noindent\textit{Transcription reads.} \(\int_0^1\log|1-t|\,dt/t=\pi^2/6\).\par
\noindent\textit{Critical reading.} \(\int_0^1\log(1-t)\,dt/t=-\pi^2/6\); put a minus sign before the positive series as well.\par
\noindent\textit{Evidence.} The scan has no minus sign; the integrand is strictly negative on \((0,1)\), and \(-\log(1-t)\) has the positive Taylor series. Faithful notebook error.\par
\subsection*{P000565 (both; faithful error): regularity class}
\noindent\textit{Transcription reads.} \(u\in W^{1,1}\)\par
\noindent\textit{Critical reading.} Use the energy class \(W^{1,2}\) for the stated uniformly elliptic system\par
\noindent\textit{Evidence.} The immediately preceding weak form and the next definition use W1,2. The standard two-dimensional regularity theorem needs its energy hypothesis.\par
\subsection*{P000565 (both; faithful error): first factor of Giusti tensor}
\noindent\textit{Transcription reads.} \(1/(n-2)\)\par
\noindent\textit{Critical reading.} Reconstruction: \(4/(n-2)\)\par
\noindent\textit{Evidence.} The scan clearly gives 1. The literal residual is \(-(3/2)(n-1)x_k/r^3\); both factors \(4/(n-2)\) make it zero and agree with Mingione (2006), equation 3.3.\par
\subsection*{P000565 (both; ambiguity): Dirichlet uniqueness}
\noindent\textit{Transcription reads.} \(\exists!\) singular solution\par
\noindent\textit{Critical reading.} Existence verified after tensor correction; uniqueness is not asserted\par
\noindent\textit{Evidence.} The homogeneous-system calculation proves one weak solution only. No matching uniqueness theorem or operator monotonicity was verified.\par
\subsection*{P000566 (both; faithful error): Periodicity in the opening display}
\noindent\textit{Transcription reads.} \(u(x+T)=u(y)\).\par
\noindent\textit{Critical reading.} \(u(x+T)=u(x)\).\par
\noindent\textit{Evidence.} The handwritten right side clearly has y. The next page explicitly writes u(x+T)=u(x); the discussion requires a periodic potential.\par
\subsection*{P000566 (both; faithful error): Extrema after gaps close}
\noindent\textit{Transcription reads.} Maxima and minima are eventually equal to 1.\par
\noindent\textit{Critical reading.} Maxima are at+1 and minima at−1; equivalently the absolute extrema are1.\par
\noindent\textit{Evidence.} The corrected scan sketch oscillates between two levels. Infinite periodic/antiperiodic eigenvalues remain; closing gaps does not make the discriminant constant.\par
\subsection*{P000567 (both; ambiguity): Scope of the opening radical}
\noindent\textit{Transcription reads.} A radical drawn above the numerator of \((\tau^2-1)/\prod(E-E_j)\).\par
\noindent\textit{Critical reading.} Interpret as \(\sqrt{(\tau^2-1)/\prod(E-E_j)}\), with the entire quotient under the radical.\par
\noindent\textit{Evidence.} The overbar ends over the numerator but the left radical stroke descends beside the denominator. Only the full-quotient reading is single-valued at each simple band endpoint. Mathematical scope is inferred, not visually certain.\par
\subsection*{P000567 (both; faithful error): Allowed and forbidden zones sentence}
\noindent\textit{Transcription reads.} Allowed and forbidden zones---where \(|\lambda|=1\).\par
\noindent\textit{Critical reading.} Allowed zones have \(|\lambda|=1\); in open forbidden intervals the multipliers are real reciprocal and have moduli different from 1.\par
\noindent\textit{Evidence.} The scan has the same abbreviated sentence; P000566 gives the correct alternative explicitly. Faithful shorthand cannot be read as a criterion for both kinds of interval.\par
\subsection*{P000569-P000570 (both; faithful error): Mapping-space target in the homology and signed-configuration displays}
\noindent\textit{Transcription reads.} \(\mathrm{Map}_k(S^n,S^n)\) with no basepoint qualification.\par
\noindent\textit{Critical reading.} Read the electric-field and Euclidean scanning target as based maps \(\Omega^n_kS^n\); for signed particles use based maps to the corrected quotient target.\par
\noindent\textit{Evidence.} The notebook explicitly fixes infinity to zero. Segal 1973, p.213, specifies base-point-preserving maps. Unbased mapping spaces are different; this is a missing hypothesis, not a changed handwritten symbol.\par
\subsection*{P000570 (both; faithful error): Tentative signed-particle sphere target}
\noindent\textit{Transcription reads.} \(C_k^{\pm}(\mathbb R^n)\simeq\mathrm{Map}_k(S^n,S^n)\).\par
\noindent\textit{Critical reading.} For the specified cancellation topology, use \(\Omega^n_k((S^n\times S^n)/\Delta S^n)\).\par
\noindent\textit{Evidence.} The scan labels the first formula a hypothesis and follows it with a \(2n\)-dimensional \(Y_n\). McDuff\textquotesingle s target is given explicitly in Kallel 1998, Example 11.5. The correction is to a provisional mathematical assertion, not a transcription error.\par
\subsection*{P000571 (both; provenance): End of Segal configuration-space record}
\noindent\textit{Transcription reads.} Catalogue record 344 stops atP000570.\par
\noindent\textit{Critical reading.} The gluing proof continues through the opening ofP000571, before the next meeting heading onP000572.\par
\noindent\textit{Evidence.} P000571 continues both restriction squares and refers to the preceding theorem; no new speaker or date intervenes.\par
\subsection*{P000573 (both; ambiguity): Arnold numerator and stated degree bound}
\noindent\textit{Transcription reads.} \(R_{n-4}\), degree at most \(n-4\).\par
\noindent\textit{Critical reading.} Retain \(n-4\) as the source reading; annotate that the classical theorem permits degree at most \(2n-3\), assuming simple affine intersections and no intersections at infinity.\par
\noindent\textit{Evidence.} The scan clearly shows n-4 both in the subscript and explanatory line. Its weaker bound is not mathematically false. This row prevents an unjustified editorial replacement with the sharper familiar formula.\par
\subsection*{P000573-P000574 (both; provenance): Date and passage boundary of the lattice-gauge discussion}
\noindent\textit{Transcription reads.} Catalogue record 345:30 September 1974, pagesP000568--P000573.\par
\noindent\textit{Critical reading.} Keep the catalogue date, but locate the detailed formulas on P000573--P000574 under the 7 October 1974 manuscript heading; only the topic announcement is on P000569 under 30 September.\par
\noindent\textit{Evidence.} Source headings and speaker transitions are visible in the scans. The final pressure and phase paragraph onP000574 continues Sinai\textquotesingle s formula at the bottom ofP000573.\par
\subsection*{P000573 (both; faithful error): Parenthetical operation in time}
\noindent\textit{Transcription reads.} Fourier transform in time.\par
\noindent\textit{Critical reading.} Imaginary-time continuation, with a transfer-matrix interpretation when justified.\par
\noindent\textit{Evidence.} The Russian handwriting says Fourier. Wilson 1974 explicitly describes imaginary rather than real times; a Fourier transform does not perform Wick rotation. This is a faithful conceptual shorthand, not a translation mistake.\par
\subsection*{P000573-P000574 (both; faithful error): Scope of the I0 partition formula}
\noindent\textit{Transcription reads.} \(\log Z=V\log I_0(\beta)\) immediately following an \(SU(2)\) example.\par
\noindent\textit{Critical reading.} For normalized Haar measure this formula requires the subsequent \(U(1)\) angular specialization, an open planar region, and \(V\) equal to the number of plaquettes. The fundamental-trace \(SU(2)\) factor is \(I_1(2\beta)/\beta\).\par
\noindent\textit{Evidence.} The next page writes \(B=e^{iA}\) and a scalar cosine of the discrete curl. Direct one-plaquette Haar integration distinguishes the two groups; periodic boundaries impose an additional global constraint.\par
\subsection*{P000574 (both; faithful error): Area and contour-length comparison}
\noindent\textit{Transcription reads.} The correlation is proportional to area in one phase and length in another.\par
\noindent\textit{Critical reading.} Interpret area/perimeter growth for \(-\log|\langle W(C)\rangle|\); specify the model before asserting two phases.\par
\noindent\textit{Evidence.} Wilson 1974, SectionsII andIV, uses exponential suppression of the loop expectation. The finite-beta planar U(1) example has an area law throughout, so the literal unqualified phase claim is not justified.\par
\subsection*{P000575 (both; faithful error): posets as categories}
\noindent\textit{Transcription reads.} category with at most one morphism between two objects\par
\noindent\textit{Critical reading.} skeletal thin category; otherwise a preorder\par
\noindent\textit{Evidence.} Scan P000575 says partially ordered. Distinct uniquely isomorphic objects show the missing antisymmetry; mathematical clarification, not a mistranscribed word.\par
\subsection*{P000575-P000578 (both; provenance): record 348 page envelope}
\noindent\textit{Transcription reads.} P000574-P000580\par
\noindent\textit{Critical reading.} substantive Segal text: P000575-P000578\par
\noindent\textit{Evidence.} The 14/X-74 Segal heading starts on P000575; P000578 ends the bundle paragraph. Adjacent pages contain other talks.\par
\subsection*{P000576 (both; ambiguity): building definition}
\noindent\textit{Transcription reads.} category of proper subspaces\par
\noindent\textit{Critical reading.} category of nonzero proper subspaces\par
\noindent\textit{Evidence.} The phrase omits nonzero, but the same scan explicitly writes 0 strictly contained in W below. Retain the shorthand and annotate its convention.\par
\subsection*{P000576 (both; ambiguity): sphere-cover components}
\noindent\textit{Transcription reads.} \(x_i\) in a connected component of \(W_1\cap W_2\)\par
\noindent\textit{Critical reading.} use contractible open components \(X_i\) in the cover reconstruction\par
\noindent\textit{Evidence.} The scan uses membership signs and small x. Interpreting them as open components is editorial, not a secure transcription correction.\par
\subsection*{P000577 (both; ambiguity): Lemmas 1 and 2}
\noindent\textit{Transcription reads.} unqualified subset and quotient formulas\par
\noindent\textit{Critical reading.} CW subcomplexes or suitable cofibration hypotheses; nonempty intersections\par
\noindent\textit{Evidence.} The scan has no point-set hypotheses. Hatcher 2002 Proposition 0.17 and Examples 0.13--0.14 specify the condition used.\par
\subsection*{P000577 (both; faithful error): Steinberg module}
\noindent\textit{Transcription reads.} \(H_{n-2}(|C|;\mathbb Z)\)\par
\noindent\textit{Critical reading.} \(\widetilde H_{n-2}(|C|;\mathbb Z)\)\par
\noindent\textit{Evidence.} Scan P000577 has no tilde. For dimension two over \(\mathbb F_q\), ordinary \(H_0\) has rank \(q+1\) but the Steinberg module has rank \(q\): a faithful mathematical omission.\par
\subsection*{P000579 (both; faithful error): Yang--Mills action display}
\noindent\textit{Transcription reads.} \(\operatorname{Sp}(\partial_\alpha b_{\alpha\beta}-\partial_\beta b_\alpha-[b_\alpha,b_\beta]^2)\).\par
\noindent\textit{Critical reading.} Use the trace or invariant inner product of the square of the whole curvature \(\partial_\alpha b_\beta-\partial_\beta b_\alpha\pm[b_\alpha,b_\beta]\), with paired indices and the chosen sign convention.\par
\noindent\textit{Evidence.} The scan has an extra alpha on the first potential and the exponent is adjacent to the commutator bracket. Yang--Mills 1954, Eq.(4) and its Lagrangian, fixes the mathematical structure; this is an emendation of an incomplete notebook formula.\par
\subsection*{P000579 (both; faithful error): Lie-algebra fibre example}
\noindent\textit{Transcription reads.} The fibre is a Lie algebra, for example SU(2).\par
\noindent\textit{Critical reading.} The Lie algebra is \(\mathfrak{su}(2)\); SU(2) is its compact Lie group.\par
\noindent\textit{Evidence.} The scan uses uppercase SU(2) in a sentence about an algebra. The surrounding connection discussion requires a Lie-algebra-valued one-form.\par
\subsection*{P000580 (both; provenance): Functional Analysis volume 2 number 1 citation}
\noindent\textit{Transcription reads.} The page cites Functional Analysis, volume 2, number 1 beside the current-group programme.\par
\noindent\textit{Critical reading.} Retain this citation as the 1968 Gelfand--Graev precursor; identify the detailed tensor and logarithmic-kernel construction with Vershik--Gelfand--Graev 1973, RMS28(5),87--132.\par
\noindent\textit{Evidence.} Math-Net faa2751 identifies the 1968 paper, received 16 October 1967. VGG1973 cites it as reference1 and says that its integral construction goes beyond the earlier finite-support examples.\par
\subsection*{P000581 (both; faithful error): Endpoint invariant-integral differential}
\noindent\textit{Transcription reads.} \(\int f(t)\,dx\).\par
\noindent\textit{Critical reading.} \(\int f(t)\,dt\).\par
\noindent\textit{Evidence.} The scan explicitly uses t in the integrand and dx on the left, but dt in the change-of-variables expression on the right. The endpoint invariant functional integrates the line coordinate t.\par
\subsection*{P000582 (both; faithful error): Arguments of the exponential-kernel factors}
\noindent\textit{Transcription reads.} Both factors are written \(\widetilde f_i(x,t)\), with differentials \(dt_1dt_2dx\).\par
\noindent\textit{Critical reading.} Use \(\widetilde f_1(x,t_1)\widetilde f_2(x,t_2)\) for real labels; conjugate the second factor in the complex-label formulation.\par
\noindent\textit{Evidence.} The scan has the repeated unsubscripted t. The matching VGG1973, Section 6.3, uses separate line variables and real functions, with complex coefficients only in the formal span.\par
\subsection*{P000583 (both; faithful error): Centering density}
\noindent\textit{Transcription reads.} \(\widetilde f(x,t)=f(x,t)-1/[\pi(1+x^2)]\).\par
\noindent\textit{Critical reading.} \(\widetilde f(x,t)=f(x,t)-1/[\pi(1+t^2)]\).\par
\noindent\textit{Evidence.} The scan unmistakably has x squared. The constant-in-t subtraction is nonintegrable and cannot make the t-integral zero. VGG1973, Section 6.3, p.128, explicitly uses the Cauchy density in t.\par
\subsection*{P000584 (both; faithful error): Frequency-projection assertion}
\noindent\textit{Transcription reads.} \(u\in\mathcal D'(\Omega)\); \(\pi'\mathrm{WF}(u)\) described using its global Fourier transform.\par
\noindent\textit{Critical reading.} Apply the frequency-projection statement to \(u\in\mathcal E'(\mathbb R^n)\), or to compactly supported localizations.\par
\noindent\textit{Evidence.} The scan confirms the unqualified sentence. \citet[Proposition 2.5.5]{hormander1971fourier} has compact support; this is a mathematical qualification, not a mistranscribed word.\par
\subsection*{P000584 (both; ambiguity): One-sided boundary-value Fourier sign}
\noindent\textit{Transcription reads.} Upper-half-plane continuation is paired with exclusion of \(\xi_0>0\), without a Fourier convention.\par
\noindent\textit{Critical reading.} With \(\widehat u(\xi)=\int e^{-ix\xi}u(x)\,dx\), upper boundary values exclude negative covectors. The written sign is consistent with the opposite Fourier convention; retain this ambiguity.\par
\noindent\textit{Evidence.} \citet[Appendix A, equation (A.4)]{bohrlefeuvrepaternain2024invariant} explicitly checks the positive wavefront of \((x+i0)^{-1}\). The notebook does not declare the Fourier sign.\par
\subsection*{P000585 (both; faithful error): Configuration count cylinder}
\noindent\textit{Transcription reads.} \(\{K:|X\cap K|=n\}\)\par
\noindent\textit{Critical reading.} \(\{X:|X\cap K|=n\}\), with \(K\) fixed.\par
\noindent\textit{Evidence.} The scan visibly binds K. The surrounding prose defines a measure on configurations X.\par
\subsection*{P000585 (both; faithful error): Finite-range independence slogan}
\noindent\textit{Transcription reads.} Finite interaction: independence.\par
\noindent\textit{Critical reading.} Conditional independence of the interior and distant exterior given the intervening collar.\par
\noindent\textit{Evidence.} The collar diagram and the next dated session identify the conditional Gibbs interpretation.\par
\subsection*{P000585-P000586 (both; faithful error): Gibbs Radon--Nikodym density}
\noindent\textit{Transcription reads.} \(dG/d\mu=e^{-H}\)\par
\noindent\textit{Critical reading.} \(dG/d\mu=Z^{-1}e^{-H}\), when \(0<Z<\infty\).\par
\noindent\textit{Evidence.} No partition factor is present in the scan; P000587 supplies it.\par
\subsection*{P000588 (both; faithful error): Compactness in the Gibbs existence sketch}
\noindent\textit{Transcription reads.} The space of configurations in a bounded region is compact.\par
\noindent\textit{Critical reading.} A compactness or tightness argument must use the hard-core bound and a specified topology; unrestricted locally finite configurations in a bounded region are not compact.\par
\noindent\textit{Evidence.} The scan states compactness without qualifying the configuration class; the next limiting step also needs boundary regularity.\par
\subsection*{P000589 (both; faithful error): Number of Gibbs distributions at coexistence}
\noindent\textit{Transcription reads.} \(k\) Gibbs distributions, with \(k-1\) and \(k-2\) along adjacent strata.\par
\noindent\textit{Critical reading.} \(k\), \(k-1\), and \(k-2\) stable pure periodic phases in the local phase diagram.\par
\noindent\textit{Evidence.} The scan says distributions. Convexity produces infinitely many Gibbs measures from any two distinct ones; \citet[Theorem 1]{pirogovsinai1975phase} is a pure-phase theorem.\par
\subsection*{P000589 (both; ambiguity): Ground-state perturbation comparison}
\noindent\textit{Transcription reads.} \(G(U_\mu)\le G(U)\)\par
\noindent\textit{Critical reading.} Possibly \(G(U_\mu)\subseteq G(U)\), or a cardinality comparison; retain uncertainty.\par
\noindent\textit{Evidence.} The scan contains a relation resembling an inequality between set-valued symbols. Its intended interpretation is not recoverable with certainty.\par
\subsection*{P000591 (both; faithful error): Phase-parametrization local diffeomorphism}
\noindent\textit{Transcription reads.} \(\lambda_\Phi\) is a local diffeomorphism into \(T^*\Omega\setminus0\).\par
\noindent\textit{Critical reading.} It is an immersion and a local diffeomorphism onto its Lagrangian image.\par
\noindent\textit{Evidence.} Dimension count: the critical set has dimension n and the cotangent target dimension 2n; the scan omits the image qualification.\par
\subsection*{P000592 (both; faithful error): Global type of Lagrangian leaves}
\noindent\textit{Transcription reads.} Leaves are \(\mathbb R^n\) or tori.\par
\noindent\textit{Critical reading.} Leaves have a natural affine structure; neither completeness nor a global topological dichotomy follows.\par
\noindent\textit{Evidence.} The scan clearly states the dichotomy. The standard cylinder foliation is a direct counterexample.\par
\subsection*{P000592 (en; ambiguity): Example after the leaf corollary}
\noindent\textit{Transcription reads.} “Kokas.” construction.\par
\noindent\textit{Critical reading.} Likely “cotangent bundle,” from the abbreviation for кокасательное расслоение; preserve uncertainty.\par
\noindent\textit{Evidence.} The scan resembles “Кокас. рассл.” The context is the standard fibrewise Lagrangian foliation of a cotangent bundle.\par
\subsection*{P000593 (both; faithful error): Normal-form hypothesis}
\noindent\textit{Transcription reads.} \(p=0,\ dp\ne0\) alone imply the homogeneous FIO normal form.\par
\noindent\textit{Critical reading.} Also require the nonradial real-principal-type hypothesis; use \(D_{X_n}=-i\partial_{X_n}\) with the stated Fourier convention.\par
\noindent\textit{Evidence.} The scan lists only \(dp\ne0\). \citet[Proposition 6.1.4]{duistermaathormander1972fourierii} explicitly requires the Hamilton field not to be tangent to the cone axis.\par
\subsection*{P000594 (both; provenance): Meaning of Bernstein’s WF}
\noindent\textit{Transcription reads.} WF of a left differential-operator ideal.\par
\noindent\textit{Critical reading.} Algebraic characteristic variety of a nonzero quotient module; distinguish it from a distributional wavefront set.\par
\noindent\textit{Evidence.} The scan defines a common symbol-zero set, and the dimension bound is Bernstein’s module inequality.\par
\subsection*{P000596 (both; faithful error): Conditioning in the Markov-chain example}
\noindent\textit{Transcription reads.} \(\Pr(x_n\mid x_k,\ k\le n)\)\par
\noindent\textit{Critical reading.} \(\Pr(x_n\mid x_k,\ k<n)\)\par
\noindent\textit{Evidence.} The scan clearly shows a non-strict inequality. Including the variable being predicted makes its conditional law degenerate; the standard Markov property conditions only on its past.\par
\subsection*{P000597 (both; provenance): Name paired with Nielsen}
\noindent\textit{Transcription reads.} Olsen / Олсен.\par
\noindent\textit{Critical reading.} Olesen / Олесен (Poul Olesen).\par
\noindent\textit{Evidence.} The scan omits the medial e. The publisher metadata for Nielsen and Olesen, Vortex-line models for dual strings (1973), identifies the intended collaborator.\par
\subsection*{P000598 (both; faithful error): Example with dz/z}
\noindent\textit{Transcription reads.} The current difference is \(d[\eta](\varphi)-[\eta](d\varphi)\); the contour has \(\varphi\,dz\), and the answer is \(\varphi(0)\).\par
\noindent\textit{Critical reading.} Use \(d[\eta]-[d_{\mathrm{reg}}\eta]\), contour \(\varphi\,dz/z\), and \(2\pi i\varphi(0)\), or normalize \(\eta=dz/(2\pi iz)\).\par
\noindent\textit{Evidence.} The scan confirms these symbols. With degree-one current signs the written difference vanishes; the positively oriented contour computes the stated corrected factor.\par
\subsection*{P000599 (both; faithful error): Curvature sign in ordinary derivative}
\noindent\textit{Transcription reads.} \([d\eta]=c(L_D)\), followed by \([D]=d[\eta]+c(L_D)\).\par
\noindent\textit{Critical reading.} For \(\eta=(2\pi i)^{-1}\partial\log\|\sigma\|^2\) and standard \(c_1=-dd^c\log h\), the ordinary derivative off \(D\) is \(-c_1\); the current identity is \(dd^c\log\|\sigma\|^2=[D]-c_1\).\par
\noindent\textit{Evidence.} Both notebook lines are visible and inconsistent together. \citet[Proposition 1.14]{griffithsking1973nevanlinna} fixes the convention.\par
\subsection*{P000599 (both; faithful error): Bézout count}
\noindent\textit{Transcription reads.} Number of intersection points equals the product of degrees.\par
\noindent\textit{Critical reading.} Sum of intersection multiplicities equals the product of degrees, assuming no common curve component.\par
\noindent\textit{Evidence.} The proof via Chern classes computes the intersection number; the omitted hypotheses are mathematical qualifications.\par
\subsection*{P000600 (both; provenance): Attribution of extension to singular sets}
\noindent\textit{Transcription reads.} Lelong, 1953.\par
\noindent\textit{Critical reading.} Checked full publication: Lelong, 1957, Theorems 7–8; the exact intended 1953 source remains unverified.\par
\noindent\textit{Evidence.} The scan clearly reads 1953. The verified paper itself lists a related 1953 contribution, so the notebook date is not declared a transcription error.\par
\subsection*{P000600 (both; faithful error): Pullback divisor identity}
\noindent\textit{Transcription reads.} \(d[\eta]-[\omega]=[D]\), and similarly after pullback.\par
\noindent\textit{Critical reading.} \([D]=d[\eta]+[\omega]\) for the convention used on P000599, equivalently \(dd^c\log\|\sigma\|^2=[D]-\omega\).\par
\noindent\textit{Evidence.} The scan has a minus sign; \citet[Proposition 1.14]{griffithsking1973nevanlinna} and the notebook's preceding page give the corrected relation.\par
\subsection*{P000600 (both; faithful error): Ball/sphere counting integral}
\noindent\textit{Transcription reads.} \(\int_{S(r)}\varphi^{n-1}\wedge f^*\omega=\int_{D_f\cap B(r)}\varphi^{n-1}\).\par
\noindent\textit{Critical reading.} The curvature integral is over \(B(r)\). The boundary term
\[\int_{S(r)}d^c\log\|f^*\sigma\|^2\wedge\varphi^{n-1}\]
must be added.\par
\noindent\textit{Evidence.} The domain letter is faint, but the displayed equality omits the logarithmic boundary term regardless. The sphere interpretation has the wrong form degree.\par
\subsection*{P000600 (both; ambiguity): Coefficient in the d-c formula}
\noindent\textit{Transcription reads.} \(d^c=(4\pi)^{-1}(\bar\partial-\partial)\), with a cancelled symbol before the parentheses.\par
\noindent\textit{Critical reading.} For the real convention used in the reconstruction, \(d^c=(i/4\pi)(\bar\partial-\partial)\). Preserve the cancelled coefficient as uncertain in the diplomatic transcription.\par
\noindent\textit{Evidence.} The scan shows a cancellation at the coefficient; the missing imaginary factor is supplied only for the consistent mathematical reconstruction.\par
\subsection*{P000601 (both; faithful error): Power of the Euclidean radius}
\noindent\textit{Transcription reads.} \(\rho^{n+1}\) in the denominator of the integrated divisor count.\par
\noindent\textit{Critical reading.} \(\rho^{2n-1}\) for ordinary Euclidean radius and raw mass \(\int_{D_f\cap B(\rho)}\varphi^{n-1}\), up to a common dimensional factor.\par
\noindent\textit{Evidence.} The scan clearly has n+1. The repaired exponent reduces to the classical dρ/ρ in dimension one, whereas the written exponent gives dρ/ρ².\par
\subsection*{P000602 (both; faithful error): Repeated first-main-theorem normalization}
\noindent\textit{Transcription reads.} \(N(D_f,r)\le T(f,r)\), and an outer \(dt\) integral with inner \(B(r)\).\par
\noindent\textit{Critical reading.} Retain \(O(1)\) unless it has been absorbed by a declared additive normalization, and use \(B(t)\) in the inner radius integral.\par
\noindent\textit{Evidence.} The preceding page retains \(O(1)\); \citet[Proposition 5.11]{griffithsking1973nevanlinna} states the normalized identity. The scan confirms the two different radius letters.\par
\subsection*{P000602 (both; ambiguity): Two Ricci normalizations}
\noindent\textit{Transcription reads.} \(\operatorname{Ric}\Omega=i\partial\bar\partial\log g\) first, then \((i/2\pi)\partial\bar\partial\log a\).\par
\noindent\textit{Critical reading.} Declare which normalization is used. With \(dd^c=(i/2\pi)\partial\bar\partial\), the one-dimensional coefficient is minus Gaussian curvature divided by \(2\pi\).\par
\noindent\textit{Evidence.} Both factors are explicit in the scan; the minus-Gaussian comment is correct for the first formula but requires the extra factor for the second.\par
\subsection*{P000602 (both; provenance): Closing Latin names and Cowen reference}
\noindent\textit{Transcription reads.} Herman → Joahim Weil; Griffiths, Cowen.\par
\noindent\textit{Critical reading.} Likely Hermann and Joachim Weyl; the Cowen–Griffiths 1976 paper provides a thematic bibliographical lead, not an identified eigenfunction theorem.\par
\noindent\textit{Evidence.} The scan was re-examined. \citet[Introduction, pp.~93--95]{cowengriffiths1976curves} explicitly discusses H. and J. Weyl in the theory of holomorphic curves. The inference remains marked as such.\par
